\documentclass[11pt,a4paper]{article}
\usepackage[T1]{fontenc}
\usepackage[utf8]{inputenc}
\usepackage{lmodern}
\ifdefined\pdfmapfile
  \pdfmapfile{=lm.map}
  \pdfmapfile{+cm.map}
  \pdfmapfile{+symbols.map}
  \pdfmapfile{+euler.map}
  \pdfmapfile{+rsfs.map}
\fi
\usepackage[a4paper,margin=27mm]{geometry}
\usepackage{amsmath,amssymb,amsthm,mathtools,mathrsfs}
\usepackage{booktabs,tabularx,longtable,array,enumitem}
\usepackage{microtype,xcolor,etoolbox}
\usepackage{hyperref}
\hypersetup{colorlinks=true,linkcolor=blue!45!black,citecolor=blue!45!black,
urlcolor=blue!45!black,pdfauthor={Research record prepared with ChatGPT},
pdftitle={The ordered potential quotient, delta tests, and relative Bedford--Taylor theory}}
\usepackage[capitalise,noabbrev]{cleveref}
\numberwithin{equation}{section}
\newtheorem{theorem}{Theorem}[section]
\newtheorem{lemma}[theorem]{Lemma}
\newtheorem{proposition}[theorem]{Proposition}
\newtheorem{corollary}[theorem]{Corollary}
\theoremstyle{definition}
\newtheorem{definition}[theorem]{Definition}
\newtheorem{example}[theorem]{Example}
\theoremstyle{remark}
\newtheorem{remark}[theorem]{Remark}
\newcommand{\Cb}{\mathcal C_b}
\newcommand{\R}{\mathbf R}
\newcommand{\Q}{\mathbf Q}
\newcommand{\Z}{\mathbf Z}
\newcommand{\PSH}{\operatorname{PSH}}
\newcommand{\BTcap}{\operatorname{Cap}}
\newcommand{\supp}{\operatorname{supp}}
\newcommand{\dd}{d'd''}
\newcommand{\an}{\mathrm{an}}
\setlist{topsep=4pt,itemsep=2pt,parsep=0pt}
\setlength{\parindent}{1.3em}
\setlength{\parskip}{2pt}
\setlength{\emergencystretch}{2em}
\makeatletter
\patchcmd{\l@section}{1.5em}{2.1em}{}{}
\renewcommand*\l@subsection{\@dottedtocline{2}{1.5em}{3.6em}}
\makeatother
\title{The ordered potential quotient,\\
delta tests, and relative\\
Bedford--Taylor theory\\[5pt]
\large Sixth research record in non-Archimedean Bedford--Taylor theory}
\author{Research record prepared with ChatGPT\\
\small Relative to the authorized manuscripts R, B, and A2}
\date{10 November 2026 (UTC)}
\begin{document}
\maketitle
\begin{abstract}
This record continues the arbitrary-background bounded Bedford--Taylor
calculus and actual bounded-classical potential quotient proved in
R4--R5. Its main further result identifies the positive cone of the
actual bounded classical difference class: if $a,b\in\mathcal C_b$,
then $D[a-b]\ge0$ in the original CLD cone if and only if $a-b$ is
itself classical psh. The proof establishes order transfer on every
complete extension and analytic test map, retains actual curve endpoint
values, and recovers ambient upper semicontinuity through proper-fiber
maxima of finite Shilov interpolants. Curvature order and continuous
coefficient relations therefore transfer to every fixed allowed
positive closed background, including discontinuous potential slots.
The resulting mixed order estimates give qualified total-variation
stability. Arbitrary assigned bounded Borel representatives remain
outside these assertions.

In the actual delta-test category, every total-degree $(1,1)$ delta-form
acts over every allowed background with ordinary locality. Constant
real cycle relations act in all degrees, and every real delta-form has
an exact finite local generation by piecewise smooth coefficients and
constant cycles. Full actual analytic delta-testing is constructed for
the local weak closure of positive closed delta-form backgrounds,
including all their bounded BT tails, smooth regular flags, singular
effective Cartier divisors, integrable extended classical potential
curvatures, and integrable diffuse barycenters. The higher-degree analytic-relation
kernel for an arbitrary background remains open.

An actual bounded classical potential below the curvature of a
continuous potential is automatically continuous. Normalized families
with a compact continuous family of upper potentials are compact in
local uniform topology. Their actual mixed images are weakly closed on
the same domain, and their represented curvature quotients have joint
weak stability even with a varying positive closed background. A local
finite Shilov normalization also gives weighted stability for the
specified representatives; it does not identify arbitrary weights by
curvature alone.

In remaining degree one, sharp Hardy and strong capacitary constants
are respectively one and four. The universal finite-energy threshold
for concave capacity-growth moduli is determined. A compact Gauss set
has zero relative BT capacity but positive energy capacity; a measure
on it has a bounded continuous energy solution but no local bounded
classical difference potential. Exact convex-envelope and diffuse
carrier capacities give further qualified Borel regularity results.
Exact continuous classical interpolation on a fast polar Cantor set
gives explicit capacities for arbitrary bounded Borel radial graphs
and filled tubes over that diffuse parameter set.

For Dirichlet theory, completed degree-one feasibility has an exact
energy-dual criterion, while higher-degree variation criteria retain
explicit feasibility hypotheses. Counterexamples rule out unrestricted
relative Hessian root concavity and show that the actual zero-boundary
classical Monge--Ampere image can be nonconvex in every remaining degree
at least two, even for sums of smooth analytic-subspace integration
currents. All new proofs, direct consequences, failed routes and
unresolved gates have separate status labels.
\end{abstract}
\section*{Reading guide}
\addcontentsline{toc}{section}{Reading guide}
Section~\ref{sec:r6-inputs} fixes the hypotheses and frozen dependencies.
The main ordered-quotient advance is in
Sections~\ref{sec:r6-potential-order-functoriality}--\ref{sec:r6-positive-dsh-representatives};
its module and stability consequences follow immediately afterward.
The delta-extension conclusions are in
Sections~\ref{sec:r6-delta-module-generation}--\ref{sec:r6-delta-barycenters}.
The sharp arbitrary-background estimates are in
Sections~\ref{sec:r6-sharp-relative-hardy}--\ref{sec:r6-strong-capacity-trace}.
The precise changes from R5 and the next proof gates are recorded at
the end. This is a continuation with complete new proofs, rather than
a replacement for the frozen primary-literature comparison or the
authorized internal manuscripts. References marked R4 or R5 refer to
the exact numbered results retained in the complete continuation archive.
\tableofcontents
\clearpage

% BEGIN canonical R6 source: sections/round6_inputs.tex
\section{Scope, conventions, and dependencies}
\label{sec:r6-inputs}

Fix a finite extension $K/\mathbf Q_p$ with its original absolute
value. An allowed space is a good Hausdorff boundaryless quasi-smooth
$K$-analytic space of pure dimension $n$. All opens are ordinary
topological opens. Put $D=d'd''$. A background of degree $(p,p)$ is
symmetric, positive in the original middle CLD cone specified by R,
and separately $d'$- and $d''$-closed. It acts on all original compact
smooth CLD forms of complementary degree. The words positive,
strongly positive, and weakly positive are not interchanged.

The class $\mathcal C_b$ consists of actual locally bounded classical
psh functions as defined in B. Values on ambient polar sets are
retained. Products with coefficient-Radon currents use actual Borel
multiplication. Local assertions are checked on all original moments
and refinements, not just on one chosen skeleton. Signed bounded
classical differences are expanded into finite signed combinations of
positive mixed currents whenever coefficient estimates are used.

\subsection{Authorized and frozen inputs}

The internal manuscripts R and B are used as authorized inputs.
R supplies continuous local regularization, the original middle-cone
coefficient-Radon estimates, and fixed-continuous-multiplier weak
continuity. B supplies the prescribed local decreasing continuous
semipositive PL rows and their original smooth positive refinements,
the class definition, and the original $[1]$ mixed calculus. Their
statements are not reproved or globally audited in this record.
A2 remains the completed geometric adelic nefness input under its
stated number-field, integrable adelic-line, and boundary assumptions.
No part of A2 is treated as a new open problem.

The following results of the preceding records are frozen inputs.
\begin{enumerate}
\item The complete bounded BT calculus for every fixed allowed
      arbitrary background: existence, positivity, separate closedness,
      recursive and regularization independence, differentiated-slot
      symmetry, multilinearity, ordinary locality and gluing, and
      weighted joint decreasing convergence on all original tests.
\item Relative CLN estimates with finitely many original positive
      control forms, bounded-classical energy measures and Green
      identities, and their weighted convergence. R4, Lemmas 71.1
      and 71.2 and Theorem 71.4 supply normal-factor variation control
      and uniform piecewise smooth test approximation.
\item R5, Theorems 10.1--10.2, identify the actual bounded classical
      curvature kernel with continuous classical pluriharmonic
      functions. In particular
      \[
           D[u]=D[v]\quad\Longrightarrow\quad D((u-v)S)=0
      \]
      for every allowed fixed $S$. Its signed differentiated-potential
      calculus is also retained. This arbitrary-background potential
      quotient is already \texttt{PROVED}; it is not an open gate
      of this round.
\item The qualified comparison, maximum-locality, quasicontinuity,
      completed bounded calculus, and faithful relative energy-form
      results retain the domains and hypotheses of R4. In particular
      an energy solution or a completed relative potential is not
      automatically an actual ambient classical psh function.
\end{enumerate}

The full eleven-direction classical-to-non-Archimedean comparison,
with primary theorem locations, is retained in R4, Section 13,
pp.~41--50. The present record gives a precise update at the gates
it actually changes. A result marked \texttt{NEW\_PROVED} has a proof
here relative to the stated inputs; the label makes no claim of
priority over every published result. \texttt{DIRECT\_COROLLARY}
is used conservatively for deductions already supplied by the frozen
calculus. External facts and conditional hypotheses are identified
where used.

\subsection{Two distinct extension questions}

Ordinary potential independence, preservation of represented curvature
order, and testing on every delta-form are different assertions. R5
settled independence. The new bounded-difference positivity criterion
settles order, and its proof uses no higher-degree delta extension.
The latter extension question retains the separate scopes below. This round uses Mihatsch's actual analytic
delta-form sheaf $B$, as defined by compatible data on all skeletons
and all refinements. Its total bidegree is used unless a tridegree is
displayed. The external map $B\to D$ is faithful, but faithfulness
of that map does not by itself extend an arbitrary ordinary current
to the larger test sheaf.

We prove the following sufficient mechanism. Let
$\mathcal W^p(V)$ be the ordinary weak closure, on one fixed ordinary
domain $V$, of positive separately closed actual delta-form
backgrounds. Full actual delta-testing is constructed for
$\mathcal W^p_{\mathrm{loc}}$. All bounded classical BT tails of
such backgrounds stay in that class. Smooth regular flags and their
integrable diffuse barycenters give additional examples. The
equality of this class with all allowed backgrounds is not assumed.

The all-background ambient cycle action is stronger in a different
direction: constant real tropical cycles and all their ambient
relations act over every fixed allowed $S$. Actual analytic
zero-presentations introduce an additional kernel. Explicit saturated
torus and matroidal chart hypotheses remove that kernel in their
fixed-chart scope. Their locality must not be upgraded to arbitrary
ordinary subopens without another argument.

\subsection{Capacity and Dirichlet notation}

On a fixed capacity domain $\Omega$, normalized testers are actual
classical functions defined on all of $\Omega$ with $0\le u\le1$.
Raw capacity is their supremum of masses, and the star denotes the
infimum over ordinary open neighborhoods in the same domain.
C1 denotes sharp compact no-gain under the relevant weak closure;
C2 denotes raw Borel equality with that ordinary-open outer capacity.
Both remain open in general unless a theorem below supplies an
explicit additional hypothesis.

When $S$ has degree $(n-1,n-1)$, the remaining degree is one and
$\mathcal E_S(f)=\int\Gamma_S(f)$ for real original compact smooth
$f$. The sharp Hardy and strong-capacity estimates below hold for
every such allowed $S$, in every ambient dimension. Faithful closed
form and Dirichlet conclusions invoke the additional frozen form
setup where stated. For a background of smaller degree one may fix
the other bounded BT slots to obtain a degree-one tail; this does
not freeze a nonlinear family in which all slots vary.

% END source: sections/round6_inputs.tex


\part{Ordinary currents and actual delta-tests}

% BEGIN canonical R6 source: sections/delta_faithfulness_comparison.tex
\section{What ordinary equality determines in the actual delta-form category}
\label{sec:r6-delta-faithfulness}

Throughout this section, $K$ is the original finite extension of
$\mathbf Q_p$, with its original absolute value.  The analytic space $X$
is good, Hausdorff, boundaryless, and purely $n$-dimensional.  When the
relative bounded-potential calculus is used, $X$ is also quasi-smooth and
has the local hypotheses of that calculus.  Write
\[
 A^{a,b}_{c}(U)=\text{the original compactly supported CLD smooth forms},
 \qquad B^{a,b}(U)=B^{a,b}_{\mathrm{Mih}}(U).
\]
Here $B$ is the sheaf of \emph{actual} delta-forms of
\cite[Definition 4.2]{Mih24}.  Its total bidegree is used: a tridegree
$(a,b,c)$ has total bidegree $(a+c,b+c)$.  We do not identify $B$ with
the algebraic sheaf $B_{\mathrm{GK}}$ without the hypotheses of the
comparison between those two sheaves.

\subsection{The external faithfulness statement}

\begin{proposition}[Faithfulness on actual delta-forms]
\label{prop:r6-delta-faithfulness}
\textnormal{\texttt{EXTERNAL\_KNOWN}.}
For every ordinary open $U\subset X$, integration defines an injective
map
\[
 \iota_U:B^{a,b}(U)\longrightarrow D^{a,b}_{\mathrm{CLD}}(U),
 \qquad
 \langle\iota_U(\omega),\eta\rangle
       =\int_U\omega\wedge\eta,
 \quad \eta\in A_c^{n-a,n-b}(U).
\]
It commutes with $d'$ and $d''$.  The source is defined using all original
toric refinements; it is not the quotient obtained from one chosen
tropicalization alone.
\end{proposition}

The precise input is \cite[Proposition 4.11, printed pp.~70--71]{Mih24}.
The standing field and analytic-space assumptions are given at the
start of its Section 3, printed p.~42; in particular its field is any
complete non-Archimedean field, and its analytic spaces are good and
Hausdorff.  Section 4 assumes pure dimension and empty analytic boundary.
Thus no algebraic closure of the fixed $K$ is needed here.  The original
smooth-test topology is recalled in the proof of that proposition.
Definition 4.2 and Lemma 4.4, printed pp.~61--62, are the relevant
all-refinement definitions and vanishing criterion.

\begin{corollary}[Ordinary equality for two actual delta-forms]
\label{cor:r6-ordinary-to-full-delta}
\textnormal{\texttt{DIRECT\_COROLLARY}.}
If $\alpha,\beta\in B^{a,b}(U)$ and
$\iota_U(\alpha)=\iota_U(\beta)$ on all original smooth tests, then
$\alpha=\beta$ as actual delta-forms.  Consequently, for every
$\gamma\in B_c^{n-a,n-b}(U)$,
\[
                         \int_U\alpha\wedge\gamma
                           =\int_U\beta\wedge\gamma.
\]
The same conclusion holds after multiplication by any further actual
delta-form for which the total degree and support make the integral
meaningful.
\end{corollary}

\begin{proof}
Injectivity gives $\alpha-\beta=0$ in the sheaf $B$.  Its wedge product
and integration are operations on sections of that sheaf.  Applying
them to this equality proves the assertions.  No density of smooth
tests in a delta-test topology is asserted or needed.
\end{proof}

\subsection{The precise BT comparison for delta-form backgrounds}

\begin{lemma}[Multiplication by a piecewise smooth function]
\label{lem:r6-pws-multiplier-delta-current}
\textnormal{\texttt{DIRECT\_COROLLARY}.}
Let $u\in PS^{0,0}(U)$ and $\omega\in B^{a,b}(U)$.  Then $u$ is a
continuous actual function, $u\omega$ is an actual delta-form, and
\[
                         \iota_U(u\omega)=u\,\iota_U(\omega)
\]
as coefficient-Radon currents.  The right side is the ordinary Borel
multiplier, with its actual point values.
\end{lemma}

\begin{proof}
By \cite[Proposition 4.7, printed p.~66]{Mih24},
$PS^{0,0}=B^{0,0,0}$.  On a common finite skeletal presentation over
the support of a test, multiplication by $u$ multiplies each smooth
coefficient of the polyhedral current by the restriction of the
continuous function $u$; this is the product described in
\cite[Remark 4.3(1), printed p.~61]{Mih24}.  The coefficient measures
of the resulting ordinary current are therefore $u$ times the
coefficient measures of $\iota_U(\omega)$.  A finite partition of unity
and the refinement compatibility in the definition of $B$ identify
the local statements.  Equality holds for every original smooth
test, including tests requiring a further moment refinement.
\end{proof}

\begin{theorem}[Comparison on an actual delta-form background]
\label{thm:r6-bt-delta-background-comparison}
\textnormal{\texttt{DIRECT\_COROLLARY}.}
Let $\sigma\in B^{p,p}(U)$ and suppose that
$S=\iota_U(\sigma)$ is symmetric, positive in the project's original
middle CLD cone, and separately $d'$- and $d''$-closed.  Let
$u_1,\ldots,u_k\in\mathcal C_b(U)\cap PS^{0,0}(U)$, with
$0\leq k\leq n-p$.  Put $D_\delta=d'd''$ in $B$.  Then
\begin{equation}
 T_{S,k}(u_1,\ldots,u_k)
   =\iota_U\bigl(D_\delta u_1\wedge\cdots\wedge
                       D_\delta u_k\wedge\sigma\bigr).
 \label{eq:r6-bt-delta-background}
\end{equation}
In particular the delta-form on the right represents a positive,
separately closed ordinary current.  This includes semipositive PL
potentials, whenever those potentials are actual piecewise smooth
functions in the indicated original moment category.

If two such delta-form expressions represent the same ordinary BT
current, then they agree as actual delta-forms and on every
complementary compact delta-test.  If $u_0\in PS^{0,0}(U)$, the
weighted current $u_0T_{S,k}$ is also represented by the delta-form
$u_0D_\delta u_1\wedge\cdots\wedge D_\delta u_k\wedge\sigma$.
\end{theorem}

\begin{proof}
The separate closedness of $S$ and Proposition
\ref{prop:r6-delta-faithfulness} imply
$d'\sigma=d''\sigma=0$ in $B$.  Suppose inductively that the product
with $u_2,\ldots,u_k$ is represented by the separately closed
delta-form
\[
                 \tau=D_\delta u_2\wedge\cdots\wedge
                                      D_\delta u_k\wedge\sigma.
\]
The ordinary recursive definition and Lemma
\ref{lem:r6-pws-multiplier-delta-current} give
\[
 \begin{split}
 T_{S,k}(u_1,\ldots,u_k)
      &=D\bigl(u_1\iota_U(\tau)\bigr)\\
      &=D\,\iota_U(u_1\tau)
       =\iota_U\bigl(D_\delta(u_1\tau)\bigr)
       =\iota_U\bigl(D_\delta u_1\wedge\tau\bigr).
 \end{split}
\]
The last equality is the Leibniz rule with both derivatives of
$\tau$ zero.  This proves the induction, starting with $\tau=\sigma$.
Positivity and closedness are supplied by the already established
ordinary BT theorem, not by an unidentified positivity convention
for arbitrary delta-forms.  The full delta-test equality follows
from Corollary \ref{cor:r6-ordinary-to-full-delta}; the weighted
assertion follows from the multiplier lemma.
\end{proof}

\begin{remark}[What this does and does not say for arbitrary backgrounds]
\label{rem:r6-arbitrary-background-delta-gate}
For every allowed fixed positive closed $S$, even when $S$ is not an
actual delta-form, the currents
$D(uS)$ and $T_{S,k}(u_1,\ldots,u_k)$ above exist by the ordinary
BT theory.  Their individual differentiated slots depend only on
$D[u_i]$, by the R5 potential quotient theorem.  This does not by
itself construct a functional on all delta-tests.

For example, a formal expression
$\sum_j\eta_j\wedge D_\delta u_{j1}\wedge\cdots\wedge
D_\delta u_{jr}$, with admissible slots
$u_{jl}\in\mathcal C_b(U)\cap PS^{0,0}(U)$ and compactly supported
original smooth coefficients $\eta_j$ of complementary degree,
can be paired with $S$ by the corresponding finite sum of ordinary
BT pairings.  To call this an action on the \emph{delta-form represented
by that expression}, one must additionally prove that every zero
relation between such expressions annihilates $S$.  Independence
of one potential in one slot is only the linear part of this
relation problem.  Proposition \ref{prop:r6-delta-faithfulness}
does not supply a product of an arbitrary current with a delta-form.

An arbitrary ordinary current is, by definition, a functional on
$A_c$, not on $B_c$.  A full delta-test identity involving it is
therefore not a statement until the enlarged action has been
defined and shown independent of all presentations and refinements.
This is distinct from the already settled comparison for two
\emph{actual delta-forms}.
\end{remark}

\subsection{The algebraic GK category is a separate input}

In \cite{GK17}, the standing base field is algebraically closed,
complete, and nontrivially valued, and the space is the analytification
of an algebraic variety; see its Notations and Terminology, printed
p.~8 of arXiv:1406.7637v3.  Definition 4.4 imposes vanishing after all
algebraic morphisms, and Definition 4.13 sheafifies the corresponding
algebras.  Remark 6.7 explicitly does not assume faithfulness of the
ordinary-current map; Proposition 6.8 gives a criterion using
\emph{all algebraic pullbacks} for its generalized delta-forms $P$.

Mihatsch states the embedding $B_{\mathrm{GK}}\hookrightarrow B$
in \cite[Section 4.2, printed p.~67]{Mih24}, under those GK hypotheses,
and explains its basis in Lemma 4.4 and Theorem 4.5.  He labels that
comparison as given without formal proof.  Once this embedding is
available for the actual datum, Corollary
\ref{cor:r6-ordinary-to-full-delta} applies to its image, and equality
there implies equality on GK delta-tests as well.  None of these
sentences replaces the original finite $K$ by its completed algebraic
closure or transfers an arbitrary current across that extension.

The theorem above accordingly uses Mihatsch's actual sheaf directly
over the original $K$.  It neither asserts that $P_{\mathrm{GK}}$
embeds into it nor that arbitrary CLD currents acquire all the
nonclosed, nonpositive, and off-diagonal operations of the delta
calculus.

% END source: sections/delta_faithfulness_comparison.tex


% BEGIN canonical R6 source: sections/delta_cycle_generation_scope.tex
\section{The local divisor generation that is actually available}
\label{sec:r6-delta-cycle-generation}

This section separates three algebras: constant-weight rational tropical
cycles, smooth-coefficient expressions in those cycles, and all actual
Mihatsch delta-forms.  The distinction matters before one asks whether a
formula involving an arbitrary background current descends through
relations.

\begin{proposition}[The fan generation input]
\label{prop:r6-fan-generation-input}
\textnormal{\texttt{EXTERNAL\_KNOWN}.}
The graded ring of rational tropical cycle fans in $\mathbf R^r$, with
rational weights and stable intersection, is isomorphic to the rational
polytope algebra.  Every codimension-$c$ such cycle is a finite rational
linear combination of $c$th powers of tropical hypersurfaces of rational
polytopes.
\end{proposition}

The precise source is \cite[Theorem 5.1 and Corollary 5.2]{JY16},
printed pp.~19--20 of arXiv:1309.7064v3.  Its convention on printed p.~3
explicitly restricts Sections 3--6 to fans.  Its Section 5 works over
$\mathbf Q$, not with arbitrary non-rational polyhedral directions.
The constructive paragraph after Corollary 5.2 uses a subdivision into
a simple polytopal fan and products of nonnegatively weighted divisor
fans.  These are equality statements for cycles, not only equality of
their supports or rational-equivalence classes.

\begin{corollary}[Constant real weights and local affine cycle germs]
\label{cor:r6-local-constant-cycle-generation}
\textnormal{\texttt{DIRECT\_COROLLARY}.}
Let $C$ be a constant-weight rational polyhedral tropical cycle of
codimension $c$ on a neighborhood of $a\in\mathbf R^r$.  Real constant
weights are allowed, and the affine offsets of its polyhedra may be
real.  After shrinking to a neighborhood $W$ of $a$, there are finitely
many real numbers $b_j$ and convex piecewise affine functions $\varphi_j$
whose gradients are integral such that
\begin{equation}
              [C]|_W=\sum_j b_j(D\varphi_j)^c|_W
              \quad\hbox{as actual real delta-forms on }W.
 \label{eq:r6-local-cycle-generation}
\end{equation}
The functions $\varphi_j$ have real additive constants; no extension of
the valued base field is required when their pullbacks are used later.
\end{corollary}

\begin{proof}
Choose a finite subordinate polyhedral complex near $a$.  On a
sufficiently small neighborhood, this cycle is the restriction of the
translate by $a$ of its local cone fan.  The directions of that fan
are rational.  For a fixed finite rational fan, its constant real
weight vectors satisfying balancing form the real kernel of a
matrix with rational entries: choose rational bases in the finitely
many normal quotient spaces and write the balancing equations.
This kernel has a rational basis.  Consequently the real-weight
fan is a finite real linear combination of rational-weight fans
on the same complex.

Apply Proposition \ref{prop:r6-fan-generation-input} to each of these
finitely many rational-weight fans.  If $P_j$ is a resulting rational
polytope, its support function
\[
                    h_{P_j}(x)=\max_{m\in P_j}\langle m,x\rangle
\]
is convex piecewise linear, and $Dh_{P_j}$ is its tropical divisor.
Translate the argument to $x-a$.  Choose an integer $d_j>0$ clearing
the denominators of the vertices of $P_j$ and set
$\varphi_j(x)=d_jh_{P_j}(x-a)$.  Then $\varphi_j$ has integral
gradients, and the factor $d_j^{-c}$ is absorbed into $b_j$.
The divisor product agrees with stable intersection by
\cite[Theorem 4.1(4)]{Mih23}, proving \eqref{eq:r6-local-cycle-generation}.
\end{proof}

\begin{corollary}[Original moment pullbacks]
\label{cor:r6-moment-cycle-generation}
\textnormal{\texttt{DIRECT\_COROLLARY}.}
Let $t_f:U\to\mathbf R^r$ be an original toric moment map over the
fixed $K$.  On a neighborhood on which a delta-form is represented
by a finite sum
\[
                        \omega=\sum_i t_f^*(\eta_i)\wedge t_f^\star[C_i],
\]
where the $\eta_i$ are smooth superforms and the $C_i$ have constant
weights and rational directions, it is locally a finite real linear
combination of expressions
\[
           t_f^*\eta\wedge(D_\delta v_1)\wedge\cdots\wedge(D_\delta v_c),
 \qquad v_j=\varphi_j\circ t_f,
\]
with continuous piecewise affine classical psh functions $v_j$.
The identities hold in the actual delta-form sheaf and hence after
every further original moment refinement.
\end{corollary}

\begin{proof}
Apply Corollary \ref{cor:r6-local-constant-cycle-generation} in a
neighborhood of $t_f(x)$ and shrink $U$ accordingly.  A convex
piecewise affine function is a finite maximum of affine functions.
Each affine pullback is a real linear combination of logarithms of
analytic units plus a real constant.  Such affine pullbacks are
pluriharmonic under all complete extensions, and their finite maxima
are classical psh.  The delta pullback commutes with $D$ and wedge
products by \cite[Definition 4.2(3)]{Mih24}; this proves the displayed
identity.  Refinement compatibility is part of the same definition,
not a separate assertion made on only the original chart.
\end{proof}

\begin{remark}[Generation is not the relation theorem]
\label{rem:r6-generation-not-relations}
The preceding corollary constructs representatives.  It does not
prove that the ordinary BT pairing of such a representative with
an arbitrary fixed $S$ is independent of the representative.  The
remaining relation problem is exactly that a zero delta-form
expression must have zero relative BT action.  The R5 potential
quotient settles individual degree-one potential changes, not
automatically all polynomial relations between several factors.
\end{remark}

\begin{example}[A false generation statement for the full Mihatsch algebra]
\label{ex:r6-smooth-generation-false}
\textnormal{\texttt{FALSE}; direct degree check.}
The full algebra $B(\mathbf R)$ is not generated by smooth forms
and the curvature factors $D\varphi$ alone.  Its degree-zero part
contains the continuous piecewise linear function
$x\mapsto\max(0,x)$, by \cite[Lemma 3.7]{Mih23}.  The degree-zero
part of the algebra generated by smooth forms and factors of
positive bidegree $(1,1)$ consists only of smooth functions.
Thus this function is absent from that proposed generated algebra.
The same example occurs over the original $K$ on an annulus whose
skeleton contains $\log|T|=0$: the function
$u=\max(0,\log|T|)$ has $D[u]=\delta_\xi$, where $\xi$ is the
unit Gauss point, by the original torus calibration.  If $u$ were
original smooth near $\xi$, its curvature would be a smooth
top-degree form there.  Such a form has Lebesgue densities on
finitely many local calibrated one-dimensional skeletal cells,
and hence has no atom at $\xi$.  This contradiction proves the
claimed non-smoothness.  A corner seen only on one fixed skeleton
would not by itself prove non-smoothness after all moment refinements.

This disproves the literal generation assertion; it does not
disprove a generation statement allowing piecewise smooth
coefficients, first derivatives, or further operations.
\end{example}

\paragraph{Scope of the two published definitions.}
The generalized delta-preforms of \cite[Definition 2.9]{GK17}
allow tropical cycles with \emph{smooth} weights.  Their differential
subalgebra $B_{\mathrm{GK}}$ is defined through the closed subspace
$Z$ and the algebra $AZ$ in \S4.6 and Definitions 4.9 and 4.13.
Consequently the constant-weight hypothesis of
Corollary \ref{cor:r6-moment-cycle-generation} must be retained;
it is not a verified description of every section of
$B_{\mathrm{GK}}$.

Mihatsch's real delta-forms are all balanced polyhedral currents,
equivalently those polyhedral currents whose $d'$ and $d''$
derivatives remain polyhedral; see \cite[Definitions 3.1--3.2 and
Theorem 3.3]{Mih23}.  His theory also permits non-rational
polyhedra.  Neither this class of variable coefficients nor
non-rational direction data is covered by the rational fan
generation theorem quoted above.  Their treatment cannot be
replaced by an invocation of the polytope algebra without an
additional proof.

\paragraph{A useful external structural formula, not a missing generation theorem.}
For real delta-forms, \cite[Theorem 4.1 and its uniqueness proof,
printed pp.~13--14]{Mih23} characterizes the product by the
Leibniz rules, multiplication by piecewise smooth forms, and
\[
 S\wedge T=p_{1,*}\bigl(\Delta\wedge(S\times T)\bigr),\qquad
 \Delta=\bigwedge_{i=1}^rD\max(x_i,y_i).
\]
Divisor intersection satisfies
\[
 (D\varphi)\cdot T
   =D(\varphi T)-\varphi DT+d''\varphi\wedge d'T
                                -d'\varphi\wedge d''T,
\]
and reduces to $D(\varphi T)$ for a separately closed $T$;
see Definition 4.6 and Remark 4.7, printed p.~15.
These formulas may guide an extension to a new class of currents.
They are proved in the delta-form category, however, and do not
assert that the required diagonal intersections, exterior products,
or presentation-independent piecewise operations exist for every
ordinary CLD current.

\paragraph{What the additional proofs in this round now supply.}
Theorem~\ref{thm:r6-real-ambient-cycle-action} treats constant cycles
with arbitrary real directions.  The exact finite construction in
Theorem~\ref{thm:r6-delta-module-generation} treats all smooth
coefficient arrays, including flat ones, and proves that every real
Mihatsch delta-form is locally a finite sum of piecewise smooth forms
times constant real cycles.  Thus the two omissions from the quoted
rational fan theorem have been addressed by separate proofs in this
round.  These proofs do not turn the quoted JY theorem into a stronger
external input.  The further actual analytic relation problem is
settled for the weak-approximation class in
Theorem~\ref{thm:r6-full-delta-test-extension}; its all-background
version remains a distinct open question.

% END source: sections/delta_cycle_generation_scope.tex


% BEGIN canonical R6 source: sections/pws_tangent_extension.tex
\section{Closed Radon currents and piecewise smooth tests}
\label{sec:r6-pws-tangent-extension}

Throughout this section $K$ is the original finite extension of
$\mathbf Q_p$, and $U$ is a good Hausdorff boundaryless quasi-smooth
$K$-analytic space of pure dimension $n$. All smooth forms are the
original CLD forms. The coefficient-Radon property and Borel
multiplication are understood on the actual Berkovich space, as in
R and B. No pushforward to one selected skeleton is used.

Write $\mathrm{PS}^{a,b}_U$ for the sheaf of piecewise smooth forms
defined using original toric moment presentations and their full
tropical images. Its equivalence relation is the one in Mihatsch,
\S4.2, equation (4.9), printed page 65. On a polyhedral complex,
the smooth pieces agree after restriction to the affine tangent
space of every common face; see \S2.1(16), printed pages 13--14.
By his Proposition 4.7, this is precisely the codimension-zero
part $B^{a,b,0}_U$ of his sheaf of actual delta-forms. It is not
the entire sheaf $B_U$.

\paragraph{Provenance before the new argument.}
R4, Theorem 71.4 (printed p.~240), already proves the
piecewise-test action and uniform smooth approximation for
arbitrary original-positive separately closed currents.
Its wall estimates even give an $O(\varepsilon^{1/2})$ bound
uniform under fixed finite mass controls. That positive
case is a \texttt{DIRECT\_COROLLARY} here, including its
identification with the codimension-zero delta-test class.
The additional assertions below concern closed
coefficient-Radon currents without a positivity assumption,
all smooth level sets rather than only affine moment walls,
and the resulting piecewise smooth module structure. They
do not replace the stronger quantitative R4 estimate.

\begin{lemma}[Tangency of the mass on a smooth level set]
\label{lem:r6-level-tangency}
Let $T$ be a coefficient-Radon current on $U$. If $d'T=0$, then
for every real original smooth function $h$ and every $c\in\mathbf R$,
\[
       d'h\wedge\mathbf 1_{\{h=c\}}T=0.
\]
If $d''T=0$, the corresponding identity with $d''h$ holds.
These are identities on all original complementary compact
smooth tests, including arbitrary further moment refinements.
No positivity hypothesis is needed.
\end{lemma}

\noindent\textbf{Status: \texttt{NEW\_PROVED}} for the stated
nonpositive closed-current extension. The positive affine-wall
case is already part of R4.

\begin{proof}
Put $g_\varepsilon(t)=\varepsilon\arctan((t-c)/\varepsilon)$.
Its absolute value is at most $\pi\varepsilon/2$, while
\[
 0\leq g_\varepsilon'(t)
    =\frac{1}{1+((t-c)/\varepsilon)^2}\leq1,
 \qquad
 g_\varepsilon'(t)\longrightarrow\mathbf1_{\{c\}}(t).
\]
Uniform convergence and the local finite coefficient variations
give $g_\varepsilon(h)T\to0$ as a current. Thus its $d'$ tends
to zero. The smooth product rule and $d'T=0$ identify that
derivative with $g_\varepsilon'(h)d'h\wedge T$.
On any fixed original test, coefficientwise bounded convergence
identifies its limit with $\mathbf1_{\{h=c\}}d'h\wedge T$.
This proves the first identity. The same proof gives the second.
Every argument is made after the chosen test has been presented
on a common refinement; hence no restriction of the test class
has occurred.
\end{proof}

\begin{definition}
Let $\mathscr R_U$ be the sheaf of coefficient-Radon currents $T$
such that, for every local real original smooth $h$ and every
$c\in\mathbf R$, both identities in
Lemma~\ref{lem:r6-level-tangency} hold. Thus separately closed
coefficient-Radon currents belong to $\mathscr R_U$.
This auxiliary sheaf records exactly the tangency needed to
multiply by piecewise smooth forms; its elements need not be closed.
\end{definition}

\begin{theorem}[Canonical piecewise smooth multiplication]
\label{thm:r6-pws-multiplication}
There is a unique local multiplication obtained from tangential
restriction of the smooth pieces,
\[
 \mathrm{PS}^{a,b}_U\times\mathscr R_U^{p,q}
       \longrightarrow\mathscr R_U^{p+a,q+b},
       \qquad(\alpha,T)\longmapsto\alpha\wedge T.
\]
It extends multiplication by original smooth forms and by actual
continuous piecewise smooth functions. It is real bilinear,
associative with the wedge product of piecewise smooth forms,
and compatible with restriction and gluing on ordinary opens.
It is independent of subdivisions, smooth extensions from
faces, toric moment presentations, and further refinements.
The resulting objects are ordinary coefficient-Radon currents.
In particular every allowed arbitrary positive separately
closed background $S$, and every frozen bounded BT product
over $S$, acts canonically on all compact piecewise smooth
complementary tests.
\end{theorem}

\noindent\textbf{Status: \texttt{NEW\_PROVED}} for the general
$\mathscr R_U$ module formulation; \texttt{DIRECT\_COROLLARY}
for the asserted action on the project's positive closed
backgrounds and BT products, by R4 and the displayed external
identification of the piecewise smooth sheaf.

\begin{proof}
Work first in one moment presentation $t:V\to\mathbf R^r$ and
on a compact stage. Choose a finite polyhedral subdivision
adapted to the piecewise smooth representative of $\alpha$.
Use all relatively open faces, so their inverse images
$E_\sigma=t^{-1}(\operatorname{relint}\sigma)$ are disjoint
Borel sets. Locally near the compact stage, choose a smooth
ambient extension $\widetilde\alpha_\sigma$ of the smooth
form on each closed face. Define
\begin{equation}
 \alpha\wedge T
   =\sum_\sigma t^*\widetilde\alpha_\sigma
                      \wedge\mathbf1_{E_\sigma}T.
 \label{eq:r6-pws-stratum-formula}
\end{equation}
There are only finitely many terms on the stage. Smooth
extensions have bounded coefficients there, so this is a
coefficient-Radon current. If several extension neighborhoods
are needed, use a finite smooth partition on the ambient
face; the same argument below proves independence of it.

Let $L_\sigma$ be the affine span of $\sigma$, written as the
common zero set of affine functions $\ell_1,\ldots,\ell_m$.
Tangency gives
\[
 d'(\ell_j\circ t)\wedge\mathbf1_{E_\sigma}T
  =d''(\ell_j\circ t)\wedge\mathbf1_{E_\sigma}T=0.
\]
Indeed $E_\sigma\subset\{\ell_j\circ t=0\}$, and restriction
of a zero coefficient-Radon current to a Borel set remains
zero. If two smooth extensions have the same tangential
restriction on $\sigma$, their pointwise difference there
belongs to the exterior ideal generated by the conormals
$d'\ell_j,d''\ell_j$. This follows by completing a basis of
$L_\sigma^*$ to an ambient dual basis. The coefficients of
that decomposition are Borel and bounded on the compact
stage. The preceding identities annihilate this difference
in \eqref{eq:r6-pws-stratum-formula}. Thus the formula does
not depend on extension choices.

Subdivision does not change the formula. On each smaller
face the old and new smooth pieces have identical tangential
restrictions, and the same conormal argument applies.
Any two finite subdivisions have a common subdivision.

It remains to check the quotient defining the analytic
sheaf, not just a fixed polyhedral presentation. Suppose
a presented piecewise smooth form is zero in that quotient.
After passing to a common moment presentation, its
restriction to the full tropical image $T'(A,t)$ is zero
for every affinoid domain $A$ in the presentation. Cover
the support of a given original test by finitely many
interiors of such affinoids, and use an original smooth
partition subordinate to those interiors. On one member
subdivide $T'(A,t)$ and the piecewise smooth presentation
simultaneously. For each resulting face, the representative
has zero tangential restriction. The conormal argument
therefore makes its product with the corresponding Borel
part of $T$ zero. Summing proves zero on the chosen original
test. This also proves invariance under the equivalence
of two moment presentations, by taking their joint moments
and subtracting the two representatives. The argument uses
the full $T'(A,t)$ and works for every affinoid $A$; it does
not replace the defining quotient by a selected skeleton.

If $h$ is any further original smooth function, apply the
same formula to $\mathbf1_{\{h=c\}}d'h\wedge(\alpha\wedge T)$.
Each summand is a smooth multiplier times a Borel restriction
of $\mathbf1_{\{h=c\}}d'h\wedge T$, hence zero. The $d''$
identity is identical. Thus the product belongs to
$\mathscr R_U$. Using one common subdivision for two
piecewise smooth forms proves associativity directly in
\eqref{eq:r6-pws-stratum-formula}. Bilinearity follows in the
same way. Restriction is manifest, and the independence
already proved glues the local constructions uniquely.
For a continuous piecewise smooth function the face values
are its actual values, so the formula is ordinary continuous
coefficient multiplication.
\end{proof}

\begin{proposition}[Coefficient control and fixed-current smoothing]
\label{prop:r6-pws-smoothing}
Fix $T\in\mathscr R_U$ and one compact stage with a finite
original moment presentation. Let $\alpha$ be a piecewise
smooth form there. There exist smooth forms $\alpha_\varepsilon$
on a slightly smaller stage such that, for every fixed
original smooth test $\eta$ on that stage,
\[
 \left|\mu_{(\alpha_\varepsilon-\alpha)\wedge T,\eta}\right|(K)
                         \longrightarrow0
\]
for each compact $K$ in the smaller stage. All these
coefficient variations admit a common finite bound by the
variations of finitely many coefficients of $T$ times a
fixed bound for the piecewise smooth coefficients and $\eta$.
The assertion is for fixed $T$; it does not assert uniformity
over an arbitrary varying family of currents.
\end{proposition}

\begin{proof}
Use a piecewise smooth representative on an open subset of
$\mathbf R^r$ containing the relevant tropical image. In
each full-dimensional polyhedron choose its smooth
coefficient array. Assign arbitrary bounded values to the
array on lower-dimensional faces and convolve locally with
a nonnegative compactly supported smooth mollifier of radius
$\varepsilon$. Work on a relatively compact smaller real
open, so no boundary extension issue occurs. The arrays
remain uniformly bounded there and define $\alpha_\varepsilon$.

For a fixed point $x$ in the relative interior of a face
$\sigma$, all adjacent pieces have the same restriction to
the tangent space of $\sigma$. Consequently, after restricting
all covectors to that tangent space, the mollified array at
$x$ tends to the common tangential array at $x$. To verify
this, discard the finitely many polyhedra whose closures do
not contain $x$, then use continuity of the smooth arrays
on each remaining closed polyhedron. The mollifier averages
only these finitely many arrays for sufficiently small
$\varepsilon$. Their common tangential limits agree.

On $E_\sigma$, all components in the conormal exterior ideal
are annihilated by $T$, as in the preceding theorem. Thus the
coefficients of $(\alpha_\varepsilon-\alpha)\wedge T$ on that
set converge pointwise, after contraction, to zero, with a
uniform bound by finitely many fixed coefficient variations
of $T$. Bounded convergence gives convergence in total
variation on $K\cap E_\sigma$. Sum the finitely many faces.
For an additional original test, make a common refinement
first; the argument and the finite coefficient domination
are unchanged. This proves the assertion on all original
tests. A partition of unity gives the same statement for
finitely many local presentations.
\end{proof}

\begin{corollary}[First derivatives of continuous piecewise smooth potentials]
\label{cor:r6-pws-first-derivative}
Let $T$ be separately closed and coefficient-Radon. For every
real continuous piecewise smooth function $f$,
\[
 d'(fT)=d'_Pf\wedge T,
 \qquad d''(fT)=d''_Pf\wedge T.
\]
In particular both first derivatives are coefficient-Radon.
These statements do not assert that $d'\alpha=d'_P\alpha$
for an arbitrary piecewise smooth form of positive degree;
corner terms can occur in that situation.
\end{corollary}

\begin{proof}
The local ambient representative of $f$ is continuous and
piecewise smooth on a finite polyhedral subdivision. It is
locally Lipschitz: the gradients on each smooth piece are
bounded on compact subsets and the pieces agree on common
faces. Its mollifications converge locally uniformly to $f$.
Their first derivatives are the mollifications of its
almost-everywhere piecewise derivatives, because continuity
eliminates a distributional jump term in the first derivative.
The tangential convergence proof of
Proposition~\ref{prop:r6-pws-smoothing} therefore applies to
these derivatives. The smooth closed product rule gives
$d'(f_\varepsilon T)=d'f_\varepsilon\wedge T$.
The left side converges distributionally to $d'(fT)$ by
local uniform convergence; the right side converges as a
coefficient-Radon current to $d'_Pf\wedge T$. Repeat for
$d''$. Local identities glue by the preceding theorem.
\end{proof}

\paragraph{Exact scope.}
The construction solves the codimension-zero delta-test
extension for every fixed allowed background and its BT
products. A positive-codimension delta-form is not merely a
piecewise smooth coefficient array. Pairing with such a form
can require an additional intersection. No arbitrary
delta-test extension, no product of two arbitrary currents,
and no varying-background monotone theorem is asserted here.

% END source: sections/pws_tangent_extension.tex


% BEGIN canonical R6 source: sections/continuous_lattice_valuation.tex
\section{A lattice valuation identity over an arbitrary background}
\label{sec:r6-continuous-lattice-valuation}

Fix the original $K$, an allowed ordinary open $U$, and one
arbitrary symmetric original-positive separately closed
$(p,p)$-current $S$. Put $q=n-p$. The signed differentiated
calculus of R5 is a frozen input: locally bounded differences
of actual classical psh functions give coefficient-Radon
mixed currents by finite polarization, with symmetry,
recursion, and dependence only on their differentiated
currents. All equalities below are ordinary current
equalities on every original compact complementary smooth
test and every further refinement.

Let $\mathscr H_c(U)$ denote the continuous actual functions
which are locally differences of two members of
$\mathcal C_b$. No continuity of those two individual
presentations is required. Their pointwise maximum is again
classical psh. Consequently $\mathscr H_c$ is a real vector
lattice: if $h=a-b$, then
\[
 h^+=\max(a,b)-b,\qquad h^-=\max(a,b)-a,
\]
and the maximum and minimum of two local differences can
be reduced to this identity. These formulas are only used
on a common local presentation; the resulting functions
are their actual pointwise positive and negative parts.

\begin{lemma}[Orthogonality of positive and negative parts]
\label{lem:r6-disjoint-parts-orthogonality}
Let $h\in\mathscr H_c(U)$ and suppose $q\geq2$. Then
\begin{equation}
       D[h^+]\wedge D[h^-]\wedge S=0.
 \label{eq:r6-disjoint-parts-orthogonality}
\end{equation}
The same conclusion holds with $S$ replaced by any signed
mixed tail generated from $S$ by the frozen bounded
classical-difference calculus, provided the total degree
is allowed.
\end{lemma}

\noindent\textbf{Status: \texttt{DIRECT\_COROLLARY}} of the
frozen signed calculus and ordinary locality. Its proof is
recorded because it supplies the contact-set identity needed
for the new cycle-relation argument below.

\begin{proof}
First let $Q$ be any one positive separately closed mixed
tail over $S$, including $Q=S$. The signed current
$C_-=D(h^-Q)$ is coefficient-Radon by finite polarization.
The set $G=\{h>0\}$ is an ordinary open because $h$ is
continuous. On $G$ the actual function $h^-$ is zero, so
the recursive definition, or ordinary restriction of
$D(h^-Q)$, gives $C_-|_G=0$.
Every coefficient Radon measure of $C_-$ therefore vanishes
on $G$. On $U\setminus G$ the actual multiplier $h^+$ is
zero. The two assertions imply
\[
                              h^+ C_-=0
\]
as a coefficient-Radon current on all of $U$. Applying $D$
and using the signed recursive calculus proves
\eqref{eq:r6-disjoint-parts-orthogonality} with $Q$ in place
of $S$. A signed mixed tail is a finite linear combination
of such positive tails on each local presentation. Apply
the preceding argument to its terms and use ordinary
restriction and gluing. No assertion about the ordinary
openness of a gap set between discontinuous psh functions
is involved, and no simultaneous limit of two singular
measures has been exchanged.
\end{proof}

\begin{theorem}[Max--min valuation for every mixed power]
\label{thm:r6-max-min-valuation}
Let $u,v\in\mathscr H_c(U)$ and set
$a=\max(u,v)$, $b=\min(u,v)$. For every $0\leq k\leq q$,
\begin{equation}
 (D[a])^k\wedge S+(D[b])^k\wedge S
       =(D[u])^k\wedge S+(D[v])^k\wedge S.
 \label{eq:r6-max-min-valuation}
\end{equation}
The degree-zero convention is $2S=2S$. A further fixed
signed mixed tail may be inserted in every term whenever
the degree permits it. The functions $a$ and $b$ are
continuous classical differences; $b$ need not be psh.
Thus the formula uses the signed calculus in its stated
scope rather than asserting positivity of every term.
\end{theorem}

\noindent\textbf{Status: \texttt{DIRECT\_COROLLARY}} of the
preceding orthogonality identity and finite polynomial algebra.
The ensuing application to ambient cycle relations is a
separate theorem, rather than a new regularization theorem.

\begin{proof}
On the module of signed mixed tails over $S$, write
$A=D[a]$, $B=D[b]$, $U_1=D[u]$, and $V_1=D[v]$ for their
commuting differentiated-slot operators. This notation
means composition of the defined recursive operations;
it does not multiply two arbitrary currents. Linearity
gives $A+B=U_1+V_1$. With $h=v-u$ one has
$a-u=h^+$ and $a-v=h^-$. Lemma
\ref{lem:r6-disjoint-parts-orthogonality} gives
\[
                   (A-U_1)(A-V_1)=0
\]
on every allowed tail. It follows that
$AB=A(U_1+V_1-A)=U_1V_1$.

For two commuting symbols $x,y$, their power sums obey
\[
 P_0=2,\quad P_1=x+y,\quad
 P_k=(x+y)P_{k-1}-xyP_{k-2}\quad(k\geq2).
\]
This is checked by expanding the right side and cancelling
the two cross terms. The pairs $(A,B)$ and $(U_1,V_1)$
have the same sum and product, hence have the same power
sums by induction. Apply that polynomial identity to $S$
or the fixed tail. All operators commute by the frozen
signed BT symmetry, and all intermediate terms have the
specified degree. Local potential presentations give the
identity locally; the original current sheaf glues it.
\end{proof}

\begin{corollary}[A finite exponential valuation]
\label{cor:r6-exponential-valuation}
For $h\in\mathscr H_c(U)$ let $L_h$ act on the module of
signed mixed tails generated over $S$ by
$L_h T=D(hT)$. The exponential
\[
                 E_h=\sum_{j=0}^{q}\frac{L_h^j}{j!}
\]
is finite, with terms beyond the available degree zero.
It satisfies
\[
 E_{h+k}=E_hE_k,\qquad E_0=I,
 \qquad E_{\max(h,k)}+E_{\min(h,k)}=E_h+E_k.
\]
If $h$ is actual pluriharmonic, then $L_h=0$ and $E_h=I$.
\end{corollary}

\noindent\textbf{Status: \texttt{DIRECT\_COROLLARY}} of
Theorem~\ref{thm:r6-max-min-valuation} and the frozen R5 quotient.

\begin{proof}
The product $hT$ and its derivative depend on the actual
current $T$, so the operator is independent of a chosen
signed-tail presentation. Expanding that presentation
shows that it stays within the stated module and agrees
with insertion of the differentiated slot. Those operators
commute and satisfy $L_{h+k}=L_h+L_k$. The ordinary finite
binomial formula gives the first assertion. The valuation
identity follows by summing the power identities with
coefficients $1/j!$. For a pluriharmonic $h$, R5 makes its
slot zero on every positive tail; finite polarization
gives the same conclusion on the entire module.
\end{proof}

\paragraph{Continuity is used at one precise place.}
The proof uses that $\{h>0\}$ is an ordinary open. For a
general discontinuous classical difference this need not
hold. The preceding argument therefore supplies no new
fine locality theorem for that larger class. Continuous
PL support functions, which are the application below,
satisfy the hypothesis exactly.

% END source: sections/continuous_lattice_valuation.tex


% BEGIN canonical R6 source: sections/ambient_polytope_action.tex
\section{Descent through ambient tropical cycle relations}
\label{sec:r6-ambient-polytope-action}

Fix the same original field and allowed space as above, one
arbitrary allowed background $S$, and original toric moments
$t=(\log|f_1|,\ldots,\log|f_r|)$ on an ordinary open $U$.
Changing all signs to additive valuations merely changes
the chosen moment convention. We use the maximum convention
for support functions and tropical divisors, consistently
with \cite[Notations and conventions, p.~3]{JY16}.

Let $\mathscr M_S$ be the graded real module of ordinary
currents obtained locally as finite linear combinations of
the signed bounded-potential mixed tails over $S$. Its
grading records the number of differentiated slots. An
operator of degree larger than the remaining dimension
acts as zero. A continuous local classical difference $h$
defines $L_h:T\mapsto D(hT)$ on this module. This is a
well-defined operator on actual currents: Borel
multiplication and $D$ do not depend on a presentation of
$T$, while the frozen signed calculus shows that its value
is another allowed mixed tail. The operators $L_h$ commute.

\begin{lemma}[Support functions under a convex union]
\label{lem:r6-support-min-convex-union}
Let $P,Q$ be nonempty compact convex polytopes such that
$P\cup Q$ is convex. Then
\[
 h_{P\cup Q}=\max(h_P,h_Q),\qquad
 h_{P\cap Q}=\min(h_P,h_Q).
\]
In particular $P\cap Q$ is nonempty.
\end{lemma}

\begin{proof}
The maximum formula is immediate from taking the supremum
over a union. Fix a linear functional $\ell$. If
$h_P(\ell)<h_Q(\ell)$, choose maximizers $p\in P$, $q\in Q$.
For every $s\in(0,1]$, the point $(1-s)p+sq$ has $\ell$-value
strictly greater than $h_P(\ell)$. The whole segment lies in
$P\cup Q$, so all these points belong to $Q$. By closedness
$p\in Q$, and the maximum on $P\cap Q$ is at least
$h_P(\ell)$. The reverse inequality is automatic. The case
of the opposite strict inequality is symmetric.
If the two support values agree, the segment of maximizers
lies in that common supporting hyperplane. Its intersections
with the closed sets $P$ and $Q$ cover the segment and contain
its respective endpoints. Connectedness forces them to
intersect. The intersection point has the common support
value. This proves the second formula and nonemptiness.
\end{proof}

\begin{theorem}[Action of the rational fan cycle ring]
\label{thm:r6-ambient-cycle-action}
Let $\mathcal T_r$ be the graded ring of constant-weight
rational tropical cycle fans in $\mathbf R^r$, allowing
real weights. There is a canonical graded real algebra
homomorphism
\[
 \mathcal A_{t,S}:\mathcal T_r
                      \longrightarrow\operatorname{End}(\mathscr M_S)
\]
characterized by
\begin{equation}
 \mathcal A_{t,S}(D h_P)T=D\bigl((h_P\circ t)T\bigr)
 \label{eq:r6-cycle-action-divisor}
\end{equation}
for every rational polytope $P$. Thus every polynomial
relation between constant rational tropical cycle fans
annihilates every allowed mixed tail over the arbitrary $S$.
The output for a homogeneous cycle of codimension $c$ is
a separately closed coefficient-Radon current of degree
$c$ above that tail. No positivity claim is made for a
general positive tropical cycle presented as a signed
combination of hypersurface products.
\end{theorem}

\noindent\textbf{Status: \texttt{NEW\_PROVED}}: the relation
descent uses the elementary frozen-calculus consequence
proved above and the external polytope algebra isomorphism.
It is not a new construction of the ordinary BT product.

\begin{proof}
First work over $\mathbf Q$. The rational polytope algebra
$\Pi_r$ is generated additively by symbols $[P]$, with
translation invariance and the relation
\[
 [P\cup Q]+[P\cap Q]=[P]+[Q]
\]
when $P\cup Q$ is convex. Its multiplication is Minkowski
sum, $[P][Q]=[P+Q]$, with $[\varnothing]=0$ and the point
class as unit; see \cite[Section 5, p.~19]{JY16}.

For $P\ne\varnothing$ assign
\[
                  [P]\longmapsto\exp L_{h_P\circ t},
 \qquad [\varnothing]\longmapsto0.
\]
The support function pulled back by $t$ is a finite maximum
of original pluriharmonic functions, hence a continuous
classical psh function. All its operators are therefore
defined. Translating $P$ adds an affine function of $t$.
That function is original smooth pluriharmonic, so its
operator is zero. Thus translation invariance holds.
The support function of a Minkowski sum is the sum of the
support functions. The finite exponential identity in
Corollary~\ref{cor:r6-exponential-valuation} proves
multiplicativity. The preceding lemma and the max--min
exponential valuation prove the additive polytope relation.
If one polytope is empty the relation is immediate. All
defining relations are consequently respected, and the
assignment gives a unital algebra homomorphism
$\rho:\Pi_r\to\operatorname{End}(\mathscr M_S)$.

By \cite[Theorem 5.1, pp.~19--20]{JY16}, the map
\[
 [P]\longmapsto
       \sum_{j=0}^{r}\frac{(D h_P)^j}{j!}
\]
is an isomorphism $\Phi:\Pi_r\to\mathcal T_{r,\mathbf Q}$.
Its logarithm sends $\log[P]$ to $D h_P$. In the nilpotent
operator algebra, $\log\exp L_{h_P\circ t}=L_{h_P\circ t}$.
Therefore $\rho\Phi^{-1}$ has property
\eqref{eq:r6-cycle-action-divisor}. The graded component of
degree $c$ of the polytope algebra is spanned by
$(\log[P])^c$, so this action is graded and its homogeneous
values are finite linear combinations of the signed BT
currents asserted in the theorem. Separate closedness and
coefficient-Radon type follow term by term.

For real weights, use scalar extension. On any fixed finite
rational fan, the real balanced weight space is the real
scalar extension of its rational balanced weight space,
because balancing is the kernel of a rational matrix.
Every finite collection of cycles has a common finite
rational subdivision. Hence
$\mathcal T_r=\mathcal T_{r,\mathbf Q}\otimes_{\mathbf Q}\mathbf R$,
with its products preserved, and the action extends
uniquely by real linearity. Hypersurface generation gives
the claimed characterization. This proof never changes
the valued field and never replaces $S$ by a geometric
or smooth background.
\end{proof}

\begin{corollary}[Local affine cycles, support, and ambient refinements]
\label{cor:r6-ambient-affine-cycle-action}
Let $C$ be a constant-weight rational polyhedral tropical
cycle on a real open $W\subset\mathbf R^r$, with real affine
offsets and real weights. There is a local action
$\mathcal A_t(C)S$ on $t^{-1}(W)$. It is linear in $C$,
separately closed, and coefficient-Radon, and it has support
contained in $t^{-1}(|C|)$. It is independent of local fan
representations and compatible with restriction, gluing,
and affine rational refinements of ambient moment maps.
Every equality of such ambient cycle germs gives an
equality of their actions on all original compact smooth
tests, including tests with additional moment coordinates.
\end{corollary}

\begin{proof}
We first verify locality of the fan action, rather than
assuming it from its global algebraic definition. Near
$a\in\mathbf R^r$, a polytope support function satisfies
\[
 h_P(x)=h_P(a)+h_{F_P(a)}(x-a),
\]
where $F_P(a)$ is the exposed maximizing face. Indeed the
finitely many other vertices have a strict support gap
at $a$, which persists in a sufficiently small neighborhood.
The equality is exact on that neighborhood. It remains
exact after moment pullback, so its differentiated
operators agree there. A finite expression for a cycle
as powers of hypersurfaces therefore localizes to the
corresponding expression for its translated local fan.
If the cycle germ is zero, that fan relation is zero,
and Theorem~\ref{thm:r6-ambient-cycle-action} annihilates it.
This proves locality and the support assertion.

Every affine polyhedral cycle is, near a given point, the
translate of its local fan. Apply the fan theorem to
that translate, using $t-a$ in place of $t$. On the
overlap of two such neighborhoods their cycle germs
agree; localizing both expressions once more at any
point of the overlap reduces their equality to the same
fan relation. Hence the currents agree on the overlap
and glue uniquely. The same argument proves linearity
and compatibility with restriction.

For an affine rational map $F$ of real moment spaces,
pullback of a hypersurface potential is literal composition
with $F$. On a sufficiently small neighborhood it is again
a convex piecewise affine function with rational gradients,
with real constants allowed. The delta pullback is
compatible with $D$ and products. Hypersurface generation
and the identity of the actual composed potentials prove
compatibility of the actions. Refinement never restricts
the final test class: the current identity can be paired
with any original test after its own common refinement.
\end{proof}

\paragraph{The remaining analytic kernel is separate.}
The preceding results kill relations in the ambient real
cycle algebra and its germs. They do not yet prove
\[
        t^\star C=0\text{ in }B(U)
             \quad\Longrightarrow\quad\mathcal A_t(C)S=0
\]
for an arbitrary original moment map on an arbitrary
analytic space. A nonzero ambient cycle may have zero
actual pullback. Annihilation of that additional kernel
is required before the construction descends to all
actual analytic delta-tests. The R5 quotient settles
the corresponding differentiated potential relation in
degree one; it does not automatically settle this
higher-degree analytic relation. The fixed-coordinate
torus case below has an independent faithfulness check.

% END source: sections/ambient_polytope_action.tex


% BEGIN canonical R6 source: sections/real_cycle_stanley_reisner_action.tex
\section{Real cycle relations and pointwise polyhedral relations}
\label{sec:r6-real-cycle-action}

The preceding rational construction has a stronger real
version. The additional argument uses the presentation of
piecewise polynomial functions on a simplicial fan. It
also identifies a class of analytic relations that can be
annihilated without approximating the background.

\begin{lemma}[A vanishing product of continuous potentials]
\label{lem:r6-zero-function-product}
Let $f_1,\ldots,f_m$ be continuous local bounded classical
differences on an ordinary open and suppose
$f_1\cdots f_m=0$ at every point. Then, over every allowed
positive separately closed background $Q$,
\[
                      D[f_1]\wedge\cdots\wedge D[f_m]\wedge Q=0
\]
whenever the degree is allowed. The conclusion also holds
on every signed mixed tail over a fixed allowed background.
\end{lemma}

\noindent\textbf{Status: \texttt{DIRECT\_COROLLARY}} of
the frozen signed recursion and ordinary locality.

\begin{proof}
Induct on $m$. The case $m=1$ is immediate. For $m>1$
let $R=D[f_2]\wedge\cdots\wedge D[f_m]\wedge Q$.
It is a coefficient-Radon current. On the ordinary open
$G=\{f_1\ne0\}$ the remaining product of actual functions
is zero. Induction on that open gives $R|_G=0$.
On its complement the multiplier $f_1$ is zero. Thus
$f_1R=0$ as an actual coefficient-Radon current, and its
$D$ is zero. Signed tails are handled by their finite
positive-tail expansions. All of these arguments are
local on ordinary opens.
\end{proof}

\begin{lemma}[Finite piecewise linear functions give allowed signed slots]
\label{lem:r6-pl-difference-convex}
Every continuous piecewise linear function on a finite
complete real fan is a difference of two finite convex
piecewise linear functions. After composition with original
unit moments it is therefore a continuous local bounded
classical difference, over the original field.
\end{lemma}

\begin{proof}
Refine the fan by the finite arrangement of all hyperplanes
containing its facets. Write their linear equations as
$\ell_j=0$ and set $\rho=\sum_j|\ell_j|$.
Continuity of the given function $f$ implies that, across
each adjacent pair of full-dimensional chambers, its
gradient jump is a scalar multiple of the corresponding
normal. There are finitely many such jumps. For sufficiently
large $M$, every oriented normal jump of $f+M\rho$ is
nonnegative. Along a generic line, its successive slopes
are then nondecreasing. By continuity, the same convexity
inequality holds on every line, including lines in walls.
Thus $f+M\rho$ and $M\rho$ are convex piecewise linear
functions. Each is a finite maximum of linear functions.
Their moment pullbacks are finite maxima of original
pluriharmonic functions, hence classical psh under the
full complete-extension curve definition. The actual
difference is the required function.
\end{proof}

\begin{lemma}[The piecewise polynomial presentation]
\label{lem:r6-stanley-reisner-presentation}
Let $\Delta$ be a finite complete simplicial fan in
$\mathbf R^r$. Choose one nonzero generator $v_\rho$ of
each ray and let $\lambda_\rho$ be its Courant function:
it is linear on each cone and takes value
$\delta_{\rho\tau}$ at $v_\tau$. Then
\[
 \mathrm{PP}(\Delta)
     =\mathbf R[x_\rho:\rho\in\Delta(1)]/I_\Delta,
 \qquad x_\rho\longmapsto\lambda_\rho,
\]
where $I_\Delta$ is generated by the squarefree monomials
whose rays are not all in one cone. Here
$\mathrm{PP}(\Delta)$ consists of continuous functions
polynomial on every cone, with the homogeneous grading.

If $C\subset\Delta$ is a subfan, the kernel of restriction
$\mathrm{PP}(\Delta)\to\mathrm{PP}(C)$ is generated by
the monomials whose ray supports lie in no cone of $C$.
\end{lemma}

\begin{proof}
A nonface monomial vanishes pointwise, since on any cone
only its own ray variables can be nonzero. Thus the
displayed map factors through the stated quotient.
In that quotient every polynomial is a linear combination
of monomials whose support is a face. If it restricts to
zero on a maximal cone $\sigma$, all monomials supported
in $\sigma$ have zero coefficients: the corresponding
$\lambda_\rho|_\sigma$ are independent linear coordinates
and a polynomial vanishing on the relative interior is
the zero polynomial. Every allowed monomial occurs in
such a cone. This proves injectivity, and the same
argument on maximal cones of a subfan proves the asserted
description of the restriction kernel.

For surjectivity take a homogeneous piecewise polynomial
$F$ and construct its expression in increasing order of
face dimension. After matching all proper faces of a
cone $\sigma$, its residual polynomial on $\sigma$
vanishes on every facet. In its independent ray
coordinates it is therefore divisible by
$\prod_{\rho\subset\sigma}\lambda_\rho$.
Write the quotient as a polynomial in those coordinates
and use the identical polynomial in the global Courant
functions, multiplied by that product. This corrects
the residual on $\sigma$ and vanishes on every other
face of dimension at most $\dim\sigma$. Processing all
faces of that dimension consequently preserves earlier
corrections. At the maximal cones the residual is zero
everywhere. The initial constant on the zero cone is
handled by a constant polynomial. This finite induction
gives surjectivity degree by degree.
\end{proof}

\begin{theorem}[Piecewise polynomial functional calculus]
\label{thm:r6-piecewise-polynomial-operator}
For original moments $t:U\to\mathbf R^r$ and an arbitrary
allowed background $S$, there is a canonical graded
algebra homomorphism
\[
 \Theta_{t,S}:\mathrm{PP}(\mathbf R^r)
                            \longrightarrow\operatorname{End}(\mathscr M_S)
\]
which sends a degree-one piecewise linear function $f$
to $L_{f\circ t}$. The source is the direct limit over
finite complete real fans and their refinements. Every
global linear function lies in its kernel, and hence so
does the ideal generated by such functions.

More strongly, if $C$ is a subfan of a finite complete
simplicial fan $\Delta$, if $t(U)\subset|C|$, and if
$F\in\mathrm{PP}(\Delta)$ vanishes on $|C|$, then
$\Theta_{t,S}(F)=0$.
\end{theorem}

\noindent\textbf{Status: \texttt{NEW\_PROVED}} for the
operator descent and the stated pointwise-image ideal.

\begin{proof}
On one simplicial fan assign
$x_\rho\mapsto L_{\lambda_\rho\circ t}$. The operators
are defined by Lemma~\ref{lem:r6-pl-difference-convex}
and commute by signed BT symmetry. A nonface monomial
has zero product of the actual Courant functions.
Lemma~\ref{lem:r6-zero-function-product} therefore
annihilates its operator. Lemma
\ref{lem:r6-stanley-reisner-presentation} gives a
homomorphism from $\mathrm{PP}(\Delta)$.

On a simplicial refinement, a coarse Courant function
is the linear combination of the fine Courant functions
given by its values at the fine ray generators. Linearity
of $h\mapsto L_h$ shows that the two homomorphisms agree.
Every two finite complete real fans have a common finite
simplicial refinement, obtained by common subdivision
and compatible triangulation of the cones. Thus these
homomorphisms give the asserted direct-limit map, with
no dependence on the ray normalizations. A global linear
function pulled back by unit moments is original smooth
pluriharmonic. Its operator on a separately closed tail
is zero, so its generated ideal is killed.

Finally the kernel of restriction to $C$ is generated
by its nonface monomials, by the preceding lemma. Their
actual products vanish on $U$ because $t(U)\subset|C|$.
The same zero-product lemma annihilates them. This proves
the stronger assertion. Notice that it uses actual
pointwise vanishing, not vanishing almost everywhere
for the ambient integration current.
\end{proof}

\begin{theorem}[Constant real tropical cycles act over every fixed background]
\label{thm:r6-real-ambient-cycle-action}
The ambient cycle action and its local affine version
extend from rational directions to all constant-weight
real tropical cycle fans and germs. It remains linear,
graded, local in the real moment domain, and compatible
with real affine changes of moments. It agrees with
the previous action on rational cycles. Every output
is an ordinary separately closed coefficient-Radon
current, and all equalities hold on all original tests.
\end{theorem}

\noindent\textbf{Status: \texttt{NEW\_PROVED}} using
Theorem~\ref{thm:r6-piecewise-polynomial-operator} and
the following external isomorphism.

\begin{proof}
Esterov's Theorem 3.1 identifies the graded ring of
continuous conewise-polynomial functions modulo the
ideal generated by global linear functions with the
ring of constant-weight real tropical cycle fans;
see \cite[Section 3, printed p.~12]{Est12}.
The identification sends a piecewise linear function
$f$ to its corner divisor. With the standard ambient
calibration, that divisor is $Df$: on the elementary
hinge $\max(0,\ell)$ both conventions assign the
normal weight $d\ell$. The ring products are the same
transverse, hence stable, intersection products. No
rationality or lattice hypothesis appears in this
real isomorphism.

The preceding theorem kills its entire defining
kernel. Composing with the inverse isomorphism gives
the required cycle action. It sends $Df$ to
$L_{f\circ t}$, including signed PL functions. This
also proves agreement with the rational action,
because both are generated by these divisors.
Homogeneous degree $c$ piecewise polynomials are finite
linear combinations of products of $c$ Courant
functions. Their values are consequently finite signed
BT currents, giving the asserted closedness and Radon
properties.

Locality and affine-germ gluing follow as in the
rational case: localize each PL generator to its exact
piecewise linear germ at the point in question and
discard its constant value. The localized polynomial
is the local cycle germ under the corner-locus
isomorphism. A zero germ is annihilated by the same
ring relation. Composition of potentials proves
compatibility with real affine moment changes.
Again, this is an ambient relation theorem; it does
not discard the remaining actual analytic pullback
kernel described above.
\end{proof}

\paragraph{An exact additional criterion for analytic relations.}
On a local moment image contained in a subfan $C$, every
piecewise polynomial representative in
\[
 \ker(\mathrm{PP}(\Delta)\to\mathrm{PP}(C))
        + (\text{global linear functions})
\]
has zero relative action over every $S$. The first ideal
is killed by actual zero products; the second by
pluriharmonicity. It remains a separate question whether
every actual zero cycle-presentation on a given analytic
chart lies in this sum of ideals. The statement here is
an independently proved sufficient criterion, not an
assumption that the local tropical image has a Euclidean
or matroidal intersection ring.

% END source: sections/real_cycle_stanley_reisner_action.tex


% BEGIN canonical R6 source: sections/delta_torus_test_action.tex
\section{An arbitrary-background action for fixed-coordinate tropical cycle tests}
\label{sec:r6-delta-torus-action}

The ambient cycle action proved in
Theorem~\ref{thm:r6-ambient-cycle-action} and
Corollary~\ref{cor:r6-ambient-affine-cycle-action} will be denoted
$C\mapsto\mathcal A_C(S)$.  We use its stated linearity, locality,
separate closedness, coefficient-Radon property, support property,
and invariance under equality of ambient constant-weight rational
cycles.  Its hypersurface action is $\mathcal A_{D\varphi}(S)=D(\varphi S)$.
For signed cycles the action is interpreted by the proved finite
polarization of positive classical-potential products.

\subsection{Coefficients and relations on all polyhedral strata}

\begin{theorem}[Piecewise smooth coefficient relations]
\label{thm:r6-cycle-coefficient-relations}
\textnormal{\texttt{NEW\_PROVED}, relative to the ambient cycle action.}
Let $t_f:U\to W\subset\mathbf R^r$ be an original moment map and
let $S$ be any allowed fixed positive separately closed background.
Let $C_1,\ldots,C_N$ be constant-weight rational cycles of the same
codimension $c$ on $W$, and let $\eta_i\in PS^{a,b}(W)$.
If
\begin{equation}
                      \sum_{i=1}^N\eta_i\wedge[C_i]=0
                         \quad\hbox{as real delta-forms on }W,
 \label{eq:r6-ambient-coefficient-zero}
\end{equation}
then the well-defined piecewise coefficient products satisfy
\begin{equation}
             \sum_{i=1}^N t_f^*\eta_i\wedge\mathcal A_{C_i}(S)=0
                        \quad\hbox{as original CLD currents on }U.
 \label{eq:r6-coefficient-action-zero}
\end{equation}
The equality is tested against all original complementary compact
smooth forms, including tests requiring further moment refinements.
Only degrees within the dimension of $U$ are asserted.
\end{theorem}

\begin{proof}
Put $T_i=\mathcal A_{C_i}(S)$.  These are separately closed
coefficient-Radon currents, not necessarily positive.  Their
piecewise smooth coefficient products are defined by
Theorem~\ref{thm:r6-pws-multiplication}.

The assertion is local over a compact moment stage.  Choose a
common finite polyhedral complex subordinate to all the $C_i$ and
$\eta_i$ on a neighborhood of that stage.  All arguments below
can be localized to slightly smaller stages, so artificial outer
faces of the chosen complex play no role.  Write $w_{i\tau}$ for
the constant weight of $C_i$ on a cell $\tau$ of dimension $r-c$,
with zero weight when that cell is absent from $C_i$.

For a face $F$ of this complex let
\[
                            E_F=t_f^{-1}(\operatorname{relint}F).
\]
These are Borel sets.  If $L_F$ is the affine span of $F$, choose
affine equations $\ell_j=0$ for $L_F$.  The original smooth
functions $\ell_j\circ t_f$ vanish on $E_F$.  Lemma
\ref{lem:r6-level-tangency}, applied separately to each $T_i$, gives
\[
 d'(\ell_j\circ t_f)\wedge\mathbf1_{E_F}T_i
   =d''(\ell_j\circ t_f)\wedge\mathbf1_{E_F}T_i=0.
\]
Thus, on this stratum, the multiplier $t_f^*\eta_i$ depends only
on the restriction of $\eta_i$ to $L_F$.  In particular all terms
vanish if $\dim F<\max(a,b)$.

Otherwise choose constant tangent coordinate forms on $L_F$ and
write, on $F$,
\[
                         \eta_i|_F=\sum_{I,J}c_{i,IJ}\,\vartheta_{IJ}.
\]
These scalar coefficients are smooth on $F$, locally bounded, and
Borel after pullback.  For every incident cell $\tau$ of dimension
$r-c$, relation \eqref{eq:r6-ambient-coefficient-zero} implies
\[
                   \sum_iw_{i\tau}\eta_i|_\tau=0.
\]
Restrict this equality to $F\subset\tau$.  For each fixed tangent
tensor component, the vector
$c_{IJ}(x)=(c_{1,IJ}(x),\ldots,c_{N,IJ}(x))$ lies in the kernel
of the fixed real matrix
\[
                         M_F=(w_{i\tau})_{\tau\supset F,\ i}.
\]
The case in which no such $\tau$ occurs causes no problem: all
$T_i$ vanish on a neighborhood of the stratum by their support
property.

Choose a fixed basis $b^1,\ldots,b^m$ of $\ker M_F$.  There are
locally bounded Borel scalar functions $q_{IJ,l}$ on $F$ such that
\[
                          c_{IJ}(x)=\sum_lq_{IJ,l}(x)b^l.
\]
For each $l$, the cycle $\sum_i b_i^lC_i$ vanishes on an ambient
ordinary neighborhood $W_F$ of $\operatorname{relint}F$: remove
the finitely many nonincident closed top cells and use the zero
weights on all incident cells.  Local ambient relation invariance
therefore gives
\[
                         \sum_i b_i^lT_i=0
                             \quad\hbox{on }t_f^{-1}(W_F).
\]
This is an equality of coefficient-Radon currents.  Restricting
it to $E_F$, multiplying by the Borel scalar
$q_{IJ,l}\circ t_f$, and then by the constant tangent form
$t_f^*\vartheta_{IJ}$ are legitimate coefficient-measure operations.
Their finite sum proves
\[
           \mathbf1_{E_F}\sum_i t_f^*\eta_i\wedge T_i=0.
\]
Finally sum over the finitely many relevant strata.  Since every
compact original test is covered by such local stages, this proves
\eqref{eq:r6-coefficient-action-zero} on every original test.
No $E_F$ was treated as an ordinary open set, and no cycle was
inserted into a nonclosed coefficient-weighted background.
\end{proof}

\subsection{A full-coordinate torus domain removes the pullback kernel}

Let $\Omega\subset\mathbf R^n$ be ordinary open and set
\[
      U=t^{-1}(\Omega)\subset\mathbf G_{m,K}^{n,\mathrm{an}},
       \qquad t=(\log|T_1|,\ldots,\log|T_n|).
\]
We keep this single coordinate moment presentation.  The use of
$\log$ instead of $-\log$ is a simultaneous real affine reflection;
all weights and differential operations are the original ones.
Define $\mathscr B_t(U)$ as the image under $t^\star$ of those real
delta-forms on $\Omega$ which are locally finite sums
$\sum_i\eta_i\wedge[C_i]$ of the indicated kind.  Its elements
are therefore locally presented over coordinate open subsets by
\[
                       \sum_i t^*\eta_i\wedge t^\star[C_i],
\]
where $\eta_i$ are piecewise smooth superforms and $C_i$ are
constant-weight rational cycles.  Different tropical codimensions
are allowed and are treated in their separate trigraded summands.
This is an explicitly specified subalgebra of actual $B(U)$.
It is not claimed to be every actual delta-form on $U$.

\begin{theorem}[A fixed-coordinate delta-test action]
\label{thm:r6-fixed-coordinate-delta-action}
\textnormal{\texttt{NEW\_PROVED}, with the stated test scope.}
For every allowed fixed background $S$ of bidegree $(p,p)$ on
$U$, the prescription
\begin{equation}
 \mathcal A_S\left(\sum_i t^*\eta_i\wedge t^\star[C_i]\right)
       :=\sum_i t^*\eta_i\wedge\mathcal A_{C_i}(S)
 \label{eq:r6-fixed-torus-action}
\end{equation}
is independent of the displayed presentation and defines a linear
coefficient-Radon current of the expected total bidegree.  It is
compatible with restriction from $\Omega$ to its real open subsets
and thus from $U$ to the corresponding saturated torus domains.
The independence is equality of original currents on all original
tests, not merely equality in the chosen tropical coordinates.

In complementary total degree $(n-p,n-p)$ it gives a well-defined
pairing with every compactly supported test in $\mathscr B_t$.
If $S$ is represented by an actual delta-form $\sigma$, this
pairing is the usual $\int_U\sigma\wedge\beta$.

For continuous convex rational PL functions $u_j=\varphi_j\circ t$,
cycle insertion recovers the relative BT products.  More precisely,
whenever degrees permit and $\beta\in\mathscr B_t$,
\[
 \mathcal A_{T_{S,k}(u_1,\ldots,u_k)}(\beta)
      =\mathcal A_S\bigl((D_\delta u_1)\wedge\cdots\wedge
                                  (D_\delta u_k)\wedge\beta\bigr).
\]
Thus identities between these actual fixed-coordinate delta-tests
remain valid after relative BT insertion over an arbitrary $S$.
\end{theorem}

\begin{proof}
The standard skeleton of $U$ is canonically identified with
$\Omega$, with calibration one.  Restricting a presented pullback
to this skeleton is the original real delta-form on $\Omega$.
Consequently, if two displayed expressions agree as actual
delta-forms on $U$, their real expressions agree on $\Omega$.
Apply Theorem~\ref{thm:r6-cycle-coefficient-relations} to their
difference, separately in each codimension.  This proves
presentation independence.  The restriction and Radon assertions
follow from those of the cycle action and piecewise multiplication.

When the test has compact support, its real expression has compact
support in $\Omega$: a point of the support of the real expression
gives a support point on the standard skeleton, lying in the
compact support of the actual test.  Choose a smooth real compact
cutoff equal to one near that real support.  Ambient locality and
the already proved coefficient relations imply that the output is
supported over the support of this cutoff.  The inverse image of
a compact subset of $\Omega$ under the torus tropicalization is
compact, being a closed subset of a compact polyradius annulus.
Thus in complementary degree the output is a compactly supported
top-degree signed Radon measure, whose mass defines the pairing.
If $S=[\sigma]$, the recursion
comparison in Theorem~\ref{thm:r6-bt-delta-background-comparison},
finite polarization, and multiplication by piecewise coefficients
identify each term with its actual delta product.  This proves the
comparison with delta integration.

Finally $D_\delta u_j$ is the pullback of the rational tropical
hypersurface $D\varphi_j$.  The ambient polytope action is
multiplicative on cycle factors and its hypersurface insertion is
$D(u_j\,\cdot)$.  The displayed compatibility follows by applying
these identities before multiplying by the piecewise coefficient
of each summand of $\beta$.  All inserted intermediate backgrounds
are separately closed signed mixed currents for which that cycle
action was constructed; no operation on an arbitrary nonclosed
current is used.
\end{proof}

\paragraph{The remaining actual-pullback gate.}
Even on the torus, the saturated-domain statement above does not
by itself prove ordinary sheaf locality relative to the
\emph{actual} support of a delta-test.  An ordinary open subset
of $U$ can miss the standard skeleton and need not be a saturated
coordinate domain.  For example $D\max(0,\log|T|)$ on a curve
annulus has actual support consisting only of its Gauss point,
whereas $t^{-1}\{0\}$ also contains residue branches.  Its relative
action does vanish on those branches, because its potential is
locally constant there; this special observation is not a general
pullback-kernel theorem.  We have proved compact saturated support
and presentation independence, not the missing ordinary sheaf
locality assertion in general.

For a general moment map $t_f:U\to\mathbf R^r$, especially with
$r>\dim U$, an actual delta-form can vanish on $U$ although its
ambient expression is nonzero.  Theorem
\ref{thm:r6-cycle-coefficient-relations} kills \emph{ambient}
relations.  To descend through this additional pullback kernel
one needs a further argument.  The standard full-coordinate
torus domain above has no such kernel because of its standard
skeleton.  A delta-form involving additional units, such as
$T_i-1$, is not automatically in $\mathscr B_t$; enlarging the
moment tuple returns to the additional pullback problem.
No extension to all actual Mihatsch tests, all GK generalized
preforms, or arbitrary nonclosed backgrounds is claimed here.

% END source: sections/delta_torus_test_action.tex


% BEGIN canonical R6 source: sections/delta_matroidal_test_action.tex
\section{Actual cycle relations on saturated matroidal tropical charts}
\label{sec:r6-matroidal-delta-action}

This section supplies a geometric condition under which an actual
analytic pullback relation is stronger than an ambient cycle relation,
but is still annihilated over every allowed background.  The condition
concerns one fixed moment chart and its calibrated tropical image.  No
existence of such a chart on an arbitrary Berkovich space is asserted.

\begin{definition}[A saturated matroidal chart]
\label{def:r6-saturated-matroidal-chart}
Let $U$ be an allowed pure $n$-dimensional space over the original $K$,
and let $t=t_f:U\to\Omega\subset\mathbf R^r$ be original unit moments.
The following data will be called a saturated matroidal chart.
\begin{enumerate}
\item $t$ is proper onto its image $C\subset\Omega$, and $C$, with its
      induced integral affine structure, is a boundaryless rational
      tropical manifold of dimension $n$.
\item There is a positive constant $m$ such that the following calibrated
      affinoid condition holds locally over every point of $C$.
      One can choose arbitrarily small compact polyhedral neighborhoods
      $P\Subset\Omega$, with the point in $P^\circ$, for which
      $A_P=t^{-1}(P)$ is affinoid,
      \[
       t\bigl(\Sigma(A_P,f)\cap\partial A_P\bigr)\subset\partial P,
       \qquad T(A_P,f)=m[C]\quad\hbox{on }P^\circ.
      \]
      Here $T(A_P,f)$ and its weights are exactly the original calibrated
      tropicalization of Mihatsch, Definition~3.4 and equation~(3.8).
\end{enumerate}
The constant $m$ can equally be allowed to be a locally constant positive
function on $C$, by making the assertions local over $\Omega$.
Restriction to a smaller real open subset of $\Omega$ always means
restriction of $U$ to its full inverse image.
\end{definition}

The hypotheses explicitly retain the proper fibers.  They do not say
that $U$ is its skeleton, or that an arbitrary ordinary subopen of $U$
has the same calibrated tropical image.

\begin{proposition}[The required tropical duality]
\label{prop:r6-matroidal-cocycle-input}
\textnormal{\texttt{EXTERNAL\_KNOWN}.}
If $L$ is a rational tropical linear space, then
\[
 \mathrm{PP}(L)/\mathrm{LPP}(L)
       \longrightarrow Z_*^{\mathrm{fan}}(L),
       \qquad P\longmapsto P\cdot[L]
\]
is an isomorphism.  Here the source is the direct limit over rational
fan subdivisions, and $\mathrm{LPP}$ is the ideal generated by global
linear functions.  The local affine version is an isomorphism of
sheaves on every boundaryless rational tropical manifold.  These
statements remain true after tensoring the integral groups with
$\mathbf R$.
\end{proposition}

\begin{proof}[Source and normalization]
These are Gross--Shokrieh, Lemma~4.4 and Theorem~4.5,
\cite[printed pp.~13--14]{GS21}.  Their Section~4.3, printed pp.~12--13,
defines the local affine piecewise polynomial sheaf and quotients by
the ideal generated by affine functions; Remark~4.2 explains why the
affine convention is needed at points with nonintegral coordinates.
Their Lemma~4.1 identifies the degree-one action with tropical divisor
intersection.  In our maximum convention this is $D\varphi$; a
simultaneous sign reflection in the real moment coordinates changes
neither the assertion nor its calibration.  Tensoring is exact over
$\mathbf Z$ with the torsion-free, hence flat, module $\mathbf R$.
Equivalently, on each fixed finite rational fan the relevant
polynomial and balanced-weight spaces are obtained from the integral
ones by scalar extension.  This proves the stated real-weight version.
\end{proof}

\begin{lemma}[From an actual zero form to a tropical zero cycle]
\label{lem:r6-matroidal-actual-detector}
\textnormal{\texttt{DIRECT\_COROLLARY}} of the actual delta-form definition.
On a saturated matroidal chart, let $\gamma$ be a real delta-form on
$\Omega$. If $t^\star\gamma=0$ as an actual analytic delta-form on
$U$, then
\[
                         \gamma\wedge[C]=0
\]
locally on $C$. This assertion applies to forms with piecewise smooth
coefficients as well as to constant-weight cycles.
\end{lemma}

\begin{proof}
Mihatsch's Lemma~4.4(5), \cite[printed p.~62]{Mih24}, applied with the
identity refinement of $f$, gives
\[
                    T(A_P,f)\wedge\gamma=0
\]
away from $t(\Sigma(A_P,f)\cap\partial A_P)$.  On $P^\circ$ the
two chart assumptions identify this equality with
$m[C]\wedge\gamma=0$.  Since $m>0$, it gives the claimed equality.
These interiors cover $C$.  This is a consequence of the full actual
delta-form vanishing; it does not test only one chosen skeleton.
\end{proof}

\begin{lemma}[Annihilation of relative cycle relations]
\label{lem:r6-matroidal-cycle-relation}
\textnormal{\texttt{NEW\_PROVED}}, using the external tropical duality.
Let $B$ be a constant-weight rational cycle of codimension $c$ on
$\Omega$, with real weights.  If $B\wedge[C]=0$ on a relative open
subset $C\cap W$, then
\[
                           \mathcal A_B(S)=0
                       \quad\hbox{on }t^{-1}(W)
\]
for every allowed fixed positive separately closed background $S$.
The same assertion holds for finite real linear relations between
such cycles, and on signed mixed tails over $S$.
\end{lemma}

\begin{proof}
Work near a point $b\in C\cap W$.  Translate $b$ to zero.  The local
cone of $C$ is integrally linearly isomorphic, on its ambient span,
to a rational tropical linear space $L$: the chart uses the induced
affine sheaf, so both directions of this identification are locally
restrictions of ambient affine maps.  We make this identification;
the germ of $B$ is a rational fan cycle $B_b$.  By the ambient cycle presentation,
choose a homogeneous degree-$c$ piecewise polynomial $P$ on a finite
complete rational fan whose corner-cycle class is $B_b$.  Refine to
a simplicial fan adapted to $L$.  Compatibility of divisor
intersection gives
\[
                        (P|_L)\cdot[L]=B_b\wedge[L]=0.
\]
Proposition~\ref{prop:r6-matroidal-cocycle-input} implies, after a
finite rational subdivision, that
\[
                  P|_L=\sum_{j=1}^N\ell_j Q_j,
\]
where $\ell_j$ are restrictions of ambient real linear functions
and $Q_j$ are homogeneous piecewise polynomials of degree $c-1$
on $L$.  For $c=0$ the assertion instead says directly that $P=0$
on $L$ and no sum is needed.

Extend the chosen finite subdivision of $L$ to a finite complete
simplicial ambient subdivision.  This can be done by first taking
a finite hyperplane arrangement refining its cones and then a
compatible simplicial subdivision.  Each $Q_j$ extends to an
ambient piecewise polynomial $\widetilde Q_j$: in the Courant
presentation, keep its face-supported monomials and assign zero
coefficients to the remaining monomials.  Thus
\[
                    P-\sum_j\ell_j\widetilde Q_j
\]
vanishes pointwise on $L$.  The pointwise-image ideal in
Theorem~\ref{thm:r6-piecewise-polynomial-operator} annihilates its
operator on the inverse image of a sufficiently small neighborhood
of $b$.  The operators of the terms
$\ell_j\widetilde Q_j$ are zero because the original moment
pullback of $\ell_j$ is pluriharmonic.  This proves
$\mathcal A_B(S)=0$ there.  Local ambient cycle invariance permits
these identities to glue over $t^{-1}(W)$.  The argument takes
place in the module of signed mixed tails, so the last assertion
and real linearity follow at the same time.
\end{proof}

\begin{theorem}[A fixed-chart actual delta-test action]
\label{thm:r6-matroidal-delta-action}
\textnormal{\texttt{NEW\_PROVED}}, on the stated chart and test class.
On a saturated matroidal chart let $\mathscr B_f^{\mathrm{rat}}(U)$
be the image of those real delta-forms on $\Omega$ which locally
are finite sums
\[
                          \gamma=\sum_i\eta_i\wedge[B_i],
\]
where $\eta_i$ are piecewise smooth forms and $B_i$ are
constant-weight rational cycles with real weights.
For every allowed fixed background $S$ the prescription
\[
             \mathcal A_S(t^\star\gamma)
                  =\sum_i t^*\eta_i\wedge\mathcal A_{B_i}(S)
\]
is independent of the displayed presentation as an \emph{actual}
delta-form on $U$.  It is linear and coefficient-Radon, and all
its identities hold on every original complementary compact
smooth test.  It is compatible with restriction to saturated
subdomains.  In complementary total degree it gives a pairing
with every compactly supported test in
$\mathscr B_f^{\mathrm{rat}}(U)$.

For a background represented by an actual delta-form this pairing
agrees with the full actual delta product and its integration.
Insertion of continuous convex rational PL potentials in the
fixed moments agrees with relative BT insertion.  No assertion
identifies $\mathscr B_f^{\mathrm{rat}}(U)$ with all of $B(U)$.
\end{theorem}

\begin{proof}
It suffices to annihilate a displayed actual zero relation, one
tropical codimension at a time.  By
Lemma~\ref{lem:r6-matroidal-actual-detector}, that relation gives
\[
                      \sum_i\eta_i\wedge Z_i=0,
                       \qquad Z_i=B_i\wedge[C].
\]
These $Z_i$ are constant-weight rational cycles of dimension
$n-c$ on $C$.  Put $R_i=\mathcal A_{B_i}(S)$.
Lemma~\ref{lem:r6-matroidal-cycle-relation} shows that $R_i$
vanishes over every real open subset on which $Z_i$ vanishes.
In particular its support is contained in $t^{-1}(|Z_i|)$.
It also shows that every constant local relation among the
$Z_i$ gives the identical relation among the $R_i$.

Choose one finite common polyhedral subdivision on a compact
moment stage, subordinate to the $Z_i$ and all coefficients.
For every face $F$, the Borel current
$\mathbf1_{t^{-1}(\operatorname{relint}F)}R_i$ is tangent to
the affine span of $F$, by Lemma~\ref{lem:r6-level-tangency}.
Consequently only the tangent restrictions $\eta_i|_F$ matter.
Expand these restrictions in constant tangent tensor components.
For each component its scalar coefficient vector is in the
kernel of the fixed matrix of weights of the $Z_i$ on the
incident $(n-c)$-dimensional cells.  This follows by restricting
the displayed zero relation from those cells to $F$.

Choose a fixed basis of that finite-dimensional kernel.  Each
basis vector defines a constant cycle relation among the $Z_i$
on a real neighborhood of $\operatorname{relint}F$, after
discarding the finitely many nonincident closed top cells.
The preceding lemma annihilates the identical combination of
the $R_i$.  Restrict that current equality to the Borel stratum
and multiply by the bounded Borel scalar coordinates of the
coefficient vector in the chosen kernel basis.  Summing tangent
components, then faces, proves
\[
                          \sum_i t^*\eta_i\wedge R_i=0.
\]
This is the same coefficient-measure argument as in
Theorem~\ref{thm:r6-cycle-coefficient-relations}, with the relative
cycles $Z_i$ in place of the ambient cycles.  It does not insert
a cycle into a nonclosed coefficient-weighted background.

For compactly supported $\beta=t^\star\gamma$, the support of
$\gamma\wedge[C]$ is contained in the compact set
$t(\operatorname{supp}\beta)$.  Indeed, outside that set the
actual form vanishes on a saturated open, and the detector gives
the claimed tropical vanishing.  The relation argument just
proved then shows that $\mathcal A_S(\beta)$ has support in
$t^{-1}(t(\operatorname{supp}\beta))$, which is compact by
properness.  Its top-degree mass therefore defines the stated
pairing.  The remaining comparison and insertion assertions
follow termwise from the ambient cycle action,
Theorem~\ref{thm:r6-bt-delta-background-comparison}, and canonical
piecewise coefficient multiplication.
\end{proof}

\begin{example}[Products of three-punctured lines]
\label{ex:r6-matroidal-pair-of-pants}
For every $n\geq1$, take
\[
       X=(\mathbf A^1_K\setminus\{0,1\})^n,
       \qquad f=(z_1,1-z_1,\ldots,z_n,1-z_n),
\]
and a real open $\Omega\subset\mathbf R^{2n}$.
The saturated domain $U=t_f^{-1}(\Omega)$, if nonempty, has the
chart data above with $m=1$.  In one factor the image is the
three-ray fan
\[
 \{(s,0):s\leq0\}\ \cup\ \{(0,s):s\leq0\}
       \ \cup\ \{(s,s):s\geq0\}.
\]
On its first two rays the parameters $z$ and $1-z$, respectively,
give degree-one annulus coordinates; on the third ray use $z$.
The primitive slopes are $(1,0)$, $(0,1)$ and $(1,1)$ up to the
outward signs, so the original calibrated weights are one.
All radii are ordinary Berkovich Gauss radii; no extension or
rescaling of $K$ is used.  Taking products gives the product fan
with product calibration one.  It is a tropical linear space,
as is also explicit in \cite[Lemma~3.1]{GS21}.

The map is the closed immersion into the torus defined by
$T_{2i-1}+T_{2i}=1$.  Inverse images of compact coordinate boxes
are therefore affinoid closed subspaces of compact polyannuli,
and the moment map is proper.  Boxes with a given point in
their interior can be chosen inside $\Omega$.  Their analytic
boundary maps to the box boundary: the closed immersion is
boundaryless, and the polyannulus boundary lies over that
boundary.  The calibrated tropicalization on the box interior
is the indicated product fan.  Thus the theorem applies to
arbitrary allowed positive separately closed $S$ on these
$n$-dimensional domains.  In dimensions greater than one this
allows genuinely non-smooth backgrounds; it is not a theorem
only about currents coming from the displayed curves.
\end{example}

\paragraph{The remaining scope.}
The theorem is about actual relations on each fixed saturated
chart.  Its proof uses the whole calibrated tropical image in
every compact stage.  An arbitrary ordinary subopen may omit
part of a tropical fiber and fail these hypotheses.  Ordinary
sheaf locality for all such subopens, arbitrary analytic moment
images, and tests outside the indicated constant-cycle-generated
class remain separate questions.  The external tropical
Poincare duality does not assert that a general analytic space
admits matroidal charts or that its actual delta-form sheaf is
generated by this one presentation.

% END source: sections/delta_matroidal_test_action.tex


% BEGIN canonical R6 source: sections/delta_module_generation.tex
\section{Every real delta-form is locally generated over piecewise smooth forms}
\label{sec:r6-delta-module-generation}

This section proves an exact finite generation statement.  Its proof
uses a decomposition of smooth coefficients by their restrictions to
coordinate faces.  It uses no polynomial density, formal Taylor series,
or convergence against singular tests.  In particular a smooth flat
remainder is retained exactly.

\begin{lemma}[Anchored face decomposition]
\label{lem:r6-anchored-face-decomposition}
Let $J$ be a finite set and let $F$ be smooth on a neighborhood of the
closed coordinate cube $[0,\varepsilon)^J$.  For $A\subset J$ put
\[
 F_A(x_A)=\sum_{B\subset A}(-1)^{|A|-|B|}F(x_B,0_{J\setminus B}).
\]
Then
\[
                         F(x)=\sum_{A\subset J}F_A(x_A).
\]
Every $F_A$ is smooth and vanishes whenever $x_a=0$ for one $a\in A$.
These formulas commute with every fixed real linear relation between
the functions involved.  If one first restricts $F$ to a coordinate
face containing $A$, its anchored $A$-component remains $F_A$.
\end{lemma}

\begin{proof}
In the sum over $A$, the coefficient of the evaluation
$F(x_B,0_{J\setminus B})$ is
$\sum_{A:B\subset A\subset J}(-1)^{|A|-|B|}$, which is zero unless
$B=J$, when it is one.  This proves the identity.  If $x_a=0$,
pair the subsets $B$ not containing $a$ with $B\cup\{a\}$;
their evaluations agree and their signs are opposite.  Smoothness,
linearity, and compatibility with the indicated restrictions follow
directly from the finite formula.
\end{proof}

\begin{theorem}[Local piecewise smooth module generation]
\label{thm:r6-delta-module-generation}
\textnormal{\texttt{NEW\_PROVED}} by the finite construction below;
no claim of external priority is made.
Let $T$ be a real delta-form of tridegree $(a,b,c)$ on an open subset
of $\mathbf R^r$, in Mihatsch's balanced-polyhedral-current sense.
Every point has a neighborhood on which
\begin{equation}
                             T=\sum_{j=1}^N\eta_j\wedge C_j,
 \label{eq:r6-full-module-generation}
\end{equation}
where $\eta_j$ is a piecewise smooth $(a,b)$-form and $C_j$ is a
constant-weight real tropical cycle of codimension $c$.
The sum is finite.  If the supporting polyhedra have rational
directions, the $C_j$ can be chosen with rational directions.
All weights and coefficients may be real.

The equality holds in the real delta algebra itself and hence
survives its pullbacks and refinements.  No positivity, closedness,
or diagonal-degree hypothesis is imposed on $T$.
\end{theorem}

\begin{proof}
Fix the point, translate it to zero, and take a finite polyhedral
presentation near it.  Discard the closed cells not incident to
zero by shrinking the neighborhood.  Each remaining cell agrees
there with its local cone.  Refine these cones to a finite
complete simplicial fan $\Delta$; cones off the original support
have coefficient zero.  This is done only in the real moment
space and changes neither the current nor its weights.  Put
$d=r-c$ and choose weights $\mu_\sigma$ for the $d$-dimensional
cones, absorbing any scalar factors into their coefficient forms.

Choose one nonzero generator $v_\rho$ for each ray and let
$\lambda_\rho$ be the corresponding Courant function.  On a
cone $\sigma$, its own ray functions are independent linear
coordinates.  Choose one ordering of all rays.  There are unique
smooth coefficient functions $F_{\sigma,I,J}$ such that
\[
       \alpha_\sigma
          =\sum_{\substack{I,J\subset\sigma(1)\\|I|=a,\ |J|=b}}
             F_{\sigma,I,J}(\lambda_{\sigma(1)})
                   d'_P\lambda_I\wedge d''_P\lambda_J\big|_\sigma.
\]
Their smoothness up to all coordinate faces follows from the
definition of a smooth form on a closed polyhedron.  There are
only finitely many functions.  Choose a common small
$\varepsilon>0$ such that all positive coordinate cubes
$0\leq\lambda_\rho<\varepsilon$ under consideration, together
with every point obtained by setting any coordinates to zero,
remain inside the original presentation neighborhood.  Each
coefficient then has a smooth extension to an open neighborhood
of its coordinate cube.

Apply Lemma~\ref{lem:r6-anchored-face-decomposition} to every
$F_{\sigma,I,J}$.  Denote its anchored $A$-component by
$F_{\sigma,I,J,A}(x_A)$.  Fix a triple $(A,I,J)$ which occurs
in at least one cone, and let $\tau$ be the face spanned by
$A\cup I\cup J$.  For each $d$-cone $\sigma\supset\tau$ put
\[
                         f_\sigma(x_A)=F_{\sigma,I,J,A}(x_A),
\]
using zero when the coefficient is absent.  If $A$ is empty
these are simply constants.

We claim that, for every $x_A\in[0,\varepsilon)^A$, the vector
$(f_\sigma(x_A))_{\sigma\supset\tau}$ is a balanced constant
weight vector on the star of $\tau$.  Let $\kappa$ be a
$(d-1)$-cone containing $\tau$.  The balancing criterion for
$T$, \cite[Definition~3.2 and Theorem~3.3]{Mih23}, says
\[
    \sum_{\sigma\supset\kappa}
         \alpha_\sigma|_\kappa\otimes n_{\sigma,\kappa}
                 \in A^{a,b}(\kappa)\otimes N_\kappa.
\]
Compare the coefficients of
$d'_P\lambda_I\wedge d''_P\lambda_J$ in the independent
coordinate basis on $\kappa$.  Indices outside $\kappa$
restrict to zero and contribute nothing.  The remaining
identity is a fixed linear relation between smooth functions
of the coordinates on $\kappa$, with values in
$\mathbf R^r/N_\kappa$.  Apply its anchored $A$-projector.
Because $A\subset\kappa$, compatibility with restriction in
Lemma~\ref{lem:r6-anchored-face-decomposition} gives exactly
\begin{equation}
           \sum_{\sigma\supset\kappa}
                      f_\sigma(x_A)n_{\sigma,\kappa}
                                \in N_\kappa.
 \label{eq:r6-anchored-star-balance}
\end{equation}
This holds throughout the indicated cube, including its faces.
It proves the claim.

Let $K_\tau$ be the finite-dimensional real kernel defined by
all equations \eqref{eq:r6-anchored-star-balance}.  Choose a
fixed basis $b^1,\ldots,b^s$.  There are smooth functions
$g_1,\ldots,g_s$ on the $A$-cube such that
\[
                    (f_\sigma(x_A))_\sigma
                                  =\sum_{l=1}^s g_l(x_A)b^l.
\]
Indeed choose a fixed linear coordinate map on $K_\tau$ and
apply it to the vector of smooth functions.  Its coordinates
extend smoothly to a neighborhood of the cube.  Since every
$f_\sigma$ vanishes when any $A$-coordinate is zero, the same
is true of every $g_l$ on those coordinate faces.

Define the constant cycle
\[
                      C_{\tau,l}
                         =\sum_{\sigma\supset\tau}
                              b^l_\sigma
                                  [\sigma+N_\tau,\mu_\sigma].
\]
The projected cones $\sigma/N_\tau$ form the star fan.
The lifted cones $\sigma+N_\tau$ have the same tangent
spaces and the same weights as the original $\sigma$.
Their facet normal equations are precisely those defining
$K_\tau$, so $C_{\tau,l}$ is balanced.  This includes the
case $\dim\tau=d$, when the quotient is zero-dimensional.
There is no lattice index or change of calibration: choosing
a weight on $N_\tau$, dividing by it in the quotient, and
multiplying by that same weight in the lift leaves
$\mu_\sigma$ unchanged.

On a smaller neighborhood of zero put
\[
                    \eta_{A,I,J,l}
                       =g_l(\lambda_A)
                             d'_P\lambda_I\wedge d''_P\lambda_J.
\]
It is piecewise smooth.  The Courant functions are continuous
PL, their piecewise differentials agree tangentially on all
faces, and composition with the displayed smooth functions
has the same compatibility.  Its tangent restriction to any
cone missing a ray of $\tau$ is zero: a missing ray in $A$
sets one coordinate of $g_l$ to zero, whereas a missing ray
in $I$ or $J$ makes the corresponding differential zero.

We now verify the exact recovery.  On the open star of $\tau$,
the lifted cone $\sigma+N_\tau$ agrees with $\sigma$ wherever
the coefficient form can contribute.  More explicitly, if
$x$ is in the relative interior of a cone $\sigma'$ containing
$\tau$ and its quotient lies in $\sigma/N_\tau$, that quotient
lies in the relative interior of $\sigma'/N_\tau$.
The star fan intersection property implies
$\sigma'\subset\sigma$.  Thus on the top-dimensional pieces
one obtains exactly the original cone $\sigma$ and its weight.
On every other cone the tangent coefficient is zero by the
preceding paragraph.  Lower-dimensional intersections do
not change a polyhedral current.  It follows that
\[
 \sum_l\eta_{A,I,J,l}\wedge C_{\tau,l}
      =\sum_{\sigma\supset\tau}
             f_\sigma(\lambda_A)
                    d'_P\lambda_I\wedge d''_P\lambda_J
                                \wedge[\sigma,\mu_\sigma].
\]
Summing over the finitely many triples $(A,I,J)$ and applying
the anchored decomposition recovers every $\alpha_\sigma$,
hence $T$.  All constructions took place in the chosen small
neighborhood, so no extension of the smooth coefficients to
an entire fan was assumed.

If the original support directions were rational, complete
and simplicial subdivisions can be chosen rational, and the
star lifts preserve rational directions.  Finally equality
of the resulting polyhedral currents is equality in the real
delta algebra, by its definition and the refinement
equivalence of \cite[Remark~2.3]{Mih23}.  This proves all the
assertions, including arbitrary tridegrees.
\end{proof}

\begin{corollary}[Exact local generation of actual analytic delta-forms]
\label{cor:r6-analytic-delta-module-generation}
\textnormal{\texttt{DIRECT\_COROLLARY}} of
Theorem~\ref{thm:r6-delta-module-generation} and the definition of
Mihatsch's actual analytic delta-forms.
On every original analytic moment neighborhood, each actual
delta-form is locally a finite sum of a piecewise smooth form
times pullbacks of constant real cycles.  By the real cycle
presentation, the latter factors can in turn be expressed as
finite real linear combinations of products of $D_\delta\varphi$,
where $\varphi$ is a continuous PL difference of classical psh
potentials in the same original moments.

Consequently, on any class of backgrounds for which the
constant-cycle-generated action has been proved to annihilate
\emph{all actual analytic relations}, that action covers the
whole actual delta-test sheaf.  The generation result does
not itself prove that relation-annihilation property for an
arbitrary background or an arbitrary analytic moment image.
\end{corollary}

\paragraph{Relation to the fixed-chart theorems.}
For a full-coordinate saturated torus domain, the ambient
real cycle action and its coefficient-relation proof apply
to every delta-form in the fixed coordinate presentation.
For the saturated matroidal theorem, the present generation
argument applies to its rational-direction test class.
The matroidal duality input used there is a rational fan
theorem; no additional nonrational relative duality is
inferred here.  Both chart statements retain their explicit
restriction to saturated domains until ordinary sheaf
locality is established by a separate argument.

% END source: sections/delta_module_generation.tex


% BEGIN canonical R6 source: sections/delta_degree_one_action.tex
\section{Degree-one delta relations over every allowed background}
\label{sec:r6-delta-degree-one}

The weak approximation hypothesis is unnecessary in total degree
$(1,1)$.  The reason is an order property specific to curvature
currents of continuous scalar potentials.  Freezing a continuous
coefficient produces an arbitrarily small \emph{order error}; it
need not produce an exact constant relation on a neighborhood.

All domains in this section are ordinary opens in the analytic
spaces allowed by B, over the original finite $K/\mathbf Q_p$.
The positive cone and tests are the original ones.  We use R,
Theorem~2.9, through B's authorized ordinary local algebraic
localization.  Thus a finite continuous potential with positive
original curvature has locally uniform original smooth positive
approximants.  No extension of this input to a new base field is
being made.

\begin{lemma}[Continuous curvature order transfers to a fixed background]
\label{lem:r6-continuous-curvature-order}
Let $S$ be an arbitrary allowed symmetric positive separately closed
$(p,p)$-current, with $p<n$.  If $h$ is finite continuous and
$D[h]\geq0$ on an ordinary open, then
\[
                              D(hS)\geq0.
\]
In particular, if $g_1,\ldots,g_m$ are continuous classical psh,
$G=\sum_i g_i$, and $c_i\in\mathbf R$ satisfy
\[
       -\varepsilon D[G]\leq D\left[\sum_i c_i g_i\right]
                                    \leq\varepsilon D[G],
\]
then the same inequalities hold after replacing $D[g]$ by $D(gS)$.
\end{lemma}

\noindent\textbf{Status: \texttt{DIRECT\_COROLLARY}} of the
authorized continuous regularization theorem R and fixed coefficient
multiplication.

\begin{proof}
Locally choose $h_j\to h$ uniformly with $Dh_j$ original-positive
smooth forms.  Since $S$ is coefficient-Radon, $h_jS\to hS$ on
every original test.  Separate closedness gives
$D(h_jS)=Dh_j\wedge S\geq0$.  Applying the continuous distributional
operator $D$ and passing to the limit gives $D(hS)\geq0$.
The assertion is local and hence glues.  Apply it to the two
continuous potentials $\varepsilon G\pm\sum_i c_i g_i$ for the
last statement.  This argument uses uniform smooth approximation
directly and does not require a new classicality equivalence.
\end{proof}

\begin{theorem}[Continuous-coefficient degree-one relation descent]
\label{thm:r6-continuous-curvature-module}
Let $g_1,\ldots,g_m$ be continuous classical psh functions and let
$f_1,\ldots,f_m$ be finite continuous real functions on the same
ordinary domain.  Set
\[
                  A=\sum_i f_iD[g_i],\qquad
                  A_S=\sum_i f_iD(g_iS).
\]
These are actual coefficient-Radon currents.  Then
\[
                       A\geq0\quad\Longrightarrow\quad A_S\geq0.
\]
Consequently $A=0$ implies $A_S=0$ on all original tests.
The last implication also holds when each $g_i$ is a continuous
local difference of two continuous classical psh functions.
\end{theorem}

\noindent\textbf{Status: \texttt{NEW\_PROVED}} by the local order
argument below.  The background is arbitrary and fixed; no membership
in $\mathcal W$ is assumed.

\begin{proof}
Put $M_i=D[g_i]$, $T_i=D(g_iS)$, and $T=\sum_iT_i$.
All $M_i,T_i$ are positive and separately closed.  Fix
$\varepsilon>0$ and a point $x$.  There is an ordinary neighborhood
$O_x$ on which
$|f_i-f_i(x)|<\varepsilon$ for every $i$.  Write $c_i=f_i(x)$.
If $A\geq0$, positivity of the coefficient measures gives on $O_x$
\[
                         \sum_i c_iM_i\geq-\varepsilon\sum_iM_i.
\]
The preceding lemma transfers this order inequality to
$\sum_i c_iT_i\geq-\varepsilon T$.  Continuous scalar multiplication
of the positive $T_i$ also gives
\[
       A_S=\sum_i c_iT_i+\sum_i(f_i-c_i)T_i
                                                 \geq-2\varepsilon T
                           \quad\hbox{on }O_x.
\]

For a fixed compact original positive test choose a finite subcover
of its support by such $O_x$ and a nonnegative original smooth
partition subordinate to that cover.  Applying the local inequality
to the partitioned tests and summing gives
$\langle A_S,\eta\rangle\geq-2\varepsilon\langle T,\eta\rangle$.
Let $\varepsilon\downarrow0$.  This proves positivity on every
original positive test, with the same cone convention as for $S$.
If $A=0$, apply the implication both to $A$ and to $-A$.
The positive test cone spans the symmetric test space; all these
currents are symmetric.  Thus $A_S=0$ on every original test.

For local differences expand each $g_i=g_i^+-g_i^-$ and replace
the corresponding coefficient by $f_i$ and $-f_i$.  On any compact
test support only finitely many local expansions and a partition
are needed.  The already proved statement applies.  No freezing
step asserted that $\sum_i c_iD[g_i]$ vanishes on a neighborhood.
Only its two-sided small order bound was needed.
\end{proof}

\begin{lemma}[Local continuous piecewise smooth potentials are differences]
\label{lem:r6-continuous-pws-dsh}
Every finite continuous piecewise smooth scalar in a finite original
moment presentation is locally a difference of two continuous
piecewise smooth classical psh functions.  The same holds for a
finite product of such scalars.
\end{lemma}

\begin{proof}
It suffices to prove a real local statement before the moment
pullback.  Work on a convex ball with compact closure in the
presentation domain, and refine the finite subdivision by all its
facet hyperplanes.  A continuous piecewise smooth function $F$ has
uniformly bounded Hessians on the finitely many closed pieces.
Across a facet its tangential derivatives agree, so its gradient
jump is normal to the facet and is uniformly bounded there.
Choose affine defining functions $\ell_H$ for the finitely many
hyperplanes, with normalized gradients.  For sufficiently large
positive constants $A,B$, the function
\[
                     Q(x)=A|x|^2+B\sum_H|\ell_H(x)|
\]
is convex, and $F+Q$ has positive semidefinite Hessian in every
open piece and nonnegative normal derivative jumps across every
facet.  Restriction to a generic segment therefore has nonnegative
second derivative on each subinterval and nonnegative first-derivative
jumps.  It is convex.  Continuity and approximation of segments
give convexity for all segments in a smaller ball.  Hence
$F=(F+Q)-Q$ is a difference of convex continuous piecewise smooth
functions there.

Convex real functions of original logarithmic unit moments are
classical psh: one may mollify the convex functions on a smaller
real ball, obtaining locally uniform original smooth positive
approximants, or apply convexity to the harmonic coordinate
pullbacks on every test curve.  This proves the asserted local
decomposition over the original field.  Products of continuous
piecewise smooth scalars are again of that same class, so the
same real argument applies to them.
\end{proof}

For $(1,1)$-forms, ``symmetric'' below means the usual CLD
prime--double-prime symmetry; in a moment coefficient matrix it
is ordinary matrix symmetry.  Write $\operatorname{Sym}$ for that
projection.  Constant codimension-one cycle factors are symmetric.

\begin{lemma}[A continuous Hessian presentation in total degree one]
\label{lem:r6-degree-one-hessian-presentation}
Every symmetric actual Mihatsch delta-form of total degree $(1,1)$
locally has a finite presentation
\[
                         \alpha=\sum_i f_iD_\delta g_i,
\]
where $f_i$ are continuous piecewise smooth scalars and $g_i$ are
continuous piecewise smooth local differences of continuous
classical psh functions.  Each $g_i$ may therefore be expanded
into two admissible positive potentials in that formula.

For a general $(1,1)$-delta-form its antisymmetric part is piecewise
smooth, so it is covered by the already canonical PS multiplication.
\end{lemma}

\begin{proof}
Module generation splits a local total-degree-one presentation into
a PS $(1,1)$-form and continuous PS scalars times constant
codimension-one cycles.  The real cycle theorem writes the latter
cycles as finite sums of $D_\delta$ of continuous PL functions,
which have the required difference decompositions.

For the symmetric PS part use the anchored construction of
Theorem~\ref{thm:r6-delta-module-generation} with codimension zero.
Every constant balanced codimension-zero cycle on a connected real
neighborhood is a constant multiple of the ambient integration
current.  The construction consequently writes the PS form as a
finite sum
\[
                   h\,\operatorname{Sym}
                         (d'_P\lambda_i\wedge d''_P\lambda_j),
\]
with continuous PS $h$ and Courant PL scalars $\lambda_i,\lambda_j$.
The actual delta Leibniz rule gives
\[
 \operatorname{Sym}(d'_P\lambda_i\wedge d''_P\lambda_j)
   =\tfrac12\bigl(D_\delta(\lambda_i\lambda_j)
                 -\lambda_iD_\delta\lambda_j
                 -\lambda_jD_\delta\lambda_i\bigr).
\]
Lemma~\ref{lem:r6-continuous-pws-dsh} applies to every displayed
potential and coefficient product.  This proves the required
presentation.  The positive-codimension part has only scalar
coefficients and symmetric cycle factors, so the antisymmetric
part of a general form comes entirely from its PS part.
\end{proof}

\begin{lemma}[Piecewise smooth product rule with a Radon derivative]
\label{lem:r6-pws-radon-leibniz}
Suppose $R$ belongs to the tangential coefficient-Radon module of
Theorem~\ref{thm:r6-pws-multiplication}, and $d'R$ is coefficient-Radon.
For a continuous PS scalar $f$,
\[
                   d'(fR)=d'_Pf\wedge R+f\,d'R.
\]
The corresponding $d''$ identity also holds.  In particular, for
continuous PS local differences $f,g$ and the fixed $S$,
\begin{align*}
 D(fgS)={}&fD(gS)+gD(fS)\\
         &+d'_Pf\wedge d''_Pg\wedge S
          +d'_Pg\wedge d''_Pf\wedge S.
\end{align*}
\end{lemma}

\begin{proof}
Mollify a continuous PS real representative of $f$ on a smaller
compact moment stage.  The functions converge uniformly, and their
first derivatives converge tangentially on each face with uniformly
bounded arrays.  The conormal annihilation and fixed coefficient
dominated convergence in Proposition~\ref{prop:r6-pws-smoothing}
give $d'f_\varepsilon\wedge R\to d'_Pf\wedge R$ in local coefficient
variation.  Uniform convergence also gives
$f_\varepsilon d'R\to f\,d'R$.  Pass to the limit in the ordinary
smooth product rule.  This proves the first assertion and its
double-prime counterpart.

For the second assertion, Corollary~\ref{cor:r6-pws-first-derivative}
gives $d''(gS)=d''_Pg\wedge S$, which belongs to the same tangential
module by its PS module property.  Its $d'$ derivative $D(gS)$ is
coefficient-Radon because $g$ is a local difference.  Apply the first
assertion to this current, and similarly with $f,g$ exchanged.
The first-derivative product rule for the scalar $fg$ then yields
the displayed formula.  This passage uses only fixed Radon measures;
it does not require a moving-measure chain rule.
\end{proof}

\begin{theorem}[Canonical degree-one action for arbitrary fixed $S$]
\label{thm:r6-arbitrary-background-degree-one-delta}
For every allowed positive separately closed $(p,p)$-current $S$,
$p<n$, every actual delta-form $\alpha$ of total degree $(1,1)$ has
a canonical ordinary coefficient-Radon current $\alpha\wedge S$.
It is linear, independent of all presentations and refinements,
and compatible with restriction and gluing on ordinary opens.
It agrees with canonical PS multiplication, with
$D_\delta g\wedge S=D(gS)$ for the continuous local differences
above, and with the actual delta product when $S$ is a delta-form.
Its support is contained in
$\operatorname{supp}\alpha\cap\operatorname{supp}S$.

If the ordinary current $[\alpha]$ is symmetric and original-positive,
then $\alpha\wedge S$ is original-positive.  On each fixed finite
moment/polyhedral presentation it has the finite coefficient $C^0$
bounds supplied by its Hessian presentation.
No hypothesis $S\in\mathcal W_{\mathrm{loc}}$ occurs.
\end{theorem}

\noindent\textbf{Status: \texttt{NEW\_PROVED}} by continuous
order descent.  This is a total-degree-one assertion, not an
all-codimension analytic relation theorem.

\begin{proof}
For symmetric $\alpha=\sum_i f_iD_\delta g_i$ define
$\alpha\wedge S=\sum_i f_iD(g_iS)$.  If a presentation represents
zero, ordinary delta integration and continuous scalar compatibility
give $\sum_i f_iD[g_i]=0$.  Theorem
\ref{thm:r6-continuous-curvature-module} annihilates it after relative
insertion.  Subtracting presentations gives independence.  The same
theorem proves the stated positivity for every positive symmetric
ordinary $[\alpha]$.

For a general form add the canonical action of its antisymmetric
PS part.  Prime--double-prime symmetry commutes with the real delta
operations and their actual pullbacks, so this decomposition is
intrinsic.  To verify agreement with PS multiplication, use the
Hessian presentation of its symmetric part.  The last identity of
Lemma~\ref{lem:r6-pws-radon-leibniz} identifies each symmetric
Hessian expression with its already defined PS wedge product.
The antisymmetric part agrees by definition.  Ordinary smooth
multiplication is a special case.

All formulas restrict ordinarily.  Independence proves their
agreement on overlaps and under all original refinements.  If
$\alpha$ vanishes on an ordinary open its presentation there is a
zero relation, so its action vanishes; if $S$ vanishes there the
formula also vanishes.  This proves the support assertion.  On an
actual delta-form background the ordinary BT--delta comparison and
continuous multiplication identify the formula with the actual
product.  If only the ordinary current $[\alpha]$ is initially
known to be symmetric, actual delta faithfulness
(Proposition~\ref{prop:r6-delta-faithfulness}) gives
$\alpha=\operatorname{Sym}\alpha$, so the antisymmetric PS action
is zero; this justifies the theorem's stated positivity hypothesis.
Finally, on a fixed finite presentation the scalar
coefficients in the anchored Hessian expression are fixed finite
linear combinations of coordinate-face evaluations of the original
arrays.  The potentials and their positive/negative decompositions
are fixed.  Their coefficient variations therefore give the claimed
$C^0$ bounds on smaller fixed compact stages.
\end{proof}

\begin{corollary}[Full complementary delta-tests in remaining degree one]
\label{cor:r6-all-background-qone-full-delta}
Let $S$ be an arbitrary allowed $(n-1,n-1)$-current.  Then
\[
                  \Lambda_S(\alpha)=\int_U\alpha\wedge S,
                  \qquad \alpha\in B_c^{1,1}(U),
\]
is a canonical full actual delta-test extension.  It is separately
closed, agrees with $S$ on all original tests, and is nonnegative
whenever the ordinary current $[\alpha]$ is original-positive.
It is determined by the ordinary current $S$, not by an auxiliary
weak approximation or a choice of representatives.
\end{corollary}

\begin{proof}
Compact actual support and the preceding theorem give the finite
pairing and its claimed independence and positivity.  To prove
closedness, every actual form of total degree $(0,1)$ is PS.
The codimension-zero anchored construction writes it locally as
$\sum_j f_jd''_P g_j$, where $f_j,g_j$ are continuous PS scalars.
For a compact test, an original smooth partition makes each $f_j$
compactly supported in its moment neighborhood.  The actual delta
Leibniz rule, agreement with PS multiplication, and Lemma
\ref{lem:r6-pws-radon-leibniz} give, term by term,
\[
        (d'(f_jd''_P g_j))\wedge S
                        =d'(f_jd''_P g_j\wedge S).
\]
The right side is the derivative of a compactly supported ordinary
current, so its integral is zero by an original smooth cutoff equal
to one near its support.  Summing proves $\Lambda_S(d'\beta)=0$.
The prime--double-prime exchanged argument proves the other identity.
No higher-codimension delta product is used in this argument.
\end{proof}

\begin{corollary}[Uniqueness in the stated positive closed class]
\label{cor:r6-qone-delta-extension-unique}
For $p=n-1$, the functional in
Corollary~\ref{cor:r6-all-background-qone-full-delta} is the unique
extension satisfying Definition~\ref{def:r6-generator-positive-extension}.
Thus existence and uniqueness on the whole complementary actual
delta-test space hold for every allowed background of remaining degree one,
under that explicit symmetry, separate closedness, and
generator-positivity convention.
\end{corollary}

\noindent\textbf{Status: \texttt{DIRECT\_COROLLARY}} of the forced-value
theorem and the finite Hessian presentation.

\begin{proof}
Our functional is symmetric and separately closed.  Each generator
in Definition~\ref{def:r6-generator-positive-extension}, now with
$k=0$ or $1$, has positive ordinary current: it is a nonnegative
continuous scalar times an original positive smooth form or a
positive curvature current.  The positivity assertion of
Theorem~\ref{thm:r6-arbitrary-background-degree-one-delta} therefore
verifies the required generator cone.

Conversely, take any extension with the stated properties.  Localize
a compact test by a finite original smooth partition.  Each symmetric
localized test has the Hessian presentation of
Lemma~\ref{lem:r6-degree-one-hessian-presentation}; expand its
differentiated potentials into admissible positive differences.
The forced-value theorem applies with the partition function as the
compact original scalar test and with the remaining continuous PS
coefficient as its weight.  It determines every term to be its
relative BT value.  The antisymmetric part is killed by symmetry of
both functionals.  Summing proves equality on the whole actual
complementary test space.  This is finite generation, not a density
argument for singular tests.
\end{proof}

\begin{corollary}[Weighted decreasing convergence on degree-one delta-tests]
\label{cor:r6-arbitrary-background-degree-one-monotone}
Let $S$ be any fixed allowed $(p,p)$-background, let $k\geq0$ satisfy
$p+k=n-1$, and let $u_{i,j}\downarrow u_i$ for $0\leq i\leq k$
be actual locally bounded classical psh rows with locally bounded
limits.  For every compact actual $\alpha\in B_c^{1,1}(U)$,
\[
 \int_Uu_{0,j}\,\alpha\wedge T_{S,k}(u_{1,j},\ldots,u_{k,j})
       \longrightarrow
 \int_Uu_0\,\alpha\wedge T_{S,k}(u_1,\ldots,u_k).
\]
All weights are the actual Borel functions.  The background is fixed
and arbitrary; no weak-approximation hypothesis is imposed.
\end{corollary}

\noindent\textbf{Status: \texttt{DIRECT\_COROLLARY}} of the frozen
ordinary weighted BT theorem and the new finite Hessian presentation.

\begin{proof}
The antisymmetric PS part of $\alpha$ pairs to zero with each scalar
multiple of a symmetric $(n-1,n-1)$-current.  For its symmetric part,
take a finite local Hessian presentation and a compact original smooth
partition.  After expanding each fixed signed potential, every term
has the form
\[
        \int \chi f\,u_{0,j}\,
                   T_{S,k+1}(u_{1,j},\ldots,u_{k,j},g),
\]
where $g$ is fixed continuous classical psh, $f$ is fixed continuous
PS, and $\chi$ is an original smooth compact scalar cutoff.  The
extra differentiated slot is a constant decreasing row, so the
ordinary weighted joint theorem applies after replacing $\chi f$
by an original smooth compact scalar approximant.  Such approximants
converge uniformly on the one fixed compact stage.  Relative CLN
and the common local bounds of all decreasing rows bound the
unweighted positive measures uniformly there; the weights $u_{0,j}$
are uniformly bounded there as well.  The scalar approximation
errors therefore tend to zero uniformly in $j$.  Pass first in $j$
and then in the scalar approximation, and sum the finite terms.
No simultaneously varying discontinuous coefficient measure has
been passed without that estimate.
\end{proof}

\paragraph{Remaining distinction.}
The degree-one result applies to every allowed background, including
those carried by ambient complete-polar sets.  In larger remaining
degree it supplies an action of all total-degree-one delta-forms,
but it does not yet annihilate every actual relation among several
positive-codimension factors.  The full-test theorem in arbitrary
remaining degree continues to use $\mathcal W_{\mathrm{loc}}$ or
another explicitly proved relation mechanism.

% END source: sections/delta_degree_one_action.tex


% BEGIN canonical R6 source: sections/delta_weak_approximation_bridge.tex
\section{A weak-approximation bridge to an actual delta-test class}
\label{sec:r6-delta-weak-approximation}

This section states a sufficient condition on the background;
it does not impose that condition on the frozen ordinary BT
theory. Its purpose is to isolate an independent analytic
condition which kills the additional pullback relations.
All fields, spaces, positivity cones, and original test
classes remain those fixed at the start of the round.

\begin{definition}[The local weak closure of delta-form backgrounds]
\label{def:r6-delta-weak-closure}
For an ordinary open $V$, let $\mathcal W^p(V)$ be the weak
closure, in the topology of all original compact smooth
tests, of the cone of currents $[\sigma]$, where
$\sigma\in B^{p,p}(V)$ and $[\sigma]$ is symmetric,
original-positive, and separately closed. A current belongs
to $\mathcal W^p_{\mathrm{loc}}$ if each point has one
ordinary neighborhood $V$ on which it belongs to
$\mathcal W^p(V)$. The closure is allowed to use nets.
The same fixed $V$ is used throughout each approximating
net; it is not a varying-domain notion.

Let $\mathscr B_{\mathrm{cc}}$ be the subsheaf of the actual
Mihatsch delta-form sheaf whose sections locally admit
finite presentations
\[
       \alpha=\sum_i t^*\eta_i\wedge t^\star C_i,
\]
with original moments $t$, piecewise smooth coefficients
$\eta_i$, and constant-weight real polyhedral tropical
cycles $C_i$. Nonrational directions, real weights, and
real affine offsets are allowed by the real cycle theorem.
This is an explicitly specified test class,
not an asserted description of all Mihatsch or GK forms.
\end{definition}

\begin{lemma}[Fixed continuous slots and varying positive backgrounds]
\label{lem:r6-fixed-continuous-slots-varying-background}
Suppose $S_\lambda\to S$ weakly on all original tests, where
all these currents are symmetric, original-positive, and
separately closed of the same degree. Fix continuous
classical psh potentials $v_1,\ldots,v_c$ on their common
ordinary domain. Then
\[
 T_{S_\lambda,c}(v_1,\ldots,v_c)
             \longrightarrow T_{S,c}(v_1,\ldots,v_c)
\]
on every original test. The same holds after wedging with
one fixed piecewise smooth form and testing the resulting
ordinary current on any original complementary smooth form.
\end{lemma}

\noindent\textbf{Status: \texttt{DIRECT\_COROLLARY}} of R's
fixed-continuous-multiplier lemma, the frozen positive BT
recursion, and R4's piecewise-test continuity. The fixed
continuity assumption on the inserted potentials is
essential to this particular argument.

\begin{proof}
For one slot, R's lemma gives $v_1S_\lambda\to v_1S$;
applying $D$ proves the assertion. Those derivatives are
again positive separately closed currents. Induct on the
number of fixed continuous slots, applying the same lemma
to the already convergent positive tails. For the last
assertion, wedge the fixed piecewise smooth coefficient
with the chosen original test on a common finite
presentation. It is a compatible piecewise smooth test
of complementary diagonal degree; replace it by its
symmetric part if needed. R4, Theorem 71.4, gives
convergence on that test, since weak convergence bounds
its finitely many positive mass controls. This proof
allows nets and does not invoke any monotone theorem
for simultaneously varying discontinuous potentials and
backgrounds.
\end{proof}

\begin{theorem}[Actual relation descent under local weak delta approximation]
\label{thm:r6-weak-delta-actual-descent}
Let $S\in\mathcal W^p_{\mathrm{loc}}$. For every section
$\alpha\in\mathscr B_{\mathrm{cc}}^{a,b}(U)$ there is a
canonical ordinary coefficient-Radon current
$\alpha\wedge S$. On a local hypersurface presentation of
the constant cycles it is computed by the corresponding
finite sum of piecewise smooth coefficients wedged with
the relative BT products over $S$.

This operation annihilates every actual analytic
zero-presentation of $\alpha$, commutes with ordinary
restriction and gluing, and has support contained in
$\operatorname{supp}\alpha\cap\operatorname{supp}S$.
If $a=b=n-p$ and $\alpha$ is compactly supported, it
therefore defines the finite canonical pairing
$\langle S,\alpha\rangle=\int_U\alpha\wedge S$.
When $S$ itself is an actual delta-form background this
agrees with its delta product and integration.
\end{theorem}

\noindent\textbf{Status: \texttt{NEW\_PROVED}} under the
explicit weak-approximation hypothesis. Its application
to every arbitrary positive closed background is
\texttt{CONDITIONAL} on the separate density assertion
$\mathcal W^p_{\mathrm{loc}}=$ all such backgrounds.
That density assertion remains \texttt{OPEN} in this round.

\begin{proof}
The real cycle theorem and its piecewise polynomial
presentation express each constant cycle as a finite
real linear combination of products of PL divisors.
Decompose each signed PL potential into two convex PL
functions by Lemma~\ref{lem:r6-pl-difference-convex},
and expand. Thus one obtains a finite real linear
combination of mixed products of continuous convex PL
potentials pulled back by the original moments.
Thus the proposed expression is a well-defined finite
sum of ordinary coefficient-Radon currents. The ambient
relation theorem above already removes ambient cycle
relations, but the following argument also handles the
additional analytic kernel.

Work on one fixed ordinary neighborhood on which
$S_\lambda=[\sigma_\lambda]\to S$ as in Definition
\ref{def:r6-delta-weak-closure}. On each $S_\lambda$, the
BT--delta comparison and piecewise smooth multiplication
identify the proposed expression with
$[\alpha\wedge\sigma_\lambda]$. In particular it is zero
whenever the represented actual delta-form $\alpha$ is
zero. Lemma
\ref{lem:r6-fixed-continuous-slots-varying-background}
passes every term of the finite expression to its
corresponding value over $S$. Hence every actual
zero-presentation also has zero value over $S$.
Subtracting two presentations proves independence.

Restriction is immediate from the formula and the fact
that restriction sends $\mathcal W^p(V)$ into
$\mathcal W^p(V')$ for an ordinary subopen $V'$.
Independence on overlaps glues the local expressions.
If $\alpha$ or $S$ is zero on an ordinary open, the local
formula and actual zero-presentation independence make
the resulting current zero there. This proves the support
assertion. In top degree the result is a signed Radon
measure; compact support makes its total integral finite.
The delta-form-background comparison was already used
term by term, and yields the final assertion as well.
\end{proof}

\begin{proposition}[Stability of the approximation class under all bounded slots]
\label{prop:r6-delta-weak-class-bt-stability}
If $S\in\mathcal W^p_{\mathrm{loc}}$ and
$u_1,\ldots,u_k\in\mathcal C_b$, then
\[
                  T_{S,k}(u_1,\ldots,u_k)
                        \in\mathcal W^{p+k}_{\mathrm{loc}}.
\]
In particular the class contains every bounded classical
BT product over $[1]$, and over every positive separately
closed actual delta-form background. These differentiated
potentials may be discontinuous.
\end{proposition}

\noindent\textbf{Status: \texttt{DIRECT\_COROLLARY}} of the
authorized local smooth rows, ordinary BT continuity, and
the definition as a weakly closed cone.

\begin{proof}
It suffices to insert one slot and work on one smaller
ordinary neighborhood where the prescribed decreasing
original smooth positive row $u_j\downarrow u$ exists.
The restrictions of $S$ remain in the corresponding weak
closure on that fixed neighborhood. For fixed $j$,
the map $R\mapsto Du_j\wedge R$ is weakly continuous.
It sends each approximating actual delta-form background
to an actual delta-form background whose ordinary current
is positive and separately closed. Consequently
$D(u_jS)$ belongs to the weakly closed cone
$\mathcal W^{p+1}$ there. The frozen monotone theorem gives
$D(u_jS)\to D(uS)$, so the latter is in the same closure.
Repeat on a common smaller neighborhood for finitely
many slots. The current $[1]$ is itself an actual
delta-form background. This proves the stated examples.
\end{proof}

\begin{lemma}[Piecewise-test approximation with a bounded Borel weight]
\label{lem:r6-weighted-pws-error}
For a fixed compact compatible piecewise smooth test
$\eta$, the smooth approximants and finite positive
controls in R4 can be chosen so that, for every allowed
positive separately closed $T$ and every bounded Borel
$w$ on that compact stage,
\[
 |\langle wT,\eta_\varepsilon-\eta\rangle|
   \leq C\|w\|_\infty\varepsilon^{1/2}
                         \sum_b\langle T,\Theta_b\rangle.
\]
In complementary degree zero use the uniform scalar
approximation error in place of $\varepsilon^{1/2}$.
\end{lemma}

\noindent\textbf{Status: \texttt{DIRECT\_COROLLARY}} of the
total-variation version of R4's normal-factor estimate.

\begin{proof}
R4 decomposes the error into a uniformly $O(\varepsilon)$
coefficient array and finitely many bounded normal-factor
arrays supported in tubes of width $O(\varepsilon)$.
Its Lemma 71.2 bounds the variations of the normal-factor
coefficient measures, not merely their signed total
pairings. Multiplication by $w$ increases each such
variation by at most $\|w\|_\infty$. The same statement
for the regular error follows from fixed coefficient
domination. Sum the finitely many terms. The proof uses
closedness only for $T$, not for the weighted current.
\end{proof}

\begin{theorem}[Weighted decreasing convergence on the specified delta-tests]
\label{thm:r6-weighted-monotone-delta-tests}
Fix $S\in\mathcal W^p_{\mathrm{loc}}$ and
$0\leq k\leq n-p$. Suppose every row
$u_{i,j}\downarrow u_i$, $0\leq i\leq k$, consists of
actual locally bounded classical psh functions and has
a locally bounded limit. For every compact
$\alpha\in\mathscr B_{\mathrm{cc}}^{n-p-k,n-p-k}(U)$,
\[
 \int_U u_{0,j}\,\alpha\wedge
             T_{S,k}(u_{1,j},\ldots,u_{k,j})
       \longrightarrow
 \int_U u_0\,\alpha\wedge T_{S,k}(u_1,\ldots,u_k).
\]
The multiplier is the actual Borel function. No
identification of weights modulo their ambient curvature
is made.
\end{theorem}

\noindent\textbf{Status: \texttt{NEW\_PROVED}} for this
actual delta-test class and these backgrounds. The
ordinary weighted BT convergence itself is frozen input.

\begin{proof}
The preceding stability proposition and actual relation
descent define all the currents in the integrals
independently of the chosen test presentation. Localize
the compact support by a finite original smooth
partition and choose local constant-cycle hypersurface
presentations. Each term is the pairing of a fixed
piecewise smooth coefficient $\eta$ with
\[
 u_{0,j}\,T_{S,k+c}
       (u_{1,j},\ldots,u_{k,j},v_1,\ldots,v_c),
\]
where the $v_i$ are fixed continuous convex PL classical
potentials. All these unweighted tails are positive
separately closed. Their finitely many mass controls
are bounded by relative CLN and the common local bounds
of the decreasing rows and fixed $v_i$.
The weights $u_{0,j}$ have one local bound as well.

Replace $\eta$ by its fixed smooth approximant
$\eta_\varepsilon$. The frozen weighted joint decreasing
theorem gives convergence for this original smooth test,
with the fixed $v_i$ treated as constant decreasing rows.
Lemma~\ref{lem:r6-weighted-pws-error} bounds both
approximation errors uniformly in $j$ by a quantity
tending to zero with $\varepsilon$. First pass to the
limit in $j$, then in $\varepsilon$. Sum the finite
presentations. Compact support and actual relation
descent make the resulting integrals independent of
all local choices. There is no varying background $S_j$
in this monotone statement.
\end{proof}

% END source: sections/delta_weak_approximation_bridge.tex


% BEGIN canonical R6 source: sections/full_delta_tests_over_weak_closure.tex
\section{Full actual delta-testing over the weak approximation class}
\label{sec:r6-full-delta-tests}

The exact module generation theorem removes the test-generation
restriction from the weak-approximation bridge. The hypothesis on
the background remains explicit and independent of that algebraic
generation statement.

\begin{theorem}[A full actual delta-test extension]
\label{thm:r6-full-delta-test-extension}
Let $S\in\mathcal W^p_{\mathrm{loc}}$ be as in Definition
\ref{def:r6-delta-weak-closure}. For every actual Mihatsch
delta-form $\alpha\in B^{a,b}(U)$ there is a canonical ordinary
coefficient-Radon current $\alpha\wedge S$ of bidegree
$(a+p,b+p)$. It has the following properties.
\begin{enumerate}
\item It is linear in $\alpha$, independent of all actual
      presentations and refinements, and commutes with
      restriction and gluing on ordinary opens.
\item Its support is contained in
      $\operatorname{supp}\alpha\cap\operatorname{supp}S$.
      Thus complementary compact delta-forms have finite
      pairings $\Lambda_S(\alpha)=\int_U\alpha\wedge S$.
\item It agrees with ordinary smooth multiplication and
      with the actual delta product on delta-form backgrounds.
      Curvature factors of continuous PL classical potentials
      act by their ordinary relative BT insertion.
\item It satisfies the distributional identities
      \[
       d'(\alpha\wedge S)=(d'\alpha)\wedge S,
       \qquad
       d''(\alpha\wedge S)=(d''\alpha)\wedge S.
      \]
      In particular $\Lambda_S$ is separately closed on
      compact actual delta-tests.
\item On every fixed finite moment, polyhedral, and support
      presentation it is continuous for smooth coefficient
      convergence. More precisely, on smaller fixed compact
      stages each original coefficient pairing is bounded by
      a constant times finitely many $C^0$ norms of the
      presented coefficient arrays. Derivative convergence
      is not needed for this individual pairing bound.
\end{enumerate}
Positivity is asserted on the explicit generator cone:
finite sums of nonnegative continuous piecewise smooth
scalars times original middle-positive smooth forms wedged
with convex PL classical curvature factors. No identification
with an unspecified stronger or weaker delta-positivity
convention is made.
\end{theorem}

\noindent\textbf{Status: \texttt{NEW\_PROVED}} under the
stated background hypothesis. It is a full delta-test
extension for this class, not a wedge product of two
arbitrary currents.

\begin{proof}
Theorem~\ref{thm:r6-delta-module-generation} and its analytic
corollary identify the actual sheaf $B$ with the locally
constant-cycle-generated sheaf $\mathscr B_{\mathrm{cc}}$
of the bridge. Theorem
\ref{thm:r6-weak-delta-actual-descent} therefore gives the
construction and assertions (1)--(3), on the whole actual
delta-test sheaf rather than just on a subalgebra.

For (4), work on one fixed ordinary neighborhood with
$S_\lambda=[\sigma_\lambda]\to S$ as in the definition.
Both $\alpha$ and $d'\alpha$ are actual delta-forms and
have finite presentations supplied by module generation.
The fixed-continuous-slot lemma and the piecewise-test
estimate consequently give
\[
 [\alpha\wedge\sigma_\lambda]\longrightarrow\alpha\wedge S,
 \qquad
 [(d'\alpha)\wedge\sigma_\lambda]
                        \longrightarrow(d'\alpha)\wedge S
\]
on every original compact smooth test. Faithfulness of
actual delta-forms makes $d'\sigma_\lambda=d''\sigma_\lambda=0$.
Their delta Leibniz identity gives
\[
 d'[\alpha\wedge\sigma_\lambda]
                     =[(d'\alpha)\wedge\sigma_\lambda].
\]
Distributional differentiation is weakly continuous;
passing to the limit proves the first claimed identity.
The second is identical. These local identities glue.
For a compact test in the degree preceding top degree,
its resulting ordinary current has compact support by (2).
Pair its derivative with an original smooth scalar cutoff
equal to one near that support. The definition of the
derivative and the vanishing of the derivative of that
cutoff there give zero. This proves separate closedness
of the full delta-test functional.

For (5), inspect the finite anchored construction on a
fixed small Courant cube. Each anchored coefficient is a
finite signed sum of evaluations of one smooth coefficient
on coordinate faces. Taking its coordinates in a fixed
star-balancing kernel is another fixed finite linear map.
The resulting piecewise smooth coefficient arrays have
$C^0$ norms bounded by a fixed constant times those of
the original arrays on finitely many enlarged compact
coordinate-face stages. All Courant covectors and cycle
factors are fixed. The output is therefore a finite sum
of these bounded coefficient multipliers applied to fixed
coefficient-Radon cycle actions over $S$. Their local
finite variation bounds give the asserted estimate on
every original test, after its common refinement. Finite
partitions handle arbitrary compact support presentations.

Finally each generator-positive term is exactly a
nonnegative actual scalar times a positive BT current
wedged with an original middle-positive smooth form.
Its top pairing is nonnegative by the frozen cone
convention and coefficient-Radon multiplication. Summing
proves the stated positivity without changing that cone.
\end{proof}

\begin{corollary}[Full delta-tests for bounded BT products and weights]
\label{cor:r6-full-delta-bounded-bt}
Fix $S\in\mathcal W^p_{\mathrm{loc}}$. Every bounded
classical mixed BT current over $S$ has the full extension
of Theorem~\ref{thm:r6-full-delta-test-extension}. The
weighted joint decreasing convergence in Theorem
\ref{thm:r6-weighted-monotone-delta-tests} holds for every
compact actual Mihatsch delta-test of complementary degree.

If two such unweighted currents are equal on all original
smooth tests, their extensions agree on all actual compact
delta-tests. They also agree after the same actual locally
bounded Borel multiplier wherever the displayed pairings
are defined. In particular this applies to all the bounded
classical products over $[1]$ from B and to the ordinary
potential-quotient identities over backgrounds in this class.
\end{corollary}

\noindent\textbf{Status: \texttt{DIRECT\_COROLLARY}} of
the full extension, its ordinary-current construction,
the class stability proposition, and the bridge's weighted
convergence theorem.

\begin{proof}
Proposition~\ref{prop:r6-delta-weak-class-bt-stability}
puts each unweighted BT current in the same kind of
weak approximation class. Module generation makes every
actual test admissible in the already proved weighted
convergence argument. The extension is determined by
the given ordinary current through fixed-potential
recursion and piecewise coefficient multiplication, not
by a choice of approximating net. Ordinary equality
therefore gives equality of the resulting coefficient
currents for each delta-form. Multiplying those identical
coefficient currents by the same actual Borel function
preserves equality. The final applications follow because
$[1]$ belongs to the approximation class and the R5
quotient gives the relevant ordinary equalities.
\end{proof}

\begin{corollary}[Weak continuity on fixed delta-tests]
\label{cor:r6-weak-continuity-full-delta-tests}
Suppose $S_j,S\in\mathcal W^p_{\mathrm{loc}}$ are positive
separately closed and $S_j\to S$ on all original compact
smooth tests. Then for every fixed actual delta-form
$\alpha$, the ordinary currents $\alpha\wedge S_j$
converge to $\alpha\wedge S$. For compact complementary
$\alpha$ their full delta-test pairings converge as well.
This is a statement for a fixed delta-test. It does not
assert simultaneous stability with varying discontinuous
potential slots and arbitrary varying backgrounds.
\end{corollary}

\begin{proof}
Choose one finite local module presentation of the fixed
$\alpha$ over each compact test support, expand its PL
slots into convex positive and negative choices, and
apply Lemma
\ref{lem:r6-fixed-continuous-slots-varying-background}
to each term. This proof needs no common approximating
delta-form net for the $S_j$. For a compact complementary
test, all the resulting measures are supported in its
one fixed actual compact support. An original smooth
cutoff equal to one near that compact converts ordinary
current convergence into convergence of their integrals.
\end{proof}

\paragraph{The precise remaining all-background question.}
Every fixed allowed $S$ already has the full ordinary
bounded BT calculus and the R5 potential quotient.
The unsolved assertion here is that every actual zero
delta-presentation has zero relative action for that
arbitrary $S$. Local weak approximation by positive
closed actual delta-forms is one sufficient mechanism,
now proved to cover all actual delta-tests. Its
availability for every $S$ has not been established.
The saturated matroidal theorem provides a different,
explicitly geometric mechanism in its fixed-chart scope.

% END source: sections/full_delta_tests_over_weak_closure.tex


% BEGIN canonical R6 source: sections/delta_smooth_flag_backgrounds.tex
\section{Integration currents obtained from smooth regular flags}
\label{sec:r6-delta-smooth-flag}

The weak-approximation hypothesis has geometric examples beyond
bounded products over $[1]$.  The approximation below is an iterated
weak-closure argument on one fixed ordinary domain.  It does not
assert a joint limit of singular logarithmic potentials.

\begin{lemma}[Exact eventual equality on one original top test]
\label{lem:r6-log-truncation-original-test}
Let $Z$ be a good Hausdorff pure $d$-dimensional analytic space over
the original $K$, and let $g\in\mathcal O(Z)$ be a non-zero-divisor.
Put
\[
             v_N=\max\{\log|g|,-N\},\qquad N\in\mathbf N,
\]
with $v_N=-N$ on $g=0$.  For every original compactly supported
smooth $(d,d)$-form $\theta$ there is $N_\theta$ such that
\[
       \int_Z v_N\theta
          =\int_{Z\setminus\{g=0\}}\log|g|\,\theta
                  \qquad(N\geq N_\theta).
\]
In particular $[v_N]\to[\log|g|]$ as original currents.  No
integrability statement against an arbitrary Radon measure is
part of this assertion.
\end{lemma}

\noindent\textbf{Status: \texttt{DIRECT\_COROLLARY}} of the
carrier argument in \cite[Lemma 4.6.1 and its proof, printed
pp.~72--73]{CLD12}.

\begin{proof}
Set $Z^*=Z\setminus\{g=0\}$.  It is a dense Zariski open because
$g$ is a non-zero-divisor.  Apply a finite original smooth partition
to the compact support of $\theta$.  The cited proof supplies, for
each resulting moment presentation $\lambda f^*\omega$, a compact
set $C$ contained in the interior of its affinoid stage and disjoint
from $Z\setminus Z^*$.  It contains the support of
$\lambda f^*\omega$ for every top-degree coefficient form in that
fixed presentation.  Explicitly, $C$ is the intersection of the
support of $\lambda$ with the union of the inverse images of the
standard $d$-dimensional torus skeleta under the standard coordinate
projections of $f$.  The original top-degree carrier has maximal
residual dimension, so it cannot meet the proper analytic zero
locus of $g$; this is exactly the disjointness proved there.

The function $\log|g|$ is continuous and finite on each such compact
$C$, hence has a finite lower bound there.  Choose $N_\theta$ larger
than the negatives of the finitely many lower bounds.  Then
$v_N=\log|g|$ on all these carriers for $N\geq N_\theta$.
The displayed integrals are consequently equal, term by term.
The logarithmic integral defines a current by the same cited lemma;
thus equality eventually on each fixed test proves the claimed
weak convergence in the topology of all original tests.  This
argument does not require a uniform $N_\theta$ over varying tests.
\end{proof}

\begin{proposition}[A smooth regular flag belongs to the approximation class]
\label{prop:r6-smooth-flag-delta-closure}
Let $V$ be one fixed ordinary open in an allowed smooth analytic
space of pure dimension $n$.  Suppose that there are analytic
functions $g_1,\ldots,g_p$ on $V$ and a flag of closed analytic
subspaces
\[
       V=Y_0\supset Y_1\supset\cdots\supset Y_p,
       \qquad
       Y_i=\{g_1=\cdots=g_i=0\},
\]
such that every $Y_i$ is good, Hausdorff, boundaryless, and
quasi-smooth of pure dimension $n-i$, and $g_i|_{Y_{i-1}}$ cuts
out the reduced effective Cartier divisor $Y_i$ in $Y_{i-1}$.
Let $[Y_i]$ denote its integration current on $V$, with the
original CLD calibration.  Then
\[
                       [Y_i]\in\mathcal W^i(V)
                       \qquad(0\leq i\leq p).
\]
More precisely, if $v_{i,N}=\max\{\log|g_i|,-N\}$ on $V$, then
\[
          D\bigl(v_{i,N}[Y_{i-1}]\bigr)
                   \longrightarrow [Y_i]
\]
on every original compact complementary smooth test, and every
current on the left belongs to $\mathcal W^i(V)$.
All approximations use this same $V$ and the original absolute value.
\end{proposition}

\noindent\textbf{Status: \texttt{DIRECT\_COROLLARY}} of
the preceding lemma, original Poincar\'e--Lelong, and the
fixed-continuous-slot weak-closure mechanism.

\begin{proof}
The initial current $[Y_0]=[1]$ is an actual delta-form background.
Assume that $[Y_{i-1}]\in\mathcal W^{i-1}(V)$.  Each $v_{i,N}$
is an actual continuous locally bounded classical psh function on
$V$.  Indeed it is constant on the ordinary neighborhood
$\{|g_i|<e^{-N}\}$ of the zero locus, and elsewhere is the maximum
of a constant and the logarithm of the norm of an analytic unit.
This also proves $v_{i,N}\in PS^{0,0}(V)$ in the original moment
category.  The same description after every complete field extension
proves the required classicality.

Choose a net $[\sigma_\lambda]\to[Y_{i-1}]$ on the fixed $V$
as in Definition~\ref{def:r6-delta-weak-closure}.  For fixed $N$,
the current
\[
 D\bigl(v_{i,N}[\sigma_\lambda]\bigr)
           =[D_\delta v_{i,N}\wedge\sigma_\lambda]
\]
is an actual delta-form background.  Its ordinary current is
positive, symmetric, and separately closed by the ordinary
relative BT theorem.  The fixed-continuous-slot lemma gives its
weak limit $D(v_{i,N}[Y_{i-1}])$.  Hence this limit belongs to
$\mathcal W^i(V)$, with no shrinking of $V$ depending on $N$.

Let $j:Y_{i-1}\hookrightarrow V$.  Pullback of every original
compact smooth test by this closed immersion is again an original
compact smooth test.  Integration and distributional differentiation
give
\[
 D\bigl(v_{i,N}[Y_{i-1}]\bigr)
               =j_*D[v_{i,N}|_{Y_{i-1}}].
\]
Lemma~\ref{lem:r6-log-truncation-original-test}, applied intrinsically
to $Y_{i-1}$, passes its right side to
$j_*D[\log|g_i|_{Y_{i-1}}]$.  Original
Poincar\'e--Lelong
\cite[Theorem 4.6.5, printed p.~74]{CLD12} identifies this with
$j_*\delta_{\operatorname{div}(g_i|_{Y_{i-1}})}=[Y_i]$.
There is no sign or degree factor: the cited convention is
$d'd''[\log|g|]=\delta_{\operatorname{div}(g)}$, and the divisor
is reduced with multiplicity one by hypothesis.

The cone $\mathcal W^i(V)$ is weakly closed by definition, so it
contains $[Y_i]$.  This completes the induction.  Only a single
logarithmic slot is sent to its singular limit at each induction
step; no simultaneous singular-product continuity has been used.
\end{proof}

\begin{corollary}[Full delta-testing for the flag and its bounded products]
\label{cor:r6-smooth-flag-full-delta}
Every integration current which locally admits the smooth regular
flags of Proposition~\ref{prop:r6-smooth-flag-delta-closure} lies in
$\mathcal W^p_{\mathrm{loc}}$.  This includes a smooth algebraic
closed subvariety of a smooth algebraic variety after restricting
to neighborhoods on which its regular immersion has an adapted
smooth flag.  All its bounded classical mixed BT products belong
to the same approximation class in the appropriate degree.
They therefore have the canonical full actual delta-test extension
of Theorem~\ref{thm:r6-full-delta-test-extension}, including its
weighted joint decreasing convergence.
\end{corollary}

\begin{proof}
The proposition proves the local membership.  For the stated
algebraic example, choose local generators of the conormal bundle
whose independent differentials give smooth successive zero loci,
and shrink the algebraic neighborhood.  Analytification supplies
the required flag on the corresponding ordinary open.  The bounded
product assertion follows from Proposition
\ref{prop:r6-delta-weak-class-bt-stability}; the full-test assertions
then follow from the cited extension theorem and its corollary.
\end{proof}

\paragraph{Scope.}
The proof covers the explicit local smooth-flag hypothesis.  It
neither asserts that arbitrary singular analytic cycles admit such
flags nor that every allowed positive closed current belongs to
$\mathcal W_{\mathrm{loc}}$.  In particular it adds examples to
the sufficient approximation class without replacing the project's
arbitrary-background ordinary theory by a geometric subclass.

% END source: sections/delta_smooth_flag_backgrounds.tex


% BEGIN canonical R6 source: sections/delta_integrable_potential_backgrounds.tex
\section{Integrable classical potentials and effective Cartier backgrounds}
\label{sec:r6-delta-integrable-potentials}

All currents in this section remain currents on the allowed original
$K$-space. An extended classical psh function takes values in
$\mathbf R\cup\{-\infty\}$. Absolute integrability below refers to
the total variations of the measures of all original compact smooth
top-degree forms. It is not an assertion of integrability against
every positive background measure.

\begin{lemma}[A fixed approximation domain for individually bounded functions]
\label{lem:r6-fixed-box-individual-rows}
In an ordinary algebraic neighborhood of a point, fix nested polynomial
norm boxes
\[
 C_1\subset\operatorname{Int}C_{\rm mid},\qquad
 C_{\rm mid}\subset\operatorname{Int}C_2\Subset O,
 \qquad V=\operatorname{Int}C_1.
\]
They can be chosen so that every classical psh function $w$ on $O$
which is bounded on a neighborhood of $C_2$ has, on this same $V$,
a decreasing row of continuous semipositive PL functions. The
bound for $w$ and the row can depend on $w$; the domain $V$ does not.
Consequently $D[w]|_V\in\mathcal W^1(V)$.
\end{lemma}

\noindent\textbf{Status: \texttt{DIRECT\_COROLLARY}} of the
explicit fixed-box collar in B, Lemma D.2, and its bounded projective
regularization theorem. The fixed-domain point is spelled out because
it is needed for an unbounded limiting potential.

\begin{proof}
Use the ordinary algebraic localization and smooth projective
compactification of B, Proposition 4.5. If its constant field is a
finite extension of $K$, retain the original extending absolute value;
the analytic functions and their moment forms are still original
$K$-analytic data on that neighborhood. The single positive
integration factor of B, Lemma D.5, appears on both sides of each
current limit and does not change weak closure or positivity.

Choose the boxes once. The constructive part of B, Lemma D.2
(printed pp.~27--28), supplies from these boxes a continuous PL collar
$b$, a pure Fubini--Study metric $q$, and $\delta>0$, with
\[
 b=0\text{ on a neighborhood of }C_1,
 \qquad b\ge\delta\text{ on }C_2\setminus\operatorname{Int}C_{\rm mid},
\]
and with $b$ the weight of $q$ in a frame near $C_2$. These objects
are chosen before the upper and lower bounds of the functions enter
that construction.

Normalize the present $w$ by a positive constant and an additive
constant to a classical function $v$ with $0\le v\le1$ on a
neighborhood of $C_2$. Fix $c>0$ in the permitted value group and
an integer $m$ with $m\delta>c+2$. The metric with local weight
\[
                         \max\{v+c,mb\}
\]
on $\operatorname{Int}C_2$, and equal to $mq$ outside
$C_{\rm mid}$, is a bounded classical metric on the same ample
line used by that collar. On the overlap the second branch dominates,
so this is the identical gluing from B. Near $C_1$ its framed weight
is exactly $v+c$. The choices of boxes, frame, $c$ and $m$ work for
every normalized $v$. No monotonicity of a family of such $v$ is
needed to construct each metric; in B it was needed only for the
additional assertion that the resulting metric family decreases.

Apply B, Theorem 4.4, separately to this bounded metric. Its decreasing
pure-FS row restricts on $V$ to a continuous semipositive PL row for
$v+c$. Undo the fixed additive shift and the positive normalization
to obtain the desired row for $w$. Real shifts and positive real
scalings are permitted for actual PS and delta-form data. Each row
curvature is an actual delta-form with positive separately closed
ordinary current. The bounded original monotone theorem, or fixed
top-test bounded convergence, gives its weak limit $D[w]$ on this
same $V$. This proves membership in $\mathcal W^1(V)$.
\end{proof}

\begin{theorem}[An integrable extended classical potential gives a weak delta background]
\label{thm:r6-integrable-classical-delta-background}
Let $u:U\to\mathbf R\cup\{-\infty\}$ be actual extended classical
psh on an allowed original-$K$ space of dimension $n\ge1$. Assume
that $u$ is locally absolutely integrable against the total variation
of every original compact smooth $(n,n)$-form. Then $[u]$ is a
degree-zero current, $D[u]$ is symmetric positive and separately
closed, and
\[
                              D[u]\in\mathcal W^1_{\rm loc}.
\]
For $u_N=\max\{u,-N\}$ one has $D[u_N]\to D[u]$ on every
original compact smooth complementary test. Near each point all
these truncation curvatures belong to $\mathcal W^1(V)$ on one
common ordinary $V$.
\end{theorem}

\noindent\textbf{Status: \texttt{DIRECT\_COROLLARY}} of the
fixed-box lemma, original integrability and weak closure. No local
potential is supplied for an arbitrary given positive current.

\begin{proof}
First the integrability hypothesis defines a current, rather than
only individual finite integrals. On a fixed finite compact moment
stage the coefficient controls for $[1]$ bound the variation of every
top test by its coefficient $C^0$ seminorm times finitely many fixed
positive compact top-form measures. Their integrals against $|u|$
are finite. This gives a continuous order-zero functional on that
test stage. Finite localization proves continuity in the original
test topology. Thus $[u]$ is well defined, with arbitrary values on
the sets on which the controlling measures vanish.

Each actual truncation $u_N$ is locally bounded classical. For
$N\ge1$, $|u_N|\le|u|$ wherever $u$ is finite, and the
integrability hypothesis makes its infinite-value set null for
every fixed top-test variation. Fixed-measure dominated convergence
therefore gives $[u_N]\to[u]$, hence $D[u_N]\to D[u]$ on
every original test. Since each $D[u_N]$ is symmetric positive
and separately closed, its limit has the same properties.

Choose the boxes in Lemma~\ref{lem:r6-fixed-box-individual-rows}
once near a prescribed point. Upper semicontinuity gives a common
finite upper bound $M$ for $u$ on a neighborhood of $C_2$. Each
$u_N$ has the finite lower bound $-N$ there. The lemma consequently
puts every $D[u_N]|_V$ in the same weakly closed cone
$\mathcal W^1(V)$. Its weak limit $D[u]|_V$ also lies there.
The geometric domain does not shrink with $N$; only the rows and
their harmless normalizations vary. This proves the local claim.
\end{proof}

\begin{corollary}[Effective Cartier divisors, with singularities and multiplicities]
\label{cor:r6-cartier-delta-background}
Let $V$ be an allowed ordinary domain and let $f\in\mathcal O(V)$
be a non-zero-divisor. Let $Z=\operatorname{div}(f)$ be its
effective Cartier divisor, with its original generic multiplicities.
No smoothness or reducedness of $Z$ is assumed. Then its ordinary
integration current satisfies
\[
 [Z]=D[\log|f|]\in\mathcal W^1(V).
\]
More precisely, on this same whole $V$,
\[
       [D_\delta\max\{\log|f|,-N\}]\longrightarrow[Z]
\]
on every original compact complementary smooth test. Hence every
locally effective Cartier divisor lies in $\mathcal W^1_{\rm loc}$,
and has the canonical full actual delta-test action already proved
for that approximation class. Its bounded classical mixed BT tails
have that action as well.
\end{corollary}

\noindent\textbf{Status: \texttt{DIRECT\_COROLLARY}} of
\cite[Lemma 4.6.1, 4.6.3 and Theorem 4.6.5, printed
pp.~72--74]{CLD12}. The cited statements require pure dimension,
topological separatedness and a non-zero-divisor; they do not
require the divisor to be smooth.

\begin{proof}
Put $v_N=\max\{\log|f|,-N\}$. It is an actual continuous
classical PS function on $V$: near a zero of $f$ the constant
branch is active, and off the zero locus it is a maximum of affine
functions of the original unit moment $\log|f|$. These assertions
hold after every complete extension. Thus $D_\delta v_N$ is an
actual delta-form, and its ordinary current is symmetric positive
and separately closed.

For every fixed original compact top test, the compact carrier in
the proof of CLD12, Lemma 4.6.1, avoids the divisor. The logarithm
is finite and bounded below on that carrier. Accordingly $v_N$
and $\log|f|$ have exactly the same pairing with that test for
all sufficiently large $N$. This is also
Lemma~\ref{lem:r6-log-truncation-original-test}; it assumes only
a pure-dimensional space and a non-zero-divisor. Apply it to the
top test $D\eta$ for any original complementary compact smooth
$\eta$. It yields
\[
     \langle D[v_N],\eta\rangle
                 \longrightarrow\langle D[\log|f|],\eta\rangle.
\]
Original Poincare--Lelong identifies the latter current with the
integration current of $Z$, including the generic Cartier
multiplicities. In the original convention its coefficient is one:
$D[\log|f|]=\delta_{\operatorname{div}(f)}$. In particular repeated
factors contribute their actual multiplicities; no smoothness or
reduced-divisor hypothesis has been substituted.

All approximants are actual delta-form backgrounds on $V$ itself,
so the fixed-domain weak closure gives $[Z]\in\mathcal W^1(V)$.
For an effective Cartier divisor with only local equations, apply
this on the equation neighborhoods; multiplication of an equation
by a unit changes its logarithm by a smooth pluriharmonic function.
The ordinary currents and the already proved full delta actions
therefore glue. Bounded BT stability of $\mathcal W_{\rm loc}$
gives the final assertion.
\end{proof}

\paragraph{Limits of this enlargement.}
The extended-potential theorem assumes original-test absolute
integrability. It does not assert that every extended classical
function has that property, nor that every positive closed current
has such a potential. The Cartier corollary is independently
unconditional because its original logarithmic carrier calculation
is exact. No assertion about general singular higher-codimension
intersections or arbitrary analytic cycles is made here.

% END source: sections/delta_integrable_potential_backgrounds.tex


% BEGIN canonical R6 source: sections/delta_barycenter_backgrounds.tex
\section{Barycenters of backgrounds in the approximation cone}
\label{sec:r6-delta-barycenters}

The weak approximation class is not restricted to a single geometric
background or to countable positive sums. The following elementary
closure mechanism is useful for diffuse currents supported on ambient
pluripolar sets. Its same-domain hypothesis is explicit.

\begin{proposition}[Integrable barycenters in one fixed weak cone]
\label{prop:r6-delta-barycenter}
Let $V$ be one fixed ordinary open and let $(A,\lambda)$ be a positive
measure space. For almost every $a\in A$ let
$S_a\in\mathcal W^p(V)$. Assume the following two conditions.
\begin{enumerate}
\item For every original compact smooth complementary test $\eta$,
      the scalar $a\mapsto\langle S_a,\eta\rangle$ is measurable.
\item On each compact test presentation the finite positive mass
      controls for the coefficient-Radon estimate are integrable in
      $a$.
\end{enumerate}
Then
\[
           \langle S,\eta\rangle
                    =\int_A\langle S_a,\eta\rangle\,d\lambda(a)
\]
defines a symmetric original-positive separately closed current, and
$S\in\mathcal W^p(V)$. In particular it admits the full actual
delta-test extension of Theorem~\ref{thm:r6-full-delta-test-extension}.
For each fixed actual delta-form $\alpha$, whenever the corresponding
finite inserted mass controls are locally integrable, one also has
\[
             \alpha\wedge S
                      =\int_A\alpha\wedge S_a\,d\lambda(a)
\]
on all original compact smooth tests.
\end{proposition}

\noindent\textbf{Status: \texttt{DIRECT\_COROLLARY}} of convex weak
closure and the finite-presentation construction. The same-domain
membership $S_a\in\mathcal W^p(V)$ is not replaced by an
$a$-dependent local membership assertion.

\begin{proof}
The positive coefficient controls imply absolute integrability of
every displayed scalar pairing and continuity on each compact
original test presentation. They therefore define a current. Symmetry,
positivity and the two closedness identities pass under the integral.

We check membership without assuming any countable collection of
tests generates the topology. Fix finitely many real original tests
$\eta_1,\ldots,\eta_N$, and consider
\[
 v(a)=\bigl(\langle S_a,\eta_1\rangle,\ldots,
                 \langle S_a,\eta_N\rangle\bigr)\in\mathbf R^N.
\]
Its components are integrable. Its integral belongs to the closed
convex cone generated by its essential values: otherwise finite
dimensional separation supplies a real linear functional nonnegative
on that cone but negative on the integral, contradicting integration
of that same nonnegative measurable scalar function. Hence the
pairings of $S$ on this finite list are approximated arbitrarily
closely by positive finite combinations of the pairings of $S_a$.
Every such combination belongs to the convex cone
$\mathcal W^p(V)$. This proves that every basic weak neighborhood of
$S$ meets that weakly closed cone, and consequently $S$ belongs to it.
The argument permits infinite total parameter mass when the stated
local controls are integrable.

For the last assertion take a finite module presentation of $\alpha$
and express its constant cycles as finite signed products of
continuous PL curvature slots. Start with one continuous slot $v$.
Uniform approximation of a fixed continuous scalar by original smooth
scalars, together with the integrable control bounds, gives
$vS=\int_AvS_a\,d\lambda(a)$.
Applying $D$ commutes with the integral on original tests. Iterate
this argument for the finitely many fixed continuous slots; relative
CLN gives the required integrable controls from the stipulated
inserted controls, or from the original controls when the fixed
potentials have common compact bounds. Finally use the uniform
piecewise-test approximation estimate. Its errors are bounded by
an integrable finite sum of the same positive control pairings,
and the smooth approximants commute with the integral. This passes
piecewise smooth coefficient multiplication through the integral
without assuming a Borel disintegration on every individual stratum.
Summing the finite expression
proves the identity. Actual relation independence holds on both
sides by the full extension theorem.
\end{proof}

\begin{corollary}[Diffuse polar-supported fiber backgrounds]
\label{cor:r6-laminar-full-delta-tests}
In the product-polydisc setting of
Proposition~\ref{prop:r6-fiber-current-disintegration}, the diffuse
fiber current
\[
                    S=\int_P[Y\times\{a\}]\,d\lambda(a)
\]
belongs to $\mathcal W^p(\Omega)$. All its bounded classical BT
products admit full actual delta-testing, separately closed extensions,
and weighted joint decreasing convergence on these tests.
This includes the nonatomic exact complete-pole parameter examples.
\end{corollary}

\begin{proof}
For every fixed $a\in P$, the equations $z_i-a_i=0$ give a smooth
regular flag on the same domain $\Omega$. Proposition
\ref{prop:r6-smooth-flag-delta-closure} therefore puts each fiber
integration current in $\mathcal W^p(\Omega)$.
The fiber-integration proof gives scalar measurability and uniform
compact control bounds on the compact parameter space $P$.
Apply Proposition~\ref{prop:r6-delta-barycenter}, followed by the
stability of $\mathcal W_{\mathrm{loc}}$ under bounded classical
slots and Corollary~\ref{cor:r6-full-delta-bounded-bt}. No assertion
that ambient pluripolar sets are null for these currents is used.
\end{proof}

\paragraph{What remains outside this closure argument.}
An arbitrary positive separately closed current has not been shown
to be an integrable barycenter of the displayed geometric examples,
nor to be a weak limit of positive separately closed delta-forms.
The proposition enlarges an established class; it does not solve
the all-background analytic relation problem by definition.

% END source: sections/delta_barycenter_backgrounds.tex


% BEGIN canonical R6 source: sections/delta_extension_uniqueness.tex
\section{Values forced by a positive closed delta-test extension}
\label{sec:r6-delta-extension-uniqueness}

An arbitrary ordinary current is not, by definition, a functional on
singular delta-tests.  The following result is a uniqueness statement
for extensions satisfying explicit positivity and Stokes requirements.
It does not assume that such an extension exists.

Fix an allowed positive symmetric separately closed current $S$ of
bidegree $(p,p)$ on $U$, and put $q=n-p$.  An admissible potential below
means a continuous function $u\in\mathcal C_b(U)\cap PS^{0,0}(U)$;
local versions on neighborhoods of the relevant compact supports are
equally allowed.  Write $D_\delta u=d'd''u$ in the actual delta algebra.

\begin{definition}[The positivity required of an extension]
\label{def:r6-generator-positive-extension}
Consider a real symmetric linear functional $\Lambda$ on the actual
compact delta-tests of complementary total bidegree to $S$.  Assume
the following three properties.
\begin{enumerate}
\item Its restriction to all original compact smooth complementary
      tests is $S$.
\item It is separately $d'$- and $d''$-closed, in the dual sense on
      the actual delta-test complex.
\item Whenever $0\leq k\leq q$, $u_1,\ldots,u_k$ are admissible,
      $a\geq0$ is a continuous piecewise smooth function near the
      support of a compact original smooth form $\theta$, and
      $\theta$ is positive in the original middle/PSD cone of
      bidegree $(q-k,q-k)$, one has
      \[
       \Lambda\!\left(a\,D_\delta u_1\wedge\cdots\wedge
                                D_\delta u_k\wedge\theta\right)\geq0.
      \]
\end{enumerate}
This is a stated generator-positivity condition.  It is not a claim
that every possible convention for positivity of a delta-current
automatically has these properties.
\end{definition}

\begin{theorem}[Forced BT values]
\label{thm:r6-forced-delta-bt-values}
\textnormal{\texttt{DIRECT\_COROLLARY}} of the frozen ordinary calculus,
under the explicitly assumed existence of $\Lambda$ above.
For every $0\leq k\leq q$, every admissible $u_1,\ldots,u_k$, every
continuous piecewise smooth scalar $u_0$, and every original compact
smooth complementary test $\eta$, one necessarily has
\[
 \Lambda\!\left(u_0\,D_\delta u_1\wedge\cdots\wedge
                              D_\delta u_k\wedge\eta\right)
       =\bigl\langle u_0 T_{S,k}(u_1,\ldots,u_k),\eta\bigr\rangle.
\]
All quantities use the actual original field, points, positivity
cone, and full original test class.  In particular two extensions
satisfying Definition~\ref{def:r6-generator-positive-extension}
agree on the linear span of these tests.  Finite signed polarization
gives the identical conclusion for differentiated slots which are
differences of admissible potentials.
\end{theorem}

\begin{proof}
First consider $k=0$.  If $\theta$ is a positive original compact
test and $h$ is a real continuous piecewise smooth scalar, then
the generator-positivity assumption applied to
$\|h\|_\infty\pm h$ gives
\[
                 |\Lambda(h\theta)|
                     \leq \|h\|_\infty\,\langle S,\theta\rangle,
\]
where the supremum is over a fixed compact neighborhood of the
support.  For an arbitrary symmetric original test, the finite
positive-form domination in the original PSD cone gives the same
bound with a finite sum of fixed positive controls and a fixed
constant.  Symmetry reduces a general test to its symmetric part.
This is the usual finite-frame coefficient domination; no strong
cone spanning assertion is being used.

On this compact stage, a continuous piecewise smooth scalar $h$
admits original smooth approximants $h_j$ converging uniformly.
In a finite moment presentation they are ordinary mollifications
of its continuous piecewise smooth representative on a slightly
larger real open; a finite original smooth partition handles
several presentations.  Positivity gives
\[
                      \Lambda(h\eta)=\lim_j\Lambda(h_j\eta).
\]
The original restriction of $\Lambda$ and the coefficient-Radon
property of $S$ identify the right side with
\[
                      \lim_j\langle h_jS,\eta\rangle
                              =\langle hS,\eta\rangle.
\]
Thus multiplication by any such scalar agrees with actual
continuous multiplication on the ordinary restriction.

For an admissible $u$, define insertion into the functional by
\[
                    \Lambda_u(\beta)
                         =\Lambda(D_\delta u\wedge\beta).
\]
It is symmetric and separately closed: $D_\delta u$ is separately
closed and the delta algebra satisfies the two graded Leibniz
rules.  The generator-positivity assumption is inherited by
$\Lambda_u$, with one fewer available differentiated slot.
The dual Leibniz identity, using separate closedness of
$\Lambda$, is
\[
                         \Lambda_u=D(u\Lambda).
\]
It is an algebraic identity of functionals on the delta-test
complex, not a convergence statement.  Restrict it to original
tests.  The preceding scalar argument identifies the ordinary
restriction of $u\Lambda$ with $uS$.  Differentiation commutes
with restriction because original smooth tests form a
subcomplex.  Hence
\[
                       \Lambda_u\big|_{\mathrm{original}}
                                  =D(uS)=T_{S,1}(u).
\]
The right side is a positive symmetric separately closed
coefficient-Radon current by the frozen ordinary BT theorem.

Repeat the argument for $\Lambda_u$.  Induction identifies the
ordinary restriction of insertion of $k$ factors with the
corresponding relative BT product.  Apply the initial scalar
argument once more to that inserted functional and $u_0$ to
obtain the displayed weighted equality.  Finally expand any
signed slots by finite real multilinearity.  Every intermediate
unsigned insertion has the proved positivity and closedness,
so no quotient theorem for an arbitrary signed background is
silently assumed.
\end{proof}

\paragraph{What this does and does not decide.}
The theorem determines the extension on a specified test span;
it is not a density argument and it is not an existence theorem.
It also shows that, if a formal relation between tests in this
span fails after the ordinary relative BT prescription, then no
extension with the three stated properties can exist.  Conversely,
annihilating those relations constructs an action only on that
span.  Generation of the whole actual delta algebra, and the
positivity or continuity needed for an extension beyond it,
remain separate questions.  In particular the unrestricted
functional counterexamples from the preceding section do not
contradict this uniqueness statement.

% END source: sections/delta_extension_uniqueness.tex


% BEGIN canonical R6 source: sections/delta_test_extension_obstructions.tex
\section{Ordinary restriction does not determine an arbitrary delta-test functional}
\label{sec:r6-delta-test-extension-obstructions}

This section concerns \emph{functionals on tests}, not equality between
two actual delta-forms.  Proposition~\ref{prop:r6-delta-faithfulness}
continues to apply to actual delta-forms.  The examples below leave its
domain and do not contradict it.

\begin{lemma}[Top delta-forms have canonical coefficient measures]
\label{lem:r6-delta-top-radon}
\textnormal{\texttt{DIRECT\_COROLLARY}.}
On the original good Hausdorff boundaryless pure $n$-dimensional
$K$-space $U$, each $\eta\in B_c^{n,n}(U)$ has a canonical finite
signed Radon measure $\mu_\eta$, supported in $\operatorname{supp}\eta$,
characterized by
\[
                  \int_U a\,d\mu_\eta=\int_U a\eta
                  \qquad(a\in A_c^{0,0}(U)).
\]
For each locally bounded Borel function $f$ on $U$,
\[
       \Lambda_f:B_c^{n,n}(U)\longrightarrow\mathbf R,
       \qquad \Lambda_f(\eta)=\int_U f\,d\mu_\eta
\]
is a well-defined linear functional.  Its restriction to original
smooth tests is the ordinary current $[f]$.

It is continuous on every fixed finite polyhedral presentation stage
with a common compact support and coefficient $C^0$ convergence.
In particular it is continuous for the stronger coefficient smooth
convergence on that stage.  This statement specifies the continuity
being used; it does not attribute an unmentioned delta-current
topology to Mihatsch.
\end{lemma}

\begin{proof}
Choose a finite family of local presentations of $\eta$ over its
compact support.  By \cite[Corollary 4.9 and Definition 4.10]{Mih24},
its support is carried by a finite union of calibrated skeletons.
On a common refinement, its restriction is a finite sum of smooth
top-degree coefficient forms on weighted polyhedra.  Their absolute
integrals on the compact stage are finite.  They therefore define a
finite signed Radon measure on that stage and then on $U$.

Different skeletal presentations give the same integrals against
every original smooth scalar test, since both compute
$\int_U a\eta$.  Original smooth scalar functions are uniformly
dense among continuous functions on a fixed compact neighborhood,
with a common compact cutoff.  Uniqueness of Radon measures then
identifies these constructions.  This also proves linearity in
$\eta$.  The support assertion follows by ordinary locality.

On a fixed stage the sum of the absolute polyhedral integrals is
bounded by a fixed constant times the maximum of the coefficient
uniform norms.  Hence
\[
 |\Lambda_f(\eta)|\leq\sup_C|f|\,|\mu_\eta|(C)
\]
has the asserted stage bound.  The Borel function $f$ is multiplied
only after the fixed signed Radon measure has been identified;
no tropical measurability condition on $f$ is needed.
For smooth $\eta$, the measure just constructed is its original
CLD integration measure, proving the restriction statement.
\end{proof}

The algebraic GK counterpart of the measure construction is
\cite[Corollary 6.15]{GK17}.  Proposition 6.16 there states the
functional result for continuous $f$.  Its displayed bound (6.16.1),
once the signed Radon measure is fixed, gives the same continuity
estimate for locally bounded Borel $f$.  This observation does not
extend the base-field hypotheses of that paper; the proof above
uses Mihatsch's integration directly over the original $K$.

\begin{example}[A positive functional invisible to all ordinary tests]
\label{ex:r6-invisible-positive-delta-functional}
\textnormal{\texttt{FALSE} for unrestricted uniqueness;
\texttt{DIRECT\_COROLLARY} construction.}
Choose a uniformizer $\pi$ of the original $K$, put
$\ell=-\log|\pi|>0$, and let
\[
       U=\{\,|\pi|<|T|<|\pi|^{-1}\,\}\subset\mathbf G_{m,K}^{\mathrm{an}},
       \quad t=\log|T|,\quad \xi=\zeta_{0,1}.
\]
The point $\xi$ is an actual $K$-Berkovich point; no algebraic closure
or rescaling of the absolute value is used.  Let $f=\mathbf1_{\{\xi\}}$.
Every original smooth compact top-degree form has zero mass at
$\xi$: locally its integration measure has one-dimensional Lebesgue
densities on finitely many calibrated skeletal cells.  Thus
\[
                     [f]_{\mathrm{ordinary}}=0,
              \qquad \Lambda_f|_{A_c^{1,1}(U)}=0.
\]
On the other hand,
\[
           b=\max(0,t)\in PS^{0,0}(U),\qquad
           \eta_0=D_\delta b\in B_c^{1,1}(U)
\]
has $\mu_{\eta_0}=\delta_\xi$, with coefficient one by the original
one-dimensional torus calibration.  Consequently
\[
                              \Lambda_f(\eta_0)=1.
\]
The zero functional and $\Lambda_f$ therefore have the same ordinary
restriction and different values on actual delta-tests.  Moreover
$\Lambda_f(\eta)\geq0$ whenever $\mu_\eta$ is a positive measure.
This is positivity for the explicitly specified top-degree test
cone.  The function $f$ is not an actual piecewise smooth function,
and $\Lambda_f$ is not induced by an actual delta-form.
\end{example}

\begin{example}[Separate closedness alone does not restore uniqueness]
\label{ex:r6-invisible-closed-delta-functional}
\textnormal{\texttt{FALSE} for closed unrestricted uniqueness.}
On the same annulus, define the dual derivative on all compact
delta-functions by
\[
                \Psi(g)=\Lambda_f(D_\delta g),
                       \qquad g\in B_c^{0,0}(U)=PS_c^{0,0}(U).
\]
It is a top-degree functional and is separately closed, since its
$d'$ and $d''$ derivatives have degrees exceeding the dimension.
Equivalently, the dual differential identities give
$d'D_\delta\Lambda_f=d''D_\delta\Lambda_f=0$.
Its ordinary restriction is zero because $Dg$ is a smooth
top-degree form for original smooth $g$.

Choose $0\leq\chi\in C_c^\infty((-\ell,\ell))$ equal to one near
zero and put $g=\chi(t)\max(0,t)$.  This is a compactly supported
nonnegative actual piecewise smooth function.  Its delta-curvature
has coefficient one at $\xi$, so $\Psi(g)=1$.  Therefore $\Psi$ is
nonzero despite its zero ordinary restriction.

This example does \emph{not} preserve positivity.  Indeed choose a
nonnegative original smooth compact cutoff $a$ equal to one on
$\operatorname{supp}g$, and $C\geq\sup g$.  Then $Ca-g\geq0$ but
\[
                       \Psi(Ca-g)=-1.
\]
It consequently does not disprove uniqueness of an extension that
is required to be both separately closed and positive on all
appropriate full Mihatsch delta-tests.
\end{example}

\begin{proposition}[The closed-cycle test subalgebra]
\label{prop:r6-closed-cycle-subalgebra-ghost}
\textnormal{\texttt{DIRECT\_COROLLARY}.}
On the same original-$K$ annulus, let $\mathscr G$ be the differential
subalgebra of $B$ generated by original smooth forms and pullbacks
of constant-weight tropical cycle forms.  Then $\Lambda_f$ restricted
to $\mathscr G_c^{1,1}$ is nonzero, is positive for the Radon
top-degree test cone, is separately closed as a functional on
$\mathscr G_c$, and has zero ordinary restriction.
\end{proposition}

\begin{proof}
A constant-weight tropical cycle generator is separately closed.
The derivatives of a product of such generators and a smooth
coefficient differentiate only that coefficient.  Thus $d'$ and
$d''$ preserve its tropical codimension.  A form of total degree
$(0,1)$ or $(1,0)$ necessarily has codimension zero.  Its derivative
of degree $(1,1)$ still has codimension zero and therefore has a
piecewise smooth Lebesgue density on one-dimensional skeletal
cells, without an atom at $\xi$.  The functional $\Lambda_f$ kills
these derivatives, proving separate closedness in this test
algebra.  It is nonzero because $\eta_0$ is the pullback of the
constant cycle consisting of the origin in $\mathbf R$ and belongs
to $\mathscr G$.  The remaining assertions were proved above.
\end{proof}

The last proposition is not a closed-positive nonuniqueness theorem
for the \emph{full} Mihatsch test algebra.  In that larger algebra,
the piecewise smooth function $g$ from the preceding example is
an admissible test and detects the failure of closedness of
$\Lambda_f$.  This is one concrete reason to state the test sheaf
before comparing delta-currents.  We have defined $\mathscr G$
over the original $K$; we have not identified it with the full
algebraic $B_{\mathrm{GK}}$.

\paragraph{The conclusion for arbitrary backgrounds.}
The ordinary coefficient-Radon data of a current alone do not
specify its values on singular delta-tests.  A proposed extension
must include its test sheaf, continuity requirements, differential
compatibility, positivity requirements, and an independence proof.
The examples establish the failures precisely stated above.
They do not rule out an additional canonical extension mechanism
for separately closed positive currents, and they do not affect
the actual delta-form comparison of
Theorem~\ref{thm:r6-bt-delta-background-comparison}.

% END source: sections/delta_test_extension_obstructions.tex


\part{The ordered potential quotient and compactness}

% BEGIN canonical R6 source: sections/potential_order_functoriality.tex
\section{Functoriality of curvature order for actual bounded potentials}
\label{sec:r6-potential-order-functoriality}

The zero-kernel transfer in R5 does not by itself transfer an order
inequality. The additional points proved here are positivity of the
piecewise traces, positivity on a carried polyhedral set, and positive
finite-coefficient descent. These preserve the original middle CLD
cone throughout. No pullback operation on arbitrary currents is defined.

Over an auxiliary complete field, a prescribed row means the fixed
ordinary-neighborhood smooth positive decreasing row of R5,
Definition 3.1. Only pullbacks of rows supplied over the original
$K$ are used in the final theorem. The field-independent coefficient
lemma and piecewise continuity in R5, Lemmas 2.1 and 3.2, remain
applicable. All complete fields carry their given extending absolute
values.

\begin{lemma}[Positivity on the polyhedral set carrying the mass]
\label{lem:r6-carried-tangent-positivity}
Let $R$ be a symmetric positive separately closed coefficient-Radon
current over any complete field in the preceding scope. On a fixed
compact original moment stage $t:V\to\mathbf R^m$, suppose its
coefficient measures are carried by a Borel set $A$ whose moment
image lies in a finite polyhedral complex $P$. Let $\eta$ be a
smooth or compatible piecewise smooth complementary form. If its
restriction to every face of $P$ is middle-positive in that face's
affine tangent space, then
\[
                         \langle R,t^*\eta\rangle\geq0
\]
whenever the displayed test has the stated compact localization.
The same conclusion holds after a nonnegative Borel restriction
or scalar multiplier supported on this stage. No positivity of an
ambient extension away from $P$ is required.
\end{lemma}

\noindent\textbf{Status: \texttt{DIRECT\_COROLLARY}} of the
coefficient matrix description and level-set tangency, with the
carried-set scope made explicit.

\begin{proof}
Use the disjoint inverse images of the relative interiors of all
faces of $P$, restricted to $A$. The arctangent proof of
Lemma~\ref{lem:r6-level-tangency} uses only smooth differentiation
and fixed Radon dominated convergence, so works over the present
complete field. On the Borel part over a face $F$, it annihilates
every primed and double-primed conormal of its affine span. This
remains true after Borel restriction. Thus only the tangential tensor
of $\eta$ can contribute on this part.

For completeness, choose a constant basis of the appropriate exterior
power of the tangent covectors of $F$ and extend it to the ambient
moment space by one fixed linear projection. The coefficient measures
of this Borel part, relative to one finite sum of diagonal control
measures, give a symmetric positive semidefinite matrix $M_F$ almost
everywhere. Indeed every constant rank-one square in the middle cone
has a positive measure; rational coefficient vectors and continuity
give positivity for all vectors. Nondecomposable squares are included.
The tangent matrix $A_F(x)$ of $\eta$ is positive semidefinite.
The pairing density is $\operatorname{tr}(A_FM_F)\geq0$ in the
usual signed identification of superforms. If the face dimension
is smaller than the form degree, that tangent form is zero. Sum over
the finite faces. Nonnegative multipliers preserve the matrix
positivity, proving the final assertion. The definition of this cone
and its restriction property are exactly
\cite[5.1.1(c), 5.1.5, 5.2.1 and 5.3.1]{CLD12}.
\end{proof}

\begin{lemma}[The finite and split projected tests preserve positivity]
\label{lem:r6-positive-piecewise-traces}
The piecewise trace of an original positive compact test under a
finite etale map, as constructed in R5, Lemma 3.3, is tangentially
middle-positive on each actual base tropical image. The same holds
for the split-disc projected complementary tests of R4, Theorem 72.1.
These statements retain every further original moment refinement.
\end{lemma}

\begin{proof}
Localize the original positive form by a nonnegative smooth partition,
using positive presentations from \cite[Definition 5.3.1]{CLD12}.
For a finite etale piece, the fixed trace formula is
\[
                        P(t)=\sum_{l=1}^d z_l^*\omega,
\]
where the joint vectors $z_l(t)$ are listed by the characteristic
polynomials of finitely many monomials and are affine on a finite
subdivision. Each vector over an actual base point is realized by
a source point, with the full positive algebraic multiplicity.
Refine the base cells by inverse images of the finitely many source
faces in a positive presentation. On each resulting base face, the
affine map $z_l$ maps its tangent space into the tangent space of
an actual source face. Pullback preserves the middle-positive cone
by \cite[Remark 5.1.5]{CLD12}; hence every summand and their sum
are positive on that face. Matching on smaller faces preserves
positivity there. The extension off the consistent moment locus in
the trace construction is immaterial by
Lemma~\ref{lem:r6-carried-tangent-positivity}. This argument also
works after adjoining extra base moments; the same actual sheets
and positive calibrated degrees are retained.

For the split disc, use the exact formula of R4, Theorem 72.1:
\[
 P(t)=\sum_i\int_{-\infty}^{L_i(t)}
            \mathcal I_s(F_i^*\omega)(t,s)\,ds,
 \qquad
 F_i(t,s)=(t,\max(s,b_{i1}),\ldots,\max(s,b_{im})).
\]
After the finite chamber subdivision, these maps are affine and
their values on actual base points and their assigned radial intervals
are actual source Gauss-ray points. Restricting further by source
faces makes the tangent maps land in their actual tangent spaces.
The form $F_i^*\omega$ is consequently middle-positive there.
Its vertical top contraction $\mathcal I_s$ is positive: on a
positive semidefinite tensor it is the compression obtained by
wedging a base exterior vector with the one vertical vector, on
both sides of the corresponding quadratic form. Its sign is the
positive product-calibration sign. Integrating this positive tensor
against $ds$, summing the assigned rays once, and multiplying by
the nonnegative localization all preserve positivity. Collapsed
intervals contribute zero. Tangential matching on base faces, proved
in the frozen formula, gives the same conclusion on their intersections.
There are no additional wall currents in that formula.
\end{proof}

\begin{proposition}[Finite etale and original product order transfer]
\label{prop:r6-finite-etale-product-order}
Let $a,b$ be actual locally bounded classical functions with
prescribed rows over a complete characteristic-zero field $F$.
If $D[a-b]\geq0$, then for a finite etale map $p$,
\[
                         D[(a-b)\circ p]\geq0.
\]
Over the original $K$, the same implication holds for projection
from an ordinary polydisc product of any finite dimension.
\end{proposition}

\begin{proof}
Put $f=a-b$ and $T=D[f]$. In the finite etale case take a positive
original compact test $\alpha$ upstairs and let $P=p_*[\alpha]$.
Localize by a nonnegative base partition so that both prescribed
rows exist on a common neighborhood of the projected support.
The frozen finite trace is a finite compatible piecewise form, and
$\nu=DP=p_*[D\alpha]$ is one fixed signed compact Radon measure.
For the smooth rows $a_j,b_j$, compact Stokes gives
\[
 \langle P,D(a_j-b_j)\rangle=\int(a_j-b_j)\,d\nu.
\]
Piecewise continuity, applied separately to the positive closed
currents $Da_j$ and $Db_j$, and fixed-measure bounded convergence
give the exact nonzero identity
\begin{equation}
             \langle D[f\circ p],\alpha\rangle
                         =\langle T,P\rangle\geq0.
 \label{eq:r6-finite-order-identity}
\end{equation}
The inequality follows from the preceding two lemmas. This does
not assume that $a_j-b_j$ has positive curvature.

For an original-$K$ disc product, follow the finite rational
localization of R4, Theorem 74.3, with nonnegative partitions.
For one positive localized test $\alpha$, take the same fixed
finite polynomial list, remove its base discriminant by the finite
etale splitting map, and set $\alpha_j=\theta_j(\log|H|)\alpha$,
where $0\leq\theta_j\leq1$ vanishes near $H=0$ and increases
pointwise to one off that fixed analytic divisor. These remain
positive tests. Finite etale order transfer gives positivity of
the pulled-back base curvature. For the split test upstairs, the
same smooth-row calculation as in
\eqref{eq:r6-finite-order-identity}, now with the split projected
test, gives a nonnegative pairing. The original finite-flat
integration formula multiplies the downstairs pairing by the
positive finite degree, so the downstairs pairing with $\alpha_j$
is nonnegative too.

The signed current $D[f\circ\pi]$ is the difference of the two
positive ambient prescribed-row curvature currents. Every coefficient
variation of these two currents vanishes on the removed analytic
divisor, by the frozen first-order nonconcentration theorem (R5,
Theorem 2.3, or its original-$K$ predecessor). Dominated convergence
against these fixed coefficient measures gives
\[
 \langle D[f\circ\pi],\alpha\rangle
   =\lim_j\langle D[f\circ\pi],\alpha_j\rangle\geq0.
\]
No uniform differentiability bound on $\alpha_j$ is needed.
Iteration gives all finite-dimensional polydisc products, always
over the original $K$ and with the pulled-back prescribed rows.
\end{proof}

\begin{lemma}[Positive finite-coefficient descent and positive fiber extraction]
\label{lem:r6-positive-descent-fiber-extraction}
The following two refinements of the frozen test-spreading arguments
hold.
\begin{enumerate}
\item In exact finite-coefficient descent from $X_L$ to
$X_E$, $E=\overline{K(c_1,\ldots,c_m)}$, a positive original
test can be descended to a positive original test. The analogous
exact descent from a completed algebraic closure to a finite
extension has the same property.
\item Let $j:X_H\hookrightarrow X_F\times_F A$ be the full
fiber over a parameter point, with $H$ its completed residue field.
Suppose a symmetric positive separately closed current $J$ on the
product is carried by this fiber and satisfies
\[
                    \langle J,\alpha\rangle
                          =\langle T_H,j^*\alpha\rangle
\]
for all original compact complementary product tests, where
$T_H$ is a coefficient-Radon current on the fiber. If exact
spreading of its original fiber tests is available as in R5,
Lemma 4.4, then $T_H$ is positive.
\end{enumerate}
The second statement does not assert that a positive fiber form
has a positive extension to the whole ambient product.
\end{lemma}

\begin{proof}
For (1), start with positive presentations from
\cite[Definition 5.3.1]{CLD12} and use nonnegative scalar
localizations. Perform the frozen rational-unit replacement on
slightly larger compact norm boxes. It preserves the actual unit
logarithms on each box, so preserves the positive tensor on its
tropical image. Include all coefficients of the rational units
\emph{and of the domains and cutoff profiles} in the finite field.
The resulting domain over $L$ is the scalar extension of the
corresponding domain over $E$. Scalar projection of an affinoid
base change is surjective; their images under these common unit
moments are therefore identical. The same real coefficient form is
positive on exactly this common polyhedral image, which is the
definition of original positivity. The descended localized tests
are consequently positive.

For descent from a completed algebraic closure, approximate the
finite coefficients by one finite algebraic extension, with errors
strictly below all unit lower bounds and all positive profile
thresholds. Closed upper and lower norm inequalities are unchanged
under these ultrametric error bounds, and profiles near possible
zeros are constant. Thus the same compact boxes and active norms
are recovered exactly after scalar extension. The identical-image
argument again proves positivity. This is stronger information
than equality of the signed integrated forms, and follows from
retaining the positive presentation itself.

For (2), take a positive compact fiber test and partition it into
nonnegative localizations of positive moment presentations on
compact fiber norm domains. Spread these particular presentations
by the frozen exact construction: it retains all unit logarithms,
norm-domain cutoffs, and zero-protected scalar profiles on the
fiber. A local spread can be written $\chi\,\widetilde t^*\eta$,
where $\chi$ is nonnegative, its restriction is the original
localizing scalar, and $\eta$ is positive on the finite tropical
image $P$ of the corresponding fiber domain.

The current $\chi J$ is positive, and its coefficient mass on the
fiber is supported inside that domain. Every level tangency of $J$
survives multiplication by $\chi$ and Borel restriction. Stratify
this mass by the faces of $P$, whether or not those faces occur in
the original ambient tropical complex. Their affine conormals are
original smooth affine combinations of the spread unit logarithms.
Lemma~\ref{lem:r6-carried-tangent-positivity} therefore gives
$\langle J,\chi\widetilde t^*\eta\rangle\geq0$. Add the finite
localizations and use the displayed restriction identity. The
arbitrary positive fiber test has a nonnegative pairing with
$T_H$, as required. Only the tangential tensor on the mass-carrying
fiber was required to be positive.
\end{proof}

\begin{theorem}[Order transfer through every complete scalar extension]
\label{thm:r6-all-complete-scalar-order}
Let $a,b\in\mathcal C_b(X)$ on an allowed original-$K$ space
and suppose $D[a-b]\geq0$ on all original tests. Then for every
complete $L/K$,
\[
                              D[a_L-b_L]\geq0
\]
on every original $L$-test.
\end{theorem}

\noindent\textbf{Status: \texttt{NEW\_PROVED}}, using the
positive trace and carried-fiber lemmas. The corresponding zero
equality was already proved in R5.

\begin{proof}
Work locally in one original algebraic neighborhood with prescribed
rows for $a,b$. Put $f=a-b$. Finite scalar extensions preserve
positivity by Proposition~\ref{prop:r6-finite-etale-product-order};
the same-space calibration factor is the positive number $[F':F]$.
Part (1) of Lemma~\ref{lem:r6-positive-descent-fiber-extraction},
and the factor-one fixed signed-measure scalar integration formula
\cite[Proposition 15.3.1]{CLD25}, then preserve positivity on
passing from an auxiliary $F$ to its completed algebraic closure.

Here is the one-parameter step, specifying the new point where order
differs from the zero-kernel argument. Suppose $D[f_F]\geq0$ on
$X_F$ and the same holds for its pullback to
$X_F\times\mathbf A_F^{1,\an}$. Given a point of the parameter
line, embed its residue field in a completed algebraic closure and
then take $C=\widehat{F^{\rm alg}}$ there, as in R5,
Proposition 4.6. The induced point $y\in\mathbf A_C^{1,\an}$
has a residue field $H$ containing the original parameter residue
field. Both assumed positive currents remain positive over $C$.

For a nonrigid $y$, R5, Proposition 4.3, gives the full-fiber
potential $b_y$: for types II and III it is
$\log\max(r,|z-c|)$, and for type IV it is the uniform limit of
the corresponding nested-disc potentials. Each is locally uniformly
approximable by original smooth positive functions: smooth the real
convex maximum, and in type IV also use the proved uniform error
$0\leq b_j-b_y\leq\log(r_j/r)$. These are explicit functions,
not an invocation of R over $C$.

On the product let $T=D[f_C\circ q]\geq0$ and put
\[
             J=D(b_yT)=D\bigl((f_C\circ q)D[b_y]\bigr).
\]
The equality is R5, Lemma 4.1, applied separately to the prescribed
rows for $a$ and $b$. Smooth positive approximation of $b_y$ and
fixed coefficient-Radon multiplication show that $J$ is positive
and separately closed. The second expression shows that it is
carried by the full parameter fiber. Its exact restriction identity
is
\[
                \langle J,\alpha\rangle
                     =\int_{X_H} f_H D(j^*\alpha),
\]
by the full-fiber formula and its fixed Borel version. Part (2) of
Lemma~\ref{lem:r6-positive-descent-fiber-extraction} now proves
$D[f_H]\geq0$ on every original fiber test. A rational parameter
point requires no new extension. Finally, pull a positive test
over the original parameter residue field to $H$; positive smooth
pullback and factor-one scalar integration descend the inequality.
This proves the one-parameter step without a positive ambient
extension of an arbitrary fiber test.

Propagate these statements simultaneously with every finite number
$r$ of affine-space factors. The initial assertions over $K$ are
Proposition~\ref{prop:r6-finite-etale-product-order}. If they hold
over $F$ for all $r$, the assertion with $r+1$ supplies the product
hypothesis for the one-parameter step applied to the assertion
with $r$. Hence they hold over the residue field of every point
of $\mathbf A_F^1$. Iterating through
$F_i=\overline{K(c_1,\ldots,c_i)}$ proves them over every finite
parameter field. These fields have the countable dense subfields
used in the frozen R5 construction; their algebraic closures and
all the prescribed rows are exactly those in that construction.

For an arbitrary complete $L/K$, descend any positive original
compact test exactly to one such finite parameter field by part (1)
of the preceding lemma. The fixed Borel scalar integration identity
identifies its pairing with the already nonnegative pairing over
that field. Nonnegative original partitions handle the initial
algebraic neighborhoods. This proves the theorem on all original
tests, without extending B's regularization hypothesis to any new
complete field.
\end{proof}

\begin{theorem}[Analytic functoriality of the represented order]
\label{thm:r6-all-extension-analytic-order}
Under the hypotheses of the preceding theorem, let $L/K$ be any
complete extension and let $F:Z\to X_L$ be an analytic morphism
from an allowed $L$-analytic space. Then
\[
                       D[(a_L-b_L)\circ F]\geq0
\]
on every original compact complementary test of $Z$.
This is an order statement for the specified actual bounded
classical functions, not a pullback of an arbitrary positive current.
\end{theorem}

\begin{proof}
First observe that the normal trace in R5, Lemma 5.2, preserves
order. If $D[c-d]\geq0$ on $B\times\mathbb D_s$, test it with
$\beta_N\wedge\pi^*\eta$ for an original positive compact
complementary $\eta$ on $B$. The normal form $\beta_N$ is smooth
positive of degree $(1,1)$, hence strongly positive; its wedge with
$\pi^*\eta$ is positive. The frozen trace identity, applied to
$c$ and $d$ separately, gives
\[
       \langle D[(c-d)|_{t=0}],\eta\rangle
          =\lim_N\langle D[c-d],\beta_N\wedge\pi^*\eta\rangle
          \geq0.
\]
The trace uses actual Gauss radial limits of the two classical
functions and the already checked normal coarea, including its full
finite-field degree. It applies over any complete field and iterates
in finitely many normal coordinates.

Use the precise geometry of R5, Theorem 5.3. Over the original $K$,
B's ordinary algebraic localization and R's tubular neighborhood
embed the target in an affine space and extend $a,b$ constantly
in normal polydisc directions. Proposition
\ref{prop:r6-finite-etale-product-order} gives the order on this
ordinary ambient neighborhood and on all its finite affine products.
Theorem~\ref{thm:r6-all-complete-scalar-order} gives those same
inequalities after extension to $L$.

At a source point use the genuine ordinary finite etale coordinates
of R5, Lemma 5.1. Their product with the ambient target is finite
etale; Proposition~\ref{prop:r6-finite-etale-product-order} transfers
the inequality to that product. Translation by the analytic map
$F$ makes its graph the zero section. Analytic isomorphisms preserve
the original forms, calibration, and their positive cone. The frozen
compact-tube argument supplies a whole normal polydisc tube around
the zero section on a smaller ordinary source neighborhood. Apply
the positive normal trace just proved in each coordinate. The
restriction is exactly the actual difference $(a_L-b_L)\circ F$.
Zero-dimensional source or target cases are immediate. Locality
proves the assertion globally on $Z$.

Every algebraic or tubular neighborhood used in this step is over
the original $K$. The new source over $L$ is handled by genuinely
etale ordinary-open coordinates and normal trace. All new-field
rows are the pullbacks of the original prescribed rows.
\end{proof}

\paragraph{Scope of the advance.}
The represented current order is now functorial for actual bounded
classical differences. A bounded Borel function with distributionally
positive curvature is not covered by this result. Nor has an arbitrary
positive current acquired a pullback, a potential, or an action on all
higher-codimension delta-tests. The carried-fiber argument specifically
avoids claiming that positivity of an arbitrary fiber test extends
to a positive ambient test.

% END source: sections/potential_order_functoriality.tex


% BEGIN canonical R6 source: sections/curve_order_domination_continuity.tex
\section{Continuity on a curve under a continuous Laplacian dominator}
\label{sec:r6-curve-order-continuity}

This section uses an order inequality of actual curve Laplacian
measures. It does not pull back an ambient current inequality.
That distinct operation is isolated in the order-transfer argument.

\begin{lemma}[Positive potential splitting]
\label{lem:r6-positive-potential-splitting}
Let $Y$ be a topological space, let $P_1,P_2:Y\to[0,\infty]$
be lower semicontinuous, and suppose $P_1+P_2$ is finite and
continuous. Then $P_1$ and $P_2$ are finite and continuous.
Consequently, for a nonnegative potential kernel $G$, if positive
measures satisfy $0\le\mu\le\tau$, their kernel potentials are
lower semicontinuous, and $P_\tau$ is finite continuous on $Y$,
then $P_\mu$ is finite continuous there.
\end{lemma}

\noindent\textbf{Status: \texttt{DIRECT\_COROLLARY}} of
semicontinuity and additivity of nonnegative integrals.

\begin{proof}
Finiteness follows from the nonnegative sum. Each $P_i$ is also
upper semicontinuous, since it is the continuous sum minus the
other lower semicontinuous function. For the potential assertion,
use $P_\tau=P_\mu+P_{\tau-\mu}$ and Tonelli's theorem.
\end{proof}

\begin{proposition}[Actual continuity under domination on a curve]
\label{prop:r6-curve-continuous-laplacian-dominator}
Let $F$ be complete, algebraically closed, nontrivially valued and
of characteristic zero, with its fixed absolute value. Let $C$ be
a good Hausdorff boundaryless smooth $F$-analytic curve. Suppose
$a$ is a finite continuous subharmonic function and $b$ is an
actual locally bounded subharmonic function on an ordinary open
$O\subset C$. If
\[
                         0\le\Delta b\le\Delta a
\]
as Radon measures on $O$, then $b$ and $a-b$ are continuous on
$O$, and $a-b$ is subharmonic there. The values at points of
all four Berkovich types are their given actual values.
\end{proposition}

\noindent\textbf{Status: \texttt{DIRECT\_COROLLARY}} of the
curve potential kernel and distributional harmonic kernel, with the
actual-value passage specified in the proof. This assertion is not
a higher-dimensional distributional positivity converse.

\begin{proof}
Local algebraization, in the precise complete-field scope of R5,
Lemma~\textup{R5--6.1}, places a neighborhood
of the point under consideration in a smooth projective algebraic
curve $\overline C$ over $F$. Choose a compact affinoid neighborhood
$Y$ of that point contained in $O$, and choose a type-II point
$\zeta\in\overline C^{\an}\setminus Y$. Shrinking $Y$ if necessary
makes this possible. Set
\[
             \tau=\mathbf1_Y\Delta a,\qquad
             \mu=\mathbf1_Y\Delta b.
\]
They are finite positive Radon measures and $\mu\le\tau$.

Use the potential kernel $g_\zeta(x,y)\ge0$ of
\cite[Definition 4.4, Lemma 4.7 and Proposition 4.8]{Wan22}.
For a finite positive measure $\rho$ supported in $Y$, write
\[
                   P_\rho(x)=\int_Y g_\zeta(x,y)\,d\rho(y).
\]
Its lower semicontinuity and its limits along paths to endpoints
are the kernel-potential properties in
\cite[Proposition 6.5]{Wan22}. Its distributional Laplacian is
$\rho(Y)\delta_\zeta-\rho$, as follows from the fixed-kernel
Laplacian identity and integration against a compact lisse test.
These facts concern the original absolute value. A positive constant
change in the convention for logarithmic lengths changes both sides
by the same factor and is not performed here.

For type-II or type-III $x$, the function $y\mapsto g_\zeta(x,y)$
is finite continuous on the compact curve, so $P_\rho(x)$ is finite.
Thus $a+P_\tau$ defines a degree-zero curve distribution on
$\operatorname{Int}Y$ with Laplacian zero. The distributional harmonic
kernel theorem (Thuillier, Lemma 3.3.12, as explicitly used in
\cite[Remarks 7.2 and 7.12]{Wan22} and in R5) gives a continuous
harmonic function $H_a$ agreeing with $a+P_\tau$ on the type-II
and type-III points of that interior.

This equality holds at the remaining actual points as well. If
$x$ has type I or IV, choose an open-ball neighborhood inside
$\operatorname{Int}Y$ and nested closed discs with boundary Gauss
points $\xi_j\to x$ on the segment toward $x$. The subharmonic
maximum principle and upper semicontinuity of $a$ give
$a(\xi_j)\to a(x)$; here continuity of $a$ would already suffice.
The path-limit property of the positive kernel potential gives
$P_\tau(\xi_j)\to P_\tau(x)$, including a possible infinite
limit before the equality is used. Since
$H_a(\xi_j)\to H_a(x)$ is finite, we obtain
\[
                  P_\tau(x)=H_a(x)-a(x)<\infty.
\]
The same equality was already known on the other two types.
It follows that $P_\tau=H_a-a$ is finite continuous throughout
$\operatorname{Int}Y$.

Apply Lemma~\ref{lem:r6-positive-potential-splitting} to
$P_\mu+P_{\tau-\mu}=P_\tau$. It shows that $P_\mu$ and
$P_{\tau-\mu}$ are continuous on this interior. Similarly,
$b+P_\mu$ has harmonic degree-zero distribution, so it agrees
on type-II and type-III points with a continuous harmonic $H_b$.
At a type-I or type-IV point, the same nested-disc argument now
uses the given subharmonicity of $b$:
\[
                  b(x)\le b(\xi_j),\qquad
                  \limsup_j b(\xi_j)\le b(x).
\]
Thus $b(\xi_j)\to b(x)$. Continuity of $P_\mu$ and $H_b$
gives $b=H_b-P_\mu$ at that point too. This proves continuity
of the actual $b$ on $\operatorname{Int}Y$.

Finally $a-b$ is continuous and its curve distribution has
Laplacian $\Delta a-\Delta b\ge0$. The continuous distributional
criterion for curve subharmonicity gives that it is subharmonic;
see \cite[Definition 3.11 and Remark 3.12]{Wan22}. The interiors
of such $Y$ cover $O$, proving all assertions locally and hence
on the entire open.
\end{proof}

\paragraph{How this enters the higher-dimensional argument.}
For a classical $u$ and a continuous classical $g$, a separately
proved order-transfer statement would give
$0\le\Delta(u\circ F)\le C\Delta(g\circ F)$ on every smooth
curve after every complete scalar extension. The proposition then
makes $u\circ F$ and $(Cg-u)\circ F$ continuous and the latter
subharmonic. This step uses actual curve functions. Ordinary
continuity on the higher-dimensional space still requires the
finite Shilov interpolation argument; it is not inferred from
continuity along a dense collection of curves alone.

% END source: sections/curve_order_domination_continuity.tex


% BEGIN canonical R6 source: sections/positive_dsh_representatives.tex
\section{Positive curvature in the actual bounded classical difference class}
\label{sec:r6-positive-dsh-representatives}

The difference class in this section consists of actual functions,
not arbitrary representatives of degree-zero currents. Its two positive
parts belong to $\mathcal C_b$ on the same ordinary neighborhood.
This restriction is essential to the curve endpoint arguments and to
the ambient semicontinuity proof.

\begin{lemma}[A positive difference of actual curve Laplacians]
\label{lem:r6-curve-positive-difference}
Let $F$ be algebraically closed, complete, nontrivially valued and of
characteristic zero. Let $a,b$ be actual locally bounded subharmonic
functions on an ordinary open in a smooth boundaryless analytic
$F$-curve. If $\Delta a-\Delta b\ge0$, then the actual function
$a-b$ is locally bounded subharmonic. Neither $a$ nor $b$ is assumed
continuous.
\end{lemma}

\noindent\textbf{Status: \texttt{DIRECT\_COROLLARY}} of the
positive curve kernel, with actual endpoint values retained.

\begin{proof}
Use the local projective algebraization and compact affinoid $Y$ of
Proposition~\ref{prop:r6-curve-continuous-laplacian-dominator}.
Choose its type-II kernel base point $\zeta$ outside $Y$, and put
$\tau=\mathbf1_Y\Delta a$, $\mu=\mathbf1_Y\Delta b$ and
$\rho=\tau-\mu\ge0$. For any of these finite positive measures
write $P_\nu(x)=\int g_\zeta(x,y)\,d\nu(y)$.
These are nonnegative lower semicontinuous potentials with the
path-limit and Laplacian properties of \cite[Propositions 6.5--6.6]{Wan22}.
They are finite at type-II and type-III points.

On $\operatorname{Int}Y$ the distributions $a+P_\tau$ and
$b+P_\mu$ have Laplacian zero. The curve distributional kernel
theorem gives continuous harmonic functions $H_a,H_b$ agreeing with
them on those two types. At a type-I or type-IV point $x$, use
nested closed discs inside an open-ball neighborhood, with boundary
points $\xi_j\to x$ along its incoming path. The actual subharmonic
maximum principle and upper semicontinuity give
$a(\xi_j)\to a(x)$ and $b(\xi_j)\to b(x)$ separately.
The kernel path limits then identify
\[
                  P_\tau=H_a-a,\qquad P_\mu=H_b-b
\]
at $x$ as well, and show that both are finite there. No continuity
of $a$ or $b$ has been used for this actual-value identification.
The same identities already hold on types II and III.

Additivity of nonnegative integrals now yields
$P_\tau=P_\mu+P_\rho$, with all terms finite on the interior.
Consequently
\[
                     a-b=H_a-H_b-P_\rho
\]
is upper semicontinuous there. We verify harmonic comparison with
these actual values, rather than use an unrestricted representative
criterion. Put $f=a-b$. Let $Z\Subset\operatorname{Int}Y$ be a
connected strictly affinoid domain and let $h$ be harmonic on $Z$
with $f\le h$ on $\partial Z$. For a type-II or type-III point
$x\in\operatorname{Int}Z$, the compact nonnegative lisse Green
function $g_{Z,x}$ of \cite[Proposition 2.12]{Wan19} satisfies
\[
 \Delta g_{Z,x}=\omega_x-\delta_x,\qquad
 \operatorname{supp}\omega_x\subset\partial Z,\qquad
 h(x)=\int_{\partial Z}h\,d\omega_x.
\]
The boundary measure $\omega_x$ is nonnegative: the lisse Green
function is nonnegative, vanishes on the boundary and outside $Z$,
and its outward slope sum at a boundary point is therefore
nonnegative. Its mass is one, as follows by taking $h=1$.
Evaluation of the positive distribution $\Delta f$ on this
nonnegative lisse test gives
\[
 0\le\langle\Delta f,g_{Z,x}\rangle
      =\int_{\partial Z}f\,d\omega_x-f(x)
      \le h(x)-f(x).
\]
At a type-I or type-IV interior point, the earlier nested-disc
argument gives the actual path limits of $a$ and $b$ separately.
Pass the last inequality along type-II/III points on that path,
using continuity of $h$, to obtain $f\le h$ there as well.
Thus $f$ satisfies the original affinoid harmonic comparison
principle, so it is subharmonic. Local boundedness is inherited
from the original difference. Cover by such interiors.
\end{proof}

\begin{lemma}[One-sided Shilov bounds without ambient semicontinuity]
\label{lem:r6-subharmonic-shilov-without-usc}
In the geometric setting of R5, Theorem
\textup{R5--7.2}, replace harmonic curve
pullbacks by actual subharmonic curve pullbacks. Then an actual
finite function $f$ satisfies
\[
                       f(x)\le\max_{\Gamma(V)}f\quad(x\in V).
\]
No upper semicontinuity of $f$ on the ambient space is assumed.
The same bound applies on each of the frozen semistable charts.
\end{lemma}

\noindent\textbf{Status: \texttt{DIRECT\_COROLLARY}} of the
one-sided argument in the frozen curve-slicing proof.

\begin{proof}
Use the finite flat map $V\to\mathbb B^n$ whose coordinates extend
regularly to the algebraic ambient space. Fixing all but one coordinate
and passing to its completed residue field gives the full affinoid
curve fiber. Its algebraic normalization is a smooth curve defined
across the whole compact fiber. Subharmonicity of the actual pullback
and the curve maximum principle bound its value at every point by its
maximum on that fiber's finite Shilov boundary. The boundary is the
preimage of the corresponding one-variable Gauss point, by the frozen
finite flat Gauss-fiber lemma.

If $M(t)$ is the maximum of $f$ on the finite nonempty fiber over
$t\in\mathbb B^n$, this proves $M(t)\le M(G_i(t))$ for each
coordinate Gauss operation $G_i$. These operations commute and their
composition is the full Gauss point. Iterating proves the displayed
inequality. All maxima of finite fibers are actual finite numbers.
The argument uses only one sign and no ambient semicontinuity.
For a semistable chart, its regular topological generators give the
required finite flat algebraic Noether coordinates by R5,
Lemma~\textup{R5--9.1}.
\end{proof}

\begin{theorem}[Upper semicontinuity of a curvewise subharmonic difference]
\label{thm:r6-dsh-curvewise-usc}
Let $a,b\in\mathcal C_b(\Omega)$ on an allowed original-$K$ space,
and put $f=a-b$ as an actual function. Suppose every pullback of $f$
to a smooth analytic curve after every complete extension is an actual
subharmonic function. Then $f$ is upper semicontinuous on $\Omega$.
In particular $f\in\mathcal C_b(\Omega)$.
\end{theorem}

\noindent\textbf{Status: \texttt{NEW\_PROVED}}. The essential
step is an upper semicontinuous majorant obtained from a proper
finite-chart cover; no ambient regularity is assumed for the difference.

\begin{proof}
Fix $x\in\Omega$ and $\epsilon>0$. Since $a$ is upper
semicontinuous, choose an ordinary neighborhood $O$ of $x$ with
$a(z)<a(x)+\epsilon$ for every $z\in O$. Use the frozen algebraic
localization and norm-box construction to obtain $x\in\operatorname{Int}V$
with $V\subset O$. Let $q:W\to V$ and its finitely many compact
semistable charts $W_\lambda$ be the cover of R5, Proposition
\textup{R5--9.3}.

For each chart, use its nonnegative unit-logarithm coordinates
$\lambda_i$, summing to one, to define the continuous pluriharmonic
interpolants
\[
 P_\lambda a=\sum_i a(q(\xi_i))\lambda_i,
 \qquad P_\lambda b=\sum_i b(q(\xi_i))\lambda_i,
 \qquad P_\lambda f=P_\lambda a-P_\lambda b.
\]
All functions are defined on genuine ordinary neighborhoods of the
charts. The Shilov values of $f\circ q-P_\lambda f$ vanish,
and this difference has actual subharmonic curve pullbacks.
Lemma~\ref{lem:r6-subharmonic-shilov-without-usc} therefore gives
\[
                         f\circ q\le P_\lambda f.
\]
The ordinary classical maximum principle applied to
$b\circ q-P_\lambda b$ gives $b\circ q\le P_\lambda b$.
Since every $q(\xi_i)$ lies in $V\subset O$ and the coefficients
are nonnegative and sum to one, $P_\lambda a<a(x)+\epsilon$ on
the entire chart. In particular, for every chart point $y$ above $x$,
\[
             P_\lambda f(y)
                \le a(x)+\epsilon-b(x)=f(x)+\epsilon.
\]

Take the disjoint union of these finitely many compact charts and
let $Q:\coprod_\lambda W_\lambda\to V$ be its continuous
surjection. The function which equals $P_\lambda f$ on each component
is continuous. Define
\[
        e(z)=\max\{P_\lambda f(y):y\in W_\lambda,\ q(y)=z\}.
\]
Every fiber is nonempty and compact. For any real $c$, the set
$\{e\ge c\}$ is the image under $Q$ of the compact set on which
this continuous componentwise function is at least $c$; it is closed
in the Hausdorff space $V$. Hence $e$ is upper semicontinuous.
The preceding inequalities give $f\le e$ on $V$ and
$e(x)\le f(x)+\epsilon$. Since $x$ is in the ordinary interior,
\[
                   \limsup_{z\to x} f(z)
                         \le e(x)\le f(x)+\epsilon.
\]
Let $\epsilon\downarrow0$. This proves upper semicontinuity at
$x$. The difference is locally bounded. Once this ambient upper
semicontinuity has been established, the smooth-curve testing reduction
in \cite[4.1.3 and Definition 3.4]{GR25} applies to the assumed
complete-extension curve property and gives membership in the actual
classical class. In particular it covers test curves with boundary or
singularities by the compositions appearing in that definition.
Zero-dimensional components are locally discrete and immediate.
\end{proof}

\begin{theorem}[Positivity determines the bounded classical difference cone]
\label{thm:r6-positive-bounded-dsh-classical}
On an allowed original-$K$ ordinary open $\Omega$, let
$f=a-b$ with $a,b\in\mathcal C_b(\Omega)$. Then
\[
             D[f]\ge0\text{ in the original CLD cone}
                         \quad\Longleftrightarrow\quad
             f\in\mathcal C_b(\Omega).
\]
The same assertion holds for a locally represented difference, with
possibly different decompositions on different ordinary neighborhoods.
No assertion is made for an arbitrary locally bounded Borel function
whose associated degree-zero distribution has positive curvature.
\end{theorem}

\noindent\textbf{Status: \texttt{NEW\_PROVED}}, relative to the
order-functoriality theorem and the preceding semicontinuity mechanism.

\begin{proof}
The reverse implication is the positive-current assertion already
contained in B. Suppose $D[a-b]\ge0$. For every complete extension
and every smooth analytic curve map, Theorem
\ref{thm:r6-all-extension-analytic-order} gives positive original
curvature for the actual pullback difference. Over an algebraically
closed complete field, the frozen curve algebraization and B's explicitly complete-field
Lemma A.2, in the scope checked in R5, Lemma
\textup{R5--6.1}, identify this as
$\Delta(a\circ F)-\Delta(b\circ F)\ge0$ as Radon measures.
Indeed the inequality holds on all nonnegative original smooth scalar
tests; their compact uniform density extends it to nonnegative
continuous tests. Lemma~\ref{lem:r6-curve-positive-difference} makes
the actual difference subharmonic on the curve.

For a curve over a non-algebraically closed complete field, perform
the argument after completed algebraic closure. Upper semicontinuity
descends through scalar projection: on a compact affinoid neighborhood
it is a continuous surjection from a compact space onto a Hausdorff
space, and the upper level sets of the fiber-constant pullback are
closed; their images give the upper level sets downstairs. Subharmonicity
then descends by \cite[Proposition 3.5(7)]{GR25}. We have proved the
actual subharmonicity assertion on every smooth complete-extension
test curve. Theorem~\ref{thm:r6-dsh-curvewise-usc} now supplies the
missing ambient upper semicontinuity and then applies the full
classical smooth-curve reduction. Local boundedness is already known.
This proves the forward implication. Ordinary restriction and locality prove the
locally represented version.
\end{proof}

\begin{corollary}[Order of represented curvatures transfers to every background]
\label{cor:r6-arbitrary-background-potential-order}
For any allowed fixed positive separately closed $(p,p)$-background
$S$, with $p<n$, and $a,b\in\mathcal C_b(\Omega)$,
\[
             D[a]\ge D[b]\quad\Longrightarrow\quad
             D(aS)\ge D(bS).
\]
All inequalities and equalities use every original complementary
compact smooth test. The same order replacement is valid in any
one differentiated slot of a mixed BT product, with its other slots
fixed in $\mathcal C_b$.
\end{corollary}

\noindent\textbf{Status: \texttt{DIRECT\_COROLLARY}}. Equality
of the represented curvatures was already handled by the R5 quotient;
the new conclusion preserves their order, including discontinuous
actual potentials.

\begin{proof}
The difference $a-b$ is actual classical by the theorem. Its product
$D((a-b)S)$ is positive by the frozen arbitrary-background calculus,
and equals $D(aS)-D(bS)$ by linearity of Borel multiplication and $D$.
For another differentiated slot, use symmetry and replace $S$ by the
positive separately closed tail containing the other fixed potentials.
No arbitrary-current wedge is introduced.
\end{proof}

\begin{theorem}[Continuous coefficient relations for bounded potential slots]
\label{thm:r6-bounded-curvature-continuous-module}
Let $g_1,\ldots,g_m\in\mathcal C_b(\Omega)$ be actual locally
bounded classical potentials, and let $f_1,\ldots,f_m$ be finite
continuous real functions. Define the coefficient-Radon currents
\[
           A=\sum_{i=1}^m f_iD[g_i],\qquad
           A_S=\sum_{i=1}^m f_iD(g_iS).
\]
For every fixed allowed background $S$, $A\ge0$ implies $A_S\ge0$.
In particular $A=0$ implies $A_S=0$ on all original tests.
Finite continuous-coefficient presentations by curvatures of locally
bounded classical differences therefore have a well-defined positive
linear action over $S$, independently of their presentation.
The differentiated potentials need not be continuous. The scalar
coefficients in this statement are continuous.
\end{theorem}

\noindent\textbf{Status: \texttt{NEW\_PROVED}} extension of the
continuous-potential module theorem to all actual bounded classical
slots. The coefficient-bearing currents are not asserted separately
closed merely because their individual curvature factors are closed.

\begin{proof}
Put $T=\sum_iD(g_iS)\ge0$. Fix $\epsilon>0$. Near any point choose
constants $c_i$ with $|f_i-c_i|<\epsilon$. On this ordinary neighborhood,
\[
 D\left(\sum_i c_i g_i+\epsilon\sum_i g_i\right)
       =A+\sum_i(c_i+\epsilon-f_i)D[g_i]\ge0.
\]
The actual function inside parentheses is a bounded classical
difference, by separating its finitely many positive and negative
real coefficients. Theorem~\ref{thm:r6-positive-bounded-dsh-classical}
therefore makes it classical, and its curvature over $S$ is positive.
Thus $\sum_i c_iD(g_iS)\ge-\epsilon T$. The coefficient error gives
\[
                A_S\ge-2\epsilon T
\]
on this neighborhood. Cover the support of a positive original test
by finitely many such neighborhoods and apply a nonnegative original
smooth partition. Positivity is local, so the same inequality holds
on that test globally. Let $\epsilon\downarrow0$, using its fixed
finite control mass against $T$. This proves $A_S\ge0$.

If $A=0$, apply the implication both to the coefficients $f_i$ and
to $-f_i$. The original middle positive cone of currents is pointed,
so $A_S=0$. To compare two signed-potential presentations, expand each
classical difference into its positive classical parts, subtract the
presentations, and use the proved zero relation. Addition and continuous
scalar multiplication are then well defined; the positive-order
assertion is exactly the first part of the theorem.
\end{proof}

\begin{corollary}[No a priori continuity is needed in a curvature order interval]
\label{cor:r6-continuity-from-continuous-curvature-domination}
Let $u\in\mathcal C_b(\Omega)$ and let $g$ be continuous classical.
If $0\le D[u]\le C D[g]$ for some $C\ge0$, then $u$ and $Cg-u$
are continuous classical functions. Consequently all compactness,
represented-quotient, actual-image and joint-continuity conclusions
above for normalized continuous curvature order families remain true
when their defining family initially allows all actual $u\in\mathcal C_b$.
This also applies to the family with $D[u]\le D[g]$ for some
$g$ in a fixed compact continuous upper-potential family $\mathscr G$.
\end{corollary}

\noindent\textbf{Status: \texttt{DIRECT\_COROLLARY}} of the
new bounded-difference positivity criterion.

\begin{proof}
The actual difference $Cg-u$ has positive curvature, so it is classical
and upper semicontinuous. Since $Cg$ is continuous, this says that
$u=Cg-(Cg-u)$ is lower semicontinuous. It is already upper
semicontinuous, hence continuous; the complement is then continuous
as well. Thus the enlarged defining family is exactly the previously
proved continuous family. In the compact-$\mathscr G$ case apply
this argument separately with the upper potential attached to each
member. All prior conclusions apply without a new approximation step.
\end{proof}

\paragraph{Dependency and remaining limits.}
The proof uses the order-transfer construction for the prescribed
functions, the curve kernel and actual endpoint values, the frozen
semistable geometric covers, and the one-sided Shilov bound. It does
not infer ambient semicontinuity from curve continuity alone, nor
invoke an unproved general distributional regularization theorem.
The bounded-difference criterion does not supply a potential for an
arbitrary positive closed current. The continuous-coefficient theorem
has one curvature factor in each summand; it does not annihilate every
higher-degree analytic delta relation or permit arbitrary Borel
coefficient replacements.

% END source: sections/positive_dsh_representatives.tex


% BEGIN canonical R6 source: logs/positive_dsh_actual_spike_scope.tex
\section{An actual upper-semicontinuous spike is outside the DSH criterion}
\label{sec:r6-actual-spike-scope}

\noindent\textbf{Status: \texttt{DIRECT\_COROLLARY}} of original
smooth top-form atomlessness and the curve maximum principle. The
proposed replacement of the actual bounded classical difference
hypothesis by bounded upper semicontinuity is \texttt{FALSE}.
This is the one-dimensional version of the already recorded R5
assigned-spike obstruction, not a new counterexample mechanism.

Fix a uniformizer $\pi$ of the original $K$, and let
\[
 \Omega=\{|\pi|^2<|T|<|\pi|^{-2}\}
                  \subset\mathbf G_{m,K}^{\mathrm{an}},
 \qquad x=\zeta_{0,1},\qquad f=\mathbf1_{\{x\}}.
\]
The actual function $f$ is bounded, Borel and upper semicontinuous,
since its nonzero level set is a closed singleton. Nevertheless
\[
                         [f]=0,\qquad D[f]=0,
\]
while $f$ is not classical psh and is not continuous.

\begin{proof}
Every original compact smooth $(1,1)$-form has an atomless signed
integration measure. The finite-moment proof is the one given in
Section~\ref{sec:r6-atomic-variation-obstruction}: positive smooth
top-form controls push forward to calibrated one-dimensional smooth
densities, which have no atom at a tropical point; positivity then
gives no mass to its entire inverse image. Rank-zero presentations
give the zero top form, and finite localization covers every test.
Thus the total variation of every such test measure gives mass zero
to $\{x\}$. Integrating the bounded assigned function $f$ yields
zero on every original test, proving both current equalities. This
does not say that arbitrary positive currents have no atoms.

If $f$ were classical on $\Omega$, it would be subharmonic on this
same original analytic curve. On the compact subannulus
\[
              A=\{|\pi|\le|T|\le|\pi|^{-1}\}\Subset\Omega
\]
the two Shilov boundary points are different from $x$ and have
value zero. The subharmonic maximum principle would give $f\le0$
throughout $A$, contradicting $f(x)=1$. The distinct Gauss points
$\zeta_{0,r}$ with $r\to1$ also show directly that $f$ is not
continuous at $x$.
\end{proof}

In particular $f$ cannot be an actual difference of two locally
bounded classical psh functions on any ordinary neighborhood of
$x$. Such a presentation would have equal original curvatures for
its two terms. The proved R5 kernel theorem would make their actual
difference pluriharmonic and continuous there, contrary to this
spike. Equivalently, the R6 ordered-DSH criterion applied to both
signs would give the same contradiction.

This check preserves the exact representative scope of
Theorem~\ref{thm:r6-positive-bounded-dsh-classical}. It does not
disprove existence of a continuous classical representative of the
zero distribution: the constant zero function is such a
representative. The failure lies in prescribing the spike's actual
point value and then treating it as if it were already in the
classical-difference domain of that theorem.

% END source: logs/positive_dsh_actual_spike_scope.tex


% BEGIN canonical R6 source: sections/continuous_curvature_module_products.tex
\section{Products in the continuous coefficient curvature module}
\label{sec:r6-continuous-curvature-module-products}

This section records a consequence of the represented-order theorem.
The input currents have a prescribed local finite potential
presentation. No product of completely arbitrary positive closed
currents is defined.

\begin{definition}
On an allowed ordinary open $U$, let $\mathscr M^1(U)$ be the
$C(U,\mathbf R)$-module of original coefficient-Radon $(1,1)$-currents
locally admitting a finite expression
\[
                         A=\sum_{\alpha}f_\alpha D[g_\alpha],
 \qquad f_\alpha\in C(U,\mathbf R),\quad g_\alpha\in\mathcal C_b(U).
\]
The expression is local; actual equality of currents is the equivalence
relation. Signed bounded classical potentials can be used after
expanding them into two classical parts. Its positive cone consists
of those represented currents which are positive in the original
middle CLD sense. Variable continuous coefficients need not preserve
separate closedness.
\end{definition}

\begin{theorem}[Positive products of continuously weighted curvatures]
\label{thm:r6-positive-continuous-module-products}
Let $S$ be any fixed allowed $(p,p)$-background and $r\le n-p$.
There is a well-defined symmetric $C(U,\mathbf R)$-multilinear map
\[
 \Phi_{S,r}:(\mathscr M^1)^r\longrightarrow
       \{\text{original coefficient-Radon $(p+r,p+r)$-currents}\}
\]
characterized on finite local presentations by
\begin{equation}
 \Phi_{S,r}(A_1,\ldots,A_r)
 =\sum_{\alpha_1,\ldots,\alpha_r}
    \left(\prod_{i=1}^r f_{i,\alpha_i}\right)
    T_{S,r}(g_{1,\alpha_1},\ldots,g_{r,\alpha_r}).
 \label{eq:r6-continuous-module-product}
\end{equation}
It commutes with ordinary restriction and gluing. If every $A_i$ is
positive, its image is positive. In particular the map is increasing
in each slot when all the other slots are positive. When every slot
is $D[g_i]$, it is the frozen bounded BT product.

No separate closedness of the output is asserted for general
continuous coefficients. The map does not factor through every
higher-degree ordinary-current relation merely from its construction.
\end{theorem}

\noindent\textbf{Status: \texttt{DIRECT\_COROLLARY}} of the new
bounded-difference positivity criterion and continuous-coefficient
relation theorem, with the positivity argument written out.

\begin{proof}
Each expression in \eqref{eq:r6-continuous-module-product} is a finite
sum of continuous multiples of positive coefficient-Radon currents.
Suppose the presentation of one slot is changed without changing its
original current. Fix one multi-index in the remaining slots and put
\[
 R=T_{S,r-1}(g_{1,\alpha_1},\ldots,
              \widehat{g_{i,\alpha_i}},\ldots,g_{r,\alpha_r}).
\]
This is an allowed positive separately closed background. The
zero-relation assertion of
Theorem~\ref{thm:r6-bounded-curvature-continuous-module}, applied to
$R$, makes the change in that slot zero. Multiplication by the
remaining continuous coefficient product and finite summation
preserve zero. Changing the slots one at a time proves presentation
independence. Finite linearity, scalar multiplication, symmetry of
the frozen BT tails, ordinary locality and gluing give the stated
module and sheaf properties.

For positivity fix a positive original compact complementary test
$\eta$. Use finitely many local presentations near its support and
nonnegative smooth localization; it suffices to treat one presentation.
All its coefficients are bounded by a fixed number $M\ge1$ on a
slightly larger compact neighborhood. For $0<\epsilon<1$, subdivide
the test support into finitely many ordinary neighborhoods on which
\[
                    |f_{i,\alpha}-c_{i,\alpha}|<\epsilon
\]
for suitable constants $c_{i,\alpha}$. On one such neighborhood put
\[
 H_i=\sum_\alpha(c_{i,\alpha}+\epsilon)g_{i,\alpha}.
\]
It is an actual bounded classical difference, and
\[
 D[H_i]=A_i+\sum_\alpha(c_{i,\alpha}+\epsilon-f_{i,\alpha})
                                      D[g_{i,\alpha}]\ge0.
\]
Theorem~\ref{thm:r6-positive-bounded-dsh-classical} therefore gives
$H_i\in\mathcal C_b$ on that neighborhood. Hence
$T_{S,r}(H_1,\ldots,H_r)\ge0$ there.

For every multi-index the scalar difference between its coefficient
in \eqref{eq:r6-continuous-module-product} and its coefficient in
$T_{S,r}(H_1,\ldots,H_r)$ has absolute value at most
$2r\epsilon(M+2)^{r-1}$. This follows by telescoping the product of
$r$ scalars, since each paired scalar differs by at most $2\epsilon$.
Let $R_0$ be the finite sum of all the positive BT tails in
\eqref{eq:r6-continuous-module-product}, without coefficients.
On this neighborhood the positive test cone therefore gives
\[
          \Phi_{S,r}(A_1,\ldots,A_r)
                   \ge -2r\epsilon(M+2)^{r-1}R_0.
\]
The same constant and the same $R_0$ work for every member of the
finite subdivision. Apply a nonnegative original smooth partition
to $\eta$, sum the inequalities, and let $\epsilon\downarrow0$.
The fixed control mass $\langle R_0,\eta\rangle$ is finite. The
pairing of $\Phi_{S,r}$ with $\eta$ is nonnegative. This proves
positivity on all original tests in the specified cone.

For monotonicity subtract the two products and telescope their
slots. Each term has one positive difference slot and all other
slots positive, so the proved positivity applies. The BT
specialization is its defining one-term presentation.
\end{proof}

\begin{example}[Borel coefficients do not preserve current relations]
\label{ex:r6-borel-curvature-module-failure}
The word continuous in the preceding module cannot be replaced by
bounded Borel, even when all potential functions are continuous.
Take the ordinary bidannulus
\[
 U=\{r<|T_1|<R,\ r<|T_2|<R\}\subset(\mathbf G_{m,K}^2)^{\an},
             \qquad 0<r<1<R,
\]
and let $Y=\{T_2=1\}\cap U$, $S=[Y]$, and
$g=\max\{0,\log|T_1|\}$. This is a continuous locally bounded
classical psh function. The frozen first-order analytic-divisor
nonconcentration theorem, applied to $D[g]$, gives
\[
                         \mathbf1_YD[g]=0
\]
as an ambient coefficient-Radon current. On the other hand, the
restriction and projection formulas for the subvariety background,
and the one-variable slope computation, give
\[
 \mathbf1_YD(g[Y])=D(g[Y])=i_*\delta_{\zeta_{0,1}}\ne0,
                  \qquad i:Y\hookrightarrow U.
\]
Here $\zeta_{0,1}$ is the radius-one Gauss point in the $T_1$ annulus;
the slope jump of $\max(0,\log|T_1|)$ is exactly one. An original
nonnegative compact smooth scalar test equal to one near this point
pairs with the displayed measure to give one. Thus the failure is
visible on original tests, not only on an enlarged test sheaf.
\end{example}

\noindent\textbf{Status: \texttt{FALSE}} for the proposed Borel
extension; the example is a \texttt{DIRECT\_COROLLARY} of the frozen
polar-carrier obstruction and first-order nonconcentration. It is not
counted as a new counterexample mechanism. R5,
Corollary~\textup{R5--23.2} and
Remark~\textup{R5--22.4}, already show why
arbitrary assigned bounded upper semicontinuous values must also be
excluded from the bounded-classical-difference positivity criterion.

\paragraph{What this construction does not settle.}
Independence of the degree-one presentations means that relations in
$\mathscr M^1$ may be substituted inside
\eqref{eq:r6-continuous-module-product}. It does not say that every
ordinary zero current formed from several curvature factors is such
a relation. The higher-degree analytic delta kernel remains a separate
problem. Nor is the output of a continuous coefficient product an
admissible new background unless its required positivity and separate
closedness are independently available.

% END source: sections/continuous_curvature_module_products.tex


% BEGIN canonical R6 source: sections/represented_order_stability.tex
\section{Quantitative stability in the represented curvature order}
\label{sec:r6-represented-order-stability}

The ordered quotient gives a quantitative convergence statement which
is stronger than weak convergence but has correspondingly stronger,
explicit hypotheses. The upper potentials here may be discontinuous.
The background is fixed throughout.

\begin{proposition}[A mixed order bound]
\label{prop:r6-mixed-represented-order-bound}
Let $S$ be any fixed allowed $(p,p)$-background and $k\le n-p$.
Let $u_i,v_i,w_i\in\mathcal C_b(U)$ and let $M_i,\epsilon_i\ge0$
be finite constants such that, on all original ambient tests,
\[
 0\le D[u_i],D[v_i]\le M_iD[w_i],\qquad
 -\epsilon_iD[w_i]\le D[u_i-v_i]\le\epsilon_iD[w_i]
                    \quad(1\le i\le k).
\]
Set
\[
 R=T_{S,k}(w_1,\ldots,w_k),\qquad
 c=\sum_{i=1}^k\epsilon_i\prod_{\ell\ne i}M_\ell.
\]
Then
\begin{equation}
 -cR\le T_{S,k}(u_1,\ldots,u_k)-T_{S,k}(v_1,\ldots,v_k)
                                      \le cR
 \label{eq:r6-mixed-order-difference}
\end{equation}
in the original positive cone. If $k=n-p$, these are Radon measures
and their signed difference $\nu$ satisfies $|\nu|\le cR$ on all
Borel sets. In particular this controls total variation on every
compact subset, without continuity of any $w_i$.
\end{proposition}

\noindent\textbf{Status: \texttt{DIRECT\_COROLLARY}} of the
new order-preserving quotient and the frozen mixed calculus.

\begin{proof}
Telescope the mixed difference, replacing the slots one at a time.
The $i$th signed term has differentiated potential $u_i-v_i$, with
all earlier slots chosen from the $v$ list and all later slots from
the $u$ list. Their product with $S$ is a positive separately closed
tail. Corollary~\ref{cor:r6-arbitrary-background-potential-order},
applied to the two inequalities for $D[u_i-v_i]$, bounds this term
above and below by $\epsilon_i$ times the same tail with slot $w_i$.
Replace each other slot in that positive tail by $M_\ell w_\ell$,
using the same order corollary and mixed symmetry. Its bound is
$\epsilon_i(\prod_{\ell\ne i}M_\ell)R$. Summation proves
\eqref{eq:r6-mixed-order-difference}. No division by $M_i$ is used,
so vanishing constants cause no difficulty.

In top degree the inequalities say $-cR\le\nu\le cR$ as signed
Radon measures. Radon--Nikodym differentiation with respect to $R$
gives $\nu=hR$ with $|h|\le c$ almost everywhere, hence
$|\nu|\le cR$. This proves the Borel and total-variation assertions.
\end{proof}

\begin{corollary}[Weighted total-variation convergence]
\label{cor:r6-order-weighted-total-variation}
Let $k=n-p$ and let $u_{i,j},u_i,w_i\in\mathcal C_b(U)$ be
\emph{sequences} with common finite constants $M_i$ such that
\[
 0\le D[u_{i,j}],D[u_i]\le M_iD[w_i],\qquad
 -\epsilon_{i,j}D[w_i]\le D[u_{i,j}-u_i]
                         \le\epsilon_{i,j}D[w_i],
              \qquad \epsilon_{i,j}\longrightarrow0.
\]
Put $R=T_{S,k}(w_1,\ldots,w_k)$, $T_j=T_{S,k}(u_{1,j},\ldots,u_{k,j})$
and $T=T_{S,k}(u_1,\ldots,u_k)$. Let $a_j,a$ be actual Borel
functions which are uniformly bounded on every compact subset and
satisfy $a_j\to a$ $R$-almost everywhere. Then
\[
                         |a_jT_j-aT|(K)\longrightarrow0
\]
for every compact $K\Subset U$.
The weights need not be continuous or psh under these hypotheses.
\end{corollary}

\noindent\textbf{Status: \texttt{DIRECT\_COROLLARY}}. This is
not a consequence of bare weak convergence with arbitrary Borel weights.

\begin{proof}
Let $c_j=\sum_i\epsilon_{i,j}\prod_{\ell\ne i}M_\ell$, and put
$M=\prod_iM_i$. The proposition gives $|T_j-T|\le c_jR$ and
order replacement gives $T\le MR$. Therefore, if $B_K$ bounds
$|a_j|$ on $K$,
\[
 |a_jT_j-aT|(K)
  \le B_Kc_jR(K)+M\int_K|a_j-a|\,dR.
\]
The first term tends to zero and the second tends to zero by fixed
finite-measure dominated convergence. The sequence hypothesis is
intentional: arbitrary pointwise convergent nets do not satisfy this
version of dominated convergence.
\end{proof}

\paragraph{Scope.}
The propositions concern finite locally represented curvature order
intervals and their explicit order errors. They neither turn ordinary
weak convergence into order convergence nor enlarge the continuous
coefficient relation module to arbitrary Borel coefficients. The Borel
weights in the last corollary are multiplied only after the relative
measures and a common relative dominating measure have been fixed.
No varying-background theorem is inferred from these estimates.

% END source: sections/represented_order_stability.tex


% BEGIN canonical R6 source: sections/continuous_joint_stability.tex
\section{A qualified varying-background continuity theorem}
\label{sec:r6-continuous-joint-stability}

The original fixed-background monotone theorem must not be replaced by
an arbitrary varying-background assertion. Local uniform convergence
to continuous potentials supplies a sufficient additional hypothesis.
It also gives weak continuity of the new degree-one delta action on
the full cone of allowed backgrounds.

\begin{proposition}[Continuous potentials and weakly varying backgrounds]
\label{prop:r6-continuous-joint-varying-background}
On one common ordinary domain, let $S_j\to S$ be a net of allowed
positive separately closed currents of degree $(p,p)$, converging on
every original compact smooth test. Suppose that, for
$1\le i\le k\le n-p$, the actual functions $u_{i,j},u_i$ are
continuous classical psh and $u_{i,j}\to u_i$ locally uniformly.
Then
\[
       T_{S_j,k}(u_{1,j},\ldots,u_{k,j})
                         \longrightarrow T_{S,k}(u_1,\ldots,u_k)
\]
on all original compact smooth tests. If $a_j,a$ are finite continuous
real functions with $a_j\to a$ locally uniformly, the same assertion
holds after multiplying the two sides by $a_j$ and $a$ respectively.
Finite signed expansions of the continuous potential slots are allowed
when their positive parts satisfy these same hypotheses.
\end{proposition}

\noindent\textbf{Status: \texttt{DIRECT\_COROLLARY}} of R's fixed
continuous multiplier lemma, coefficient control and the frozen BT
recursion. No monotonicity is required under these stronger uniform
and continuity hypotheses.

\begin{proof}
First record a multiplier fact. If $R_j\to R$ are positive currents
of one degree and $b_j\to b$ locally uniformly with $b$ continuous,
then for each original compact test $\eta$,
\[
 \langle b_jR_j-bR,\eta\rangle
    =\langle(b_j-b)R_j,\eta\rangle
                         +\langle b(R_j-R),\eta\rangle\longrightarrow0.
\]
The second term converges by R. The first is bounded by
$\|b_j-b\|_\infty$ on a fixed compact stage times finitely many
positive control pairings of $R_j$. Those controls are eventually
bounded by weak convergence. This proves the fact for nets, without
a global mass bound or a topology determined by countably many tests.

Use induction on $k$. The positive separately closed tails with
$k-1$ slots converge by induction. Apply the multiplier fact with
$b_j=u_{1,j}$ and then the continuous distributional operator $D$.
The frozen recursion identifies the two derivatives with the stated
mixed currents. The weighted assertion is one further application
of the same multiplier fact to those positive currents. Finite
signed expansions are handled term by term.
\end{proof}

\begin{corollary}[All-background weak continuity on a fixed degree-one delta-test]
\label{cor:r6-all-background-degree-one-weak}
Let $S_j,S$ be as above and let $\alpha$ be a fixed actual delta-form
of total degree $(1,1)$. Then the canonical ordinary currents
\[
                         \alpha\wedge S_j\longrightarrow\alpha\wedge S
\]
converge on every original compact smooth test. If $p=n-1$ and
$\alpha$ has compact actual support, the full delta pairings converge.
Thus, on the positive separately closed cone in remaining degree one,
the ordinary weak topology already controls every fixed compact actual
delta-test.
\end{corollary}

\noindent\textbf{Status: \texttt{DIRECT\_COROLLARY}} of the proved
degree-one relation theorem and continuous multiplier stability.

\begin{proof}
For the symmetric part use a fixed local Hessian presentation
$\alpha=\sum_i f_iD_\delta g_i$, with continuous PS coefficients
and continuous PS classical-difference potentials. Expand each $g_i$
into its fixed continuous positive parts. The proposition, first
with one constant potential row and then with the fixed continuous
coefficient $f_i$, proves convergence of each term on every original
test. A finite partition covers a compact test support.
The antisymmetric part is PS; the frozen piecewise-test convergence
theorem applied to $S_j$ handles it on any original complementary
test. Actual relation independence makes these local formulas the
canonical action for all backgrounds in the statement.

For a compact complementary $\alpha$, the ordinary supports of all
outputs are contained in the same actual compact support of $\alpha$.
An original smooth cutoff equal to one near that compact converts the
current convergence into convergence of their total integrals.
Original smooth tests are themselves actual delta-tests, so the
assertion about the two weak topologies follows in this degree.
\end{proof}

\paragraph{Excluded extension.}
Neither assertion allows a simultaneously varying discontinuous Borel
weight under mere weak background convergence. Likewise locally bounded
psh rows converging only pointwise cannot be substituted for the local
uniform continuous rows in the proposition. The frozen weak-current
counterexamples and the main fixed-$S$ monotone hypothesis retain
their force.

% END source: sections/continuous_joint_stability.tex


% BEGIN canonical R6 source: sections/curvature_order_compactness.tex
\section{Compactness of continuous curvature order intervals}
\label{sec:r6-curvature-order-compactness}

The order bound in this section is on the original ambient curvature.
It holds on the whole domain on which the functions are normalized.
No finite-dimensional assumption on these curvatures is made.
Fix an allowed ordinary open $\Omega$, a continuous
$g\in\mathcal C_b(\Omega)$, and a real constant $C\geq0$. Put
\begin{equation}
 \mathscr P(g,C;\Omega)=
 \{u\in C(\Omega)\cap\mathcal C_b(\Omega):
       0\leq u\leq1,\quad 0\leq D[u]\leq C D[g]\}.
 \label{eq:r6-curvature-order-family}
\end{equation}
The inequalities between currents use every original positive
complementary test. All scalar extensions below retain the given
absolute value.

\begin{lemma}[The continuous curvature converse used here]
\label{lem:r6-continuous-positive-classical}
Let $h$ be a finite continuous function on an allowed ordinary open.
If $D[h]$ is original CLD-positive, then $h$ is locally bounded
classical psh. In particular, both $u$ and $Cg-u$ are continuous
classical psh for every $u\in\mathscr P(g,C;\Omega)$.
\end{lemma}

\noindent\textbf{Status: \texttt{DIRECT\_COROLLARY}} of the
authorized continuous regularization in R and the external smooth
curve criterion. This is not a converse for discontinuous
distributional potentials.

\begin{proof}
After B's ordinary algebraic localization, R, Theorem 2.9, supplies
on a fixed neighborhood smooth original functions $h_j$ with positive
curvature forms and $h_j\to h$ locally uniformly. An original smooth
function with positive curvature form is classical psh: original
smooth forms pull back along analytic maps and complete scalar
extension, and positive $(1,1)$ forms remain positive. On a smooth
curve the resulting smooth function is subharmonic by the smooth case
of Wanner's Theorem 4.7 \cite{Wan19}. The passage to all analytic
curves and all complete extensions is exactly the smooth-function
argument recorded in \cite[\S1.2, p.~3; \S4.1.3]{GR25}.
It concerns pullback of smooth forms, not pullback of an arbitrary
current. Wanner's auxiliary normalization of the logarithmic value
group is harmless here: multiplying all logarithmic lengths by a
positive constant preserves positivity and the subharmonic comparison
inequality, and the absolute value in the present construction is
never changed.

The continuous pointwise limit $h$ is classical psh by
\cite[Theorem 6.5]{GR25}. This assertion can also be checked directly
by passing a local uniform limit through the curve comparison
inequalities. Thus no regularization theorem over a new complete
field has been used. Continuity gives local boundedness. Apply this
to $h=Cg-u$ for the last assertion.
\end{proof}

\begin{lemma}[Positive Shilov interpolation on the fixed charts]
\label{lem:r6-positive-shilov-interpolation}
Let $V$ be an original algebraic norm box as in R5, Lemma 9.2,
contained in an ordinary algebraic neighborhood. Use the compact
surjective semistable cover $q:W\to V$ and one of its finitely many
affinoid charts $W_\lambda$ from R5, Proposition 9.3. Write its
Shilov boundary as $\{\xi_0,\ldots,\xi_r\}$ and its logarithmic
coordinates as
\[
 \lambda_i=-\frac{\log|x_i|}{-\log|\varpi|},
 \qquad P_\lambda a=\sum_{i=0}^r a(q(\xi_i))\lambda_i.
\]
Then $\lambda_i\geq0$ and $\sum_i\lambda_i=1$ on $W_\lambda$.
If $a$ is continuous classical psh on an ordinary neighborhood of
$V$, then
\[
                  a\circ q\leq P_\lambda a\quad\hbox{on }W_\lambda.
\]
For every $u\in\mathscr P(g,C;\Omega)$ and every such chart over
$V\subset\Omega$,
\begin{equation}
 0\leq P_\lambda u-u\circ q
       \leq C(P_\lambda g-g\circ q)
       \leq C\,\operatorname{osc}_V(g).
 \label{eq:r6-shilov-order-error}
\end{equation}
\end{lemma}

\noindent\textbf{Status: \texttt{DIRECT\_COROLLARY}} of the
frozen semistable charts and the affinoid maximum principle.

\begin{proof}
The regular functions $x_i$ are power bounded on their formal generic
fiber, and $x_0\cdots x_r=\varpi$. Hence the logarithmic coordinates
are nonnegative and sum to one. R5, Lemma 8.2, identifies their
values at the Shilov points as $\lambda_i(\xi_j)=\delta_{ij}$.
Each $x_i$ is an analytic unit on a genuine ordinary algebraic
neighborhood of the chart. Therefore $P_\lambda a$ is a real linear
combination of logarithms of units and is pluriharmonic there.

The function $a\circ q-P_\lambda a$ is classical psh near the entire
chart and vanishes at each of its Shilov points. The affinoid maximum
principle \cite[Theorem 6.4]{GR25} proves the asserted inequality.
Apply it first to $a=u$ and then to $a=Cg-u$, which is classical by
Lemma~\ref{lem:r6-continuous-positive-classical}. Linearity of
$P_\lambda$ gives the middle inequality in
\eqref{eq:r6-shilov-order-error}. Since $P_\lambda g$ is a convex
combination of values of $g$ in $V$, the final inequality follows.
No arbitrary vertex interpolant is assumed harmonic: these are
precisely the unit-logarithm interpolants on the frozen charts.
\end{proof}

\begin{theorem}[Finite evaluation and local uniform compactness]
\label{thm:r6-curvature-order-compactness}
For each compact $K_0\Subset\Omega$ and each $\epsilon>0$ there is
a finite set $F_\epsilon\subset\Omega$, depending on
$g,K_0,\epsilon$ and the chosen geometric neighborhoods, such that
\begin{equation}
 \|u-v\|_{K_0}
   \leq \max_{z\in F_\epsilon}|u(z)-v(z)|+2C\epsilon
 \qquad(u,v\in\mathscr P(g,C;\Omega)).
 \label{eq:r6-curvature-order-finite-evaluation}
\end{equation}
The family $\mathscr P(g,C;\Omega)$ is compact in the topology of
uniform convergence on compact subsets of $\Omega$. In particular,
every net in this family has a subnet converging locally uniformly
to an actual member of the same family.

The sets $F_\epsilon$ need not be independent of $\epsilon$.
Neither finite-dimensionality of the family nor convergence in
total variation of its curvature measures is asserted.
\end{theorem}

\noindent\textbf{Status: \texttt{NEW\_PROVED}}, using R's
continuous regularization, the external affinoid maximum principle,
and the frozen R5 geometric construction. This label makes no
claim about priority over other compactness theorems.

\begin{proof}
By continuity of $g$, cover $K_0$ by finitely many ordinary
neighborhoods on each of which $\operatorname{osc}(g)<\epsilon$.
Use B's ordinary local algebraization and R5, Lemma 9.2, to choose
finitely many norm boxes $V_a$ within these neighborhoods whose
ordinary interiors cover $K_0$. Apply the frozen semistable cover
to each $V_a$. It gives finitely many charts and finitely many
Shilov points on each chart. Define $F_\epsilon$ to be all their
images in $\Omega$. These points and covers depend on $g$ and the
accuracy, but not on $u$ or $v$.

For $x\in K_0$, choose a box containing it and a preimage
$y\in W_\lambda$ in one of its charts. The preceding lemma and
the nonnegative coordinates give
\begin{align*}
 |u(x)-v(x)|
 &\leq |P_\lambda u(y)-P_\lambda v(y)|+2C\epsilon\\
 &\leq \max_{z\in F_\epsilon}|u(z)-v(z)|+2C\epsilon.
\end{align*}
This proves \eqref{eq:r6-curvature-order-finite-evaluation}. In
particular, a finite grid in $[0,1]^{F_\epsilon}$ gives finite
uniform nets on $K_0$: choose one actual family member for each
occupied grid cell and apply the estimate to any two members in
that cell. No interpolant has to descend through the alteration.

For full compactness, take any net $(u_j)$ in the family. Tychonoff
compactness of $[0,1]^\Omega$ gives a pointwise convergent subnet;
write its actual limit as $u$. Given a compact $K_0$ and a positive
tolerance, choose $\epsilon$ so that $2C\epsilon$ is smaller than
half the tolerance (when $C=0$ choose any $\epsilon$). The subnet
is Cauchy on the finite set $F_\epsilon$, so the displayed estimate
makes it uniformly Cauchy on $K_0$. Completeness of $C(K_0)$ gives
uniform convergence there, and that uniform limit is the already
chosen pointwise function $u$.

The space $\Omega$ is locally compact Hausdorff and has compact
neighborhoods, so these compact restrictions show that $u$ is
continuous and that convergence is locally uniform. The inequalities
$0\leq u\leq1$ persist. Every fixed original compact test pairs
continuously with $D[u_j]$ under local uniform convergence: apply
the fixed signed top-form measure to $u_j-u$. Consequently
$0\leq D[u]\leq C D[g]$. Finally, $u$ is classical by
\cite[Theorem 6.5]{GR25}, or by
Lemma~\ref{lem:r6-continuous-positive-classical}. Hence $u$ belongs
to the same family. The net characterization of compactness now
proves the assertion without a countability assumption on $\Omega$.
\end{proof}

\begin{lemma}[Uniform continuity of the fixed-background BT map]
\label{lem:r6-uniform-bt-map}
Fix an allowed positive separately closed background $S$. Suppose
$u_{i,j},u_i$ are continuous classical psh, for $1\leq i\leq k$,
and $u_{i,j}\to u_i$ locally uniformly. Then, on every original
compact complementary smooth test,
\[
 T_{S,k}(u_{1,j},\ldots,u_{k,j})
                  \longrightarrow T_{S,k}(u_1,\ldots,u_k).
\]
The assertion remains true after multiplication by continuous
$w_j\to w$ locally uniformly, with locally bounded $w_j$.
\end{lemma}

\noindent\textbf{Status: \texttt{DIRECT\_COROLLARY}} of the
fixed-continuous-multiplier result in R and relative CLN.

\begin{proof}
Proceed by induction through the ordered recursion. If the positive
tails $R_j$ converge to $R$, relative CLN gives uniform coefficient
variation bounds on each fixed compact test stage. Thus
$(u_{1,j}-u_1)R_j\to0$ by local uniform convergence. R's fixed
continuous multiplication gives $u_1R_j\to u_1R$, so applying $D$
proves the next induction step. The same two-term decomposition
with $w_j-w$ proves the weighted assertion. This proof works for
nets, uses all original test forms, and never varies the initial
background $S$.
\end{proof}

\begin{corollary}[Actual representatives of weak limits on the same domain]
\label{cor:r6-dominated-family-actual-weak-closure}
Fix $k\leq n-p$ and, for each $i$, continuous classical $g_i$ on
the whole $\Omega$ and constants $C_i\geq0$. For every fixed
allowed background $S$, the set
\[
 \left\{T_{S,k}(u_1,\ldots,u_k):
               u_i\in\mathscr P(g_i,C_i;\Omega)\right\}
\]
is compact and closed in the ordinary weak current topology.
Every member of its weak closure therefore has a representation
with actual normalized continuous classical potentials on the
same $\Omega$, satisfying the same ambient curvature bounds.
The corresponding diagonal image is also compact and closed.

For $k=n-p$, the restricted diagonal raw capacity has no compact
mass gain under weak closure. This conclusion concerns the displayed
family; no common continuous curvature dominator for the unrestricted
normalized family is asserted.
\end{corollary}

\noindent\textbf{Status: \texttt{DIRECT\_COROLLARY}} of
Theorem~\ref{thm:r6-curvature-order-compactness}.

\begin{proof}
The finite product of the potential families is compact, and the
BT map is continuous by Lemma~\ref{lem:r6-uniform-bt-map}. Its
image is compact, and the weak current topology is Hausdorff because
all original tests are retained. Thus the image is closed. Restrict
the same argument to the diagonal when all slots use one family.
A measure in that diagonal image is itself one of the normalized
test measures defining the restricted raw capacity, which proves
the final compact-mass assertion with constant one. This argument
does not infer compact-mass continuity from weak convergence.
\end{proof}

\begin{corollary}[Ambient weak curvature stability within an order interval]
\label{cor:r6-order-interval-quotient-stability}
For each $i$, let $u_{i,j}\in\mathscr P(g_i,C_i;\Omega)$ and
suppose that the ambient currents $D[u_{i,j}]$ converge weakly on
all original tests to $A_i$. There are actual
$u_i\in\mathscr P(g_i,C_i;\Omega)$ with $D[u_i]=A_i$, and for
every fixed allowed background $S$,
\[
       T_{S,k}(u_{1,j},\ldots,u_{k,j})
                       \longrightarrow T_{S,k}(u_1,\ldots,u_k).
\]
The right side is independent of the chosen cluster potentials.
No statement for varying backgrounds or for an undominated family
is made.
\end{corollary}

\noindent\textbf{Status: \texttt{DIRECT\_COROLLARY}} using the
frozen R5 arbitrary-background potential quotient, Theorem 10.2.

\begin{proof}
Compactness gives a simultaneous locally uniform cluster tuple,
and pairing with original tests identifies its curvatures with
the $A_i$. Every other cluster tuple has those same curvatures.
The R5 quotient and differentiated-slot symmetry identify their
mixed products over $S$, slot by slot with a positive closed tail.
Lemma~\ref{lem:r6-uniform-bt-map} identifies every cluster of the
mixed-current net with this same current. Since its image lies in
a compact Hausdorff set, a unique cluster point implies convergence
of the whole net. No convergence of a chosen pluriharmonic gauge is
needed.
\end{proof}

\paragraph{Limits of the conclusion.}
The local uniform compactness is a property of continuous actual
potentials under a fixed continuous ambient order bound. It does not
follow from a relative energy bound alone. In particular the sine
example in Section~\ref{sec:r6-capacity-dominated-curvatures} is
compatible with this theorem: its potentials converge uniformly,
while the curvature measures fail to converge in total variation.
If the order bound is known only on a smaller $W\subset\Omega$,
the argument gives compactness on compact subsets of $W$; it does
not produce a normalized representative on all of $\Omega$.

% END source: sections/curvature_order_compactness.tex


% BEGIN canonical R6 source: sections/curvature_interval_topology.tex
\section{The weak topology on represented curvature order intervals}
\label{sec:r6-curvature-interval-topology}

This section states the topological consequence on ambient curvature
currents themselves. Its compactness input concerns actual normalized
continuous potentials on one common ordinary domain. It is not a
compactness assertion for the unrestricted bounded psh class.

For continuous classical $g$ on $\Omega$ and $C\geq0$, define
\[
 \mathscr A(g,C;\Omega)
     =\{D[u]:u\in\mathscr P(g,C;\Omega)\},
\]
where $\mathscr P$ is defined in
\eqref{eq:r6-curvature-order-family}. Give $\mathscr A$ the subspace
topology of ordinary weak current convergence on every original test.
Write $\mathscr S_+^p(\Omega)$ for the cone of allowed symmetric
original-positive separately closed $(p,p)$ currents, with the same
weak topology.

\begin{proposition}[The curvature quotient has its ordinary weak topology]
\label{prop:r6-curvature-interval-quotient-topology}
The set $\mathscr A(g,C;\Omega)$ is compact, convex, and Hausdorff.
The map
\[
 \mathscr P(g,C;\Omega)\longrightarrow\mathscr A(g,C;\Omega),
                          \qquad u\longmapsto D[u],
\]
is a closed quotient map. Its fibers are exactly the equivalence
classes obtained by allowing pluriharmonic differences while retaining
membership in the normalized potential family.
\end{proposition}

\noindent\textbf{Status: \texttt{DIRECT\_COROLLARY}} of the
compactness theorem and the frozen R5 actual potential quotient.

\begin{proof}
The potential family is compact by
Theorem~\ref{thm:r6-curvature-order-compactness}. Its map to currents
is continuous, since a fixed original pairing with $D[u]$ is an
integral of $u$ against a fixed signed top-form measure. The target
weak current space is Hausdorff. A continuous surjection from a compact
space to a Hausdorff space is closed and is a quotient map. Both the
potential family and its curvature image are convex, by the defining
linear order inequalities. R5, Theorems 10.1--10.2, identify equality
of the actual bounded-classical curvatures with a continuous
pluriharmonic difference. This identifies the fibers without selecting
a gauge or asserting the existence of a continuous section.
\end{proof}

\begin{theorem}[Joint weak continuity in background and ambient curvatures]
\label{thm:r6-curvature-interval-joint-topology}
Fix $k\leq n-p$. For each $i$, fix a continuous classical $g_i$ on
the whole $\Omega$ and $C_i\geq0$, and put
$\mathscr A_i=\mathscr A(g_i,C_i;\Omega)$. There is a well-defined
continuous map
\begin{equation}
 \begin{split}
 \mathscr S_+^p(\Omega)\times\prod_{i=1}^k\mathscr A_i
     &\longrightarrow\mathscr S_+^{p+k}(\Omega),\\
 (S,A_1,\ldots,A_k)&\longmapsto
             \mathsf T(S;A_1,\ldots,A_k)
              :=T_{S,k}(u_1,\ldots,u_k),\qquad D[u_i]=A_i.
 \end{split}
 \label{eq:r6-curvature-interval-joint-map}
\end{equation}
All topologies in this display are the ordinary weak topologies on
all original compact smooth tests. In particular, for nets
$S_j\to S$ and $A_{i,j}\to A_i$ in these specified spaces,
\[
       \mathsf T(S_j;A_{1,j},\ldots,A_{k,j})
              \longrightarrow\mathsf T(S;A_1,\ldots,A_k).
\]
For each fixed actual delta-form $\alpha$ of total degree $(1,1)$,
with $p+k+1\leq n$, the same map followed by the canonical operation
$R\mapsto\alpha\wedge R$ is continuous on original tests. If
$p+k=n-1$, this includes convergence against every fixed compact
complementary actual delta-test.
\end{theorem}

\noindent\textbf{Status: \texttt{DIRECT\_COROLLARY}} of
Theorem~\ref{thm:r6-curvature-order-compactness},
Proposition~\ref{prop:r6-continuous-joint-varying-background}, and
the frozen arbitrary-background potential quotient. The hypotheses
on the represented order intervals are essential to the deduction.

\begin{proof}
Choose representatives $u_i$ in the normalized continuous potential
families. The R5 quotient, with positive separately closed mixed tails,
shows that replacing any representative by another with the same
ambient curvature leaves the product unchanged. This proves that
\eqref{eq:r6-curvature-interval-joint-map} is well defined.

To check continuity, consider a convergent net of input tuples and
choose arbitrary normalized representatives $u_{i,j}$ for its curvature
slots. Every subnet of the resulting tuple of representatives has a
further subnet converging locally uniformly, simultaneously in its
finitely many slots, by compactness of their finite product. Let
$(u_1,\ldots,u_k)$ be such a limit. Its ambient curvatures are the
specified limits $A_i$, by local uniform convergence on original tests.
The backgrounds still converge to $S$ along this further subnet.
Proposition~\ref{prop:r6-continuous-joint-varying-background} gives
\[
 T_{S_j,k}(u_{1,j},\ldots,u_{k,j})
                  \longrightarrow T_{S,k}(u_1,\ldots,u_k).
\]
The right side is the value prescribed in
\eqref{eq:r6-curvature-interval-joint-map}, independently of this
choice of cluster representatives. Thus every subnet admits a further
subnet converging to the same output. This is precisely convergence
of the original output net: if an output neighborhood failed to contain
a tail, one could take a subnet lying in its complement and obtain
a contradiction. This proves joint continuity without metrizability,
a gauge, or a continuous choice of representatives.

The outputs are positive separately closed currents.
Corollary~\ref{cor:r6-all-background-degree-one-weak} therefore applies
to them and proves the final assertion. Its actual delta-test is fixed;
no convergence of varying delta-test presentations is asserted.
\end{proof}

\paragraph{Weights are not curvature classes.}
The theorem concerns differentiated slots. A weight cannot be passed
to the same quotient: the constants $0$ and $1$ have the same ambient
curvature, while their products with any nonzero background are
different. If continuous actual weights converge locally uniformly,
the weighted conclusion follows instead from
Proposition~\ref{prop:r6-continuous-joint-varying-background}.
Likewise, the theorem does not concern arbitrary weakly convergent
curvature currents outside the displayed compact represented intervals,
nor does it assert total variation convergence.

% END source: sections/curvature_interval_topology.tex


% BEGIN canonical R6 source: sections/compact_dominator_families.tex
\section{Compact families of continuous upper potentials}
\label{sec:r6-compact-dominator-families}

A curvature upper bound supplied by one continuous potential is a useful
sufficient condition, but it can be replaced by compactness of the upper
potentials themselves.
All families and normalizations in this section live on one common
ordinary domain $\Omega$.

Let $\mathscr G$ be a nonempty compact subset of
$C(\Omega)\cap\mathcal C_b(\Omega)$ for the topology of uniform
convergence on compact subsets. Define
\begin{align*}
 \mathscr P_{\mathscr G}^{\,c}
  &=\{u\in C(\Omega)\cap\mathcal C_b(\Omega):0\le u\le1,
                   \quad D[u]\le D[g]\text{ for some }g\in\mathscr G\},\\
 \mathscr A_{\mathscr G}&=\{D[u]:u\in\mathscr P_{\mathscr G}^{\,c}\}.
\end{align*}
The curvatures $D[u]$ are already nonnegative by membership in the
classical class. No convexity of $\mathscr G$ or of the displayed
unions is assumed.

\begin{proposition}[Compactness with varying upper potentials]
\label{prop:r6-compact-upper-potential-family}
The family $\mathscr P_{\mathscr G}^{\,c}$ is compact in local
uniform topology. Its curvature image $\mathscr A_{\mathscr G}$
is compact Hausdorff in the ordinary weak topology, and the curvature
map is a closed quotient map with the frozen pluriharmonic fibers.
For every compact $K_0\Subset\Omega$ and $\epsilon>0$, there is
a finite set $F_\epsilon\subset\Omega$ such that
\[
 \|u-v\|_{K_0}
     \le\max_{z\in F_\epsilon}|u(z)-v(z)|+2\epsilon
             \quad(u,v\in\mathscr P_{\mathscr G}^{\,c}).
\]
\end{proposition}

\noindent\textbf{Status: \texttt{DIRECT\_COROLLARY}} of positive
Shilov interpolation, with a larger compactness scope than one fixed
curvature dominator.

\begin{proof}
First, compactness of $\mathscr G$ gives uniform small oscillation
near each point. Choose a compact neighborhood $L$ of a given point
$x$. The restrictions of $\mathscr G$ to $L$ form a compact subset
of $C(L)$, so choose finitely many functions giving a uniform
$\epsilon/3$ net. Continuity of these finitely many functions gives
an ordinary neighborhood of $x$ inside $L$ on which all their
oscillations are less than $\epsilon/3$. Every $g\in\mathscr G$
then has oscillation less than $\epsilon$ on that neighborhood.

Cover $K_0$ by finitely many such neighborhoods and choose the
frozen norm boxes and finite semistable Shilov covers inside them.
For any $u$ choose an upper potential $g\in\mathscr G$. Both
$u$ and $g-u$ are continuous classical by
Lemma~\ref{lem:r6-continuous-positive-classical}, so positive
Shilov interpolation gives on each of the fixed charts
\[
                   0\le P_\lambda u-u\circ q
                         \le P_\lambda g-g\circ q<\epsilon.
\]
The covers and their finite set of Shilov image points depend only
on $\mathscr G,K_0,\epsilon$, not on either chosen potential.
Subtract the inequalities for $u,v$ and use the positive barycentric
coordinates. This proves the finite-evaluation estimate.

For a net $u_j$ choose an associated $g_j\in\mathscr G$. Compactness
of $[0,1]^\Omega\times\mathscr G$ gives a simultaneous subnet
with pointwise $u_j\to u$ and locally uniform $g_j\to g$.
The finite-evaluation estimate makes the same subnet uniformly Cauchy
on each compact set, hence locally uniformly convergent to $u$.
The actual limit is continuous classical, lies between zero and one,
and satisfies $0\le D[u]\le D[g]$ by original smooth testing.
Thus it belongs to the displayed family. This proves compactness
without a countable exhaustion. The quotient and fiber assertions
follow exactly as in Proposition
\ref{prop:r6-curvature-interval-quotient-topology} from the continuous
curvature map, Hausdorff target, and frozen R5 quotient.
\end{proof}

\begin{corollary}[Joint stability and actual weak-image closure]
\label{cor:r6-compact-upper-potential-stability}
For finitely many compact upper-potential families $\mathscr G_i$,
the operation
\[
 (S,A_1,\ldots,A_k)\longmapsto T_{S,k}(u_1,\ldots,u_k),
                  \qquad D[u_i]=A_i\in\mathscr A_{\mathscr G_i},
\]
is well defined and jointly continuous for ordinary weak convergence
of the positive separately closed background $S$ and all ambient
curvature slots. For fixed $S$ its image is compact and closed,
with actual normalized representatives on the whole $\Omega$.
The same is true for the diagonal image of one such family.
Consequently, in top degree this diagonal family has sharp compact
no-gain under weak closure. These conclusions do not assert general Borel outer
regularity, setwise convergence, or a common measure dominator.
\end{corollary}

\noindent\textbf{Status: \texttt{DIRECT\_COROLLARY}}.

\begin{proof}
For every subnet of representatives choose a further subnet converging
locally uniformly in the compact finite product of their potential
families. The qualified varying-background theorem identifies its
mixed-current limit, and the frozen quotient makes the answer
independent of the cluster representatives. The every-subnet argument
of Theorem~\ref{thm:r6-curvature-interval-joint-topology} proves
convergence of the whole net. For fixed $S$ the image of the compact
family is compact in the Hausdorff current space, hence closed.
In top degree a weak-limit diagonal measure is itself one of the
actual normalized testers on $\Omega$, so its mass on any compact set
is bounded by
the raw capacity. This uses actual representatives, not an invalid
compact-set Portmanteau inequality in the opposite direction.
\end{proof}

\begin{example}[Compact upper potentials without a common measure upper bound]
\label{ex:r6-moving-hinge-upper-family}
On the original annulus with $t=-\log|T|\in(-2,2)$, put
\[
              g_r(t)=\frac14|t-r|,\qquad -\frac12\le r\le\frac12.
\]
These are continuous classical functions, lie in $[0,1]$, and depend
uniformly continuously on $r$. Their family is compact even in the
uniform norm, and
\[
                         Dg_r=\frac12\delta_{\zeta_{e^{-r}}}
\]
with the original logarithmic calibration. Thus they are included
in the preceding compactness statement by choosing each potential
as its own upper potential.

There is no single locally finite positive Radon measure $\rho$ on
this annulus for which every $Dg_r$ is at most $C_r\rho$ with a finite
constant $C_r$. Otherwise $\rho$ would have a positive atom at every
one of the uncountably many displayed Gauss points in a fixed compact
subannulus. A finite measure has only countably many positive atoms:
for each integer $m$ only finitely many can have mass at least $1/m$.
In particular no single continuous psh potential has curvature
dominating this whole family, even with member-dependent finite
multipliers.

Also, if $r_j\to0$ with $r_j\ne0$, then $g_{r_j}\to g_0$ uniformly
and their measures converge weakly, whereas
\[
       Dg_{r_j}(\{\zeta_1\})=0,
              \qquad Dg_0(\{\zeta_1\})=\tfrac12.
\]
So the compact upper-potential theorem cannot be upgraded to setwise
convergence. It genuinely differs from the common-Radon-dominator
mechanism in Section~\ref{sec:r6-capacity-dominated-curvatures}.
\end{example}

\noindent\textbf{Status of the example: \texttt{DIRECT\_COROLLARY}}
of the original annular hinge calibration and elementary measure
countability. It limits this new sufficient criterion; it does not
supply a counterexample to general compact C1.

% END source: sections/compact_dominator_families.tex


% BEGIN canonical R6 source: sections/capacity_compact_dominator_outer_failure.tex
\section{Compact continuous dominators do not give Borel outer regularity}
\label{sec:r6-compact-dominator-outer-failure}

This example concerns the full normalized family with a curvature
dominated by one member of a specified compact continuous family.
It does not replace that family by just the displayed extremizers.
It realizes the frozen abstract vertical-fiber obstruction on an
actual smooth Berkovich torus, with its original tests and calibration.

\begin{theorem}[A restricted capacity with different Borel and outer values]
\label{thm:r6-compact-dominator-outer-failure}
Work over the original $K$, with its original absolute value. Put
\[
 \Omega=\{(T_1,T_2)\in(\mathbf G_{m,K}^{2})^{\rm an}:
                         -1<t_1,t_2<2\},
 \qquad t_i=\log|T_i|,
 \qquad S=3[d't_2\wedge d''t_2].
\]
Let
\[
 A=\{0\}\cup\{1/j:j\ge1\},\qquad
 g_s=\tfrac13\max\{t_1-s,0\}\quad(s\in A),
 \qquad \mathscr G=\{g_s:s\in A\}.
\]
The family $\mathscr G$ is compact in $C_{\rm loc}(\Omega)$ and
consists of continuous classical psh functions with $0\le g_s\le1$.
For the entire family
\[
 \mathscr P(\mathscr G)=
 \{u\in\mathcal C_b(\Omega):0\le u\le1,
                \ D[u]\le D[g_s]\text{ for some }s\in A\},
\]
define
\[
 C_{\mathscr G}(E)=\sup_{u\in\mathscr P(\mathscr G)}
                              D(uS)(E),\qquad
 C_{\mathscr G}^{*}(E)=
       \inf_{\substack{G\supset E\\G\text{ ordinary open in }\Omega}}
                              C_{\mathscr G}(G).
\]
There is an $F_\sigma$ subset $E$ of a fixed compact subset of
$\Omega$ such that
\[
                       C_{\mathscr G}(E)=\tfrac12,
                  \qquad C_{\mathscr G}^{*}(E)=1.
\]
Thus Borel outer regularity fails for this restricted capacity even
though the whole potential family $\mathscr P(\mathscr G)$ is
compact in the local uniform topology and its represented measure
image is weakly compact.
\end{theorem}

\noindent\textbf{Status: \texttt{NEW\_PROVED}} realization in
the actual normalized curvature-order family, using the frozen
abstract vertical-fiber example. No new general abstract capacity
counterexample or failure of unrestricted BT capacity is claimed.

\begin{proof}
The background $S$ is an original smooth, symmetric, positive,
separately closed $(1,1)$-current. The affine unit moment $t_1-s$
is pluriharmonic after every complete extension, so its maximum
with zero is classical psh. Since $t_1<2$ and $s\ge0$, one has
$0\le g_s<2/3$. Also
\[
                  \|g_s-g_r\|_{\infty,\Omega}
                           \le\tfrac13|s-r|.
\]
The parameter space $A$ is compact, hence $\mathscr G$ is compact
even for the global uniform norm. By
Proposition~\ref{prop:r6-compact-upper-potential-family} and
Corollary~\ref{cor:r6-continuity-from-continuous-curvature-domination},
the full displayed family $\mathscr P(\mathscr G)$ is compact in
$C_{\rm loc}$; it is not necessary to assume continuity in its
definition.

Let $\sigma:(-1,2)^2\to\Omega$ be the original two-variable
Gauss section, with radii $(e^{t_1},e^{t_2})$. Original torus
calibration gives the equality of Radon measures
\begin{equation}
                D(g_sS)=\nu_s
                    :=\sigma_*(\delta_s\otimes dy),
             \qquad -1<y<2.
 \label{eq:r6-vertical-hinge-measure}
\end{equation}
The slope jump of $g_s$ is $1/3$, and the factor $3$ in $S$
cancels it. There is no change of the field or its absolute value.
This is an identity on every original smooth compact scalar test:
on a common moment refinement, torus top integration is supported
on the original Gauss section and has degree one. Equivalently,
smooth convex approximations of the one-variable hinge give the
same product-calibrated measure in the limit. The identity is
therefore not a prescription of pairings only against functions
which factor through the two radius coordinates.

If $u\in\mathscr P(\mathscr G)$, choose its parameter $s$.
The arbitrary-background order theorem gives
\[
                       0\le D(uS)\le D(g_sS)=\nu_s.
\]
Conversely $g_s\in\mathscr P(\mathscr G)$ for every $s$.
It follows for every Borel $B\subset\Omega$ that
\begin{equation}
                  C_{\mathscr G}(B)=\sup_{s\in A}\nu_s(B).
 \label{eq:r6-full-family-vertical-capacity}
\end{equation}
This equality establishes the upper bound for all normalized
competitors, not only for the hinges. Local uniform BT continuity
also shows that the full represented measure image is weakly compact.

In the real radius box set
\[
 E_0=(\{0\}\times[0,\tfrac12])
            \ \cup\ \bigcup_{j\ge1}
                    (\{1/j\}\times(\tfrac12,1]),
                 \qquad E=\sigma(E_0).
\]
The Gauss map is a topological embedding, and
$\sigma(A\times[0,1])$ is compact in $\Omega$. Every half-open
vertical segment in this formula is a countable union of compact
segments. Thus $E$ is $F_\sigma$ in the actual Berkovich space.
Each $\nu_s$ assigns it mass $1/2$, so
\eqref{eq:r6-full-family-vertical-capacity} gives
$C_{\mathscr G}(E)=1/2$.

Let $G$ be any ordinary open set containing $E$. The inverse
image $O=\sigma^{-1}(G)$ is an ordinary real open set containing
$E_0$. Compactness of $\{0\}\times[0,1/2]$ gives a $\delta>0$
such that
\[
                (-\delta,\delta)\times[0,1/2]\subset O.
\]
For every sufficiently large $j$, this includes the lower half of
the vertical segment at $1/j$, while its upper half already belongs
to $E_0$. Hence $\{1/j\}\times[0,1]\subset O$ and
$\nu_{1/j}(G)\ge1$. The full-family identity proves
$C_{\mathscr G}(G)\ge1$ for every such ordinary open $G$.

For the opposite bound, take $0<\epsilon<1$ and the ordinary open
set
\[
                 G_\epsilon=\{-\epsilon<t_2<1+\epsilon\}
                                         \cap\Omega.
\]
It contains $E$, and every $\nu_s$ gives it mass $1+2\epsilon$.
Thus $C_{\mathscr G}(G_\epsilon)=1+2\epsilon$ by the same
identity. Taking the infimum and then $\epsilon\downarrow0$
proves $C_{\mathscr G}^{*}(E)=1$.
\end{proof}

\begin{corollary}[The restricted outer capacity fails continuity from below]
\label{cor:r6-compact-dominator-outer-from-below}
There are increasing compact sets $K_m\uparrow E$ in the preceding
example such that
\[
 C_{\mathscr G}^{*}(K_m)=C_{\mathscr G}(K_m)=\tfrac12
       \quad\text{for every }m,\qquad
 C_{\mathscr G}^{*}(E)=1.
\]
\end{corollary}

\begin{proof}
For $m\ge3$ let
\[
 K_m=\sigma\left((\{0\}\times[0,\tfrac12])
       \cup\bigcup_{j=1}^{m}
             (\{1/j\}\times[\tfrac12+1/m,1])\right).
\]
These are compact, increase, and have union $E$. Their raw
capacity is $1/2$ by \eqref{eq:r6-full-family-vertical-capacity}.
To see the outer value directly, separate the finitely many
parameters $1,\ldots,1/m$ by pairwise disjoint small real open
intervals, each containing no other point of $A$. Around $0$ use
one disjoint interval containing only $0$ and a tail of $A$.
Take their products with slightly enlarged versions of the respective
vertical segments, and pull back the resulting real open union by
$(t_1,t_2)$. For any $\epsilon>0$ the enlargements can be chosen
so that every $\nu_s$ gives this ordinary open neighborhood of
$K_m$ mass at most $1/2+\epsilon$. The outer value is consequently
at most $1/2$, and the reverse inequality follows from monotonicity.
The theorem supplies the strict increase at the union.
\end{proof}

\paragraph{What the example excludes.}
Compactness of an actual normalized potential family, compactness of
its weak measure image, and sharp compact-set outer regularity do
not imply Borel outer regularity for its upper-envelope capacity.
This example does not contradict the fixed-Radon domination theorem:
the measures $\nu_s$ charge distinct vertical fibers, and no locally
finite Radon measure can dominate all of them with one fixed density
bound on this compact rectangle. Nor is this a counterexample for
the unrestricted capacity formed from every normalized classical psh
function on $\Omega$. Additional unrestricted competitors may change
both displayed capacity values; no value for that larger capacity
is asserted here.

% END source: sections/capacity_compact_dominator_outer_failure.tex


% BEGIN canonical R6 source: sections/capacity_compact_dominator_all_degrees.tex
\section{The restricted outer-regularity failure in every positive relative degree}
\label{sec:r6-compact-dominator-all-degrees}

\noindent\textbf{Status: \texttt{DIRECT\_COROLLARY}} of
Theorem~\ref{thm:r6-compact-dominator-outer-failure}, mixed order
replacement, and original product calibration. This is a dimensional
propagation of the same restricted-family obstruction.

\begin{proposition}
\label{prop:r6-compact-dominator-all-degrees}
For every pair $p,q\ge1$ there is an allowed smooth background of
bidegree $(p,p)$ on a smooth $(p+q)$-dimensional original-$K$ torus
domain, and a compact continuous normalized upper-potential family,
whose restricted relative degree-$q$ capacity satisfies
\[
                            C(E)=\tfrac12,
                      \qquad C^*(E)=1
\]
on one $F_\sigma$ set $E$ contained in a fixed compact subset.
The family defining $C$ includes every normalized actual classical
potential whose ambient curvature is bounded by a member of the
specified upper family.
\end{proposition}

\begin{proof}
Use coordinates $T_1,\ldots,T_q,Z_1,\ldots,Z_p$ and put
$x_i=\log|T_i|$, $y_a=\log|Z_a|$. Let
\[
 \Omega_{p,q}=\{-1<x_i,y_a<2\}\subset
                  (\mathbf G_{m,K}^{p+q})^{\rm an},
 \qquad A=\{0\}\cup\{1/j:j\ge1\}.
\]
Take
\[
 g_s=\frac{(x_1-s)_++\sum_{i=2}^{q}(x_i)_+}{3q},
 \qquad
 S=\frac{(3q)^q}{q!}
              \left[\bigwedge_{a=1}^{p}d'y_a\wedge d''y_a\right].
\]
An empty hinge sum for $q=1$ is zero. The background is smooth,
symmetric, strongly positive and separately closed. Every $g_s$
is continuous classical psh, $0\le g_s<2/3$, and
\[
                  \|g_s-g_r\|_\infty\le\frac{|s-r|}{3q}.
\]
Thus $\mathscr G=\{g_s:s\in A\}$ is a compact normalized
continuous family on this same whole domain.

Let $\sigma$ be the original $(p+q)$-variable Gauss section. The
exact ordinary top measure is
\begin{equation}
 T_{S,q}(g_s,\ldots,g_s)
       =\nu_s:=\sigma_*
           (\delta_s\otimes\delta_0^{\otimes(q-1)}
                                      \otimes dy_1\cdots dy_p).
 \label{eq:r6-all-degree-vertical-measure}
\end{equation}
To check the coefficient, each one-variable hinge curvature has
calibration one and squares to zero. In the $q$th power of their
sum, only terms using all $q$ distinct coordinate curvatures survive;
there are $q!$ such terms. The remaining scalar is
\[
                         \frac{(3q)^q}{q!}
                              \frac{q!}{(3q)^q}=1.
\]
Wedging with the $p$ smooth passive-coordinate top factors gives
their product Lebesgue measure. This verifies the entire product
calibration, rather than multiplying $p$ or $q$ independent field
degrees. As in the two-dimensional case, smooth convex approximants
of each hinge and the original torus top-form measure identify
\eqref{eq:r6-all-degree-vertical-measure} on all original compact
scalar tests, including further moment refinements.

Define the whole competitor family by
\[
 \mathscr P(\mathscr G)=\{u\in\mathcal C_b(\Omega_{p,q}):
       0\le u\le1,\ D[u]\le D[g_s]\text{ for some }s\in A\}.
\]
For a competitor with parameter $s$, replace the differentiated
slots one at a time using
Corollary~\ref{cor:r6-arbitrary-background-potential-order}.
Each remaining mixed tail is a fixed positive separately closed
background, so this yields
\[
       0\le T_{S,q}(u,\ldots,u)
                          \le T_{S,q}(g_s,\ldots,g_s)=\nu_s.
\]
Each $g_s$ is itself a competitor. Consequently the restricted
capacity is exactly $C(B)=\sup_{s\in A}\nu_s(B)$ for every
Borel $B$, with no restriction to just the displayed hinge testers.

Write $Q_0=[0,1]^{p-1}$, interpreted as a point if $p=1$, and set
\[
 \begin{split}
 E_0={}&\{(0,0,\ldots,0)\}\times([0,\tfrac12]\times Q_0)\\
 &{}\cup\bigcup_{j\ge1}
       \{(1/j,0,\ldots,0)\}\times((\tfrac12,1]\times Q_0),
 \qquad E=\sigma(E_0).
 \end{split}
\]
It is $F_\sigma$ and lies in the compact Gauss image of
$A\times\{0\}^{q-1}\times[0,1]^p$. Every $\nu_s$ assigns
$E$ mass $1/2$, proving the raw value.

For any ordinary open $G\supset E$, its inverse image under
$\sigma$ contains the compact lower plate
$\{0\}^q\times[0,1/2]\times Q_0$. A uniform small displacement
in the $x_1$-direction stays in this real open set. Therefore, for
all sufficiently large $j$, it contains the lower plate over
$x_1=1/j$. The upper plate there already belongs to $E_0$.
The entire passive cube $[0,1]^p$ at this parameter is contained
in $\sigma^{-1}(G)$, so $\nu_{1/j}(G)\ge1$ and $C^*(E)\ge1$.
Conversely the ordinary open sets
\[
        G_\epsilon=\{-\epsilon<y_a<1+\epsilon
                          \text{ for }1\le a\le p\}
                                  \cap\Omega_{p,q},
       \qquad 0<\epsilon<1,
\]
contain $E$ and have $\nu_s$-mass $(1+2\epsilon)^p$ for every
$s$. Let $\epsilon\downarrow0$ to obtain $C^*(E)\le1$.
This proves the assertion for all $p,q\ge1$, including the requested
$p=1$, $n=q+1$ case.
\end{proof}

The ambient background in this construction is smooth and belongs
to the full delta approximation class. The failure is caused by
the restricted normalized family, not by a lack of delta tests or
by singularity of the background. It neither disproves compact-set
outer regularity for that compact represented family nor determines
the unrestricted relative Bedford--Taylor capacity of the same set.

% END source: sections/capacity_compact_dominator_all_degrees.tex


% BEGIN canonical R6 source: sections/local_pluriharmonic_gauge.tex
\section{A local finite-point normalization of the potential quotient}
\label{sec:r6-local-ph-gauge}

A differentiated potential quotient does not itself determine a
multiplier. On the compact represented families constructed above,
one can nevertheless make an explicit local choice of representatives.
The choice depends on a fixed affinoid and finitely many normalization
points; it is not a canonical global gauge for all psh functions.

Fix a compact affinoid norm box $V\Subset\Omega$ with nonempty
ordinary interior, and let $\Gamma(V)$ be its finite Shilov boundary.
Let $\mathscr H_V$ be the vector space of restrictions to $V$ of
continuous pluriharmonic functions defined on some ordinary neighborhood
of $V$ inside $\Omega$. The neighborhood may depend on the function.
Finite sums are defined on intersections of these neighborhoods, so
this is a vector space of actual continuous functions on $V$.

\begin{lemma}[Finite Shilov normalization]
\label{lem:r6-finite-shilov-ph-normalization}
The space $\mathscr H_V$ has dimension $d\le\#\Gamma(V)$.
There are points $\xi_1,\ldots,\xi_d\in\Gamma(V)$, a basis
$h_1,\ldots,h_d$ of $\mathscr H_V$ represented by pluriharmonic
functions on one fixed ordinary neighborhood $W$ of $V$, and a
bounded linear projection
\[
        L_V:C(V)\longrightarrow\mathscr H_V,
       \qquad (L_Vf)(\xi_j)=f(\xi_j)\quad(1\le j\le d).
\]
The projection is the identity on $\mathscr H_V$. Its construction
uses actual values at the displayed points.
\end{lemma}

\noindent\textbf{Status: \texttt{DIRECT\_COROLLARY}} of the
classical Shilov maximum principle and finite-dimensional linear
algebra.

\begin{proof}
Apply the affinoid maximum principle to both signs of any representing
pluriharmonic function, as in \cite[Theorem 6.4]{GR25}. Its restriction
to $V$ is determined by its values on $\Gamma(V)$: zero values there
force the function to vanish on all of $V$. Evaluation therefore embeds
$\mathscr H_V$ in $\mathbf R^{\Gamma(V)}$, proving the dimension
bound. Choose a basis, intersect its finitely many defining ordinary
neighborhoods, and call the intersection $W$.

The evaluation matrix of this basis on $\Gamma(V)$ has rank $d$.
Choose $d$ rows with invertible matrix $M=(h_i(\xi_j))_{j,i}$ and set
\[
    L_Vf=\sum_{i=1}^d c_i h_i,
                   \qquad (c_i)_i=M^{-1}(f(\xi_j))_j.
\]
This is a bounded linear map in the uniform norm, interpolates the
selected values, and reproduces every member of $\mathscr H_V$.
In the vacuous zero-dimensional vector-space case use the zero map.
\end{proof}

\begin{theorem}[Continuous local representatives from weak curvature data]
\label{thm:r6-local-continuous-potential-section}
Let $\mathscr P$ be any compact family in
$C(\Omega)\cap\mathcal C_b(\Omega)$ for local uniform topology,
with $0\le u\le1$ for all its members. Fix $V,W,L_V$ as above,
shrinking $W$ inside $\Omega$ if necessary, and give
\[
             \mathscr A_W=\{D[u]|_W:u\in\mathscr P\}
\]
the ordinary weak topology on every original compact test in $W$.
For $A=D[u]|_W$ put
\begin{equation}
                    s_V(A)=u|_V-L_V(u|_V).
 \label{eq:r6-local-potential-section}
\end{equation}
This is independent of the chosen $u\in\mathscr P$, and defines
a continuous map $s_V:\mathscr A_W\to C(V)$. Its values are actual
continuous classical potentials on $\operatorname{Int}V$, with
\[
 D[s_V(A)]=A|_{\operatorname{Int}V},\qquad
 s_V(A)(\xi_j)=0,\qquad
 \|s_V(A)\|_V\le1+\|L_V\|.
\]
Here the curvature equality refers to the unique restriction to the
ordinary interior. Each representative is also defined classically
on a neighborhood of $V$, although no common choice outside $V$ is
part of the section. In particular, weak convergence $A_j\to A$
within $\mathscr A_W$ implies uniform convergence of these specified
representatives on $V$.
\end{theorem}

\noindent\textbf{Status: \texttt{DIRECT\_COROLLARY}} of the
frozen actual potential quotient and compactness, with the local
normalization explicitly constructed.

\begin{proof}
If $D[u]|_W=D[v]|_W$, the frozen R5 quotient makes $u-v$ a continuous
pluriharmonic function on $W$. Its restriction belongs to $\mathscr H_V$,
and hence $L_V(u-v)=u-v$ on $V$. Subtracting proves independence in
\eqref{eq:r6-local-potential-section}.

The curvature map from $\mathscr P$ onto $\mathscr A_W$ is continuous,
as seen on each original compact test. Its domain is compact and its
target is Hausdorff, so it is a closed quotient map. The map
$u\mapsto u|_V-L_V(u|_V)$ is continuous into $C(V)$ and constant on
its fibers. The defining property of a quotient topology therefore
gives continuity of $s_V$. This argument uses no metrizability or
choice of convergent lifts.

For any lift $u$, subtract the fixed finite linear combination of the
pluriharmonic basis functions on $W$. This produces an actual continuous
classical potential near $V$, has unchanged curvature, and restricts
to $s_V(A)$. Reproduction and the interpolation identities give the
normalizations; the operator norm gives the uniform bound. Different
lifts give the same actual function on $V$, hence the same potential
on its ordinary interior. Nothing is asserted about their values
outside $V$.
\end{proof}

\begin{corollary}[Weighted stability after the specified local normalization]
\label{cor:r6-gauged-weighted-weak-stability}
For $0\le i\le k$, let $A_{i,j}\to A_i$ be nets in the restricted
curvature images of compact normalized families as in the theorem,
and put $u_{i,j}=s_V(A_{i,j})$, $u_i=s_V(A_i)$. Let $S_j\to S$
be arbitrary allowed positive separately closed backgrounds on
$\operatorname{Int}V$. Then
\[
 u_{0,j}T_{S_j,k}(u_{1,j},\ldots,u_{k,j})
             \longrightarrow u_0T_{S,k}(u_1,\ldots,u_k)
\]
on all original compact complementary smooth tests in that interior.
The same conclusion holds against a fixed actual degree-one delta-test
when it is complementary to the displayed mixed currents.
\end{corollary}

\noindent\textbf{Status: \texttt{DIRECT\_COROLLARY}}.

\begin{proof}
All displayed potentials converge uniformly on $V$ by the theorem.
Their restrictions are continuous classical on its interior.
Proposition~\ref{prop:r6-continuous-joint-varying-background} proves
the original-test assertion. For the degree-one delta assertion, insert
its fixed continuous-coefficient Hessian presentation, expand its fixed
continuous classical-difference slots, and use the same proposition
term by term with the locally uniformly convergent continuous weight.
For its PS part the fixed PWS test approximation has the uniform normal
coefficient bound, and multiplication by the uniformly bounded weights
preserves that bound. A compact cutoff converts each top-degree current
convergence into the desired pairing.
\end{proof}

\paragraph{Scope of the normalization.}
This section does not identify arbitrary actual multipliers with their
curvatures. For example $0$ and $1$ have equal curvature but are distinct
multipliers; the specified normalization sends both to zero. Its weighted
conclusion applies only after this explicit local choice. The compact
families of Theorem~\ref{thm:r6-curvature-order-compactness} and
Proposition~\ref{prop:r6-compact-upper-potential-family} are concrete
inputs. There is no continuous-section claim for the unrestricted weak
space of bounded classical curvatures, and no global compatibility of
different choices of $V$, normalization points, or projections is asserted.

% END source: sections/local_pluriharmonic_gauge.tex


\part{Sharp relative capacity estimates}

% BEGIN canonical R6 source: sections/dirichlet_sharp_hardy.tex
\section{A sharp Hardy estimate for every normalized relative tester}
\label{sec:r6-sharp-relative-hardy}

Fix an arbitrary original-positive symmetric separately closed
$(n-1,n-1)$-current $S$ on an allowed ordinary open $\Omega$.
For original real compact smooth $f$ write
\[
                    \mathcal E_S(f)=\int_\Omega\Gamma_S(f).
\]
The first estimate below needs only the frozen bounded BT calculus and
its local coefficient-Radon energy measures. The faithful closed form
is used subsequently for the capacity and Dirichlet consequences.

\begin{lemma}[Polynomially weighted local energy limits]
\label{lem:r6-polynomial-weighted-energy-limit}
On a fixed local compact regularization stage let
$a_{i,j}\downarrow a_i$ be finitely many bounded classical psh rows.
Let $P,Q,R$ be real polynomials in these functions and in finitely many
fixed smooth local classical differences. Put $P_j=P(a_{1,j},\ldots)$,
and similarly for $Q_j,R_j$. Then
\[
            R_j\,\Gamma_S(P_j,Q_j)
                    \longrightarrow R\,\Gamma_S(P,Q)
\]
on original compact smooth scalar tests. All functions and measures
use their actual representatives.
\end{lemma}

\noindent\textbf{Status: \texttt{DIRECT\_COROLLARY}} of the weighted
one-slot theorem and the definition of $\Gamma_S$.

\begin{proof}
After shifting the finitely many psh functions by constants on the
compact stage, they are nonnegative. Every polynomial in them is a
finite real linear combination of powers of their positive sums, by
the ordinary polarization formula for monomials. Each such power is
classical psh and decreases along the given rows. Fixed smooth local
classical differences have the same description after choosing their
finite psh decompositions; these functions are constant rows.
Thus every polynomial belongs to the local bounded classical-difference
algebra. Expand
\[
 R_j\Gamma_S(P_j,Q_j)
  =\tfrac12\{R_jD(P_jQ_jS)-R_jP_jD(Q_jS)-R_jQ_jD(P_jS)\}.
\]
Every term is a finite linear combination of weighted one-slot products
of decreasing bounded psh rows. The frozen weighted theorem passes
each term. The limiting identity is the Borel multiplier of the
coefficient-Radon measure defining $\Gamma_S(P,Q)$.
\end{proof}

\begin{theorem}[Sharp normalized-tester Hardy inequality]
\label{thm:r6-sharp-relative-hardy}
Let $w\in\mathcal C_b(\Omega)$ satisfy $0\le w\le1$ everywhere.
For every original real $f\in A_c^{0,0}(\Omega)$ and every
$\varepsilon>0$,
\begin{equation}
 \int_\Omega\frac{f^2}{1+\varepsilon-w}\,D(wS)
                                      \le\mathcal E_S(f).
 \label{eq:r6-relative-picone}
\end{equation}
Consequently
\begin{equation}
                       \int_\Omega f^2\,D(wS)
                                      \le\mathcal E_S(f).
 \label{eq:r6-sharp-hardy}
\end{equation}
The coefficient one in \eqref{eq:r6-sharp-hardy} cannot be decreased
uniformly over all allowed domains and backgrounds.
\end{theorem}

\noindent\textbf{Status: \texttt{NEW\_PROVED}} for the sharp uniform
estimate. It strengthens the earlier square-root estimate; that weaker
estimate is not used in its proof.

\begin{proof}
First suppose $w$ is smooth near the support of $f$. Put
$v=1+\varepsilon-w>0$. The positive covector square gives
\[
 \Gamma_S(f)-\Gamma_S(v,f^2/v)
  =\Gamma_S(f)-2(f/v)\Gamma_S(f,v)+(f^2/v^2)\Gamma_S(v)\ge0.
\]
Compact Green and $D(vS)=-D(wS)$ therefore yield
\[
 \mathcal E_S(f)\ge\int\Gamma_S(v,f^2/v)
                         =\int(f^2/v)D(wS).
\]

For bounded possibly discontinuous $w$, a global smooth row near the
entire support is not assumed. Choose finitely many nonnegative
original smooth functions $\chi_i$, supported in local compact
regularization stages, with $\sum_i\chi_i=1$ on a neighborhood of
$\operatorname{supp}f$. On the stage for $\chi_i$, choose a decreasing
original smooth classical row $w_j\downarrow w$. After discarding a
prefix, $0\le w_j<1+\varepsilon/2$ on a fixed neighborhood of the
relevant compact support. Indeed the increasing open sublevel sets
below $1+\varepsilon/2$ cover that compact set. Set
\[
       v_j=1+\varepsilon-w_j,\qquad r_j=1/v_j.
\]
The $r_j$ form a bounded decreasing classical psh row on that fixed
neighborhood: the scalar reciprocal function of $w_j$ is convex and
increasing on its range. The limit is $r=(1+\varepsilon-w)^{-1}$.

Multiply the smooth positive-square inequality by $\chi_i$ before
integrating. The localized Green identity is
\[
 \int\chi_i\Gamma_S(v_j,f^2r_j)
    =\int\chi_i f^2r_jD(w_jS)
                              +\int f^2r_j\Gamma_S(w_j,\chi_i).
\]
The first term passes by weighted one-slot continuity for the two rows
$r_j,w_j$. The second passes by
Lemma~\ref{lem:r6-polynomial-weighted-energy-limit}, treating $r_j$ as
an additional psh row. Fixed original smooth $f$ and $\chi_i$ are local
classical differences by the frozen smooth compensator construction.
Consequently
\[
 \int\chi_i\Gamma_S(f)
   \ge\int\chi_i f^2rD(wS)+\int f^2r\Gamma_S(w,\chi_i).
\]
Sum these finitely many inequalities. Their cross terms cancel
exactly: by bilinearity and ordinary locality,
$\Gamma_S(w,\sum_i\chi_i)=0$ near the support of $f$.
Also $\Gamma_S(f)$ is supported there. This proves
\eqref{eq:r6-relative-picone} without a partition error or a global
regularization assumption. Since $1+\varepsilon-w\le1+\varepsilon$,
it gives
$\int f^2D(wS)\le(1+\varepsilon)\mathcal E_S(f)$.
Let $\varepsilon\downarrow0$ to obtain \eqref{eq:r6-sharp-hardy}.

For sharpness let $\lambda=-\log|\pi|>0$, take the original annulus
$|\pi|<|T|<|\pi|^{-1}$, put $s=\log|T|/\lambda\in(-1,1)$,
and set $S=[1]$, $w=|s|$. Then
\[
                 D(wS)=(2/\lambda)\delta_{\zeta_{0,1}},
       \qquad \mathcal E_S(f(s))
                   =\lambda^{-1}\int_{-1}^1|f'(s)|^2\,ds.
\]
Choose smooth compact functions on $(-1,1)$ converging to
$1-|s|$ in the usual $H^1_0$ norm and with values at zero tending
to one; smoothing slightly truncated tents gives such a sequence.
The numerator in \eqref{eq:r6-sharp-hardy} tends to $2/\lambda$,
as does the energy. All functions used as tests are original smooth
functions of the unit moment. Thus a universal smaller coefficient
is impossible.
\end{proof}

\begin{corollary}[Weak limits and the energy-capacity bound]
\label{cor:r6-sharp-energy-capacity}
Assume the faithful form exists on $F=\operatorname{supp}S$ as in R4.
Every positive weak limit of normalized measures $D(w_jS)$ on all
original tests satisfies \eqref{eq:r6-sharp-hardy}. This includes the
normalized completed image defined by global normalized classical
rows. Moreover, for every subset $A\subset\Omega$,
\begin{equation}
            \operatorname{Cap}_S^*(A;\Omega)
                  \le\operatorname{cap}_{\mathcal E,1}(A\cap F).
 \label{eq:r6-linear-energy-capacity}
\end{equation}
For original smooth compact $f$ and $t>0$ there is the sharper direct
level-set bound
\begin{equation}
        \operatorname{Cap}_S^*(\{|f|>t\};\Omega)
                          \le\mathcal E_S(f)/t^2.
 \label{eq:r6-sharp-capacity-chebyshev}
\end{equation}
These are one-way bounds. Neither the reverse capacity inequality nor
the general raw-versus-completed compact equality is asserted.
\end{corollary}

\noindent\textbf{Status: \texttt{DIRECT\_COROLLARY}} of the sharp estimate
and the already verified regular-form compact-core capacity formula.

\begin{proof}
For weak limits, $f^2$ is one fixed original compact smooth scalar
test, so its inequality passes directly, also for nets. For a compact
$K\subset F$ and original smooth core $f>1$ on $K$, every normalized
tester satisfies
\[
 D(wS)(K)\le\int f^2D(wS)\le\mathcal E_S(f)
                                      \le\mathcal E_1(f,f).
\]
The verified compact-core formula gives
$D(wS)(K)\le\operatorname{cap}_{\mathcal E,1}(K)$, with the same
bound for all testers. For a relative ordinary open $G\subset F$,
inner regularity of each fixed Radon measure therefore gives
$D(wS)(G)\le\operatorname{cap}_{\mathcal E,1}(G)$.
Choose an ordinary open $O\subset\Omega$ with $O\cap F=G$.
All these measures are supported on $F$. If $G$ contains $A\cap F$,
then $O\cup(\Omega\setminus F)$ contains $A$, and taking the
supremum over testers bounds its raw BT capacity by the energy
capacity of $G$. The two outer infima give
\eqref{eq:r6-linear-energy-capacity}. Only a uniform compact estimate
was used before Radon inner regularity; no infimum/supremum exchange
has been made.

Finally, Chebyshev and \eqref{eq:r6-sharp-hardy} give
$D(wS)\{|f|>t\}\le\mathcal E_S(f)/t^2$ for every normalized
tester. The level set is an ordinary open, so its raw capacity equals
its outer capacity by definition. Taking the supremum proves
\eqref{eq:r6-sharp-capacity-chebyshev}.
\end{proof}

\begin{corollary}[Linear capacity domination is enough for finite energy]
\label{cor:r6-linear-capacity-energy-data}
Let $\mu$ be a positive Radon measure supported on a compact subset of
$V\Subset\Omega$, with mass $M>0$, and suppose
\[
              \mu(E)\le A\operatorname{Cap}_S(E;\Omega)
\]
for every Borel subset of that support. Then
\[
                      \|\mu\|_{-1,V}\le2\sqrt{AM}.
\]
It has the unique energy solution of
Theorem~\ref{thm:r6-measure-energy-dirichlet}. More generally, if
$\mu(E)\le A\operatorname{Cap}_S(E;\Omega)^\alpha$ with
$\alpha>1/2$, then
\[
 \|\mu\|_{-1,V}
     \le\frac{2\alpha}{2\alpha-1}
             A^{1/(2\alpha)}M^{1-1/(2\alpha)}.
\]
No boundedness or ambient psh property of that solution follows from
these assertions alone.
\end{corollary}

\noindent\textbf{Status: \texttt{DIRECT\_COROLLARY}} by layer cake.

\begin{proof}
For a core function of energy $e^2$, the preceding level-set bound
gives
\[
              \mu\{|f|>t\}\le\min\{M,Ae^{2\alpha}t^{-2\alpha}\}.
\]
Split the integral of this tail at
$t_0=(A/M)^{1/(2\alpha)}e$. Since $2\alpha>1$, its two pieces sum
to the stated constant times $e$. The case $\alpha=1$ gives
$2\sqrt{AM}$. A zero-energy core function has zero integral by the
same level-set bound. Taking the core supremum completes the proof.
\end{proof}

\paragraph{Frozen higher-degree tails.}
For an arbitrary allowed $(p,p)$-background $S_0$ and
$q=n-p\ge1$, fix bounded classical potentials $v_2,\ldots,v_q$ and
put $R=T_{S_0,q-1}(v_2,\ldots,v_q)$. The theorem applies verbatim
to this fixed degree-one background. Its capacity is
$\operatorname{Cap}_R(E;\Omega)=\sup_{0\le w\le1}D(wR)(E)$ and
its energy is $\mathcal E_R(f)=\int\Gamma_R(f)$. No single quadratic
form is thereby asserted to control the nonlinear diagonal family
$T_{S_0,q}(w^q)$ when all $q$ slots vary.

% END source: sections/dirichlet_sharp_hardy.tex


% BEGIN canonical R6 source: sections/dirichlet_strong_capacity_trace.tex
\section{A sharp strong-capacity inequality and energy traces}
\label{sec:r6-strong-capacity-trace}

Retain the arbitrary fixed degree-one background $S$ and the notation
of Section~\ref{sec:r6-sharp-relative-hardy}. Capacities below use
normalized actual classical testers on the entire fixed $\Omega$.
All smooth functions and tests are original CLD objects.

\begin{theorem}[Strong capacitary inequality with optimal constant]
\label{thm:r6-strong-capacity-four}
For every real original $f\in A_c^{0,0}(\Omega)$,
\begin{equation}
 \int_0^\infty
        \operatorname{Cap}_S(\{|f|>t\};\Omega)\,2t\,dt
                              \le4\mathcal E_S(f).
 \label{eq:r6-strong-capacity-four}
\end{equation}
The level sets are ordinary open, so the same formula holds with outer
capacity. If a positive Radon measure $\mu$ satisfies
\begin{equation}
                  \mu(E)\le A\operatorname{Cap}_S(E;\Omega)
 \label{eq:r6-linear-capacity-domination}
\end{equation}
for every Borel $E\subset\Omega$, then
\begin{equation}
                        \int f^2\,d\mu\le4A\mathcal E_S(f).
 \label{eq:r6-capacity-energy-trace}
\end{equation}
Both universal constants four are optimal, already on one original
analytic annulus with $S=[1]$.
\end{theorem}

\noindent\textbf{Status: \texttt{NEW\_PROVED}} for the sharp capacitary
and trace estimates. The proof uses only the preceding Hardy theorem
and explicit scalar truncations; no raw-capacity outer regularity is
assumed.

\begin{proof}
Fix $t>0$ and put $g_t=\min(f^2/t^2,1)$. To apply the smooth-core
Hardy inequality, choose smooth functions $P_\delta$ on the real line
such that $P_\delta(0)=0$, $0\le P_\delta'\le1$,
$P_\delta(s)\to\min(s,1)$ for $s\ge0$, and
$P_\delta'(s)\to0$ for $s>1$. They may be obtained by integrating
a smooth transition from derivative one below $1-\delta$ to zero
above $1+\delta$, chosen symmetrically so the final constant is one.
For $\delta<1/2$ the functions are exactly $s$ near zero and satisfy
$0\le P_\delta(s)\le1$ for $s\ge0$.

Then $P_\delta(f^2/t^2)$ is an original smooth compact function.
For each fixed normalized tester $w$, Hardy and the ordinary chain
rule give
\[
 \int P_\delta(f^2/t^2)^2D(wS)
       \le\frac4{t^4}\int
              f^2P_\delta'(f^2/t^2)^2\,\Gamma_S(f).
\]
On the fixed compact support both measures are finite. Bounded
convergence on the left and dominated upper limits on the right yield
\[
 D(wS)\{|f|>t\}
    \le\int g_t^2D(wS)
    \le\frac4{t^4}\int_{\{|f|\le t\}}f^2\,\Gamma_S(f).
\]
The inclusion of the equality level makes this argument independent
of any assertion about energy mass on a level set. The right side is
independent of $w$, so the same bound holds for the raw capacity.

Multiply by $2t$ and integrate. Tonelli's theorem for the nonnegative
measure $\Gamma_S(f)$ gives
\[
 \begin{split}
 \int_0^\infty\operatorname{Cap}_S\{|f|>t\}\,2t\,dt
 &\le8\int f^2
          \left(\int_{|f|}^{\infty}t^{-3}\,dt\right)\Gamma_S(f)\\
 &=4\int_{\{f\ne0\}}\Gamma_S(f)
 \le4\mathcal E_S(f).
 \end{split}
\]
At $f=0$, the inner-integrand contribution is interpreted before
integration as identically zero; no $0\cdot\infty$ expression is used.
The capacity of these level sets equals their outer capacity because
they are ordinary open. Finally layer cake and
\eqref{eq:r6-linear-capacity-domination} imply
\[
 \int f^2\,d\mu
       =\int_0^\infty\mu\{|f|>t\}\,2t\,dt
       \le A\int_0^\infty\operatorname{Cap}_S\{|f|>t\}\,2t\,dt,
\]
which proves \eqref{eq:r6-capacity-energy-trace}.
\end{proof}

\begin{example}[Sharpness on an original annulus]
\label{ex:r6-four-is-sharp}
Put $\lambda=-\log|\pi|>0$ and take
\[
       \Omega=\{1<|T|<|\pi|^{-1}\},\qquad
       s=\log|T|/\lambda\in(0,1),\qquad S=[1].
\]
Let $\sigma_s$ be the Gauss point of radius $\exp(\lambda s)$.
For $0<\varepsilon<1/4$ define the finite compactly supported measure
\[
 \mu_\varepsilon
    =\lambda^{-1}\sigma_*
                \bigl(\mathbf1_{[\varepsilon,1/2]}s^{-2}ds\bigr).
\]
It satisfies $\mu_\varepsilon(E)\le\operatorname{Cap}_{[1]}(E;\Omega)$
for every Borel $E\subset\Omega$.

Indeed, if $E_s=\{s\in[\varepsilon,1/2]:\sigma_s\in E\}$ is nonempty,
put $a=\inf E_s>0$. Then
\[
        \mu_\varepsilon(E)\le\frac1\lambda\int_a^\infty s^{-2}ds
                                      =\frac1{\lambda a}.
\]
For each $r\in E_s$ the globally normalized classical hinge
\[
                     w_r(s)=\max(0,1-s/r)
\]
has original curvature $D w_r=(\lambda r)^{-1}\delta_{\sigma_r}$.
Thus the capacity is at least
$\sup_{r\in E_s}(\lambda r)^{-1}=(\lambda a)^{-1}$.
If $E_s$ is empty the assertion is immediate. No formula for the
capacity of a general set is being assumed.

Choose fixed smooth scalar cutoffs $\alpha,\beta$ with
$\alpha=0$ on $(-\infty,1]$, $\alpha=1$ on $[2,\infty)$,
$\beta=1$ on $(0,1/2]$, and $\beta=0$ on $[3/4,1)$.
Set
\[
                f_\varepsilon(s)
                       =\alpha(s/\varepsilon)\beta(s)\sqrt{s}.
\]
These are original smooth compact tests. Direct integration gives
\[
 \int f_\varepsilon^2\,d\mu_\varepsilon
                   =\lambda^{-1}\log(1/\varepsilon)+O(1),
 \qquad
 \mathcal E_S(f_\varepsilon)
                   =(4\lambda)^{-1}\log(1/\varepsilon)+O(1).
\]
For the energy, the derivative on $[2\varepsilon,1/2]$ is
$1/(2\sqrt{s})$. Its two cutoff collars have uniformly bounded
energy: the lower collar has length $O(\varepsilon)$ and derivative
$O(\varepsilon^{-1/2})$, and the upper collar is fixed. The mass
calculation has the same bounded collar errors. Hence the trace ratio
tends to four. Since each $\mu_\varepsilon$ is dominated by the raw
capacity with coefficient one, its layer-cake integral also gives a
lower bound for the left side of \eqref{eq:r6-strong-capacity-four}.
Both constants are therefore optimal. The field and all tests remain
the original ones, including Gauss points of non-strict radii.
\end{example}

\begin{corollary}[Canonical energy traces]
\label{cor:r6-capacity-energy-trace-map}
Assume the faithful form exists and $\mu$ satisfies
\eqref{eq:r6-linear-capacity-domination}. The original smooth restriction
map extends uniquely to a bounded linear map
\[
             J_\mu:\mathscr D\longrightarrow L^2(\mu),
           \qquad \|J_\mu u\|_{L^2(\mu)}
                          \le2\sqrt A\,\mathcal E(u,u)^{1/2},
\]
It is represented by the canonical relative
capacitary representative of $u$, $\mu$-almost everywhere. The same
assertion holds on every part space $\mathscr D_0(V)$.
\end{corollary}

\noindent\textbf{Status: \texttt{DIRECT\_COROLLARY}} of the trace bound
and the frozen canonical energy representative.

\begin{proof}
The estimate makes any energy-Cauchy core sequence Cauchy in
$L^2(\mu)$, proving existence, independence and the norm bound.
On a fixed compact stage, the same core sequence converges locally in
outer relative capacity to the canonical representative by R4.
For each Borel gap, its $\mu$-mass is at most $A$ times its raw capacity
and hence at most $A$ times its outer capacity. Thus it converges in
$\mu$-measure on every compact stage. Its $L^2(\mu)$ limit must be
that representative almost everywhere. This argument selects the
canonical representative; it does not identify arbitrary modifications
on $\lambda$-null sets.
\end{proof}

\begin{corollary}[A controlled negative measure perturbation]
\label{cor:r6-form-small-measure-potential}
Fix $V\Subset\Omega$, put $H=\mathscr D_0(V)$, and let $\mu$ be
carried by $V$ and satisfy \eqref{eq:r6-linear-capacity-domination}.
For $c\ge0$ with $4cA<1$, the bilinear form
\[
        \mathcal B_c(u,v)=\mathcal E(u,v)
                              -c\int (J_\mu u)(J_\mu v)\,d\mu
\]
is continuous and coercive on $H$, with
\[
              \mathcal B_c(u,u)\ge(1-4cA)\mathcal E(u,u).
\]
For every bounded energy functional $\ell\in H^*$ there is a unique
$u\in H$ with $\mathcal B_c(u,\phi)=\ell(\phi)$ for all $\phi\in H$,
and
\[
                      \|u\|_H\le\frac{\|\ell\|_{H^*}}{1-4cA}.
\]
If $\ell(\phi)=-\int\phi\,d\sigma$ for a finite-energy Radon datum,
this is the original-current equation
\[
                         D(uS)+c(J_\mu u)\mu=\sigma
                                                    \quad\hbox{on }V.
\]
The boundary condition is the stated energy part space. No ambient
psh, nonlinear Monge--Ampere, or Markov property of the perturbed form
is claimed.
\end{corollary}

\begin{proof}
Cauchy--Schwarz in $L^2(\mu)$ and the preceding bound give
$|\int J_\mu uJ_\mu v\,d\mu|\le4A\|u\|_H\|v\|_H$.
The displayed lower bound follows. Thus $\mathcal B_c$ is an
equivalent Hilbert inner product, and its Riesz theorem gives the
solution and estimate. For an original compact test, the ordinary
Green identity identifies its energy term with $-D(uS)$.
The trace factor is locally $\mu$-integrable by Cauchy--Schwarz and
local finiteness, so it defines the indicated Radon measure. Rearranging
the two terms gives the claimed sign in the current equation.

The coefficient threshold is a uniform sufficient threshold. The
annular family in Example~\ref{ex:r6-four-is-sharp}, scaled by $A$,
with part domains containing the respective compact test supports,
shows that for $4cA>1$ these forms need not even be nonnegative.
At the equality threshold it supplies no uniform positive coercivity
constant. This does not decide critical behavior for one fixed measure.
\end{proof}

As throughout this energy discussion, a fixed higher mixed tail can
replace $S$. The resulting capacity and energy belong to that fixed
degree-one tail; this does not control all slots of a nonlinear
diagonal family by one quadratic form.

% END source: sections/dirichlet_strong_capacity_trace.tex


% BEGIN canonical R6 source: sections/dirichlet_energy_growth_modulus.tex
\section{A square-integrable capacity modulus for finite-energy data}
\label{sec:r6-energy-growth-modulus}

Throughout the first statement, $S$ is an arbitrary allowed original
positive separately closed current of bidegree $(n-1,n-1)$ on the
fixed ordinary domain $\Omega$. The energy and part spaces are the
frozen faithful ones. All capacities use normalized actual classical
testers on the whole $\Omega$.

\begin{theorem}[A modulus criterion from the strong capacity inequality]
\label{thm:r6-square-integrable-capacity-modulus}
Let $\mu$ be a finite positive Radon measure, carried by an ordinary
$V\subset\Omega$, with total mass $M$. Suppose that $F:[0,\infty)\to
[0,\infty)$ is nondecreasing, finite, locally absolutely continuous,
and $F(0)=0$, and that
\[
            \mu(E)\le F(\operatorname{Cap}_S(E;\Omega))
                          \qquad(E\subset\Omega\text{ Borel}).
\]
For infinite capacity, $F(\infty)$ means its monotone extended limit.
If $r_0>0$ satisfies $F(r_0)\ge M$ and
$F'\in L^2(0,r_0)$, then $\mu$ has finite energy on $V$ and
\begin{equation}
 \|\mu\|_{-1,V}
          \le 2\left(\int_0^{r_0}|F'(r)|^2\,dr\right)^{1/2}.
 \label{eq:r6-square-modulus-energy-bound}
\end{equation}
Consequently the energy-class Dirichlet equation has its unique
solution in $\mathscr D_0(V)$ whenever this coercive part space is
defined as in Theorem~\ref{thm:r6-measure-energy-dirichlet}.
\end{theorem}

\noindent\textbf{Status: \texttt{DIRECT\_COROLLARY}} of the new sharp
strong-capacity inequality, with the displayed modulus criterion
proved explicitly. It is an energy-class solvability statement;
no classical or completed-psh membership is asserted.

\begin{proof}
Let $f$ be one original real compact smooth core function in $V$,
extended by zero to $\Omega$, and put
\[
 c(t)=\operatorname{Cap}_S(\{|f|>t\};\Omega),\qquad
 L(r)=\big|\{t>0:c(t)>r\}\big|.
\]
The function $c$ is nonincreasing. Its superlevel sets in the variable
$t$ are therefore initial intervals, up to endpoints. Tonelli's
theorem and Theorem~\ref{thm:r6-strong-capacity-four} give
\[
 \int_0^\infty L(r)^2\,dr
    =\int_0^\infty\int_0^{c(t)}2t\,dr\,dt
    =\int_0^\infty2t c(t)\,dt\le4\mathcal E_S(f).
\]
All analytic level sets here are ordinary open, so neither a raw/outer
identification on general Borel sets nor Choquet regularity is used.
Since $F(r_0)\ge M$, the datum bound gives
\[
 \begin{split}
 \int|f|\,d\mu
  &=\int_0^\infty\mu\{|f|>t\}\,dt\\
  &\le\int_0^\infty F(\min\{c(t),r_0\})\,dt\\
  &=\int_0^{r_0}F'(r)L(r)\,dr\\
  &\le2\left(\int_0^{r_0}|F'(r)|^2\,dr\right)^{1/2}
                                 \mathcal E_S(f)^{1/2}.
 \end{split}
\]
Absolute continuity and $F(0)=0$ justify the second Tonelli identity;
its derivative is nonnegative almost everywhere. This proves the
core dual bound, and the frozen energy Riesz theorem supplies the
stated solution. If $M=0$, the same conclusion is immediate.
\end{proof}

\begin{corollary}[The critical logarithmic family]
\label{cor:r6-critical-log-energy-modulus}
For
\[
       F_\beta(r)=A\,r^{1/2}
                         [\log(e+r^{-1})]^{-\beta},\qquad A>0,
       \quad F_\beta(0)=0,
\]
every finite measure satisfying the preceding capacity-growth bound
has finite energy if $\beta>1/2$. In contrast, for every
$0\le\beta\le1/2$ there is a finite positive measure on the original
unit disc, compactly supported inside one fixed ordinary smaller
disc, satisfying that same bound with $A=1$, whose energy dual norm
is infinite.
\end{corollary}

\noindent\textbf{Status: \texttt{NEW\_PROVED}} for the exact universal
logarithmic threshold within this specified family. The sufficient
direction is the preceding direct criterion; the counterexample uses
the frozen original radial calibration. This is not an if-and-only-if
criterion for an individual measure.

\begin{proof}
Write $L(r)=\log(e+r^{-1})$. For $r>0$,
\[
 F_\beta'(r)
    = A r^{-1/2}L(r)^{-\beta}
           \left(\frac12+\frac{\beta}{(er+1)L(r)}\right).
\]
For $\beta\ge0$ this is positive, and its square is integrable near
zero exactly when $2\beta>1$. The function is unbounded at infinity,
so a finite $r_0$ with $F_\beta(r_0)\ge M$ always exists. This proves
the sufficient direction.

For the opposite direction fix $\lambda=-\log|\pi|$, let
$t_j=2^j\lambda$, $z_j=\zeta_{0,e^{-t_j}}$, and put
\[
 W_j=F_\beta(t_j^{-1})
       =t_j^{-1/2}[\log(e+t_j)]^{-\beta},\qquad
 m_j=W_j-W_{j+1}>0,\qquad
 \mu=\sum_{j\ge1}m_j\delta_{z_j}.
\]
The total mass is $W_1$, the support is contained in
$\{0\}\cup\{z_j\}$, and $\mu\{0\}=0$. All nonzero-mass support
points are original divisorial Gauss points. As in
Proposition~\ref{prop:r6-half-power-infinite-energy}, take the fixed
$V=\{|T|<e^{-\lambda/2}\}\Subset\Omega=\{|T|<1\}$.
The normalized hinge with curvature $t_j^{-1}\delta_{z_j}$ proves
for every Borel $E$ containing a least-index support atom $z_j$ that
\[
        \mu(E)\le W_j
             \le F_\beta(\operatorname{Cap}_{[1]}(E;\Omega)).
\]
The inequality is immediate for sets containing no support atom.

Use the finite test functions from that proposition with these new
coefficients,
\[
          f_N(t)=\sum_{j=1}^N m_j
                           \min\{(t-\lambda)_+,t_j-\lambda\}.
\]
Their energy and partial-measure pairing are again equal:
\[
 E_N=\sum_{j=1}^N(t_j-t_{j-1})
                            (W_j-W_{N+1})^2
        =\int f_N\,d\mu_N\le\int f_N\,d\mu,
          \qquad t_0=\lambda.
\]
The quantities $E_N$ increase to
\[
 \sum_{j\ge1}(t_j-t_{j-1})W_j^2
       =\frac12\sum_{j\ge1}[\log(e+t_j)]^{-2\beta}=\infty
                                     \quad(0\le\beta\le1/2).
\]
Here monotone convergence is applied to the nonnegative discrete
sum; no rate estimate for $m_j/W_j$ is needed. The same original
smooth compact approximations inside the fixed $V$ as in
Proposition~\ref{prop:r6-half-power-infinite-energy} pass both energy
and integrals against the finite $\mu$. A finite core dual norm
would imply $E_N\le C\sqrt{E_N}$, a contradiction.
\end{proof}

\paragraph{Power bounds and scope.}
For $F(r)=Ar^\alpha$, $\alpha>1/2$, taking
$r_0=(M/A)^{1/\alpha}$ gives
\[
 \|\mu\|_{-1,V}
   \le\frac{2\alpha}{\sqrt{2\alpha-1}}
                  A^{1/(2\alpha)}M^{1-1/(2\alpha)}.
\]
One may use the smaller of this bound and the earlier direct
layer-cake bound. The new improvement is at the half-power endpoint;
the linear case still gives $2\sqrt{AM}$. A fixed higher-degree
background with frozen mixed slots may be replaced by its resulting
degree-one current, with that current's own capacity and energy.
No bound for a nonlinear capacity with all slots varying is inferred.

% END source: sections/dirichlet_energy_growth_modulus.tex


% BEGIN canonical R6 source: sections/dirichlet_energy_threshold_obstruction.tex
\section{Sharpness of the power-growth threshold for finite energy}
\label{sec:r6-energy-threshold-obstruction}

The following construction reuses the frozen radial Gauss-point
calibration and first-atom tail comparison from
Section~\textup{R4--60}. Its new conclusion
concerns finite faithful energy, rather than bounded classical
Monge--Ampere solvability.

\begin{proposition}[The exponent one half cannot be included]
\label{prop:r6-half-power-infinite-energy}
For every $0<\alpha\le1/2$ there is a finite positive Radon measure
$\mu$ on the original open unit disc $\Omega$, supported on a fixed
compact subset of an ordinary $V\Subset\Omega$, such that
\begin{equation}
                   \mu(E)\le\operatorname{Cap}_{[1]}(E;\Omega)^\alpha
                         \qquad(E\subset\Omega\text{ Borel}),
 \label{eq:r6-low-power-capacity-bound}
\end{equation}
but $\|\mu\|_{-1,V}=\infty$ for the faithful degree-one energy of
$S=[1]$. Therefore it has no solution in $\mathscr D_0(V)$ of
$D(uS)=\mu$ on all original scalar tests. The measure is carried by
divisorial Gauss points and gives zero mass to their rigid limit point.
\end{proposition}

\noindent\textbf{Status: \texttt{NEW\_PROVED}} as the sharpness
application for the R6 finite-energy criterion. The radial curvature
and original integration normalization are frozen inputs, not new
results of this proposition. No additional bounded-classical
Dirichlet conclusion is claimed.

\begin{proof}
Fix $\lambda=-\log|\pi|>0$ and put
\[
 t_j=2^j\lambda,\qquad z_j=\zeta_{0,e^{-t_j}},\qquad
 W_j=t_j^{-\alpha},\qquad
 m_j=W_j-W_{j+1}=(1-2^{-\alpha})t_j^{-\alpha}.
\]
Set $\mu=\sum_{j\ge1}m_j\delta_{z_j}$. Its total mass is
$W_1<\infty$. Since $e^{-t_j}=|\pi|^{2^j}$, all $z_j$ are points
of the original $K$-analytic disc and are divisorial. They converge
to its rigid origin, so
$K_0=\{0\}\cup\{z_j:j\ge1\}$ is compact. The measure gives no
mass to zero. Take the fixed ordinary disc
\[
                         V=\{|T|<e^{-\lambda/2}\}\Subset\Omega.
\]
It contains $K_0$.

For each $j$, the actual normalized classical function on the whole
$\Omega$
\[
             w_j(T)=\max\{0,1+t_j^{-1}\log|T|\}
\]
is assigned value zero near the origin and has
$D w_j=t_j^{-1}\delta_{z_j}$, by the original radial calibration.
If $E$ contains some support atom, choose its least index $j$.
Then
\[
       \mu(E)\le\sum_{k\ge j}m_k=W_j=t_j^{-\alpha}
                       \le\operatorname{Cap}_{[1]}(E;\Omega)^\alpha.
\]
If it contains none, its mass is zero. This proves the bound for
every Borel set, including sets containing the limit point. It uses
neither compact outer regularity nor an optimizing envelope.

On the punctured unit disc put $t=-\log|T|$. For $N\ge1$ define
the nonnegative continuous radial function
\[
 f_N(t)=\sum_{j=1}^N m_j\min\{(t-\lambda)_+,t_j-\lambda\},
\]
with its constant terminal value at the origin. It is zero for
$t\le\lambda$, and constant for $t\ge t_N$. For
$t_{j-1}<t<t_j$, where $t_0=\lambda$, its derivative is
\[
                         M_{j,N}=\sum_{k=j}^N m_k.
\]
The one-dimensional original energy is consequently
\begin{equation}
 E_N=\int_\lambda^\infty|f_N'(t)|^2dt
       =\sum_{j=1}^N(t_j-t_{j-1})M_{j,N}^2
       \ge\frac{(1-2^{-\alpha})^2}{2}
                       \sum_{j=1}^N t_j^{1-2\alpha}.
 \label{eq:r6-radial-energy-divergence}
\end{equation}
This tends to infinity for every $0<\alpha\le1/2$.
No factor of $\lambda$ is missing: $t$ is the unnormalized unit
moment $-\log|T|$, so its energy calibration is $dt$ itself.

Let $\mu_N=\sum_{j\le N}m_j\delta_{z_j}$. A direct finite sum gives
\[
 \begin{split}
 \int f_N\,d\mu_N
  &=\sum_{i,j\le N}m_im_j
                       \min(t_i-\lambda,t_j-\lambda)\\
  &=\int_\lambda^\infty
                  \left(\sum_{j\le N}m_j\mathbf1_{\{t<t_j\}}\right)^2dt
    =E_N.
 \end{split}
\]
Since $\mu\ge\mu_N$ and $f_N\ge0$, one has
$\int f_N\,d\mu\ge E_N$.

We now realize these comparisons by original smooth core functions,
rather than testing a current with a merely PL function. Extend $f_N$
to the real line by zero to the left and its terminal constant to the
right, and convolve with smooth nonnegative kernels of radius
$\delta<\lambda/4$. The resulting $f_{N,\delta}(t)$ is zero for
$t\le\lambda-\delta$ and constant for $t\ge t_N+\delta$.
It therefore extends as an original smooth function across the rigid
origin and has compact support inside the fixed $V$. For each fixed
$N$ the approximation is uniform on the entire analytic disc, and
its derivative converges in $L^2(dt)$ to $f_N'$.
The original calibration gives
\[
 \mathcal E_S(f_{N,\delta})\longrightarrow E_N,
       \qquad \int f_{N,\delta}\,d\mu\longrightarrow\int f_N\,d\mu.
\]
The second assertion uses the finite total mass of $\mu$ and uniform
approximation, not a singular-test convergence theorem.

If $\|\mu\|_{-1,V}=C<\infty$, apply its core bound to
$f_{N,\delta}$ and let $\delta\downarrow0$. It would give
\[
                  E_N\le\int f_N\,d\mu\le C\sqrt{E_N}
\]
for every $N$, contradicting \eqref{eq:r6-radial-energy-divergence}.
This proves infinite energy. Any energy solution in the specified
part space would make the core datum bounded by its own energy norm
through the original Green identity, so such a solution cannot exist.
\end{proof}

\paragraph{Ambient dimension and scope.}
The same remaining-degree-one example can be embedded in every ambient
dimension by taking $S$ to be the integration current of the coordinate
disc in a product polydisc. The radial normalized testers pull back
from its first coordinate, while the core functions are multiplied by
fixed smooth cutoffs equal to one on the other coordinate origins.
Their restriction and energy are exactly the displayed ones.
This preserves arbitrary original-test meanings. It does not turn this
linear energy obstruction into a claim about a nonlinear operator with
several moving energy slots.

% END source: sections/dirichlet_energy_threshold_obstruction.tex


% BEGIN canonical R6 source: sections/dirichlet_concave_modulus_threshold.tex
\section{The universal finite-energy criterion for concave moduli}
\label{sec:r6-concave-modulus-threshold}

The word ``universal'' in the following statement quantifies over all
measures obeying a fixed growth modulus. It does not characterize the
energy of an individual measure and does not characterize classical
Monge--Ampere solvability.

\begin{theorem}[An exact criterion within the concave class]
\label{thm:r6-concave-modulus-threshold}
Let $F:[0,\infty)\to[0,\infty)$ be finite, nondecreasing and
concave, with $F(0)=0$ and $F$ continuous at zero. The following
universal finite-energy assertion holds if and only if
\begin{equation}
                         \int_0^a |F'(r)|^2\,dr<\infty
                         \quad\text{for some }a>0:
 \label{eq:r6-concave-modulus-threshold}
\end{equation}
for every allowed original degree-one background $S$ on a fixed
ordinary domain $\Omega$ carrying the faithful energy form, and
every finite positive Radon measure $\mu$ supported on a compact
$K\Subset V\Subset\Omega$ satisfying
\[
        \mu(E)\le F(\operatorname{Cap}_S(E;\Omega))
                       \qquad(E\subset\Omega\text{ Borel}),
\]
the core dual norm $\|\mu\|_{-1,V}$ is finite.

If \eqref{eq:r6-concave-modulus-threshold} fails, the counterexample
already occurs for $S=[1]$ on the original open unit disc. Its support
is a fixed compact sequence of divisorial Gauss points and their
rigid limit, inside one fixed smaller ordinary disc; the limit
point has mass zero. Thus failure of the criterion does not require
a varying domain, an atomic rigid pole, or a complicated background.
\end{theorem}

\noindent\textbf{Status: \texttt{NEW\_PROVED}} as the exact universal
threshold for this specified concave class, obtained by combining the
new strong-capacity criterion with the explicit radial construction.
All original Gauss and energy calibrations are frozen inputs.

\begin{proof}
A finite concave function is locally absolutely continuous away
from zero. Its nonnegative decreasing almost-everywhere derivative
integrates to $F$ also from zero, by continuity there. In particular,
it is absolutely continuous on every finite closed interval. On
intervals bounded away from zero its derivative is bounded. Thus
\eqref{eq:r6-concave-modulus-threshold} for one $a$ implies the same
condition on every finite interval.

Suppose first that the condition holds. Relative CLN gives
$C_K=\operatorname{Cap}_S(K;\Omega)<\infty$. If $C_K=0$, the
growth bound gives $\mu=0$. Otherwise
$\mu(\Omega)=\mu(K)\le F(C_K)$. Apply
Theorem~\ref{thm:r6-square-integrable-capacity-modulus} with
$r_0=C_K$. It gives, explicitly,
\[
             \|\mu\|_{-1,V}
                  \le2\left(\int_0^{C_K}|F'(r)|^2\,dr\right)^{1/2}
                  <\infty.
\]
This proves the full arbitrary-fixed-background direction.

Conversely, assume that the integral of $|F'|^2$ diverges near zero.
Concavity and $F(0)=0$ imply, almost everywhere,
\[
                    0\le F'(r)\le\frac{F(r)}r.
\]
Indeed the right-hand side is the average of the decreasing
derivative on $(0,r)$. Hence
$\int_0^a F(r)^2r^{-2}\,dr=\infty$ for every $a>0$.
Fix $\lambda=-\log|\pi|$, let $t_j=2^j\lambda$ and
$r_j=t_j^{-1}$. Monotonicity of $F$ gives
\[
 \int_{r_{j+1}}^{r_j}\frac{F(r)^2}{r^2}\,dr
       \le F(r_j)^2\left(\frac1{r_{j+1}}-\frac1{r_j}\right)
       =t_jF(t_j^{-1})^2.
\]
It follows that $\sum_j t_jF(t_j^{-1})^2=\infty$.

Set $W_j=F(t_j^{-1})$, $m_j=W_j-W_{j+1}\ge0$ and
\[
             z_j=\zeta_{0,e^{-t_j}},\qquad
             \mu=\sum_{j\ge1}m_j\delta_{z_j}.
\]
Continuity at zero implies $W_j\downarrow0$, so $\mu$ is finite
with total mass $W_1$. It is supported in
$\{0\}\cup\{z_j\}\Subset V=\{|T|<e^{-\lambda/2}\}$ and gives
zero mass to zero. The normalized original hinge of curvature
$t_j^{-1}\delta_{z_j}$ proves the growth bound for every Borel set:
if $j$ is the least index of an atom of positive mass in that set,
its mass is at most $W_j$, and its capacity is at least $t_j^{-1}$.
If no such index exists, its mass is zero. No capacity regularity
is used.

The finite radial test functions
\[
          f_N(t)=\sum_{j=1}^Nm_j
                       \min\{(t-\lambda)_+,t_j-\lambda\},
                       \qquad t=-\log|T|,
\]
have, with $t_0=\lambda$,
\[
 \mathcal E_{[1]}(f_N)
   =\sum_{j=1}^N(t_j-t_{j-1})(W_j-W_{N+1})^2
   =\int f_N\,d\mu_N\le\int f_N\,d\mu.
\]
Their energies increase to
\[
       \sum_{j\ge1}(t_j-t_{j-1})W_j^2
                =\frac12\sum_{j\ge1}t_jF(t_j^{-1})^2=\infty.
\]
The original smooth compact approximations inside this same fixed
$V$ are exactly those justified in
Proposition~\ref{prop:r6-half-power-infinite-energy}: mollify each
finite real PL profile, retaining its zero initial region and its
constant terminal region. Uniform convergence and $L^2(dt)$
derivative convergence pass both the finite-measure pairing and
the energy. A finite core dual norm would therefore bound the
displayed diverging energy by a constant times its square root,
which is impossible. This completes the necessity assertion.
\end{proof}

\paragraph{Qualification.}
Concavity is used only for the sharp converse. For a general
nondecreasing absolutely continuous modulus, the square-integrable
derivative condition remains sufficient, but this proof asserts
no necessity. Neither direction identifies the energy capacity
with relative BT capacity. The conclusion supplies the energy-class
linear equation through the Riesz theorem; all psh feasibility and
trace constraints remain separate.

% END source: sections/dirichlet_concave_modulus_threshold.tex


% BEGIN canonical R6 source: sections/dirichlet_normalized_image_obstruction.tex
\section{Capacity domination does not characterize the normalized image}
\label{sec:r6-normalized-image-obstruction}

This is a scope consequence of the two sharp estimates just proved.
It concerns a fixed normalization on a fixed ordinary domain; it is
neither a counterexample to C1 nor an assertion of Dirichlet
nonexistence after removing the normalization.

\begin{corollary}[A strict distinction on an original annulus]
\label{cor:r6-capacity-not-normalized-image}
There are an ordinary annulus $\Omega$ over the original field $K$
and a finite positive Radon measure $\mu$, compactly supported in
$\Omega$, such that
\[
             \mu(E)\le\operatorname{Cap}_{[1]}(E;\Omega)
                    \qquad(E\subset\Omega\text{ Borel}),
\]
but $\mu$ does not belong to the weak closure, on all original compact
smooth scalar tests, of
\[
 \mathcal M=\{D w: w\in\mathcal C_b(\Omega),\ 0\le w\le1
                                      \text{ on the whole }\Omega\}.
\]
In particular it is not the curvature of a normalized completed
potential defined by a same-domain row in this normalized classical
family.
\end{corollary}

\noindent\textbf{Status: \texttt{DIRECT\_COROLLARY}} of
Theorem~\ref{thm:r6-sharp-relative-hardy} and
Example~\ref{ex:r6-four-is-sharp}. No new compact-regularity assertion
is used.

\begin{proof}
Put $\lambda=-\log|\pi|$ and
$\Omega=\{1<|T|<e^\lambda\}$. Write
$s=\log|T|/\lambda\in(0,1)$, and let $\sigma(s)$ be its Gauss section.
For $0<\varepsilon<1/4$ take
\[
 \mu_\varepsilon
    =\lambda^{-1}\sigma_*
                    (\mathbf1_{[\varepsilon,1/2]}s^{-2}\,ds).
\]
These measures have compact support and finite mass. The normalized
hinges $w_r(s)=\max(0,1-s/r)$ have
$D w_r=(\lambda r)^{-1}\delta_{\sigma(r)}$ on this same annulus.
For a Borel set $E$ meeting the support of $\mu_\varepsilon$, put
$a=\inf\{s\in[\varepsilon,1/2]:\sigma(s)\in E\}$. Then
\[
 \mu_\varepsilon(E)\le\frac1{\lambda a}
   =\sup_{\sigma(r)\in E,\ r\in[\varepsilon,1/2]}
                                      \frac1{\lambda r}
   \le\operatorname{Cap}_{[1]}(E;\Omega).
\]
The bound is immediate when $E$ misses this support. Attainment of
the infimum is unnecessary.

The original smooth compact tests constructed in
Example~\ref{ex:r6-four-is-sharp} satisfy
\[
 \frac{\int f_\varepsilon^2\,d\mu_\varepsilon}
      {\mathcal E_{[1]}(f_\varepsilon)}\longrightarrow4.
\]
Choose and then fix one sufficiently small $\varepsilon$ for which
this ratio is greater than one. On the other hand, the sharp Hardy
estimate gives
\[
        \int f_\varepsilon^2\,d\nu
                  \le\mathcal E_{[1]}(f_\varepsilon)
                         \qquad(\nu\in\mathcal M).
\]
The same inequality holds in the weak closure of $\mathcal M$ because
$f_\varepsilon^2$ is one fixed original compact smooth scalar test.
Thus the displayed strict inequality separates
$\mu_\varepsilon$ from that weak closure, for sequences and nets
alike. The completed-image conclusion follows from the fixed-bound
completed continuity in Corollary~\ref{cor:r6-sharp-energy-capacity}.
\end{proof}

\paragraph{Direction of the implication.}
C1 asks whether a weak-limit or completed measure is dominated, on
compact sets, by the raw classical-test capacity. The present measure
already has that domination and fails to be a weak limit. It disproves
the converse characterization, not C1. In particular a capacitary
upper bound cannot replace the support-functional constraints in
Theorem~\ref{thm:r6-energy-dirichlet-alternative}. No assertion is made
here about unnormalized classical potentials or arbitrary boundary
values.

% END source: sections/dirichlet_normalized_image_obstruction.tex


% BEGIN canonical R6 source: sections/capacity_energy_null_gap.tex
\section{Relative BT and energy capacities can have different null sets}
\label{sec:r6-capacity-energy-null-gap}

The sharp one-way energy bound cannot be reversed, even at the level of
null sets. The counterexample below uses a smooth original-positive
background on a full torus polyannulus. It is not a counterexample to
compact no-gain or to ordinary-open outer regularity of BT capacity.

Fix $n\ge3$, let $Q=(-2,2)^n$, and put
\[
 \Omega=\operatorname{trop}^{-1}(Q)\subset(\mathbf G_{m,K}^n)^{\an},
 \qquad t_i=\log|T_i|,\qquad
 \beta=\sum_{i=1}^n d't_i\wedge d''t_i,
 \qquad S=\frac{\beta^{n-1}}{(n-1)!}\wedge[1].
\]
The background $S$ has bidegree $(n-1,n-1)$, is symmetric,
original-positive, and separately closed. Write
$\sigma:Q\to\Omega$ for the original Gauss section. The original
torus calibration gives, for $\varphi\in C_c^\infty(Q)$,
\begin{equation}
 D(\varphi\circ\operatorname{trop})\wedge S
                  =\sigma_*((\Delta\varphi)\,dx),\qquad
 \mathcal E_S(\varphi\circ\operatorname{trop})
                  =\int_Q|\nabla\varphi|^2\,dx.
 \label{eq:r6-torus-laplace-calibration}
\end{equation}
The factorial in $S$ is chosen precisely for these identities. In
particular no normalization of the original absolute value is made.

\begin{lemma}[Uniform Gauss-section growth for all classical testers]
\label{lem:r6-gauss-lipschitz-capacity}
There is a constant $C_n$ such that, for every actual
$u\in\mathcal C_b(\Omega)$ with $0\le u\le1$, and every Euclidean
ball $B(x,2r)\subset(-1,1)^n$,
\begin{equation}
              D(uS)(\operatorname{trop}^{-1}B(x,r))
                                      \le C_n r^{n-1}.
 \label{eq:r6-uniform-gauss-ball-mass}
\end{equation}
Consequently, if $A\subset(-1/2,1/2)^n$ is compact and
$\mathcal H^{n-1}(A)=0$, then
\[
                    \operatorname{Cap}_S^*(\sigma(A);\Omega)=0.
\]
\end{lemma}

\begin{proof}
Put $g=u\circ\sigma$. We first justify a restriction property of
$g$ rather than assume it from the word ``skeleton.'' Fix all but one
of its logarithmic coordinates. Let $F$ be the completed residue field
of the Gauss point of the corresponding $(n-1)$-torus over $K$.
Its coordinate elements have the prescribed radii. The analytic map
from an annulus over $F$ that varies the remaining coordinate and
keeps these elements fixed sends its Gauss point of each radius to
the corresponding original $K$-Gauss point in $\Omega$. This is the
iterated Gauss seminorm identity on Laurent polynomials; it concerns
actual points, not replacement of full analytic fibers in integration.

By the defining complete-extension curve test for classical psh, the
pullback of $u$ is subharmonic. On a closed subannulus, subtract the
affine harmonic function of the logarithmic radius that agrees with
its two Gauss endpoint values. The subharmonic Shilov maximum
principle on that annulus gives a nonpositive difference. If necessary
one first extends $F$ to realize the two endpoint radii in its value
group; the same actual Gauss identities and classical curve test
remain valid. Thus $g$ is convex in the one varying coordinate.
Repeating for each coordinate proves separate convexity.

The bound $0\le g\le1$ on $Q$ now gives a uniform Lipschitz bound
on $(-1,1)^n$: the one-dimensional convex slope bounds hold uniformly
in all the fixed other coordinates, and one joins two points by
coordinate segments. In particular $g$ has bounded weak first
derivatives there, by a constant depending only on $n$.

Let $\mu_u=D(uS)$, a positive Radon measure by the frozen calculus.
For every $\varphi\in C_c^\infty(Q)$, recursion and
\eqref{eq:r6-torus-laplace-calibration} give
\[
 \langle\operatorname{trop}_*\mu_u,\varphi\rangle
   =\langle uS,D(\varphi\circ\operatorname{trop})\rangle
   =\int_Q g\Delta\varphi\,dx.
\]
Tropicalization is proper on compact subsets of $Q$, so this is an
identity of Radon measures: $\operatorname{trop}_*\mu_u=\Delta g$.
Choose a nonnegative scalar cutoff equal to one on $B(x,r)$,
supported in $B(x,2r)$, with gradient bounded by a fixed multiple
of $r^{-1}$. Integration by parts against the bounded weak gradient
of $g$ gives
\[
 \Delta g(B(x,r))
    \le\int\varphi_r\,d\Delta g
    =-\int\nabla g\cdot\nabla\varphi_r\,dx
    \le C_n r^{n-1}.
\]
The constant is independent of $u$; this proves
\eqref{eq:r6-uniform-gauss-ball-mass}.

Cover $A$ by finitely many sufficiently small balls $B(x_i,r_i)$
with $\sum_i r_i^{n-1}<\epsilon$, keeping the doubled balls in
$(-1,1)^n$. Their tropical inverse images form one ordinary open
neighborhood $G$ of $\sigma(A)$ in the same $\Omega$. The uniform
estimate and finite subadditivity give
$\operatorname{Cap}_S(G;\Omega)\le C_n\epsilon$.
Let $\epsilon\downarrow0$. This proves outer capacity zero directly;
no compact outer-regularity theorem has been used.
\end{proof}

\begin{lemma}[Every original energy dominates the Gauss restriction energy]
\label{lem:r6-energy-gauss-restriction}
For every original real compact smooth $f$ on $\Omega$, the actual
function $g=f\circ\sigma$, extended by zero outside $Q$, belongs to
$H_0^1(Q)$ and satisfies
\begin{equation}
                         \int_Q|\nabla g|^2\,dx\le\mathcal E_S(f).
 \label{eq:r6-energy-restriction-bound}
\end{equation}
\end{lemma}

\begin{proof}
The function $g$ is continuous and compactly supported in $Q$.
For $\varphi\in C_c^\infty(Q)$, compact Green, the torus
calibration, and the original energy Cauchy--Schwarz inequality give
\begin{equation}
 \left|\int_Qg\Delta\varphi\,dx\right|^2
   =|\mathcal E_S(f,\varphi\circ\operatorname{trop})|^2
   \le\mathcal E_S(f)\int_Q|\nabla\varphi|^2\,dx.
 \label{eq:r6-restriction-dual-bound}
\end{equation}
This uses every original $f$, not just a torus-invariant subcore.

Choose a nonnegative even smooth probability mollifier $\kappa_\epsilon$.
For small enough $\epsilon$, both
$g_\epsilon=g*\kappa_\epsilon$ and
$\varphi_\epsilon=g*\kappa_\epsilon*\kappa_\epsilon$ are smooth
and compactly supported in $Q$. Symmetry, Fubini and ordinary
integration by parts give
\[
 -\int g\Delta\varphi_\epsilon\,dx
        =\int|\nabla g_\epsilon|^2\,dx,
 \qquad
 \|\nabla\varphi_\epsilon\|_2
        \le\|\nabla g_\epsilon\|_2.
\]
Substitute into \eqref{eq:r6-restriction-dual-bound}. After cancellation,
or trivially when this norm is zero, one obtains
$\|\nabla g_\epsilon\|_2^2\le\mathcal E_S(f)$.
As $g_\epsilon\to g$ in $L^2$, weak compactness in $H_0^1(Q)$
and lower semicontinuity prove
\eqref{eq:r6-energy-restriction-bound}.
This argument requires no assertion that a smooth original test has
a smooth restriction in one fixed real moment chart.
\end{proof}

\begin{lemma}[A compact real set with the two needed size properties]
\label{lem:r6-cantor-newton-set}
There is a compact $A\subset(-1/2,1/2)^n$ with
$\mathcal H^{n-1}(A)=0$ and a probability measure $\vartheta$
on $A$ such that
\[
                  I(\vartheta)=\iint|x-y|^{2-n}
                                  \,d\vartheta(x)d\vartheta(y)<\infty.
\]
Moreover there is a constant $c_A>0$ such that every continuous
compactly supported $g\in H_0^1(Q)$ with $g\ge1$ on $A$ satisfies
$\int|\nabla g|^2\ge c_A$.
\end{lemma}

\begin{proof}
Take a small affine copy $C$ of the middle-third Cantor set, a small
$(n-2)$-cube $J$, and put $A=C\times J\times\{0\}$. At scale
$3^{-j}$ it is covered by at most
$C_0 2^j3^{j(n-2)}$ cubes of that side length. The sum of their
$(n-1)$st diameter powers is at most $C_1(2/3)^j$, proving the
Hausdorff measure assertion.

Let $d=\log2/\log3>0$, and let $\vartheta$ be the product of
the Cantor Bernoulli probability, normalized Lebesgue measure on
$J$, and the point mass at zero. The usual interval count in the
Cantor construction gives
\[
                  \vartheta(B(x,r))\le C_2 r^{n-2+d}
                                                   \quad(0<r<1).
\]
Layer cake for $|x-y|^{2-n}$ now bounds its near-diagonal integral
by a constant times $\int_0^1r^{d-1}\,dr<\infty$; the remainder
is bounded. This proves the finite energy assertion.

For completeness, this energy supplies the required Sobolev lower
bound directly. For a real $g\in C_c^\infty(\mathbf R^n)$, the
divergence formula for the Newton kernel gives
\[
 g(x)=c_n\int_{\mathbf R^n}
                  \nabla g(y)\cdot\frac{x-y}{|x-y|^n}\,dy
\]
with a fixed dimensional constant (its sign is irrelevant below).
The convolution of the two nonnegative kernels $|x|^{1-n}$ is
$C_n|x|^{2-n}$: translation, rotation and scaling give this form,
and the defining integral is finite away from zero since $n\ge3$.
Tonelli and Cauchy--Schwarz therefore imply
\begin{equation}
           \left|\int g\,d\vartheta\right|^2
              \le C_n' I(\vartheta)\int|\nabla g|^2\,dx.
 \label{eq:r6-newton-trace-bound}
\end{equation}
For continuous compactly supported $g\in H_0^1(Q)$, convolution
approximates uniformly as well as in $H^1$; thus the same inequality
holds. If $g\ge1$ on $A$, its left side is at least one. Take
$c_A=(C_n'I(\vartheta))^{-1}>0$.
\end{proof}

\begin{theorem}[Strict separation of the two capacity null ideals]
\label{thm:r6-energy-BT-null-gap}
For the smooth background $S$ above and the compact set
$K=\sigma(A)\Subset\Omega$ of
Lemma~\ref{lem:r6-cantor-newton-set},
\begin{equation}
       \operatorname{Cap}_S^*(K;\Omega)=0,
       \qquad \operatorname{cap}_{\mathcal E,1}(K)>0.
 \label{eq:r6-energy-BT-null-gap}
\end{equation}
Here the second capacity is that of the frozen faithful regular
Dirichlet form on $F=\operatorname{supp}S$.
Thus neither a reverse comparison with a finite constant nor equality
of the two null ideals holds in general.
\end{theorem}

\noindent\textbf{Status: \texttt{NEW\_PROVED}} explicit obstruction
to the reverse energy-capacity route, in every $n\ge3$. It makes no
negative claim about the general relative C1 or C2 conjectures.

\begin{proof}
The first equality follows from
Lemma~\ref{lem:r6-gauss-lipschitz-capacity}. Every Gauss point of
$Q$ lies in $F$: a nonnegative original cutoff supported in any
neighborhood of it, paired with $\beta\wedge S$, has positive
integral on a nonempty real neighborhood of that point. Hence
$K\subset F$.

In the verified compact-core formula for
$\operatorname{cap}_{\mathcal E,1}(K)$, let $f$ be any original
smooth compact admissible function, so $f>1$ on $K$. Its continuous
Gauss restriction $g=f\circ\sigma$ is at least one on $A$.
The preceding two lemmas give
\[
           \mathcal E_1(f,f)\ge\mathcal E_S(f)
                    \ge\int|\nabla g|^2\,dx\ge c_A.
\]
Taking the infimum proves the second assertion. In particular the
failure cannot be repaired by passing to all original smooth tests:
the energy lower bound was proved for that full test class.
\end{proof}

\paragraph{Exact scope of the obstruction.}
This example has a fixed smooth background and remaining degree one.
Its BT outer capacity is already zero, so it exhibits no compact
weak-closure gain and no raw/outer discrepancy. It shows why the
new sharp Hardy bound, although valid for every normalized tester and
its weak limits, cannot close C1 merely by identifying the controlling
energy capacity with BT capacity. The actual nonlinear psh constraint
is visible in the uniform Gauss-section Lipschitz bound.

% END source: sections/capacity_energy_null_gap.tex


% BEGIN canonical R6 source: sections/capacity_gauss_convex_reduction.tex
\section{Exact capacity reduction on a Gauss section}
\label{sec:r6-gauss-convex-capacity}

Let $n\ge1$, let $Q\subset\mathbf R^n$ be a nonempty bounded convex
ordinary open set, and set
\[
 \Omega=\operatorname{trop}^{-1}(Q)\subset(\mathbf G_{m,K}^n)^{\an},
 \qquad S=\frac{\beta^{n-1}}{(n-1)!}\wedge[1],
 \qquad \beta=\sum_i d't_i\wedge d''t_i.
\]
Write $\sigma:Q\to\Omega$ for the original Gauss section and define
\[
 \mathcal G(Q)=\{g:Q\to\mathbf R:\ g\text{ is convex},\ 0\le g\le1\},
 \qquad c_Q(E)=\sup_{g\in\mathcal G(Q)}\Delta g(E)
\]
for Borel $E\subset Q$. The Laplacian is the positive distributional
Radon measure on the real domain. All noncompact sets considered below
are contained in a fixed compact subset of $Q$. Let $c_Q^*$ be the
outer capacity obtained using ordinary real open neighborhoods in the
same $Q$.

\begin{lemma}[Joint convexity from complete-extension annuli]
\label{lem:r6-joint-gauss-convexity}
Every actual $u\in\mathcal C_b(\Omega)$ with $0\le u\le1$ has
$u\circ\sigma\in\mathcal G(Q)$. Conversely every
$g\in\mathcal G(Q)$ gives a normalized continuous classical psh
function $g\circ\operatorname{trop}$ on the same $\Omega$.
\end{lemma}

\begin{proof}
The coordinate-annulus argument in
Lemma~\ref{lem:r6-gauss-lipschitz-capacity} applies on every real box
compactly contained in $Q$. It proves separate convexity and local
Lipschitz continuity of $g=u\circ\sigma$.

Fix a primitive integral vector $v\in\mathbf Z^n$. Complete $v$ to
the first column of a matrix $B\in\operatorname{GL}_n(\mathbf Z)$.
The associated monomial change of torus coordinates has logarithmic
coordinates $y=B^{-1}t$. Fix $y_2,\ldots,y_n$ and form the completed
residue field of their Gauss point. The analytic annulus over that
field varying $y_1$ maps into $\Omega$ over exactly the interval
cut out by $B^{-1}Q$. Its Gauss point at each radius maps to the
original Gauss point of $Q$ with those coordinates. This is the
monomial Gauss-seminorm identity, not a fiber-integration assertion.

The complete-extension classical curve test and the subharmonic
maximum principle on each closed subannulus therefore show that
$g$ is convex along every segment parallel to $v$. Endpoint radii
can first be realized after a further complete extension, as in the
coordinate case. Convexity follows along every rational direction.
For $x,y\in Q$, approximate $y-x$ by rational vectors $v_j$ with
$x+v_j\in Q$. Convexity of $Q$ keeps all the corresponding segments
inside $Q$. The already proved continuity of $g$ passes their convex
inequalities to the segment $[x,y]$. Thus $g$ is jointly convex.

For the converse, on each relatively compact real box one may
approximate the continuous convex $g$ locally uniformly by smooth
convex functions. Each such function of the unit-log coordinates
is classical psh: after any complete extension and analytic curve
map, the coordinates are harmonic and the convex Hessian chain rule
is nonnegative. Locally uniform limits preserve the subharmonic
curve inequalities. Hence $g\circ\operatorname{trop}$ is classical
on $\Omega$, with precisely the original normalization $0\le g\le1$.
\end{proof}

\begin{theorem}[Raw and outer capacities reduce to real convex traces]
\label{thm:r6-gauss-capacity-reduction}
For every Borel $E$ contained in a fixed compact core of $Q$,
\begin{align}
 \operatorname{Cap}_S(\sigma(E);\Omega)
 &=\operatorname{Cap}_S(\operatorname{trop}^{-1}(E);\Omega)
   =c_Q(E), \label{eq:r6-gauss-raw-reduction}\\
 \operatorname{Cap}_S^*(\sigma(E);\Omega)
 &=\operatorname{Cap}_S^*(\operatorname{trop}^{-1}(E);\Omega)
   =c_Q^*(E). \label{eq:r6-gauss-outer-reduction}
\end{align}
Every competitor in these identities is normalized on the entire
displayed capacity domain; no smaller-domain tester is substituted.
\end{theorem}

\noindent\textbf{Status: \texttt{DIRECT\_COROLLARY}} of the
actual complete-extension curve test, ordinary torus calibration, and
the frozen fixed-background calculus. It identifies a concrete real
capacity, rather than proving general relative Borel outer regularity.

\begin{proof}
For $u$ as in the preceding lemma, write $\mu_u=D(uS)$ and
$g=u\circ\sigma$. The calculation against original toric scalar
tests in Lemma~\ref{lem:r6-gauss-lipschitz-capacity} gives
\begin{equation}
                  \operatorname{trop}_*\mu_u=\Delta g.
 \label{eq:r6-gauss-push-laplacian}
\end{equation}
It is an identity of Radon measures because tropicalization is proper
over compact subsets of $Q$. Conversely, for the toric competitor
$u=g\circ\operatorname{trop}$ one has the stronger identity
\begin{equation}
                         D(uS)=\sigma_*(\Delta g).
 \label{eq:r6-toric-convex-laplacian}
\end{equation}
For smooth real $g$ this is the original torus calibration on all
original scalar tests. Local uniform smooth convex approximation,
fixed coefficient-Radon multiplication, and positivity pass it to
continuous convex $g$. The restriction of an arbitrary original test
to $\sigma(Q)$ is used only as a continuous compact function.

Equation~\eqref{eq:r6-gauss-push-laplacian} bounds both raw capacities
in \eqref{eq:r6-gauss-raw-reduction} by $c_Q(E)$, while
\eqref{eq:r6-toric-convex-laplacian} supplies the reverse bounds.
The sets are Borel: $\sigma$ is a closed embedding over every compact
real core, and tropical inverse images preserve Borel sets.

If $G\subset\Omega$ is any ordinary open set containing $\sigma(E)$,
then $O=\sigma^{-1}(G)$ is a real open neighborhood of $E$. Toric
competitors alone give
\[
       \operatorname{Cap}_S(G;\Omega)
           \ge\sup_{g\in\mathcal G(Q)}\Delta g(O)=c_Q(O)
           \ge c_Q^*(E).
\]
No saturated-neighborhood assumption is made in this lower bound.
Conversely, for every real open $O\supset E$, the ordinary saturated
open set $\operatorname{trop}^{-1}(O)$ has capacity $c_Q(O)$ by the
raw identity. These two bounds prove the outer formula for
$\sigma(E)$. Since
$\sigma(E)\subset\operatorname{trop}^{-1}(E)$, the same saturated
upper bound and monotonicity prove it also for the full inverse image.
All infima can be restricted to neighborhoods inside one larger
relatively compact core by intersection, so local mass bounds suffice.
\end{proof}

\begin{corollary}[Compact no-gain on the Gauss section]
\label{cor:r6-gauss-compact-no-gain}
The family $\mathcal G(Q)$ is compact in the topology of local uniform
convergence, and $g\mapsto\Delta g$ is continuous for local weak
Radon convergence. Consequently, for compact $A\Subset Q$,
\[
                         c_Q^*(A)=c_Q(A).
\]
If a net of normalized original relative measures $D(u_jS)$ converges
weakly to $\nu$ on $\Omega$, then
\[
            \nu(\sigma(A))\le\operatorname{Cap}_S(\sigma(A);\Omega).
\]
These statements give no Borel identity $c_Q^*=c_Q$ for arbitrary
Borel sets.
\end{corollary}

\noindent\textbf{Status: \texttt{DIRECT\_COROLLARY}}. The C1
consequence for this smooth nonpolar remaining-degree-one background
is already within the frozen R3 nonpolar C1 theorem. The new explicit
description here is the reduction to $\mathcal G(Q)$ and its exact
ordinary-open outer counterpart.

\begin{proof}
Normalized convex functions have one uniform Lipschitz bound on
each compact subset of $Q$. Arzela--Ascoli on a countable real compact
exhaustion gives compactness of $\mathcal G(Q)$; local uniform limits
retain convexity and the same bounds on all of $Q$. Distributional
convergence of their Laplacians, together with uniform local mass
bounds, gives weak Radon convergence. This is compactness of the
real normalized family, not a claim that arbitrary relative currents
have actual classical representatives.

If $A_j\downarrow A$ are compact sets in one core, choose
$g_j\in\mathcal G(Q)$ with
$\Delta g_j(A_j)\ge c_Q(A_j)-j^{-1}$. Pass to a locally uniformly
convergent subsequence with limit $g$. For every fixed $l$, compact
Portmanteau gives
\[
 \limsup_j\Delta g_j(A_j)
      \le\limsup_j\Delta g_j(A_l)\le\Delta g(A_l).
\]
Let $l\to\infty$ and use finite Radon continuity above. This proves
$c_Q(A_j)\downarrow c_Q(A)$. Taking shrinking compact neighborhoods
of $A$ now yields $c_Q^*(A)=c_Q(A)$.

For the final assertion, compactness gives a subnet
$u_j\circ\sigma\to g\in\mathcal G(Q)$ locally uniformly.
Equation~\eqref{eq:r6-gauss-push-laplacian}, tested first by compact
real smooth functions, gives $\operatorname{trop}_*\nu=\Delta g$.
Positivity therefore implies
\[
 \nu(\sigma(A))\le\nu(\operatorname{trop}^{-1}A)
      =\Delta g(A)\le c_Q(A),
\]
as required. The weak limit itself need not be toric. Nothing in this
proof exchanges an infimum over open neighborhoods with a supremum
over measures on a general Borel set.
\end{proof}

\paragraph{Remaining real question.}
The equality $c_Q(E)=c_Q^*(E)$ for all Borel $E\Subset Q$ would,
by the exact reduction, settle C2 on these Gauss-section sets.
Compactness of the convex measure image alone does not prove that
equality, as the frozen abstract weak-compact-family counterexample
already demonstrates. This restricted real question is distinct from
the full arbitrary-background C2 problem.

% END source: sections/capacity_gauss_convex_reduction.tex


% BEGIN canonical R6 source: sections/capacity_gauss_borel_regularity.tex
\section{Borel outer regularity on Gauss sets of finite codimension-one size}
\label{sec:r6-gauss-finite-size-regularity}

Retain the full convex torus domain and the smooth remaining-degree-one
background of Section~\ref{sec:r6-gauss-convex-capacity}, and assume
$n\ge2$. Fix a compact real core $Q_0\Subset Q$. Hausdorff measure
below refers to the ordinary Euclidean metric on $Q$, not to a metric
on the Berkovich space.

\begin{proposition}[An outer bound by codimension-one Hausdorff measure]
\label{prop:r6-gauss-hausdorff-outer-bound}
There is a constant $C=C(Q_0,Q,n)<\infty$ such that every subset
$E\subset Q_0$ satisfies
\[
                 c_Q^*(E)\le C\mathcal H^{n-1}(E).
\]
Equivalently the same bound holds for the ordinary-open outer relative
capacities of $\sigma(E)$ and $\operatorname{trop}^{-1}(E)$.
For non-Borel sets only these outer capacities are asserted.
\end{proposition}

\begin{proof}
The normalized convex family has a uniform Lipschitz bound on a fixed
larger real core. The cutoff calculation in
Lemma~\ref{lem:r6-gauss-lipschitz-capacity} therefore gives
\[
                   c_Q(B(x,r))\le C_0r^{n-1}
\]
for all sufficiently small balls centered near $Q_0$, with their
doubled balls in the larger core. Cover $E$ by countably many such
balls. Their union is one real open neighborhood; countable
subadditivity, valid directly for the supremum of the positive
measures $\Delta g$, bounds its capacity by
$C_0\sum_i r_i^{n-1}$. Taking the infimum over covers of arbitrarily
small diameter proves the bound, up to the fixed normalization of
Hausdorff measure.

The exact outer comparison proof of
Theorem~\ref{thm:r6-gauss-capacity-reduction} used only open
neighborhoods and monotonicity. It therefore applies to arbitrary
subsets for the outer capacities, and gives the stated equivalent
formulations. Alternatively the saturated inverse images of the
covering balls give the Berkovich upper bounds directly.
\end{proof}

\begin{corollary}[A concrete Borel C2 class]
\label{cor:r6-gauss-finite-hausdorff-C2}
If $E\subset Q_0$ is Borel and
$\mathcal H^{n-1}(E)<\infty$, then
\[
                         c_Q(E)=c_Q^*(E).
\]
Consequently raw and ordinary-open outer relative capacities agree
for both $\sigma(E)$ and $\operatorname{trop}^{-1}(E)$, with capacity
domain the same whole $\Omega$. This includes every Borel subset of
a compact smooth real hypersurface in $Q$ of finite area.

If instead $E$ is a countable union of Borel sets of finite
$\mathcal H^{n-1}$ measure, then
\[
                  c_Q(E)=0\quad\Longleftrightarrow\quad c_Q^*(E)=0.
\]
The second assertion is an equality of null classes; it does not
claim equality of positive capacity values for every such countable
union.
\end{corollary}

\noindent\textbf{Status: \texttt{DIRECT\_COROLLARY}} of the
uniform growth bound and the already proved compact C1 theorem.
This is a qualified Borel outer-regularity result, not general C2.

\begin{proof}
The finite Borel measure
$\mathcal H^{n-1}|_E$ on the real core is Radon. For every
$\epsilon>0$ choose compact $K\subset E$ with
$\mathcal H^{n-1}(E\setminus K)<\epsilon$. Compact C1 and the
preceding proposition give
\[
 c_Q^*(E)
   \le c_Q^*(K)+c_Q^*(E\setminus K)
   \le c_Q(K)+C\epsilon
   \le c_Q(E)+C\epsilon.
\]
Let $\epsilon\downarrow0$ and combine with the always valid reverse
inequality. The exact raw and outer reductions then yield the
Berkovich statements. No uniform setwise continuity of individual
tester measures is asserted or needed.

For the last assertion write $E=\bigcup_iE_i$ with each $E_i$ Borel
and of finite Hausdorff measure. If $c_Q(E)=0$, then $c_Q(E_i)=0$;
the first part gives $c_Q^*(E_i)=0$. Countable subadditivity gives
$c_Q^*(E)=0$. The converse follows from $c_Q\le c_Q^*$.
\end{proof}

\paragraph{What remains.}
The proof uses a finite measure on the particular Borel set to
approximate it by a compact subset, together with an independent
outer estimate on the residual set. Merely knowing compactness of
the potential family supplies neither that residual estimate nor
ordinary-open outer continuity for all Borel sets. In particular the
argument leaves the general real Borel C2 gate, and a fortiori the
arbitrary-background Berkovich C2 gate, open.

% END source: sections/capacity_gauss_borel_regularity.tex


% BEGIN canonical R6 source: sections/capacity_gauss_projection_criterion.tex
\section{A coordinate-projection criterion for real and Gauss Borel capacities}
\label{sec:r6-gauss-projection-C2}

Retain $c_Q$, the Gauss section, and the smooth remaining-degree-one
background from Section~\ref{sec:r6-gauss-convex-capacity}, with $n\ge2$.
Fix a compact rectangular core $Q_0$ and open real boxes
\[
                  Q_0\Subset W\Subset W^+\Subset Q.
\]
Let $\operatorname{pr}_i:\mathbf R^n\to\mathbf R^{n-1}$ delete
the $i$th coordinate. Lebesgue outer measure on this target is denoted
by $\mathcal L^{n-1,*}$.

\begin{lemma}[Uniform projected-size bound]
\label{lem:r6-gauss-projection-outer-bound}
There is a finite constant $C=C(W,W^+,Q,n)$ such that every subset
$E\subset Q_0$ satisfies
\[
        c_Q^*(E)\le C\sum_{i=1}^n
                          \mathcal L^{n-1,*}(\operatorname{pr}_i E).
\]
The same inequality holds for the ordinary-open outer relative
capacities of $\sigma(E)$ and $\operatorname{trop}^{-1}(E)$.
\end{lemma}

\noindent\textbf{Status: \texttt{DIRECT\_COROLLARY}} of the
uniform convex slope bound, directional Fubini, and the exact outer
Gauss reduction. It is an upper bound, not an identification of
capacity with projected size.

\begin{proof}
All $g\in\mathcal G(Q)$ have one common Lipschitz bound $L$ on
$W^+$. Fix $i$. On each coordinate line with transverse parameter
$y\in\operatorname{pr}_iW$, the restriction $g_y$ is convex, and
its one-dimensional second derivative on the interval cut out by
$W$ is a positive measure of mass at most $2L$.

Directional distributional Fubini gives
\[
                  (\partial_{ii}g)|_W
                     =\int_{\operatorname{pr}_iW}g_y''\,dy.
\]
For completeness, for a compact smooth test $\varphi$ this follows
by writing $\int_Wg\,\partial_{ii}\varphi$ as the iterated integral
on these lines and applying the one-dimensional distributional
identity. The kernels are weakly Borel: their compact smooth pairings
are integrals of the continuous $g$ against second derivatives, or
limits of such pairings after restriction to compact line intervals.
The uniform mass bound makes the resulting positive Radon kernel
integral well defined and identifies the measures. In particular,
for every real open $G\subset W$,
\[
            \Delta g(G)\le2L\sum_i
                                   \mathcal L^{n-1}(\operatorname{pr}_iG).
\]

For $\varepsilon>0$ choose open $O_i\supset\operatorname{pr}_iE$
with
\[
 \mathcal L^{n-1}(O_i)
       \le\mathcal L^{n-1,*}(\operatorname{pr}_iE)+\varepsilon.
\]
The real open set $G=W\cap\bigcap_i\operatorname{pr}_i^{-1}(O_i)$
contains $E$. The preceding inequality is uniform in all normalized
convex testers and bounds $c_Q(G)$ by $2L\sum_i\mathcal L^{n-1}(O_i)$.
Let $\varepsilon\downarrow0$. The outer comparison proof in
Theorem~\ref{thm:r6-gauss-capacity-reduction} applies to arbitrary
subsets, so it gives both Berkovich formulations as well.
\end{proof}

\begin{proposition}[Borel coordinate sheets]
\label{prop:r6-borel-coordinate-sheet-C2}
Let $E\subset Q_0$ be Borel. Suppose that for each $i$ the restriction
$\operatorname{pr}_i|_E$ is one-to-one, its image $A_i$ is Borel,
and its inverse $s_i:A_i\to E$ is Borel. Then
\[
                              c_Q(E)=c_Q^*(E).
\]
The corresponding same-domain raw and ordinary-open outer relative
capacities agree on both $\sigma(E)$ and $\operatorname{trop}^{-1}(E)$.
The same conclusion holds for a finite union of such Borel sheets.
No Hausdorff-measure bound on $E$ is assumed.
\end{proposition}

\noindent\textbf{Status: \texttt{DIRECT\_COROLLARY}} of the
projection bound and compact C1. All Borel inverse hypotheses are
explicit, so no unproved measurable selection is being invoked.

\begin{proof}
For each $i$, push the finite measure $\mathcal L^{n-1}|_{A_i}$
forward by $s_i$, and call the resulting finite Borel measure on
$\mathbf R^n$ $\lambda_i$. Finite Borel measures on Euclidean space
are Radon. Put $\lambda_E=\sum_i\lambda_i$; it is carried by $E$.
For every compact $K\subset E$, injectivity gives
\[
 \operatorname{pr}_i(E\setminus K)=A_i\setminus\operatorname{pr}_iK,
 \qquad
 \mathcal L^{n-1}(\operatorname{pr}_i(E\setminus K))
                             =\lambda_i(E\setminus K).
\]
These projections are Borel, since $\operatorname{pr}_iK$ is compact.
Inner Radon regularity on the Borel set $E$ supplies $K\subset E$
compact with $\lambda_E(E\setminus K)<\varepsilon$. The projection
bound and compact C1 imply
\[
 c_Q^*(E)\le c_Q^*(K)+c_Q^*(E\setminus K)
       \le c_Q(K)+C\varepsilon\le c_Q(E)+C\varepsilon.
\]
Let $\varepsilon\downarrow0$ and use $c_Q\le c_Q^*$.

For a finite union of sheets, add their finitely many projection
pullback measures. For every residual $E\setminus K$, its coordinate
projections lie in the unions of the corresponding sheet projections,
so the same upper bound by $\lambda_E(E\setminus K)$ holds. The same
compact approximation proves the result. Finally apply the exact raw
and outer Gauss reductions.
\end{proof}

\paragraph{A directly specified class in dimension two.}
Let $A,B$ be Borel subsets of bounded real intervals, and let
$F:A\to B$ be a bijection such that both $F$ and $F^{-1}$ are Borel.
If its graph lies in $Q_0$, it satisfies the proposition: the inverse
coordinate maps are $x\mapsto(x,F(x))$ and
$y\mapsto(F^{-1}(y),y)$. Thus these Borel Gauss graphs and their full
tropical inverse images have exact ordinary-open outer regularity.
Continuity or a finite area assumption is not used.

\paragraph{Scope of the remaining gate.}
The finite Radon measure constructed here lives on the particular
Borel set and controls the outer capacity of its residual subsets.
It is not a common neighborhood dominator for all tester measures.
The proposition gives no such residual control for an arbitrary
Borel set, and leaves general real Borel C2, and general relative
Berkovich C2, open.

% END source: sections/capacity_gauss_projection_criterion.tex


% BEGIN canonical R6 source: sections/capacity_gauss_open_compact.tex
\section{A finite gluing mechanism for open sets with compact additions}
\label{sec:r6-gauss-open-compact}

Retain the bounded convex real domain $Q$, normalized convex class
$\mathcal G(Q)$, and Laplacian capacity $c_Q$ of
Section~\ref{sec:r6-gauss-convex-capacity}. All sets below lie in one
fixed compact core of $Q$. The normalization is on the whole $Q$.

\begin{lemma}[Pasting one local-uniform row on two fixed pieces]
\label{lem:r6-convex-local-uniform-pasting}
Let $g_j,g\in\mathcal G(Q)$ with $g_j\to g$ locally uniformly, and
let $K_1,K_2\Subset Q$ be disjoint compact sets. There are
$v_j\in\mathcal G(Q)$ and numbers $\epsilon_j\downarrow0$ such that,
for every Borel $A_i\subset K_i$,
\[
 \Delta v_j(A_1\cup A_2)
    \ge(1+\epsilon_j)^{-1}
                        \bigl(\Delta g(A_1)+\Delta g_j(A_2)\bigr).
\]
The same $v_j$ works for all these Borel subsets.
\end{lemma}

\noindent\textbf{Status: \texttt{DIRECT\_COROLLARY}} of the
explicit maximum-and-compensator method in frozen
Theorem~\textup{R4--79.4}. We give the real proof
because only local uniform convergence is assumed here, while the
normalization must remain global.

\begin{proof}
Choose smooth $0\le\chi_i\le1$ with disjoint compact supports in
$Q$, each identically one near $K_i$. A sufficiently large positive
quadratic $\rho$ on $\mathbf R^n$ is nonnegative on $Q$ and makes
$\rho+2\chi_i$ convex on the whole $Q$. Since $Q$ is bounded,
\[
                         A=\sup_Q\rho+2<\infty.
\]
Fix a compact neighborhood $H\Subset Q$ of the two regions where
the cutoffs equal one. Choose positive numbers $\delta_j\downarrow0$
such that $|g_j-g|\le\delta_j$ on $H$; this is possible by local
uniform convergence, enlarging the errors to decreasing tail suprema
if necessary. Set
\[
 v_j=\frac{\max\{g+\delta_j(\rho+2\chi_1),
                       g_j+\delta_j(\rho+2\chi_2)\}}
                  {1+A\delta_j}.
\]
Both branches are convex on the entire $Q$, and their maximum is
convex. The global bounds $0\le g,g_j\le1$ give $0\le v_j\le1$.
Closeness outside $H$ is neither assumed nor needed.

Near $K_1$ the first branch strictly exceeds the second by at least
$\delta_j$, and near $K_2$ the second strictly exceeds the first.
Ordinary locality of distributional Hessians therefore gives
\[
 \Delta v_j|_{K_1}
       \ge(1+A\delta_j)^{-1}\Delta g|_{K_1},\qquad
 \Delta v_j|_{K_2}
       \ge(1+A\delta_j)^{-1}\Delta g_j|_{K_2}.
\]
Here the positive added quadratic curvature is simply discarded.
The assertions first hold on neighborhoods and hence on every Borel
subset of the compact pieces. Their disjointness proves the lemma
with $\epsilon_j=A\delta_j$.
\end{proof}

\begin{theorem}[Borel outer regularity for an open set plus a compact set]
\label{thm:r6-convex-open-compact-C2}
Let $O\Subset Q$ be real open and let $K\Subset Q$ be compact.
Then
\[
                         c_Q(O\cup K)=c_Q^*(O\cup K).
\]
Consequently raw and ordinary-open outer relative capacities agree
on both $\sigma(O\cup K)$ and
$\operatorname{trop}^{-1}(O\cup K)$ in the same whole torus domain.
\end{theorem}

\noindent\textbf{Status: \texttt{NEW\_PROVED}} for this exact
set class, using the displayed finite gluing argument. Weak compactness
of an arbitrary measure family would not prove this assertion.

\begin{proof}
The case $K=\varnothing$ is immediate. Otherwise choose decreasing
open metric neighborhoods $N_j$ of $K$, with compact closures in
one fixed core, such that $\bigcap_j\overline{N_j}=K$.
They are cofinal among open neighborhoods of the compact $K$.
Thus $G_j=O\cup N_j$ are cofinal among open neighborhoods of
$E=O\cup K$ that contain $O$, and
\[
                         c_Q(G_j)\downarrow c_Q^*(E).
\]
Choose $g_j\in\mathcal G(Q)$ with
$\Delta g_j(G_j)\ge c_Q(G_j)-j^{-1}$.
The already proved local-uniform compactness of $\mathcal G(Q)$
gives a subsequence converging locally uniformly to
$g\in\mathcal G(Q)$. Keep this subsequence fixed for the rest of
the proof. Its Laplacians converge weakly on the fixed core, and
its displayed open-set masses still tend to $c_Q^*(E)$.

Now fix an arbitrary open neighborhood $N$ of $K$ with compact
closure in that core. Eventually $N_j\subset N$, so
\[
 \Delta g_j(G_j)
      \le\Delta g_j(O\setminus N)+\Delta g_j(\overline N).
\]
Compact Portmanteau implies
\begin{equation}
 c_Q^*(E)\le\limsup_j\Delta g_j(O\setminus N)
                                        +\Delta g(\overline N).
 \label{eq:r6-open-compact-upper}
\end{equation}

The sets $K_1=K$ and $K_2=\overline{O\setminus N}$ are fixed
disjoint compacts: $K$ lies inside the open $N$, while $K_2$ lies
in its closed complement. Apply the preceding lemma to these
pieces, with the Borel subsets $A_1=K$ and $A_2=O\setminus N$.
Since both subsets lie in $E$, it gives
\[
 c_Q(E)\ge\Delta g(K)+\limsup_j\Delta g_j(O\setminus N).
\]
This does not discard mass on $\partial N$: the second cutoff is
one on a neighborhood of the entire compact closure $K_2$.
Combining with \eqref{eq:r6-open-compact-upper} yields
\[
                 c_Q^*(E)\le c_Q(E)
                                +\Delta g(\overline N)-\Delta g(K).
\]
Only after taking the row limit do we shrink $N$ through nested
neighborhoods of $K$. The fixed finite Radon measure $\Delta g$
then makes the last difference tend to zero. This order of choices
is essential: the cutoff cost may grow as $N$ shrinks, but for each
fixed $N$ the local-uniform error $\delta_j$ already tends to zero.
The opposite inequality is always true. The exact raw and outer
Gauss reductions finish the Berkovich statement.
\end{proof}

\begin{corollary}[Adding a Borel set with controlled compact residuals]
\label{cor:r6-convex-open-residual-C2}
Let $O\Subset Q$ be open and let $B$ be Borel in one compact core.
Suppose that for every $\varepsilon>0$ there is compact $K\subset B$
with $c_Q^*(B\setminus K)<\varepsilon$. Then
\[
                            c_Q(O\cup B)=c_Q^*(O\cup B).
\]
This applies when $B$ has finite $\mathcal H^{n-1}$ measure, or
when $B$ is a finite union of the Borel coordinate sheets in
Proposition~\ref{prop:r6-borel-coordinate-sheet-C2}.
\end{corollary}

\noindent\textbf{Status: \texttt{DIRECT\_COROLLARY}}.

\begin{proof}
Outer subadditivity and the theorem give
\[
 c_Q^*(O\cup B)\le c_Q^*(O\cup K)+c_Q^*(B\setminus K)
                 <c_Q(O\cup B)+\varepsilon.
\]
Let $\varepsilon\downarrow0$. The two specified residual estimates
were proved in the Hausdorff and projection criteria, respectively.
\end{proof}

\paragraph{The unresolved general Borel step.}
For a general Borel $E$, its open neighborhoods need not differ from
$E$ only near one compact piece, and no compact-residual estimate has
been proved. The argument therefore does not establish continuity of
outer capacity from below on arbitrary increasing Borel sets. General
real Borel C2 and arbitrary-background C2 remain open.

% END source: sections/capacity_gauss_open_compact.tex


% BEGIN canonical R6 source: sections/capacity_energy_boundary.tex
\section{A boundary criterion for qualified C1 and C2 statements}
\label{sec:r6-energy-boundary-criterion}

\noindent\textbf{Status: \texttt{DIRECT\_COROLLARY}.}
Let $S$ be an arbitrary fixed original-positive symmetric separately
closed $(n-1,n-1)$-current on an ordinary norm-box domain $\Omega$
with the frozen faithful energy form. Put $F=\operatorname{supp}S$.
All capacities use normalized actual classical testers on all of
$\Omega$. Fix a compact core $K_0\Subset\Omega$.

For a Borel $E\subset K_0$, set
\[
              B_E=\partial_F(E\cap F),\qquad
              e_E=\operatorname{cap}_{\mathcal E,1}(B_E).
\]
Then
\begin{equation}
 0\le\operatorname{Cap}_S^*(E;\Omega)-\operatorname{Cap}_S(E;\Omega)
                                                     \le e_E.
 \tag{B1}
\end{equation}
The quantities on the left are finite by local CLN. In particular
$e_E=0$ gives Borel ordinary-open outer regularity for this $E$;
for compact $E$ it gives sharp compact C1.

\begin{proof}
Every normalized BT measure is carried by the closed $F$. An ordinary
open set $O\subset\Omega$ may be chosen with
$O\cap F=\operatorname{int}_F(E\cap F)$, and may be intersected
with a fixed relatively compact neighborhood of $K_0$. Its raw
capacity equals that of this relative interior and is at most the
raw capacity of $E$. The part of $E$ outside $F$ is contained in an
ordinary open set of zero capacity. Finally
\[
 E\cap F\subset\operatorname{int}_F(E\cap F)\cup B_E.
\]
Finite outer subadditivity and the sharp one-way inequality
\(\operatorname{Cap}_S^*(B_E;\Omega)\le e_E\)
from Corollary~\ref{cor:r6-sharp-energy-capacity} prove (B1).
Only relative topology on $F$ is used; the ambient boundary would
give an unnecessarily strong condition for polar-carried backgrounds.
\end{proof}

There is also a quantitative continuity-set statement. Suppose a net
of normalized measures $\mu_j=D(u_jS)$ converges weakly on all original
tests to a positive Radon measure $\nu$. Then
\begin{equation}
 \nu(E)-e_E\le\liminf_j\mu_j(E)
               \le\limsup_j\mu_j(E)\le\nu(E)+e_E.
 \tag{B2}
\end{equation}
Hence $e_E=0$ makes $E$ a simultaneous weak continuity set.

\begin{proof}
The fixed-core mass bounds and smooth scalar density turn the stated
weak convergence into ordinary local vague convergence. All measures
are carried by $F$, so this may be considered on $F$; compact continuous
functions there extend locally to the ambient space. The sharp Hardy
inequality passes to $\nu$ on each fixed original square test. Its
compact-core and ordinary-open argument therefore also gives
$\nu(B_E)\le e_E$. Apply the open and compact inequalities for vague
convergence to the relative interior and closure of $E\cap F$.
Their differences from $E\cap F$ lie in $B_E$, proving (B2).
One may insert a fixed original cutoff equal to one near $K_0$ to
reduce every step to finite measures on one compact set.
\end{proof}

These statements do not close general C1: the energy capacity of a
boundary can be positive even when its BT outer capacity is zero,
as Theorem~\ref{thm:r6-energy-BT-null-gap} demonstrates. They also
do not assert a converse boundary criterion. A fixed higher mixed
tail can replace $S$, with its own faithful energy form. For a family
with all higher slots moving, one would need a uniform estimate for
those varying tails; no such estimate is implicit here.

% END source: sections/capacity_energy_boundary.tex


% BEGIN canonical R6 source: sections/capacity_dominated_curvatures.tex
\section{Sharp regularity under continuous ambient-curvature domination}
\label{sec:r6-capacity-dominated-curvatures}

This criterion applies in every remaining degree and to an arbitrary
fixed background. The finite-dimensional curvature hypothesis from R5
is replaced by an order bound in an infinite-dimensional family. The
conclusion concerns that explicitly dominated family, or a weakly
convergent net lying in it; no uniform domination of all normalized
psh testers is assumed to exist.

Let $\Omega$ be an allowed ordinary open, let $S$ be symmetric,
original-positive and separately closed of bidegree $(p,p)$, and put
$q=n-p\ge1$. Fix an ordinary open $W\subset\Omega$ and a compact
core $K_0\Subset W$. All actual competitors below are defined and
normalized on the \emph{whole} domain $\Omega$.

\begin{lemma}[Ambient order gives mixed relative order]
\label{lem:r6-mixed-curvature-domination}
Let $1\le k\le q$. Suppose $u_1,\ldots,u_k\in\mathcal C_b(\Omega)$
are continuous on $W$, and $g_1,\ldots,g_k\in\mathcal C_b(W)\cap C(W)$.
If, for constants $M_i\ge0$,
\[
                     0\le D[u_i]|_W\le M_iD[g_i],
\]
then on every original complementary test in $W$ one has
\begin{equation}
 0\le T_{S,k}(u_1,\ldots,u_k)|_W
          \le\Bigl(\prod_{i=1}^kM_i\Bigr)
                    T_{S|_W,k}(g_1,\ldots,g_k).
 \label{eq:r6-mixed-curvature-domination}
\end{equation}
\end{lemma}

\noindent\textbf{Status: \texttt{DIRECT\_COROLLARY}} of R's continuous
regularization, expressed in Lemma~\ref{lem:r6-continuous-curvature-order},
and the frozen mixed BT calculus.

\begin{proof}
For any fixed positive separately closed tail $R$, the continuous
function $h_i=M_i g_i-u_i$ has positive original curvature on $W$.
The cited order-transfer lemma gives
\[
                       D(u_iR)\le M_iD(g_iR).
\]
This uses uniform smooth positive approximation of $h_i$, not a
claim that an arbitrary discontinuous difference with positive
distributional curvature is classical psh. Symmetry permits each
slot to be placed first, while the other bounded psh slots form a
positive separately closed tail. Replace the slots one at a time.
All intermediate tails retain positivity and separate closedness,
which proves \eqref{eq:r6-mixed-curvature-domination}, including
zero values of the $M_i$.
\end{proof}

\begin{theorem}[An all-degree dominated-family capacity criterion]
\label{thm:r6-dominated-curvature-capacity}
Fix continuous classical $g_1,\ldots,g_q$ on $W$ and constants
$M_1,\ldots,M_q\ge0$. For each $i$, let $\mathscr P_i$ be any
subfamily of actual $u\in\mathcal C_b(\Omega)$ satisfying
\[
 0\le u\le1\text{ on }\Omega,\qquad u|_W\in C(W),\qquad
                              D[u]|_W\le M_iD[g_i].
\]
Let $\mathscr M$ be the family of positive measures
$T_{S,q}(u_1,\ldots,u_q)$ with $u_i\in\mathscr P_i$, or any
subfamily of these measures. Define its restricted raw capacity by
\[
             c_{\mathscr M}(E)=\sup_{\mu\in\mathscr M}\mu(E),
             \qquad\sup\varnothing:=0.
\]
Its outer regularization uses ordinary open neighborhoods in the
same $\Omega$. Then:
\begin{enumerate}
\item For every Borel $E\subset K_0$,
      $c_{\mathscr M}(E)=c_{\mathscr M}^*(E)$.
\item If a net from $\mathscr M$ converges on all original compact
      smooth scalar tests in $W$ to a Radon measure $\mu$, then
      $\int b\,d\mu_j\to\int b\,d\mu$ for every bounded Borel
      $b$ supported in $K_0$. In particular
      $\mu(E)\le c_{\mathscr M}(E)$ for every such Borel $E$.
\item The same restricted capacities have continuity from above on
      decreasing Borel sequences in $K_0$. Together with their
      usual increasing-sequence continuity and Borel inner
      regularity, this supplies these finite compact-core regularity
      properties for the restricted family.
\end{enumerate}
No conclusion of convergence in total variation is made.
\end{theorem}

\noindent\textbf{Status: \texttt{DIRECT\_COROLLARY}} with an explicit
infinite-dimensional sufficient hypothesis. The transfer from
\emph{ambient} curvature bounds to a uniform measure bound for an
arbitrary $S$ is the substantive input.

\begin{proof}
On $W$ put
\[
                 \rho=\Bigl(\prod_{i=1}^qM_i\Bigr)
                          T_{S|_W,q}(g_1,\ldots,g_q).
\]
The preceding lemma proves $0\le\mu\le\rho$ on $W$ for every
$\mu\in\mathscr M$. Choose an ordinary $V$ with
$K_0\subset V\Subset W$. The positive Radon measure $\rho$ is
finite on $\overline V$.

For Borel $E\subset K_0$, Radon outer regularity supplies an ordinary
open $G$ with $E\subset G\subset V$ and
$\rho(G\setminus E)<\epsilon$. For every family member,
\[
            \mu(G)\le\mu(E)+\rho(G\setminus E)
                                      \le c_{\mathscr M}(E)+\epsilon.
\]
Take the supremum and then $\epsilon\downarrow0$. Since these are
ordinary neighborhoods in $\Omega$, this proves the first assertion.
The local nature of the dominator does not change the capacity domain.

For the second assertion, fixed nonnegative smooth tests first imply
$0\le\mu\le\rho$ on $W$. A bounded Borel $b$ supported in $K_0$
can be approximated in $L^1(\rho)$ by original smooth compact scalars
supported in $V$: use Radon regularity, compact continuous
approximation, and the original smooth scalar density on that compact
stage. For such an approximant $a$,
\[
 \left|\int(b-a)\,d\mu_j\right|
       +\left|\int(b-a)\,d\mu\right|
                      \le2\int|b-a|\,d\rho.
\]
The middle smooth pairing converges by hypothesis. Let the
$L^1(\rho)$ error tend to zero. This works for nets as well as
sequences and proves the asserted fixed-Borel convergence. It is
setwise convergence on the core, not uniform convergence over all
Borel indicators.

If $E_j\downarrow E$ are Borel subsets of $K_0$, domination gives
\[
       0\le c_{\mathscr M}(E_j)-c_{\mathscr M}(E)
                                 \le\rho(E_j\setminus E)\to0.
\]
For $E_j\uparrow E$, the two suprema over the index and over family
members commute, giving increasing continuity. For a Borel $E$,
inner regularity of each fixed member similarly gives
$c_{\mathscr M}(E)=\sup_{K\subset E\text{ compact}}c_{\mathscr M}(K)$.
These statements, with the preceding ordinary-open regularity and
finite core bound, prove the stated restricted regularity properties.
No continuity assertion for arbitrary decreasing Borel nets is used.
\end{proof}

\begin{corollary}[A maximal dominated subfamily and a sharp C1 test]
\label{cor:r6-maximal-dominated-capacity-family}
For fixed $g\in\mathcal C_b(W)\cap C(W)$ and $M\ge0$, set
\[
 \mathscr P(g,M;W)=\{u\in\mathcal C_b(\Omega):
         0\le u\le1\text{ on }\Omega,\ u|_W\in C(W),\
                                      D[u]|_W\le M D[g]\}.
\]
The restricted \emph{diagonal} capacity using this entire subfamily
is Borel outer regular on $K_0$, and its weak closure has no compact
mass gain. In particular, if normalized $u_j$ lie in this family and
$(Du_j)^qS\to\nu$ weakly on $W$, then for every compact
$K\subset K_0$,
\[
                  \nu(K)=\lim_j(Du_j)^qS(K)
                                      \le\operatorname{Cap}_S(K;\Omega).
\]
The mixed restricted capacities in the theorem are not identified
with diagonal capacities. If the unrestricted normalized family has
no such common ambient bound, no full C1 or C2 conclusion follows.
\end{corollary}

\begin{proof}
Use the diagonal subfamily of the preceding theorem, with all
$g_i=g$ and $M_i=M$. The last upper bound holds because every
member is still a normalized actual tester on all of $\Omega$.
\end{proof}

\begin{example}[Infinite-dimensional curvature and failure of TV convergence]
\label{ex:r6-dominated-curvature-not-tv}
On the original annulus $e^{-2}<|T|<e^2$ put $t=\log|T|$, take
$S=[1]$, and define
\[
 u_j(t)=\frac12+\frac{t^2}{16}-\frac{\sin(jt)}{32j^2},
 \qquad u(t)=\frac12+\frac{t^2}{16},
 \qquad g(t)=\frac{5t^2}{64}.
\]
These are smooth classical functions, $0<u_j<1$ throughout the
annulus, and
\[
 0\le D u_j
   =\left(\frac18+\frac{\sin(jt)}{32}\right)d't\wedge d''t
   \le Dg=\frac5{32}d't\wedge d''t.
\]
Their curvature currents span an infinite-dimensional real space.
They converge weakly, and setwise on every fixed compact core, to
$Du$. On the compact skeleton segment $-1\le t\le1$, however,
the total variation of the difference is
\[
             \frac1{32}\int_{-1}^1|\sin(jt)|\,dt
                                          \longrightarrow\frac1{8\pi}>0.
\]
All calibrations use $t=\log|T|$ with the original absolute value.
Thus the criterion extends beyond exact finite-dimensional curvature
families and does not falsely upgrade setwise convergence to TV.
\end{example}

% END source: sections/capacity_dominated_curvatures.tex


% BEGIN canonical R6 source: sections/variable_curvature_dominators.tex
\section{Variable curvature bounds and integrable relative density control}
\label{sec:r6-variable-curvature-dominators}

The continuous coefficient relation theorem allows upper bounds with
variable scalar coefficients. Its measure consequences require a
uniform integrability hypothesis, which is stated explicitly. The
result concerns the specified family of normalized testers, not all
normalized psh functions without that hypothesis.

\begin{proposition}[Variable continuous curvature bounds]
\label{prop:r6-variable-curvature-domination}
Fix any allowed $(p,p)$-background $S$ and $k\le n-p$. Suppose
$u_i,w_i\in\mathcal C_b(U)$ and $f_i\ge0$ are finite continuous
functions on $U$, with
\[
                         D[u_i]\le f_iD[w_i]
                              \quad(1\le i\le k)
\]
as original coefficient-Radon current inequalities. Then
\[
 0\le T_{S,k}(u_1,\ldots,u_k)
       \le \left(\prod_{i=1}^k f_i\right)
                                T_{S,k}(w_1,\ldots,w_k).
\]
Neither the $w_i$ nor the $u_i$ are assumed continuous. No closedness
of the continuously weighted upper current is required or asserted.
\end{proposition}

\noindent\textbf{Status: \texttt{DIRECT\_COROLLARY}} of
Theorem~\ref{thm:r6-bounded-curvature-continuous-module}.

\begin{proof}
For the positive separately closed tail containing all slots other
than $u_1$, apply the module theorem to
$f_1D[w_1]-D[u_1]\ge0$. This gives the displayed inequality with
$u_1$ replaced by $w_1$ and scalar factor $f_1$ outside the product.
Use symmetry, apply the same argument to the next remaining $u_i$
inside the still positive separately closed tail, and multiply the
result by the already extracted nonnegative scalar factors. Iterating
proves the assertion. At no step is a variable-coefficient current
used as a closed background.
\end{proof}

Hereafter let $q=n-p$ and fix an ordinary $V\Subset U$ on which
\[
                          \rho=T_{S,q}(w_1,\ldots,w_q)
\]
has finite mass. Let $\mathscr F$ be a nonempty family of actual
normalized tuples $(u_1,\ldots,u_q)$, each defined on the same whole
normalized domain, whose restrictions to $V$ satisfy the preceding
bounds with coefficients $f_{i,u}$. For some $1<r<\infty$ assume
\begin{equation}
 \sup_{u\in\mathscr F}
      \left\|\prod_{i=1}^qf_{i,u}\right\|_{L^r(V,\rho)}\le C<\infty.
 \label{eq:r6-variable-density-bound}
\end{equation}
Put $\mu_u=T_{S,q}(u_1,\ldots,u_q)$, and define the restricted raw
capacity $c_{\mathscr F}(E)=\sup_{u\in\mathscr F}\mu_u(E)$.
The corresponding star is the infimum over ordinary open neighborhoods
in the same normalized domain.

\begin{theorem}[Qualified C1, C2 and Borel-weighted convergence]
\label{thm:r6-integrable-density-capacity-stability}
Under \eqref{eq:r6-variable-density-bound}, the following hold.
\begin{enumerate}
\item For every Borel $E\subset V$,
\[
                    c_{\mathscr F}(E)
                               \le C\rho(E)^{1-1/r}.
\]
For every Borel $E\Subset V$, the restricted raw capacity equals
its ordinary-open outer capacity.
\item If a net of measures from this family converges weakly on every
original compact scalar test in $V$ to a Radon measure $\mu$, then
it converges setwise on every Borel $E\Subset V$. In particular
\[
                       \mu(E)\le c_{\mathscr F}(E)
\]
there, giving sharp compact no-gain without asserting that $\mu$ has
an actual normalized representative.
\item Write $r'=r/(r-1)$. On compact localizations in $V$, the
Radon--Nikodym densities of the net and its limit converge weakly
in $L^r(\rho)$. Consequently Borel weights converging strongly in
$L^{r'}(\rho)$ may vary jointly with the measures in their pairings.
Locally uniformly bounded weights converging locally in $\rho$-measure
satisfy this condition. This last assertion is valid for nets with
convergence in measure; mere pointwise convergence of nets is not
substituted for it.
\end{enumerate}
The analogous conclusions hold with a uniform $L^\infty$ density
bound, using $L^1$ weights and weak-star density convergence. No total-
variation convergence follows solely from these hypotheses.
\end{theorem}

\noindent\textbf{Status: \texttt{DIRECT\_COROLLARY}} of variable
curvature order, finite Radon regularity and elementary $L^r$ duality.
It enlarges the stated common bounded-density criterion; general
unrestricted arbitrary-background C1 and C2 are not asserted.

\begin{proof}
The preceding proposition gives
$0\le\mu_u\le F_u\rho$, where $F_u=\prod_i f_{i,u}$.
Thus its density $h_u$ satisfies $0\le h_u\le F_u$ and
$\|h_u\|_{L^r(V,\rho)}\le C$. Holder's inequality proves (1)'s
mass bound. If $E\Subset V$ is Borel, finite Radon outer regularity
gives ordinary open $G$ with $E\subset G\Subset V$ and
$\rho(G\setminus E)$ arbitrarily small. For every family member,
\[
      \mu_u(G)\le\mu_u(E)+C\rho(G\setminus E)^{1-1/r}.
\]
Take the supremum and then shrink this error. Ordinary outer
monotonicity gives the reverse inequality, proving the asserted C2.
The same argument allows opens in the original normalized domain,
since the displayed open neighborhoods already belong to it.

For (2) and (3), the closed positive $L^r(V,\rho)$ ball of radius
$C$ is weakly compact. Every subnet of densities has a weakly
convergent further subnet, with some nonnegative density $h$ of
norm at most $C$. Its integrals against every original compact
smooth scalar test equal the assumed current limit. Original smooth
scalar tests are uniformly dense on compact stages; the uniform
finite-mass bound and cutoffs therefore identify $h\rho=\mu$ on
$V$. Radon uniqueness also makes this density unique up to
$\rho$-null sets. The same identification on every convergent
subnet proves weak $L^r$ convergence of the whole density net.
This argument does not require a countable exhaustion or a selected
convergent lifting of the potentials.

In particular testing with $\mathbf1_E\in L^{r'}(V,\rho)$ gives
setwise convergence and the sharp capacity bound. If $a_j\to a$
strongly in $L^{r'}(V,\rho)$, then
\[
 \left|\int a_jh_j\,d\rho-\int ah\,d\rho\right|
 \le C\|a_j-a\|_{L^{r'}(V,\rho)}
                 +\left|\int a(h_j-h)\,d\rho\right|\longrightarrow0.
\]
For local statements multiply weights by a fixed compact cutoff
and apply the same estimate on its support. Locally bounded
convergence in measure gives strong $L^{r'}$ convergence on finite-
measure compact localizations: split into $|a_j-a|\le\epsilon$
and its complement, use the common bound on the latter, then let
$\epsilon\downarrow0$. This proof works for nets. The $L^\infty$
case uses the weak-star compact ball and the same argument with
$L^1$ tests and weights.
\end{proof}

\paragraph{Interpretation and limitations.}
The scalar functions $f_{i,u}$ enter through original ambient current
inequalities; the dominating Radon measure $\rho$ is the fixed
relative BT product on the actual space. Thus the hypothesis supplies
a common relative density control, including on carriers that may be
ambient pluripolar. It is not an assertion of nonconcentration on all
ambient polar sets. The existence of such coefficients with uniform
integrable products must be established for the intended tester
family; normalizing $0\le u_i\le1$ alone does not supply it.

% END source: sections/variable_curvature_dominators.tex


% BEGIN canonical R6 source: sections/capacity_laminar_rectangles.tex
\section{Diffuse fiber backgrounds and exact capacity on rectangles}
\label{sec:r6-laminar-rectangles}

This section gives an all-dimensional class of relative capacities on
which sharp ordinary-open outer approximation can be computed. The
parameter measure may be diffuse and may be carried by an ambient
complete pluripolar set. The set being tested is a rectangle; arbitrary
nonrectangular compact sets are not included in the conclusion.

Fix integers $p,q\geq1$, put $n=p+q$, and let $Y$ and $Z$ be the open
unit polydiscs of dimensions $q$ and $p$ over the original $K$.
Write $\Omega=Y\times_K Z$ and let $P\subset Z(K)$ be compact.
Here $Z(K)$ means actual $K$-rational analytic points, not all Berkovich
points. Let $\lambda$ be an arbitrary finite positive Radon measure on
$P$. For $a\in P$, denote the closed fiber embedding by
$i_a:Y\hookrightarrow\Omega$, $t\mapsto(t,a)$.

The set $\pi_Z^{-1}(P)$ is naturally homeomorphic to the ordinary
topological product $Y\times P$. Indeed every residue field at a point
of $P$ is exactly $K$, so each full analytic fiber is $Y$. On compact
affinoid stages the map $(t,a)\mapsto i_a(t)$ is continuous by uniform
polynomial approximation of analytic functions; for a polynomial the
assertion follows from the ultrametric estimate when its coefficients
vary with $a$. It is a continuous bijection from a compact space to a
Hausdorff space and hence a homeomorphism there. These stages cover
$Y\times P$. We use this identification for the Borel rectangles below.
It does not replace a general nonrigid analytic fiber by an independent
section.

\begin{proposition}[Original-current disintegration over rigid parameters]
\label{prop:r6-fiber-current-disintegration}
The formula
\begin{equation}
       \langle S,\eta\rangle
          =\int_P\left(\int_Y i_a^*\eta\right)d\lambda(a),
       \qquad \eta\in A_c^{q,q}(\Omega),
\label{eq:r6-laminar-current}
\end{equation}
defines a symmetric original-positive separately $d'$- and $d''$-closed
$(p,p)$-current, carried by $\pi_Z^{-1}(P)$.
For any $0\leq k\leq q$ and $u_1,\ldots,u_k\in\mathcal C_b(\Omega)$,
\begin{equation}
 \left\langle T_{S,k}(u_1,\ldots,u_k),\eta\right\rangle
   =\int_P
       \left\langle
          T_{[1]_Y,k}(u_1\circ i_a,\ldots,u_k\circ i_a),
                       i_a^*\eta\right\rangle d\lambda(a)
\label{eq:r6-fiber-bt-disintegration}
\end{equation}
for every original compact smooth complementary test $\eta$.
At $k=q$ this is an equality of positive Radon measure kernels on the
actual product $Y\times P$. In particular it holds on every Borel set,
not only on smooth tests or on a selected skeleton.
\end{proposition}

\noindent\textbf{Status: \texttt{DIRECT\_COROLLARY}} of the frozen
fixed-background calculus and CLN estimates, together with the original
fiber-integration theorem. The explicit disintegration proof is given
to fix the test class, measurability and calibration.

\begin{proof}
The projection $\pi_Z$ is separated, flat, boundaryless, and of pure
relative dimension $q$. Original fiber integration
\cite[Theorem 16.3.3(d), pp.~447--450]{CLD25} shows that
$a\mapsto\int_Y i_a^*\eta$ is continuous, in particular bounded on
the compact set $P$. Its integration in
\eqref{eq:r6-laminar-current} is therefore defined.
On a fixed compact original moment presentation, dominate a smooth
test by finitely many fixed positive controls as in R. Their fiber
integrals are continuous and bounded on $P$. This gives the requisite
finite coefficient seminorm bound for the displayed functional and
proves continuity in the original test topology. Pullback preserves the
original positive cone and symmetry, so the functional is symmetric
and positive. For a compact test of one lower degree, pullback commutes
with $d'$ and $d''$, and Stokes on the boundaryless fiber gives zero.
Thus $S$ is separately closed. A test supported off $\pi_Z^{-1}(P)$
has zero restriction on every fiber, proving the carrier assertion.
The fibers are the original $K$-space $Y$; there is no residue-degree
factor or absolute-value renormalization in this formula.

Here are the measurability and uniform bounds needed for the higher
formula. A compact test is covered by finitely many original local
regularization stages for the bounded potentials, with a smooth
partition. On each stage take the frozen simultaneous smooth
majorizing rows. After restriction to $i_a^{-1}$ of that stage they
are legitimate rows on the same original $K$-space $Y$.
The smooth-row fiber pairings are continuous in $a$ by the same
fiber-integration theorem. Their pointwise limits are the fiber BT
pairings by the already proved calculus with background $[1]_Y$.
They are consequently Borel functions of $a$.

The relative CLN proof on a finite original presentation is uniform
after restriction to these fibers: the real coefficient bounds and
cutoff derivatives are the same, while the finitely many background
control pairings are fiber integrals of fixed original smooth forms.
Those pairings are bounded on $P$ by fiber integration. The local
bounds of the $u_i$ on the fixed compact stage give one integrable bound
for all pairings under discussion. This argument permits degenerate
restricted moment frames; no strictly positive frame on every fiber
is assumed.

One can therefore define the right side of
\eqref{eq:r6-fiber-bt-disintegration} as a current $R_k$. It is positive
and separately closed by the corresponding fiber properties, and has
the finite coefficient controls just established. Thus it has the
actual coefficient-Radon property. For fixed tests its coefficient
measures are the integrals of the fiber coefficient measures: first
test this against original smooth scalar multipliers, then use
uniqueness of Radon measures and their bounded Borel multiplication.
The same statement is valid with a bounded Borel scalar on a compact
stage, including the actual $u_i$.

Induct on $k$. The assertion at zero is
\eqref{eq:r6-laminar-current}. Assuming it for the tail product, the
coefficient-measure statement and commutation of pullback with $D$ give
\[
 \begin{split}
 \langle D(u_1R_{k-1}),\eta\rangle
  &=\int_P\left\langle
       u_1\circ i_a\,
       T_{[1]_Y,k-1}(u_2\circ i_a,\ldots,u_k\circ i_a),
                              D(i_a^*\eta)\right\rangle d\lambda(a)\\
  &=\int_P\left\langle
       T_{[1]_Y,k}(u_1\circ i_a,\ldots,u_k\circ i_a),
                                      i_a^*\eta\right\rangle d\lambda(a).
 \end{split}
\]
This is precisely the fixed-background recursion. It proves the formula
on all original tests and hence equality with $T_{S,k}$.
For $k=q$, the fiber measures are weakly Borel measurable and locally
uniformly finite, as proved above. The usual Radon kernel integral on
$Y\times P$ is therefore defined. Its compact continuous pairings agree
with the current formula; original smooth scalar density identifies
the Radon measures, giving the Borel assertion.
\end{proof}

Write $C_Y(A)=\operatorname{Cap}_{[1]_Y}(A;Y)$ and
$C_S(E)=\operatorname{Cap}_S(E;\Omega)$, with normalized actual classical
testers on the indicated \emph{whole} domains. Stars denote infima over
ordinary open neighborhoods in those same domains.

\begin{theorem}[Sharp rectangular capacity and ordinary-open regularity]
\label{thm:r6-laminar-rectangle-capacity}
Fix an ordinary core $W\Subset Y$. For every Borel $A\subset W$ and
every Borel $B\subset P$, the rectangle $E=i(A\times B)$ satisfies
\begin{equation}
                  C_S(E)=\lambda(B)C_Y(A)=C_S^*(E).
\label{eq:r6-laminar-rectangle-capacity}
\end{equation}
In particular compact rectangles have sharp compact outer regularity,
in every dimension and every remaining degree $q\geq1$.
\end{theorem}

\noindent\textbf{Status: \texttt{DIRECT\_COROLLARY}} of the preceding
disintegration and the frozen \emph{ambient} Borel outer-regularity
theorem: R2, Corollary~\textup{R4--36.9} (Corollary 36.9,
printed p.~126 of the frozen R2 report). That theorem applies to $Y$,
an exhaustion polydisc, and to $A$ in its fixed relatively compact core.
It is not an assumption of general relative Borel outer regularity.

\begin{proof}
For any $0\leq u\leq1$ classical on all of $\Omega$, every restriction
$u\circ i_a$ is normalized classical on all of $Y$. Disintegration gives
\[
          \int_E(Du)^qS
              =\int_B\left(\int_A(D(u\circ i_a))^q\right)d\lambda(a)
              \leq\lambda(B)C_Y(A).
\]
Conversely a normalized classical $v$ on $Y$ pulls back to the normalized
classical $v\circ\pi_Y$ on all of $\Omega$. Its integral over $E$ is
$\lambda(B)\int_A(Dv)^q$. Taking suprema proves the raw equality,
including the case $\lambda(B)=0$. All quantities are finite here by
CLN on the fixed core and finiteness of $\lambda$.

For $\varepsilon>0$, the cited ambient theorem supplies an ordinary
open $O_Y$ with
\[
        A\subset O_Y\Subset Y,\qquad C_Y(O_Y)\leq C_Y(A)+\varepsilon.
\]
The finite measure $\lambda$ extended by zero from the compact $P$ to
$Z$ is Radon. Choose an ordinary open $O_Z\subset Z$ containing $B$
such that $\lambda(P\cap O_Z)\leq\lambda(B)+\varepsilon$.
The ordinary open set
\[
            G=\pi_Y^{-1}(O_Y)\cap\pi_Z^{-1}(O_Z)
\]
contains $E$. The same raw rectangle formula, and the carrier of $S$,
give
\[
 C_S^*(E)\leq C_S(G)
     =\lambda(P\cap O_Z)C_Y(O_Y)
     \leq(\lambda(B)+\varepsilon)(C_Y(A)+\varepsilon).
\]
Let $\varepsilon$ decrease to zero. The general inequality
$C_S(E)\leq C_S^*(E)$ gives the reverse bound. This proof does not
assert that an arbitrary nonsaturated open neighborhood contains one
fixed product neighborhood of a noncompact Borel rectangle.
\end{proof}

\begin{example}[A diffuse complete-pole parameter measure]
\label{ex:r6-diffuse-polar-parameters}
These backgrounds need not be nonpolar or countable sums of fiber
integration currents. For $p=1$, choose a uniformizer $\pi$ and put
\[
 P=\left\{\sum_{j=1}^{\infty}\epsilon_j\pi^{4^j}:
                          \epsilon_j\in\{0,1\}\right\}\subset Z(K).
\]
The coding by the binary product space is continuous and injective;
thus $P$ is compact. Give it the nonatomic Bernoulli probability
measure $\lambda$. With $\ell=-\log|\pi|$, set
\[
 \phi_N(z)=\int_P\max(\log|z-a|,-N)\,d\lambda(a),
               \qquad \phi=\lim_{N\to\infty}\phi_N.
\]
Every $\phi_N$ is continuous classical and nonpositive on $Z$.
Indeed the integrand is continuous jointly on compact stages, so
finite positive atomic quadratures converge uniformly on each such
stage; each quadrature is continuous classical. Decreasing closure
then makes $\phi$ extended classical.

For $z\in P$, the parameters agreeing with its first $j$ digits have
measure $2^{-j}$ and satisfy $|z-a|\leq|\pi|^{4^{j+1}}$. Hence
\[
       \int_P-\log|z-a|\,d\lambda(a)
                         \geq\ell\,4^{j+1}2^{-j}\longrightarrow\infty.
\]
Thus $\phi=-\infty$ on $P$. If an analytic point $z\notin P$, the
continuous function $a\mapsto|T-a|(z)$ is everywhere positive on the
compact set $P$ and has a positive minimum. Therefore $\phi(z)$ is
finite. This proves that $P$ is an exact complete pole set. For larger
$p$ embed $P$ as $P\times\{0\}^{p-1}$ and use the witness
$\max(\phi(z_1),\log|z_2|,\ldots,\log|z_p|)$.

The current in \eqref{eq:r6-laminar-current} is carried by the complete
pole set $Y\times P$, yet gives zero coefficient mass to each
individual fiber because $\lambda$ has no atoms. It is consequently
not a positive order limit of finite sums of nonzero subcurrents
carried by individual fibers. Formula
\eqref{eq:r6-laminar-rectangle-capacity} nevertheless applies exactly.
\end{example}

\paragraph{The next laminar gate.}
For a general compact $K\subset Y\times P$, disintegration gives only
the fiberwise upper bound by the integral of intrinsic capacities,
when that upper envelope is measurable. Choosing an optimizing tester
independently on uncountably many fibers does not supply one ambient
classical function normalized on $\Omega$. In particular a compact
graph of varying fiberwise Gauss points is not a rectangle, and no
claim of C1 or C2 for that graph follows from this section. This is a
concrete diffuse polar-supported test case for the remaining sharp
selection and extension mechanism.

% END source: sections/capacity_laminar_rectangles.tex


% BEGIN canonical R6 source: sections/capacity_laminar_graphs.tex
\section{Exact capacity on moving Gauss graphs over diffuse pole carriers}
\label{sec:r6-laminar-gauss-graphs}

Retain the original field, polydiscs, compact rigid parameter set $P$,
finite parameter measure $\lambda$, and disintegrated background $S$
of Section~\ref{sec:r6-laminar-rectangles}. The following computation
handles a nonrectangular family. Its moving atoms are followed by an
actual ambient psh threshold; no countable dense family of fixed fiber
testers is substituted for that construction.

For $L>0$ let $\zeta_{e^{-L}}\in Y$ be the diagonal Gauss point with
all $q$ coordinate radii $e^{-L}$. On the whole unit polydisc $Y$,
\begin{equation}
 \operatorname{Cap}_{[1]}(\{\zeta_{e^{-L}}\};Y)
 =\operatorname{Cap}_{[1]}(\{|t_1|,\ldots,|t_q|\le e^{-L}\};Y)
 =L^{-q}.
 \label{eq:r6-arbitrary-radius-capacity}
\end{equation}
This is the same exact computation as
Lemma~\textup{R4--60.1}, now with any real $L>0$:
$U=\max(-1,L^{-1}\max_i\log|t_i|)$ has
$(DU)^q=L^{-q}\delta_{\zeta_{e^{-L}}}$ by the original slope-simplex
calibration. The normalized tester $1+U$ gives the lower bound.
Comparison of $(1+\epsilon)U$ with an arbitrary normalized tester
minus one gives the upper bound on the closed polydisc; let
$\epsilon\downarrow0$. No change of the original absolute value,
strict-radius assumption, or outer-regularity interchange is involved.

\begin{theorem}[A sharp formula for moving Gauss graphs]
\label{thm:r6-laminar-gauss-graph-capacity}
Let $f\in\mathcal C_b(Z)\cap C(Z)$ satisfy $f\le0$ on all of $Z$
and
\[
                        0<m\le-f(a)\le M<\infty\qquad(a\in P).
\]
Suppose that $P$ admits finite partitions
$P=\bigsqcup_iP_{r,i}$ into compact sets such that each $P_{r,i}$
is an exact complete pole set of a nonpositive extended classical psh
function on $Z$, and
\[
              \max_i\operatorname{osc}_{P_{r,i}}f\longrightarrow0.
\]
For any Borel $B\subset P$ put
\[
 E_B=\{i_a(\zeta_{\exp f(a)}):a\in B\}\subset\Omega.
\]
Then, with normalized actual classical testers on the whole $\Omega$,
\begin{equation}
 C_S(E_B;\Omega)=C_S^*(E_B;\Omega)
                   =\int_B(-f(a))^{-q}\,d\lambda(a).
 \label{eq:r6-gauss-graph-capacity}
\end{equation}
In particular $E_P$ is compact and has sharp compact outer regularity.
The theorem holds for every $p,q\ge1$.
\end{theorem}

\noindent\textbf{Status: \texttt{DIRECT\_COROLLARY}} of the proved
disintegration, exact radial calibration, and frozen finite pole-carrier
additivity, with the explicit same-domain testers below. This computes
a new nonrectangular class; it does not prove general relative C1 or C2.

\begin{proof}
The map $a\mapsto\zeta_{\exp f(a)}$ is continuous: on every polynomial,
the Gauss seminorm is a finite maximum of continuous functions of the
radius. Together with the product identification in
Section~\ref{sec:r6-laminar-rectangles}, this makes $E_P$ a compact
graph homeomorphic to $P$ and $E_B$ a Borel subset of it. All its
$Y$-coordinates lie in the one fixed compact polydisc of radius
$e^{-m}<1$.

For a normalized ambient tester $u$, the fiber function $u\circ i_a$
is normalized on all of $Y$. The Borel kernel identity and
\eqref{eq:r6-arbitrary-radius-capacity} therefore give
\[
 \int_{E_B}(Du)^qS
  =\int_B (D(u\circ i_a))^q(\{\zeta_{\exp f(a)}\})\,d\lambda(a)
  \le\int_B(-f(a))^{-q}\,d\lambda(a).
\]
This proves the raw upper bound.

Fix one of the finite partitions and abbreviate its parts by $P_i$.
Let $S_i$ be the background obtained by integrating the same full
fibers against $\lambda|_{P_i}$. Proposition
\ref{prop:r6-fiber-current-disintegration} proves directly that each
$S_i$ is an allowed positive separately closed current, and
$S=\sum_iS_i$. The pullback of the pole witness of $P_i$ has exact
pole set $Y\times P_i$ and carries every coefficient of $S_i$.
These closed carriers are pairwise disjoint. Ordinary restriction and
locality make every BT measure for $S_i$ vanish on the carrier of
$S_l$ for $l\ne i$. Thus the hypotheses of the frozen
Theorem~\textup{R4--79.1} hold for these actual
diffuse currents, not merely for individual fiber integration currents.
It follows that
\begin{equation}
                       C_S(E_B;\Omega)=\sum_i C_{S_i}(E_B;\Omega).
 \label{eq:r6-graph-finite-additivity}
\end{equation}

Put $M_i=\max_{P_i}(-f)$ and $m_i=\min_{P_i}(-f)$, omitting empty
parts. The globally defined function
\[
 u_i(t,z)=\max\left\{0,\,1+M_i^{-1}
               \max\big(\log|t_1|,\ldots,\log|t_q|,f(z)\big)\right\}
\]
is continuous classical psh and satisfies $0\le u_i\le1$ on all of
$\Omega$. The inner maximum is finite because $f$ is finite; the
outer maximum supplies the global lower bound even away from $P_i$.
For $a\in P_i$ the outer floor is redundant, and exact radial
calibration on the fiber gives
\[
                   (D(u_i\circ i_a))^q
                              =M_i^{-q}\delta_{\zeta_{\exp f(a)}}.
\]
Consequently \eqref{eq:r6-graph-finite-additivity} implies
\begin{equation}
                    C_S(E_B;\Omega)
                         \ge\sum_i M_i^{-q}\lambda(B\cap P_i).
 \label{eq:r6-graph-lower-sum}
\end{equation}
The carrier theorem takes the independent finite suprema in one
same-domain normalized class; no existence of a single exactly
optimizing tester is asserted here.

For the outer upper bound set
$A_i=\{\zeta_{\exp f(a)}:a\in P_i\}\subset Y$.
It is compact, contains the Gauss point of the largest radius
$e^{-m_i}$, and is contained in the closed polydisc of that radius.
Equation~\eqref{eq:r6-arbitrary-radius-capacity} implies
$C_Y(A_i)=m_i^{-q}$. The Borel rectangles
$i(A_i\times(B\cap P_i))$ cover $E_B$. The already proved
rectangle theorem and finite subadditivity of ordinary-open outer
capacity therefore give
\begin{equation}
             C_S^*(E_B;\Omega)
               \le\sum_i C_S^*(i(A_i\times(B\cap P_i));\Omega)
               =\sum_i m_i^{-q}\lambda(B\cap P_i).
 \label{eq:r6-graph-upper-sum}
\end{equation}
The function $x\mapsto x^{-q}$ is uniformly continuous on $[m,M]$.
The partition oscillation condition makes the difference of the upper
and lower sums tend to zero, bounded by $\lambda(P)$ times that
uniform modulus. Both sums converge to the displayed integral in
\eqref{eq:r6-gauss-graph-capacity}. The inequalities
\eqref{eq:r6-graph-lower-sum}, $C_S\le C_S^*$, and
\eqref{eq:r6-graph-upper-sum} prove the theorem.
\end{proof}

\begin{corollary}[Diffuse complete-polar Cantor parameters]
\label{cor:r6-cantor-gauss-graphs}
For the explicit parameter set $P$ in
Example~\ref{ex:r6-diffuse-polar-parameters}, the preceding theorem
applies to every continuous classical $f\le0$ satisfying the stated
strict negative bounds on $P$, and to every finite positive Radon
measure $\lambda$ on $P$, including nonatomic ones.
\end{corollary}

\begin{proof}
Partition $P$ by its first $r$ binary digits. Each part is compact,
and the diameters of the parts tend to zero, so continuity of $f|_P$
gives the required uniform oscillation condition. Each part has an
exact complete-pole witness: use the same logarithmic integral as in
Example~\ref{ex:r6-diffuse-polar-parameters}, with the Bernoulli
probability conditioned on that cylinder. The divergent estimate is
unchanged up to its fixed positive normalization, and compactness
gives finiteness at every analytic point outside that cylinder.
For higher $p$, include the other zero coordinates by the same finite
maximum. This construction depends on the parameter set, not on the
chosen measure $\lambda$.
\end{proof}

\paragraph{Why this does not give arbitrary fiberwise optimization.}
A weakly dense or countable family of fixed fiber testers can miss all
the atoms of a moving graph. The proof above instead follows its
radius through the actual joint psh function $f(z)$, and only the
normalizing slopes are discretized. For a general compact
$K\subset Y\times P$, no such joint psh threshold is supplied.
The existence of measurable fiberwise optimizers, even if established,
would still not produce normalized ambient classical testers. Neither
general compact no-gain nor actual representability of every completed
relative measure follows from the graph formula.

% END source: sections/capacity_laminar_graphs.tex


% BEGIN canonical R6 source: sections/capacity_cantor_borel_graphs.tex
\section{Exact interpolation on a pole Cantor set and Borel Gauss graphs}
\label{sec:r6-cantor-borel-graphs}

Let $Z_1$ be the ordinary open unit disc over the original field $K$,
and retain the compact set
\[
 P=\left\{\sum_{j\ge1}\epsilon_j\pi^{4^j}:
                         \epsilon_j\in\{0,1\}\right\}\subset Z_1(K)
\]
from Example~\ref{ex:r6-diffuse-polar-parameters}. A cylinder means
the subset obtained by prescribing finitely many initial digits.
The following interpolation statement concerns this specified set.
It is not an extension theorem for arbitrary pluripolar sets.

\begin{lemma}[Continuous pole wells on cylinders]
\label{lem:r6-cantor-continuous-wells}
Let $P_1,\ldots,P_s$ be the nonempty cylinders of one finite partition
of $P$. There are pairwise disjoint ordinary open discs $N_i\supset P_i$
with the following property. For every $\kappa>0$ there are continuous
classical functions $w_i$ on the whole $Z_1$ such that
\[
 -1\le w_i\le0,\qquad w_i|_{P_i}=-1,
 \qquad -\kappa\le w_i\le0\quad\hbox{on }Z_1\setminus N_i.
\]
\end{lemma}

\noindent\textbf{Status: \texttt{DIRECT\_COROLLARY}} of the
explicit logarithmic pole witnesses, with the uniform bound recorded.

\begin{proof}
The finite cylinders have strictly smaller internal diameters than
their mutual distances. Choose a disc around each cylinder with radius
strictly between these two quantities. The discs are ordinary open,
pairwise disjoint, and contain the respective cylinders.

Let $\tau_i$ be the Bernoulli probability conditioned on $P_i$ and put
\[
              \phi_i(z)=\int_{P_i}\log|T-a|(z)\,d\tau_i(a).
\]
The construction in Example~\ref{ex:r6-diffuse-polar-parameters}
shows that $\phi_i\le0$ is extended classical and has exact pole set
$P_i$. It is continuous and finite off $P_i$: on a compact neighborhood
of any such point the continuous positive function $|T-a|(z)$ has a
uniform positive lower bound for $a\in P_i$, so the logarithmic
integral is continuous there. Since $\phi_i$ is upper semicontinuous,
\[
                         w_{i,N}=\max\{\phi_i/N,-1\}
\]
is locally constant equal to $-1$ near every point of $P_i$, and is
continuous everywhere. It is a bounded classical function by the
maximum and decreasing-closure properties.

If $N_i$ has center $b_i$ and radius $r_i$, then for $z\notin N_i$
and $a\in P_i$ one has $|T-a|(z)\ge r_i$, by the ultrametric
inequality and $|a-b_i|<r_i$. Thus
$\phi_i(z)\ge\log r_i$ uniformly on the whole complement of $N_i$.
Taking $N$ sufficiently large makes $w_{i,N}\ge-\kappa$ there.
This proves the claimed uniform, same-domain wells.
\end{proof}

\begin{theorem}[Continuous interpolation with arbitrary range slack]
\label{thm:r6-cantor-classical-interpolation}
For every $F\in C(P,\mathbf R)$ and every $\varepsilon>0$, there is
an actual bounded continuous classical function $f$ on the whole
$Z_1$ such that
\[
 f|_P=F,\qquad
 \min_P F-\varepsilon\le f\le\max_P F+\varepsilon
                              \quad\hbox{on }Z_1.
\]
For $P\subset Z_1\times\{0\}^{p-1}\subset Z$, the same assertion
holds on the whole $p$-dimensional unit polydisc $Z$, by pullback.
\end{theorem}

\noindent\textbf{Status: \texttt{NEW\_PROVED}} for this explicit
pole Cantor set. The proof uses the displayed continuous wells and a
uniformly convergent decreasing construction; it does not infer an
ambient representative from relative capacity convergence.

\begin{proof}
Write $a=\min_PF$, $b=\max_PF$, and start with the constant classical
function $v_0=b+\varepsilon/4$. Put
\[
                    \eta_j=\frac{\varepsilon}{64}\,8^{-j}.
\]
We construct continuous classical $v_j$ decreasing on $Z_1$, with
continuous residuals $e_j=v_j|_P-F$. Initially
$e_0\ge\varepsilon/4=16\eta_0$.
At step $j$ assume $e_j\ge2\eta_j$, and let
$H_j=\max_Pe_j$. Choose a finite cylinder partition such that the
oscillation of $e_j$ on each part $P_i$ is at most $\eta_j$.
Define
\[
                    a_i=\min_{P_i}e_j-\eta_j>0.
\]
Use the preceding lemma with pairwise disjoint discs $N_i$, choosing
$\kappa>0$ so small that $sH_j\kappa\le\eta_j/2$. Set
\[
                       v_{j+1}=v_j+\sum_{i=1}^s a_iw_i.
\]
This is continuous classical and at most $v_j$.

At a point of $P_i$, the $i$th well is exactly $-1$ and every other
well lies in $[-\kappa,0]$. Therefore
\[
                    \eta_j/2\le e_{j+1}\le2\eta_j.
\]
In particular $e_{j+1}\ge4\eta_{j+1}$, so the induction continues.
At any point of $Z_1$ at most one of the disjoint $N_i$ occurs.
The corresponding well contributes at most $H_j$ to the absolute
decrement; all others contribute together at most $\eta_j/2$.
Consequently
\[
       0\le v_j-v_{j+1}\le H_j+\eta_j/2\quad\hbox{on }Z_1.
\]
For $j\ge1$ one has $H_j\le2\eta_{j-1}$. The tail decrements are
uniformly summable, and $v_j$ converges uniformly on the entire disc
to a bounded continuous function $f$. Decreasing classical closure
with bounded limit makes this actual function classical. The residual
bounds give $f|_P=F$ exactly.

The upper bound is $f\le v_0=b+\varepsilon/4$. Also
\[
 \sum_{j\ge0}\|v_j-v_{j+1}\|_\infty
 \le H_0+\frac{20}{7}\eta_0,
 \qquad H_0=b-a+\varepsilon/4.
\]
Hence $f\ge a-(20/7)\eta_0>a-\varepsilon$. This proves the range
claim. Pullback by the first coordinate projection gives the final
assertion and preserves the original-field classical curve condition.
\end{proof}

\begin{corollary}[Every bounded Borel radial graph in this lamination]
\label{cor:r6-borel-cantor-gauss-graphs}
Let $Y,Z,\Omega=Y\times_KZ$ and
$S=\int_P[Y\times\{a\}]\,d\lambda(a)$ be as in
Section~\ref{sec:r6-laminar-rectangles}, with this explicit embedded
Cantor set $P$ and any finite positive Radon measure $\lambda$ on $P$.
Let $F:P\to[-M,-m]$ be Borel, where $0<m\le M<\infty$.
For every Borel $B\subset P$, define
\[
 E_B=\{i_a(\zeta_{\exp F(a)}):a\in B\}.
\]
Then, in every $p,q\ge1$, with normalized actual classical testers
on the same whole $\Omega$,
\[
              C_S(E_B;\Omega)=C_S^*(E_B;\Omega)
                     =\int_B(-F(a))^{-q}\,d\lambda(a).
\]
\end{corollary}

\noindent\textbf{Status: \texttt{DIRECT\_COROLLARY}} of the
new interpolation theorem, the proved continuous-threshold graph
formula, Lusin approximation, and the rectangle formula. The datum
$F$ need not be continuous or the trace of a preassigned psh function.

\begin{proof}
The real radius interval $[e^{-M},e^{-m}]$ parametrizes a compact
Gauss arc in $Y$. The graph of the Borel map $a\mapsto\exp F(a)$
is Borel in its product with the compact metrizable $P$. The original
rigid-fiber product identification makes $E_B$ a Borel subset of a
fixed compact core of $\Omega$.

For any normalized ambient tester, its fiber restrictions are
normalized on the entire $Y$. The Borel disintegration and exact
singleton capacity in \eqref{eq:r6-arbitrary-radius-capacity} give
the upper bound
\[
                       C_S(E_B)\le I_B,
       \qquad I_B=\int_B(-F(a))^{-q}\,d\lambda(a).
\]
Fix $\delta>0$. Finite Radon inner regularity and Lusin's theorem
give a compact $K\subset B$ with
$\lambda(B\setminus K)<\delta$ such that $F|_K$ is continuous.
Extend it continuously to $\widetilde F:P\to[-M,-m]$ by the
Tietze extension theorem on the compact metrizable space $P$.
Apply Theorem~\ref{thm:r6-cantor-classical-interpolation} with
range slack smaller than $m/2$. It gives a continuous classical
$f\le-m/2<0$ on the whole $Z$ with $f|_P=\widetilde F$.
Every prefix partition consists of exact complete pole sets and its
$f$-oscillation tends to zero. Thus the already proved
Theorem~\ref{thm:r6-laminar-gauss-graph-capacity}, applied to the
Borel parameter subset $K$, gives
\[
              C_S(E_K)=C_S^*(E_K)=I_K.
\]

Let $A$ be the compact Gauss arc of radii in
$[e^{-M},e^{-m}]$. Its intrinsic capacity is $m^{-q}$: it contains
the largest-radius Gauss point and is contained in the closed
polydisc of that radius. The remaining graph is contained in the
Borel rectangle $i(A\times(B\setminus K))$. The rectangle theorem
and outer subadditivity therefore imply
\[
 C_S^*(E_B)\le I_K+m^{-q}\lambda(B\setminus K)
                          \le I_B+m^{-q}\delta.
\]
On the other hand $C_S(E_B)\ge C_S(E_K)=I_K\ge I_B-m^{-q}\delta$.
Let $\delta\downarrow0$ and use $C_S\le C_S^*$. Every comparison
uses the original whole capacity domain; no fiberwise supremum is
interchanged with an integral.
\end{proof}

\paragraph{The remaining selection problem.}
This result treats one scalar moving radial threshold over the explicit
pole Cantor set. It does not extend an arbitrary family of intrinsic
psh functions to a jointly psh ambient function, identify the capacity
of arbitrary Borel subsets of $Y\times P$, or settle unrestricted C1
or C2 for general backgrounds. The successful mechanism is exact
scalar interpolation followed by finite-measure Lusin localization;
weakly dense lists of fixed fiber testers would not follow these
moving atoms.

% END source: sections/capacity_cantor_borel_graphs.tex


% BEGIN canonical R6 source: sections/capacity_laminar_borel_tubes.tex
\section{Filled Borel radial tubes and their intermediate sets}
\label{sec:r6-laminar-borel-tubes}

Retain the explicit pole Cantor lamination, the finite measure $\lambda$,
the background $S$, and the same whole capacity domain $\Omega$ from
Corollary~\ref{cor:r6-borel-cantor-gauss-graphs}. Let
$F:P\to[-M,-m]$ be Borel, $0<m\le M<\infty$, and let $B\subset P$
be Borel. Write $E_B$ for its diagonal Gauss graph and set
\[
 H_B=\{i_a(t):a\in B,\ \max_{1\le l\le q}|t_l|
                                      \le\exp F(a)\}.
\]
The latter is the filled moving closed polydisc on each parameter
fiber. Both sets are Borel and lie in the fixed compact core consisting
of the radius-$e^{-m}$ closed polydisc times $P$.

\begin{corollary}[Exact capacity of filled tubes]
\label{cor:r6-filled-borel-radial-tubes}
Every Borel set $E$ with $E_B\subset E\subset H_B$ satisfies
\[
                 C_S(E;\Omega)=C_S^*(E;\Omega)
                       =\int_B(-F(a))^{-q}\,d\lambda(a).
\]
In particular this holds for the entire filled tube $H_B$, in every
$p,q\ge1$.
\end{corollary}

\noindent\textbf{Status: \texttt{DIRECT\_COROLLARY}} of the
Borel graph formula and the same-domain rectangle formula. There is
no new fiberwise selection or extremal-function assumption.

\begin{proof}
For each mesh size $\delta>0$, partition the compact interval $[m,M]$
into finitely many disjoint Borel intervals $I_i$ of diameter at most
$\delta$. Put
\[
 B_i=\{a\in B:-F(a)\in I_i\},\qquad s_i=\inf I_i\ge m,
 \qquad A_i=\{t\in Y:\max_l|t_l|\le e^{-s_i}\}.
\]
These are Borel parameter pieces and fixed compact intrinsic
polydiscs. Their rectangles cover $H_B$. The rectangle theorem and
the exact radial calibration give
\[
 C_S^*(H_B;\Omega)
       \le\sum_i C_S^*(i(A_i\times B_i);\Omega)
       =\sum_i s_i^{-q}\lambda(B_i).
\]
Since $x\mapsto x^{-q}$ is uniformly continuous on $[m,M]$, the
right side differs from $\int_B(-F)^{-q}\,d\lambda$ by at most
$\lambda(P)$ times its modulus of continuity at $\delta$. Let
$\delta\downarrow0$. The graph inclusion and the proved graph
formula give the matching raw lower bound. For any displayed $E$,
\[
 C_S(E_B)\le C_S(E)\le C_S^*(E)\le C_S^*(H_B),
\]
so every term equals the claimed integral.

The Borel assertion follows from the actual rigid-fiber product
identification and the Borel inequality between the continuous
function $\max_l|t_l|$ and the parameter function $\exp F(a)$.
Every tester used in the rectangle and graph inputs is normalized on
the original whole domain $\Omega$.
\end{proof}

This computation includes nonsaturated Borel intermediate sets whenever
they contain the Gauss graph. It does not identify the capacity of an
arbitrary Borel subset of the lamination or settle general relative C2.

% END source: sections/capacity_laminar_borel_tubes.tex


% BEGIN canonical R6 source: sections/capacity_upper_envelope_obstruction.tex
\section{Why an ambient upper contact cannot recover relative compact mass}
\label{sec:r6-capacity-upper-envelope-obstruction}

The following example eliminates a specific proposed route to compact
outer regularity. It is not a counterexample to compact outer regularity
itself. The completed representative and its ambient upper envelope must
be kept separate even after the potential-quotient theorem of R5.

\begin{proposition}[Loss of all relative mass under an ambient upper envelope]
\label{prop:r6-upper-envelope-relative-mass-loss}
Let $p,q\geq1$, $n=p+q$, and let $\Omega$ be the open unit polydisc
over the original $K$, with coordinates $(t_1,\ldots,t_q,z_1,\ldots,z_p)$.
Put $Y=\{z_1=\cdots=z_p=0\}$ and $S=[Y]$.
There are continuous classical $a_j\in\mathcal C_b(\Omega)$ with
$0\leq a_j\leq1$ and $a_j\uparrow a$ such that:
\begin{enumerate}
\item $a$ belongs to the bounded relative completion for $S$, and
      $\mu_{a_j}=(Da_j)^qS=\mu_a$ is a fixed nonzero positive measure;
\item the ambient upper semicontinuous regularization is $a^*=1$,
      so $(Da^*)^qS=0$;
\item at every point of $Y$, no normalized upper semicontinuous
      majorant $b\leq1$ of $a$ on an ordinary neighborhood can contact
      $a$. Such a $b$ is identically one on that part of $Y$.
\end{enumerate}
In particular, if the majorant is classical, its relative top measure
vanishes there. Nevertheless, $\mu_a$ itself is realized by a globally
normalized continuous classical potential on $\Omega$.
\end{proposition}

\noindent\textbf{Status: \texttt{DIRECT\_COROLLARY}} of the frozen
increasing completion and integration-current compatibility, with an
explicit example. The upper-contact strategy it disproves is recorded
as \texttt{FAILED\_ATTEMPT}.

\begin{proof}
Choose a uniformizer $\pi$, write $\lambda=-\log|\pi|>0$, and fix
$0<\delta<1/2$. Set
\[
 b_i(t_i)=\max\{0,1+\lambda^{-1}\log|t_i|\},\qquad
 h(t)=\delta+\frac{1-2\delta}{q}\sum_{i=1}^q b_i(t_i).
\]
These are actual continuous classical functions, with
$\delta\leq h\leq1-\delta$. Define
\[
 \psi(z)=\max_{1\leq l\leq p}\log|z_l|,\qquad
 a_j(t,z)=\max\{h(t),1+j^{-1}\psi(z)\},
\]
with $\psi=-\infty$ on $Y$. The function $\psi$ is extended classical
and nonpositive. Its logarithmic singularity is removed by the finite
continuous lower branch $h$ in $a_j$. Hence $a_j$ is continuous
classical, $\delta\leq a_j\leq1$, and $a_j$ increases. Its actual limit is
\[
                a=h\quad\hbox{on }Y,\qquad
                a=1\quad\hbox{on }\Omega\setminus Y.
\]
The fixed-background increasing-limit theorem puts this actual function
in the relative completion. For each classical row, restriction to the
smooth analytic subspace gives
\[
               (Da_j)^q[Y]=i_*\bigl((Dh)^q\bigr),
\]
where $i:Y\hookrightarrow\Omega$. Completed convergence gives exactly
the same measure for $a$. On the original product calibration each
$Db_i$ has mass $\lambda^{-1}$ at the radius-$|\pi|$ Gauss point in its
coordinate. Thus, if $y\in Y$ is the Gauss point with all these radii,
\[
       \mu_a=q!\left(\frac{1-2\delta}{q\lambda}\right)^q\delta_y.
\]
Only nonvanishing is needed below; the formula also makes it explicit.

Every point of $Y$ is an ordinary limit of points outside $Y$.
For example, after passage to its completed residue field, take the
normal Gauss polydisc of radius $r>0$ centered at zero and let $r$ tend
to zero; its images in $\Omega$ approach the original point and lie
outside $Y$. Therefore $a^*=1$ at every point of $Y$, and it is already
one off $Y$. If $b\leq1$ is an upper semicontinuous majorant of $a$
on an ordinary neighborhood, upper semicontinuity along the same points
forces $b=1$ on its intersection with $Y$. Since $h\leq1-\delta$,
there is no contact with $a$ there. When $b$ is classical,
integration-current compatibility and ordinary locality imply
$(Db)^q[Y]=0$ on that neighborhood.

Finally the function $h(t)$, pulled back by the coordinate projection,
is globally normalized and continuous classical on $\Omega$, and its
relative measure is precisely the same $\mu_a$. The example therefore
invalidates upper-contact recovery for this representative, without
invalidating the raw-capacity inequality for the measure.
\end{proof}

\paragraph{The precise excluded inference.}
One cannot start from a completed potential $a$, replace it by $a^*$ or
by normalized ambient classical majorants, and use a contact inequality
to conclude that the mass of $(Da)^qS$ on a compact set survives. In the
example the contact set is empty near its positive atom. Any successful
proof of general compact outer regularity must instead choose different
testers, or supply an additional mechanism that transfers their compact
mass while retaining the normalization on the same whole domain.

% END source: sections/capacity_upper_envelope_obstruction.tex


\part{Relative Dirichlet problems and their constraints}

% BEGIN canonical R6 source: sections/dirichlet_energy_separation.tex
\section{A dual solvability criterion for the completed degree-one Dirichlet problem}
\label{sec:r6-dirichlet-energy-separation}

Throughout this section $K/\mathbf Q_p$ is the original finite extension,
with its prescribed absolute value. Let $\Omega$ be an ordinary algebraic
norm-box domain relatively compact in a smooth affine analytification of
pure dimension $n\ge1$. Fix an arbitrary symmetric original-positive,
separately $d'$- and $d''$-closed $(n-1,n-1)$-current $S$.
All tests are original compact smooth CLD tests. In particular, no
smoothness, nonpolarity or geometric origin of $S$ is assumed.

Use the faithful measure $\lambda$ and closed coercive form
$(\mathcal E_S,\mathscr D_S)$ of
Theorem~\textup{R4--84.3}. Choose
\[
       L\Subset W\Subset\Omega,
\]
where $L$ is compact and lies in a smaller norm box and $W$ is ordinary
open. Write $H=\mathscr D_{S,0}(W)$, with Hilbert norm
$\|x\|_H^2=\mathcal E_S(x,x)$.
Fix the smooth classical outer potential $h$ and common bound $B$ from
Section~\textup{R4--80}. Let $\mathscr P$ be any
nonempty convex subfamily of
\[
 \mathscr P_{h,L}^B
   =\{u\in\mathcal C_b(\Omega): |u|\le B,\ u=h
                                      \text{ on }\Omega\setminus L\}.
\]
Let $\mathscr A$ be its bounded global closure for all-domain local
outer $S$-capacity convergence, modulo local relative-null equality.
Thus membership means approximation by this specified classical family,
with its one bound and its actual exterior values. For $a\in\mathscr A$
put $\mu_a=D(aS)$ and $\mu_h=D(hS)$.

\subsection{Identification of the constraint in the faithful energy space}

\begin{proposition}[The completed collar constraint is an energy-closed convex set]
\label{prop:r6-completed-energy-convex-set}
The map
\[
         a\longmapsto x_a=a-h\quad\text{in }L^2(\lambda)
\]
identifies $\mathscr A$ with the nonempty closed bounded convex set
\[
          C=\overline{\{u-h:u\in\mathscr P\}}^{\,H}\subset H.
\]
It is injective on the indicated relative classes and satisfies
\begin{equation}
 \|x_a-x_b\|_H^2=I_S(a,b)
      =\int(a-b)\,d(\mu_b-\mu_a).
 \label{eq:r6-energy-completion-isometry}
\end{equation}
For every original compact smooth scalar $\phi$ supported in $W$,
\begin{equation}
 \mathcal E_S(x_a,\phi)
        =-\int\phi\,d(\mu_a-\mu_h).
 \label{eq:r6-energy-current-pairing}
\end{equation}
Each class retains an actual representative bounded by $B$ and equal to
$h$ on $\Omega\setminus L$.
\end{proposition}

\noindent\textbf{Status: \texttt{NEW\_PROVED}} for the identification
of this completed constraint with a closed convex subset of the faithful
form. Its ingredients are the previously proved energy and completion
theorems.

\begin{proof}
For $u\in\mathscr P$, the actual difference $u-h$ vanishes off $L$.
Choose an original smooth cutoff equal to one near $L$ and supported in
$W$. Proposition~\textup{R4--86.1}, including its
support-preserving core approximations, puts $u-h$ in $H$.
The compact Green identity gives
\[
       \|u-v\|_H^2=I_S(u,v)\qquad(u,v\in\mathscr P).
\]
Relative CLN and the common bound give one uniform finite bound on
$I_S(u,h)$, so these vectors form a bounded subset of $H$.

Suppose a defining row $u_j\in\mathscr P$ converges in all-domain local
outer capacity to a bounded completed potential $a$. The completed
weighted calculus, applied to the four terms defining $I_S(u_j,u_k)$,
shows that $I_S(u_j,u_k)\to0$ as $j,k\to\infty$. In each of these
four weighted terms insert the fixed original compact smooth cutoff
equal to one near $L$ and supported in $W$. The full signed energy is
unchanged because $u_j-u_k$ vanishes off $L$; the individual terms are
thereby legitimate compact test pairings. No finite total mass on all
of $\Omega$ or weak convergence against a Borel indicator is used.
For complete detail,
failure would give indices $j_l,k_l\to\infty$ with a fixed positive
lower bound; all four weighted terms along those two rows have the same
completed limits with cancelling signs, a contradiction. Thus
$u_j-h$ is Cauchy in $H$.

Its $L^2(\lambda)$ limit is $a-h$. Indeed the faithful reference
potential is normalized and classical on $\Omega$, so on the fixed
compact $L$ its measure satisfies
$\lambda(E)\le\operatorname{Cap}_S(E;\Omega)$ for Borel $E$.
Capacity convergence, the common bound, and a positive-gap splitting
therefore give $L^2(\lambda)$ convergence. Off $L$ all these differences
are zero. Passing the energy identity through two defining rows proves
\eqref{eq:r6-energy-completion-isometry}.

Conversely, an $H$-Cauchy sequence $u_j-h$ with $u_j\in\mathscr P$
is capacity Cauchy by the fixed-collar energy-capacity estimate. Reverse
framed domain comparison gives the all-domain local condition.
The bounded completion theorem supplies a limit in $\mathscr A$;
the preceding paragraph identifies its $H$ limit. Hence the energy
closure is precisely the displayed set $C$. Convexity follows either
from the defining family or from its energy closure. If two classes
have the same energy vector, their energy difference is zero;
the same energy-capacity estimate makes every positive gap outer-null.
This proves injectivity. The actual limsup representative from the
bounded completion retains the common bound and exterior equality.

Finally Green's identity holds for each classical defining row by
Proposition~\textup{R4--86.1}. Its left side passes by
strong form convergence, and its right side by completed weak
convergence on the fixed original test. This proves
\eqref{eq:r6-energy-current-pairing}.
\end{proof}

\subsection{An explicit family of linear inequalities on the datum}

Let $V\Subset W$ be ordinary open, or let $V=W$.
Set $M=\mathscr D_{S,0}(V)$, regarded as a closed subspace of $H$.
The choice $V=W$ is allowed, with $M=H$.
For an original compact smooth scalar $\phi$ supported in $V$, define
the support functional of the prescribed classical constraints by
\begin{equation}
 \Sigma_{\mathscr P}(\phi)
    =\sup_{u\in\mathscr P}
          \left[-\int_\Omega (u-h)\,D(\phi S)\right].
 \label{eq:r6-classical-support-functional}
\end{equation}
This is finite. It involves only the prescribed classical family and
the fixed coefficient measures $D(\phi S)$; it contains no unknown
solution and no unknown completed Monge--Ampere measure.

A finite positive Radon datum $\nu$ on $V$ is called finite-energy here
if $\delta=\nu-\mu_h|_V$ satisfies
\begin{equation}
 \|\delta\|_{-1,V}
  :=\sup_{\substack{\phi\in A_c^{0,0}(V;\mathbf R)\\
                         \mathcal E_S(\phi,\phi)\le1}}
             \left|\int_V\phi\,d\delta\right|<\infty.
 \label{eq:r6-negative-energy-norm}
\end{equation}
This condition includes annihilation of energy-zero core functions.
For example, $\delta=f\lambda|_V$ with $f\in L^2(V,\lambda)$ satisfies
it, with norm at most $\sqrt2\|f\|_2$, by coercivity.
Neither $D(aS)\ll\lambda$ for general bounded potentials nor an equality
of energy-null and BT-null ideals is assumed.

\begin{theorem}[A smooth-test alternative for completed solvability]
\label{thm:r6-energy-dirichlet-alternative}
For a finite-energy datum $\nu$, the following are equivalent.
\begin{enumerate}
\item There is $a\in\mathscr A$ with $D(aS)|_V=\nu$ on all original
compact smooth scalar tests, equivalently as Radon measures on $V$.
\item For every original real compact smooth $\phi$ supported in $V$,
\begin{equation}
       -\int_V\phi\,d(\nu-\mu_h)
                         \le\Sigma_{\mathscr P}(\phi).
 \label{eq:r6-smooth-test-feasibility}
\end{equation}
\end{enumerate}
If these conditions fail, one such original smooth scalar test gives a
strict reverse inequality. Thus nonsolvability has a single smooth-test
separation certificate.

If the defining family additionally has one fixed actual exterior
$u=e$ on $\Omega\setminus V$, then the solution class is unique.
If its defining rows are trapped between fixed continuous barriers on
$\overline V$ having the same boundary value $g$, the selected solution
representative has the ordinary boundary limit $g$. These exterior and
barrier conditions are hypotheses on the nonempty classical family;
their availability for arbitrary boundary data is not asserted.
\end{theorem}

\noindent\textbf{Status: \texttt{NEW\_PROVED}} for this exact
finite-energy criterion over arbitrary fixed $S$. The separating-space
argument is elementary Hilbert-space duality. The criterion retains the
full prescribed classical constraints through the explicit functional
\eqref{eq:r6-classical-support-functional}; it is not a bounded-subsolution
criterion and does not assert that every positive datum is feasible.

\begin{proof}
The datum functional
$\phi\mapsto-\int\phi\,d\delta$ extends continuously to $M$.
Let $y\in M$ be its Riesz vector:
\[
          \mathcal E_S(y,\phi)=-\int\phi\,d\delta
                                    \quad(\phi\in A_c^{0,0}(V)).
\]
Let $P_M:H\to M$ be orthogonal projection and put $C_V=P_M C$.
The set $C$ is closed, bounded and convex in a Hilbert space, hence
weakly compact. Its continuous linear image $C_V$ is weakly compact
and convex, hence norm closed. In particular it is not necessary to
assume that the image of an arbitrary unbounded closed convex set
under a linear map is closed.

For every $\phi\in M$ let
$\sigma_{C_V}(\phi)=\sup_{z\in C_V}\mathcal E_S(z,\phi)$.
On the original smooth core this equals
$\Sigma_{\mathscr P}(\phi)$, by the energy closure defining $C$ and
Green's identity. Both support functionals are Lipschitz in the energy
norm, with constant $\sup_{x\in C}\|x\|_H$. Core density therefore
extends \eqref{eq:r6-smooth-test-feasibility} to all $\phi\in M$.

If $y\in C_V$, all these inequalities hold. If $y\notin C_V$, let
$z$ be its nearest point in this closed convex set. The elementary
projection inequality is
\[
       \mathcal E_S(y-z,b-z)\le0\qquad(b\in C_V).
\]
For $\phi=y-z$ it gives
\[
 \mathcal E_S(y,\phi)-\sigma_{C_V}(\phi)
                              \ge\|y-z\|_H^2>0.
\]
Approximating $\phi$ in the energy norm by an original smooth core
function preserves a strict positive gap. This proves the stated
separation alternative and shows that the inequalities are equivalent
to $y\in C_V$.

If $y=P_Mx_a$ for some $a\in\mathscr A$, equation
\eqref{eq:r6-energy-current-pairing} gives the required equality on all
original tests in $V$. Conversely such equality gives $P_Mx_a=y$ by
core density. This proves the equivalence. Original smooth scalars
determine the Radon measures by the same original compact
Stone--Weierstrass argument already used for the strong data metric.

For uniqueness under fixed actual exterior values, choose the actual
representatives supplied by the common defining rows. Then $a-b=0$
at every point of $\Omega\setminus V$, including $\partial V$.
The completed weighted product is therefore exactly
\[
 (a-b)(\mu_b-\mu_a)=\mathbf1_V(a-b)(\mu_b-\mu_a)=0
\]
when the two interior equations agree. Thus $I_S(a,b)=0$, and
\eqref{eq:r6-energy-completion-isometry} and the capacity estimate give
equality of the relative classes. No equality of their boundary
measures is used. The actual limsup construction preserves the exterior
values and both barriers, so continuity of the barriers at each ordinary
boundary point gives the asserted boundary limit.
\end{proof}

\subsection{Best feasible data and a quantitative certificate}

\begin{corollary}[The energy-distance defect is exactly a support gap]
\label{cor:r6-energy-data-projection}
With the preceding notation, set $d=\operatorname{dist}_H(y,C_V)$.
Then
\begin{equation}
 d=\sup_{\substack{\phi\in A_c^{0,0}(V;\mathbf R)\\
                            \|\phi\|_H\le1}}
       \left[-\int_V\phi\,d(\nu-\mu_h)
                         -\Sigma_{\mathscr P}(\phi)\right]
   =\min_{a\in\mathscr A}\|\nu-\mu_a|_V\|_{-1,V}.
 \label{eq:r6-exact-separation-distance}
\end{equation}
The minimum is attained. Its image measure is unique; under the fixed
actual exterior hypothesis its completed potential class is unique.
The retraction onto this feasible image is firmly nonexpansive in the
negative-energy norm. More precisely, if $y,y'$ are the Riesz vectors
of two data and $z,z'$ their nearest points in $C_V$, then
\[
             \|z-z'\|_H^2\le\mathcal E_S(z-z',y-y').
\]
For $V=W$, the feasible potential vectors themselves are the projections
$z,z'$, so their energy distance is at most the data distance.
For a proper interior domain the assertion concerns the projected
vectors and the interior data; no unproved energy bound for the inverse
of $P_M|_C$ is asserted.
\end{corollary}

\begin{proof}
For any $z\in C_V$ and $\|\phi\|_H\le1$, the bracket in
\eqref{eq:r6-exact-separation-distance} is at most
$\mathcal E_S(y-z,\phi)\le\|y-z\|_H$.
For the nearest point $z$, if $y\ne z$, the direction
$\phi=(y-z)/\|y-z\|_H$ attains equality in the Hilbert space.
Approximation by the smooth core, with rescaling to the unit ball,
gives the same supremum over the displayed original tests. If $y=z$,
every bracket is nonpositive and the zero test gives zero.
The current pairing identifies the negative-energy norm of a data
difference with $\|y-P_Mx_a\|_H$, proving the minimum formula.
Every point of $C_V$ has a lift in $C$ by its definition, so attainment
gives an actual completed potential. Hilbert projection is unique,
and original scalar tests distinguish the corresponding measures.
Adding the two projection inequalities gives the displayed firm
nonexpansiveness. The uniqueness and full-domain claims then follow
from the theorem and from $P_M=\mathrm{id}$ when $V=W$.
\end{proof}

When the datum also lies in the admissible class of R4 and $V=W$, this
projection agrees with the previous variational retraction. Indeed for
$x=a-h$ the degree-one Aubin--Mabuchi energy is
\[
 \mathcal E_{\rm AM}(a)=\int x\,d\mu_h-\tfrac12\|x\|_H^2,
\]
so its datum objective equals
$-\tfrac12\|x-y\|_H^2+\tfrac12\|y\|_H^2$.
For completed potentials, this identity follows by classical defining
rows and their strong energy limits; when actual datum integrals are
used, R4's capacity absolute continuity permits the same limit.
The present finite-energy construction does not require that extra
absolute continuity merely to define the Hilbert projection.

\paragraph{Frozen mixed slots.}
For an arbitrary original-positive closed $(p,p)$-background $S_0$,
put $q=n-p\ge1$ and fix $v_2,\ldots,v_q\in\mathcal C_b(\Omega)$.
Apply this entire section to
$R=T_{S_0,q-1}(v_2,\ldots,v_q)$. It gives the same exact criterion for
the linear mixed equation
\[
          T_{S_0,q}(a,v_2,\ldots,v_q)|_V=\nu.
\]
Here $a$ belongs to the completed constraint constructed for this fixed
degree-one tail $R$, and the displayed left side means $D(aR)$.
This does not assert that $a\in\mathcal A_{S_0,b}$ or that all of its
slots may instead be varied in a jointly completed product over $S_0$.
This is a qualified all-dimensional mixed result with one free slot.
It does not turn the diagonal equation with $q>1$ into a linear problem,
and it does not establish an unrestricted bounded-subsolution theorem.

% END source: sections/dirichlet_energy_separation.tex


% BEGIN canonical R6 source: sections/dirichlet_general_data_alternative.tex
\section{The feasibility inequalities themselves force finite energy}
\label{sec:r6-general-data-alternative}

\noindent\textbf{Status: \texttt{DIRECT\_COROLLARY}} of
Theorem~\ref{thm:r6-energy-dirichlet-alternative}. This removes an a
priori data assumption from the solvability alternative. Finite energy
is still required for the quantitative Hilbert projection of arbitrary
infeasible data.

Retain its fixed degree-one background, prescribed nonempty convex
classical family $\mathscr P$, completed constraint $\mathscr A$, and
ordinary interior domain $V$. Let $\nu$ be any positive Radon measure
on $V$, initially only locally finite. Define
$\delta=\nu-\mu_h|_V$ on compact subsets and use the same support
functional $\Sigma_{\mathscr P}$ on original smooth compact scalars.
Then the following are equivalent:
\begin{enumerate}
\item $D(aS)|_V=\nu$ for some $a\in\mathscr A$;
\item for every real original $\phi\in A_c^{0,0}(V)$,
\[
             -\int_V\phi\,d\delta\le\Sigma_{\mathscr P}(\phi).
\]
\end{enumerate}
If they hold, $\nu$ is in fact finite on $V$, $\delta$ has finite
negative-energy norm, and all boundary and uniqueness conclusions of
the preceding theorem apply. If they fail, there is an original compact
smooth scalar giving a strict reverse inequality.

\begin{proof}
Set $R=\sup_{x\in C}\|x\|_H<\infty$. By the support-functional
identity and Cauchy--Schwarz,
\[
                   \Sigma_{\mathscr P}(\pm\phi)\le R\|\phi\|_H.
\]
Applying the assumed feasibility inequality to both signs therefore
gives
\[
                    \left|\int_V\phi\,d\delta\right|
                                        \le R\|\phi\|_H.
\]
Thus the datum defines a bounded functional on the original form core,
annihilates its null vectors, and extends to the part space. The Riesz
and projection proof of Theorem~\ref{thm:r6-energy-dirichlet-alternative}
now applies verbatim: its construction uses the datum only on these
compact tests until it produces a completed potential. It gives
$D(aS)|_V=\nu$ as Radon measures. The left side has finite mass on
$V$, since $\overline V$ lies in the fixed compact collar stage and
relative CLN applies to the common defining bound. Hence $\nu(V)<\infty$.

Conversely, a completed solution has its energy vector in $C$ and hence
satisfies every displayed support inequality. The strict failure
alternative is the logical negation of these inequalities; for a
finite-energy infeasible datum it is also supplied by the explicit
Hilbert separation proof. If the datum has infinite negative-energy
norm, the uniform bound $R$ shows directly that at least one signed
original smooth test violates the inequality.
\end{proof}

This is not a replacement of the condition by positivity or by a bounded
subsolution. The right side retains the full prescribed ambient
classical constraints, and the toric certificate shows that those
constraints exclude some smooth positive dominated data. The same
corollary applies to a fixed mixed tail exactly in the qualified sense
stated at the end of Section~\ref{sec:r6-dirichlet-energy-separation}.

% END source: sections/dirichlet_general_data_alternative.tex


% BEGIN canonical R6 source: sections/dirichlet_energy_data.tex
\section{Measure data and bounded subsolutions in the relative energy class}
\label{sec:r6-relative-energy-data}

This section deliberately changes the admissible solution class. The
solutions belong to the faithful relative energy space; they need not
be ambient classical psh or belong to a bounded classical capacity
completion. This distinction is essential even for smooth nonpolar
backgrounds.

Fix an arbitrary original-positive symmetric separately closed
$(n-1,n-1)$-current $S$ on an ordinary norm-box neighborhood $\Omega$
in the allowed space. Let $F=\operatorname{supp}S$, let $\lambda$ and
$(\mathcal E,\mathscr D)$ be the faithful measure and coercive regular
form of Theorem~\textup{R4--84.3}, and fix an
ordinary open $V\Subset\Omega$. Put
\[
                  H=\mathscr D_0(V),\qquad
                  \|u\|_H=\mathcal E(u,u)^{1/2}.
\]
The part-space closure is the energy boundary condition. No pointwise
boundary trace is built into that definition.

For a signed Radon measure $\sigma$ on $V$, define its possibly infinite
core dual norm by
\[
 \|\sigma\|_{-1,V}
   =\sup\left\{\left|\int_V\phi\,d\sigma\right|:
       \phi\in A_c^{0,0}(V;\mathbf R),\ \|\phi\|_H\le1\right\}.
\]
Finiteness includes vanishing on zero-energy core functions. Total mass
on $V$ need not be finite for this definition.

\begin{lemma}[Positive finite-energy functionals preserve order]
\label{lem:r6-positive-energy-data}
Suppose a positive Radon measure $\tau$ has finite core dual norm.
Its core integral extends to a bounded functional $\ell_\tau$ on $H$,
and $\ell_\tau(v)\ge0$ for every $v\in H$ with $v\ge0$
$\lambda$-almost everywhere. If $\tau_1\le\tau_2$ and both have
finite core dual norm, their extended functionals have the same order.

For every real $v\in H$,
\[
                    \mathcal E(v,v^+)\ge\mathcal E(v^+,v^+).
\]
\end{lemma}

\begin{proof}
Core density gives the unique bounded extension. To verify positivity
without assuming an integral formula for arbitrary energy vectors,
take core $v_j\to v$ in $H$. Smooth nonnegative normal contractions
approximating $t^+$ give nonnegative core functions $w_j$ whose energies
are bounded by those of $v_j$, and with $w_j-v_j^+\to0$ in $L^2$.
For example choose the smoothing error small enough on the finite
$\lambda$-mass of each individual core support. Since $v_j^+\to v$
in $L^2$, every weakly convergent subsequence of the bounded $w_j$ in
$H$ has limit $v$. Their nonnegative functional values therefore give
$\ell_\tau(v)\ge0$. Apply this argument to $\tau_2-\tau_1$ for
the order assertion.

The normal-contraction property from the frozen form theorem gives
$v^+,v^-\in H$ and
$\mathcal E(|v|,|v|)\le\mathcal E(v,v)$. Expanding the two quadratic
forms yields $\mathcal E(v^+,v^-)\le0$. Since $v=v^+-v^-$, the
displayed inequality follows. Thus no extra strong-local chain rule is
needed for the order argument.
\end{proof}

\begin{theorem}[Finite-energy measure Dirichlet problem]
\label{thm:r6-measure-energy-dirichlet}
Every signed Radon measure $\sigma$ with finite $\|\sigma\|_{-1,V}$
has a unique $u_\sigma\in H$ satisfying
\[
                 \mathcal E(u_\sigma,\phi)=-\ell_\sigma(\phi)
                                      \qquad(\phi\in H).
\]
Its coefficient-defined original current satisfies
\[
                       D(u_\sigma S)=\sigma\quad\hbox{on }V
\]
on every original compact smooth scalar test. The solution map is
linear and
\[
 \|u_\sigma\|_H=\|\sigma\|_{-1,V},\qquad
 \|u_\sigma-u_{\sigma'}\|_H=\|\sigma-\sigma'\|_{-1,V}.
\]
For positive data the solution is nonpositive almost everywhere.
If $0\le\sigma\le\sigma'$ and both norms are finite, then
$u_{\sigma'}\le u_\sigma\le0$ almost everywhere and for the canonical
energy-quasicontinuous representatives outside an energy-polar set.
\end{theorem}

\noindent\textbf{Status: \texttt{DIRECT\_COROLLARY}} of the R4 faithful
coercive form, now applied to finite-energy Radon data instead of only
$L^2(\lambda)$ densities. It is not an ambient psh existence theorem.

\begin{proof}
Riesz representation in the energy Hilbert space gives the solution,
uniqueness, and both equalities. Coefficient domination and coercivity
make $u_\sigma S$ a well-defined current: on every compact stage its
coefficient variations are bounded by a fixed multiple of $\lambda$,
and $u_\sigma\in L^2(\lambda)\subset L^1_{\rm loc}(\lambda)$.
Approximation by the original core passes compact Green's identity to
\[
 \mathcal E(u_\sigma,\phi)
          =-\langle D(u_\sigma S),\phi\rangle
                         \qquad(\phi\in A_c^{0,0}(V;\mathbf R)).
\]
This proves the asserted full original-test equation.

For positive $\sigma$, test the weak equation with $u_\sigma^+$.
The preceding lemma gives
\[
 0\le\mathcal E(u_\sigma^+,u_\sigma^+)
 \le\mathcal E(u_\sigma,u_\sigma^+)
                  =-\ell_\sigma(u_\sigma^+)\le0.
\]
Coercivity implies $u_\sigma\le0$. Apply the identical argument to
$u_{\sigma'}-u_\sigma$, whose datum is positive, for order comparison.
A zero vector has zero canonical quasicontinuous representative outside
an energy-polar set; normal contraction gives the stated representative
version of these inequalities. This does not authorize arbitrary
changes on merely $\lambda$-null sets.
\end{proof}

\begin{theorem}[Bounded subsolution domination in the energy class]
\label{thm:r6-bounded-energy-subsolution}
Suppose $w\in H\cap L^\infty(\lambda)$ and its original curvature
$\tau=D(wS)$ on $V$ is a positive Radon measure. Then every positive
Radon measure $0\le\mu\le\tau$ has finite energy and admits the unique
solution from the preceding theorem, with
\[
                w\le u_\mu\le0,\qquad
                \|u_\mu\|_H\le\|w\|_H,\qquad
                \|u_\mu\|_\infty\le\|w\|_\infty.
\]
The order and bound first hold almost everywhere, and hold for canonical
energy representatives outside an energy-polar set.

If in addition $w$ has a continuous energy representative $w_0$ on
$V\cap F$ extending to a bounded continuous function on $\overline V$,
with $w_0\le0$ on $V$ and $w_0=0$ on the ordinary boundary, then
$u_\mu$ has a bounded actual Borel representative $u_0$ on $V$ such
that
\[
                       w_0\le u_0\le0
\]
at every point. Consequently $u_0$ has ordinary boundary limit zero.
It represents the same current and canonical relative energy class.
No interior continuity or classical psh property is asserted.
\end{theorem}

\noindent\textbf{Status: \texttt{DIRECT\_COROLLARY}} with a complete
proof of the datum bound. This is a bounded-subsolution criterion in
the larger relative energy class; the corresponding unrestricted claim
in the bounded classical completion remains false.

\begin{proof}
For a real core function $\phi$ set
$P_\epsilon(t)=\sqrt{t^2+\epsilon^2}-\epsilon$.
It is nonnegative, smooth, vanishes at zero, and has derivative of
absolute value at most one. Thus $P_\epsilon(\phi)$ is a nonnegative
original core function and its energy is at most that of $\phi$.
The original current equation for $w$ and compact Green give
\[
 \int P_\epsilon(\phi)\,d\tau
       =-\mathcal E(w,P_\epsilon(\phi))
                         \le\|w\|_H\|\phi\|_H.
\]
Let $\epsilon\downarrow0$ using monotone convergence on the compact
support. Since $\mu\le\tau$,
\[
       \left|\int\phi\,d\mu\right|
          \le\int|\phi|\,d\mu
          \le\int|\phi|\,d\tau
                         \le\|w\|_H\|\phi\|_H.
\]
This proves the required dual bound without assuming finite total mass
near $\partial V$. The same argument shows that $\tau$ has finite
energy. Its energy solution is exactly $w$, by uniqueness and the
original core equation. The order comparison gives $w\le u_\mu\le0$
and the stated bounds.

For the last assertion take the canonical energy representative on
$F$. Outside one energy-polar exceptional set it is finite and lies
between $w_0$ and zero. Assign value zero on that exceptional set and
off $F$. This gives a Borel representative with the displayed pointwise
bounds. The exceptional set is $\lambda$-null, hence these changes
preserve the coefficient-defined current; by the one-way energy-to-BT
capacity theorem they also preserve the canonical relative capacitary
class. The squeeze between $w_0$ and zero proves the ordinary boundary
limit. The role of $w_0$ is explicit; a part-space boundary condition
alone was not turned into an ordinary pointwise boundary assertion.
\end{proof}

\begin{proposition}[Relative Dini-capacity growth supplies finite energy]
\label{prop:r6-dini-relative-energy-data}
Let $H_0\Subset V$ be compact and let $\mu$ be a positive Radon measure
supported in $H_0$, of mass $M>0$. Let $\omega:(0,\infty)\to[0,\infty)$
be finite and nondecreasing, and assume
\[
                    \int_0^r\omega(s)\,\frac{ds}{s}<\infty
                                      \quad\hbox{for some }r>0.
\]
Suppose for every Borel $E\subset H_0$,
\[
 \mu(E)\le A\,\operatorname{Cap}_S(E;\Omega)
                \omega(\operatorname{Cap}_S(E;\Omega)),
\]
where the right side is zero when the capacity is zero. Then $\mu$ has
finite energy and hence a unique relative energy Dirichlet solution.
If $C_{H_0}$ is the constant in the frozen one-way capacity comparison,
then
\begin{equation}
 \|\mu\|_{-1,V}
 \le\sqrt3\,C_{H_0}
       \inf_{r>0}\left\{\frac Mr+A\int_0^r\omega(s)\,\frac{ds}{s}\right\}.
 \label{eq:r6-dini-energy-norm}
\end{equation}
For $\omega(s)=s^\alpha$, $\alpha>0$, the right side is at most
\[
 \sqrt3\,C_{H_0}(1+\alpha^{-1})
                  A^{1/(1+\alpha)}M^{\alpha/(1+\alpha)}.
\]
These hypotheses alone do not assert boundedness, ambient psh, or
membership in the bounded classical completion of the solution.
\end{proposition}

\noindent\textbf{Status: \texttt{DIRECT\_COROLLARY}} of
Theorem~\textup{R4--85.2}, by the following layer-cake
calculation. The two capacities are not identified.

\begin{proof}
For an original real core function $\phi$, put $e=\|\phi\|_H$.
The energy capacity definition, applied to $|\phi|/t$, and coercivity
give
\[
 \operatorname{cap}_{\mathcal E,1}(\{|\phi|>t\}\cap F)
       \le\frac{\mathcal E_1(\phi,\phi)}{t^2}
       \le\frac{3e^2}{t^2}.
\]
Consequently, with $C=\sqrt3\,C_{H_0}$,
\[
 \operatorname{Cap}_S(H_0\cap\{|\phi|>t\};\Omega)
       \le \operatorname{Cap}_S^*(H_0\cap\{|\phi|>t\};\Omega)
       \le Ce/t.
\]
The growth hypothesis and monotonicity of $s\omega(s)$ imply
\[
          \mu\{|\phi|>t\}
                 \le\min\{M,A(Ce/t)\omega(Ce/t)\}.
\]
If $e=0$, the same capacity estimate makes the integral zero. Otherwise
split its layer-cake integral at $t=Ce/r$, and substitute $s=Ce/t$
in the upper part. One obtains
\[
 \int|\phi|\,d\mu
       \le Ce\left[\frac Mr+A\int_0^r\omega(s)\,\frac{ds}{s}\right].
\]
Take the core supremum and then the infimum over $r$. The Dini
assumption makes these quantities finite (also for every other finite
$r$, by local boundedness of $\omega$). For a power choose
$r=(M/A)^{1/(1+\alpha)}$ when $A>0$; the zero-datum case is trivial.
\end{proof}

\paragraph{Frozen mixed slots.}
All statements apply with
$R=T_{S_0,q-1}(v_2,\ldots,v_q)$ in place of $S$, for any arbitrary
allowed original background $S_0$ and fixed bounded classical slots.
The solution then lies in the faithful part energy space for $R$ and
satisfies $D(uR)=\mu$. This is a linear equation in its one unfrozen
potential, not a diagonal nonlinear equation with $q$ moving energy
potentials. In particular no new arbitrary-current wedge or unproved
completion inclusion is being introduced.

% END source: sections/dirichlet_energy_data.tex


% BEGIN canonical R6 source: sections/dirichlet_energy_subsolution_hardy.tex
\section{Sharp Hardy control from a bounded energy subsolution}
\label{sec:r6-energy-subsolution-hardy}

Retain the arbitrary degree-one background, faithful measure $\lambda$,
and part Hilbert space $H=\mathscr D_0(V)$ of
Section~\ref{sec:r6-relative-energy-data}. The potential in this section
is an energy vector, not an assumed classical or completed psh
function. In particular its values on an arbitrary $\lambda$-null set
are never used to define a measure multiplier.

\begin{lemma}[The bounded chain test in the closed core]
\label{lem:r6-energy-reciprocal-test}
Let $h\in H$ satisfy $-B\le h\le0$ $\lambda$-almost everywhere,
where $B<\infty$. For an original real compact smooth $f$ in $V$
and $\varepsilon>0$, the function class
\[
                        \psi=\frac{f^2}{\varepsilon-h}
\]
belongs to $H$ and has support in the compact support of $f$.
Its form pairing satisfies
\begin{equation}
                       -\mathcal E(h,\psi)\le\mathcal E(f,f).
 \label{eq:r6-closed-core-picone}
\end{equation}
No integral of $\psi$ against $D(hS)$ is required to formulate or
prove this assertion.
\end{lemma}

\noindent\textbf{Status: \texttt{DIRECT\_COROLLARY}} of the
closed-core gradient calculus. The following construction records the
approximation and admissibility explicitly.

\begin{proof}
For original core functions $a,b$, their energy measures
$\Gamma_S(a,b)$ are absolutely continuous with respect to $\lambda$.
The original positive-square inequality gives their measure-valued
Cauchy--Schwarz inequality. One may therefore form the Hilbert module
generated by symbols $da$ with bounded Borel coefficients, with
\[
 \left\|\sum_i c_i\,da_i\right\|^2
       =\int\sum_{i,j}c_ic_j\,d\Gamma_S(a_i,a_j),
\]
first for finite sums and then by quotient and completion. Positivity
holds pointwise after taking densities; for any finite family the
matrix is positive semidefinite almost everywhere, as can be tested
on rational coefficient vectors. Bounded Borel multiplication is
well defined in this module.

If $h_j$ is an original core sequence converging to $h$ in $H$, then
$dh_j$ is Cauchy and defines $dh$. The energy-measure inequality also
shows that $\Gamma_S(h_j,h_j)$ is Cauchy in total variation, with
limit absolutely continuous with respect to $\lambda$. In particular
\[
                 \|dh\|^2=\mathcal E(h,h),
           \qquad \|c\,dh\|^2=\int c^2\,d\Gamma_S(h).
\]
These are closed-form objects, independent of any actual representative
of $h$ on a null set.

For a smooth function $\Phi:\mathbf R\to\mathbf R$ with
$\Phi(0)=0$ and bounded derivative, this construction proves the
chain rule directly. Pass to a subsequence with $h_j\to h$
$\lambda$-almost everywhere. The smooth core rule and
\[
 \begin{split}
 \Phi'(h_j)dh_j-\Phi'(h)dh
   ={}&\Phi'(h_j)(dh_j-dh)\\
      &+(\Phi'(h_j)-\Phi'(h))dh
 \end{split}
\]
show convergence in the gradient module. The second term tends to
zero by bounded convergence against the fixed
$\Gamma_S(h)\ll\lambda$. Lipschitz continuity gives the $L^2(\lambda)$
limit $\Phi(h)$, so closability yields
$\Phi(h)\in H$ and $d\Phi(h)=\Phi'(h)dh$. The approximants are still
original core functions in $V$, since $\Phi(0)=0$.

Choose such a smooth global $\Phi$ agreeing with
$(\varepsilon-s)^{-1}-\varepsilon^{-1}$ on $[-B,0]$.
Multiplication by a fixed original compact smooth function is bounded
in the form norm, by the frozen coefficient estimate. Hence
\[
           \psi=\varepsilon^{-1}f^2+f^2\Phi(h)\in H,
\]
with original core approximations supported inside
$\operatorname{supp}f$. The product and chain rules give
\[
            d\psi=\frac{2f}{\varepsilon-h}\,df
                       +\frac{f^2}{(\varepsilon-h)^2}\,dh.
\]
The nonnegative module square is now exactly
\[
 \left\|df+\frac{f}{\varepsilon-h}dh\right\|^2
      =\mathcal E(f,f)+\mathcal E(h,\psi)\ge0.
\]
This proves \eqref{eq:r6-closed-core-picone}. Writing
$v=\varepsilon-h$ would give the usual Picone identity, but no
global form membership of the constant $\varepsilon$ is needed.
\end{proof}

\begin{theorem}[The energy-subsolution Hardy constant]
\label{thm:r6-energy-subsolution-hardy}
Let $h\in H\cap L^\infty(\lambda)$ and suppose that
$\tau=D(hS)$ is a positive Radon measure on $V$. Put
$B=\|h\|_{L^\infty(\lambda)}$. For every positive Radon measure
$0\le\mu\le\tau$ and every original real compact smooth $f$ in $V$,
\begin{equation}
                        \int f^2\,d\mu\le B\mathcal E(f,f).
 \label{eq:r6-energy-subsolution-hardy}
\end{equation}
The coefficient one multiplying $B$ is optimal.

The core restriction map extends uniquely to
\[
                   J_\mu:H\longrightarrow L^2(\mu),
           \qquad \|J_\mu u\|_2\le\sqrt B\,\|u\|_H.
\]
It is the trace of the canonical part-form energy-quasicontinuous representative,
not of an arbitrary $\lambda$-almost-everywhere representative.
The bounded energy solution from
Theorem~\ref{thm:r6-bounded-energy-subsolution} still satisfies
$h\le u_\mu\le0$. If $M=\mu(V)<\infty$, its norm also obeys
\begin{equation}
             \|u_\mu\|_H=\|\mu\|_{-1,V}
                  \le\min\{\|h\|_H,\sqrt{BM}\}.
 \label{eq:r6-bounded-energy-data-height}
\end{equation}
No membership of $h$ or $u_\mu$ in a psh class is asserted.
\end{theorem}

\noindent\textbf{Status: \texttt{DIRECT\_COROLLARY}} of the
closed-core Picone identity and the already proved positive-functional
and bounded-energy-subsolution theorems. The sharp constant and the
measure-trace consequence are included in full.

\begin{proof}
The original core Green identity identifies
$\ell_\tau=-\mathcal E(h,\cdot)$, so $\tau$ has finite core dual
norm. The bounded-energy-subsolution theorem gives finite energy
for $\mu$, and the positive-functional lemma gives their extensions
with $0\le\ell_\mu\le\ell_\tau$ on the positive cone of $H$.
That same theorem gives $-B\le h\le0$ in the energy order.

For the test $\psi$ of the preceding lemma, the almost-everywhere
inequality
$\psi\ge f^2/(\varepsilon+B)$ is an inequality between vectors
of $H$. Consequently
\[
 \frac1{\varepsilon+B}\int f^2\,d\mu
     \le\ell_\mu(\psi)
     \le\ell_\tau(\psi)
     =-\mathcal E(h,\psi)
     \le\mathcal E(f,f).
\]
Only the first integral is an integral of an original smooth core
function. In particular this argument has not presumed an actual
integral formula for $\psi$, nor used the trace inequality it is
proving. Let $\varepsilon\downarrow0$ to obtain
\eqref{eq:r6-energy-subsolution-hardy}; the case $B=0$ follows too.

The inequality makes original core approximations Cauchy in
$L^2(\mu)$ and gives the unique bounded map $J_\mu$. It also shows
that $\mu$ charges no part-form energy-polar Borel set: for a compact set,
apply the compact-core formula for the part energy capacity to functions at
least one on it, and then use Radon inner regularity. A form-convergent
core sequence has a subsequence converging outside an energy-polar
set to the canonical energy-quasicontinuous representative. Its
$L^2(\mu)$ limit is therefore that representative $\mu$-almost
everywhere. This identifies the trace without authorizing arbitrary
changes on a merely $\lambda$-null set.

The already proved bounded-energy-subsolution theorem supplies the
order bounds and $\|u_\mu\|_H\le\|h\|_H$. If $M<\infty$, ordinary
Cauchy--Schwarz for a core $f$ and
\eqref{eq:r6-energy-subsolution-hardy} give
\[
 \left|\int f\,d\mu\right|
          \le\sqrt M\left(\int f^2\,d\mu\right)^{1/2}
          \le\sqrt{BM}\,\|f\|_H.
\]
Riesz isometry proves \eqref{eq:r6-bounded-energy-data-height}.
\end{proof}

\begin{example}[Sharpness with an original annular tent]
\label{ex:r6-energy-subsolution-hardy-sharp}
Let $S=[1]$ on an annulus containing
$V=\{|\pi|<|T|<|\pi|^{-1}\}$, and put
$s=\log|T|/\lambda$, $\lambda=-\log|\pi|$.
For any $B>0$ the function
\[
                      h=B(|s|-1)\quad(-1<s<1)
\]
belongs to $\mathscr D_0(V)$, has essential supremum norm $B$,
and satisfies $D h=(2B/\lambda)\delta_{\sigma(0)}$ on $V$.
Its part-space membership follows by cutting off its zero-boundary
linear tails and smoothing; both the real one-dimensional energy
and the bounded actual functions converge. Original smooth compact
approximations of $f(s)=1-|s|$ have
\[
             \int f^2\,d(Dh)\longrightarrow2B/\lambda,
                  \qquad\mathcal E(f,f)\longrightarrow2/\lambda.
\]
Thus the constant $B$ in the trace inequality cannot be reduced.
The same example attains the height-mass norm bound in the limit.
\end{example}

\begin{corollary}[Form-small perturbations from a bounded energy barrier]
\label{cor:r6-energy-barrier-negative-potential}
Under the theorem's assumptions, if $c\ge0$ and $cB<1$, then
\[
       \mathcal B_c(u,v)=\mathcal E(u,v)
                              -c\int(J_\mu u)(J_\mu v)\,d\mu
\]
is a bounded coercive form on $H$, with
$\mathcal B_c(u,u)\ge(1-cB)\|u\|_H^2$.
Every functional $L\in H^*$ therefore has a unique weak solution
$\mathcal B_c(u,\phi)=L(\phi)$, and
$\|u\|_H\le\|L\|/(1-cB)$.
For $L=-\ell_\sigma$ this is, on all original compact tests,
\[
                         D(uS)+c(J_\mu u)\mu=\sigma.
\]
The product is locally integrable by the trace bound and local Radon
finiteness. This is a linear energy equation; no Markov property of
the perturbed form and no classical psh solution are asserted.
\end{corollary}

\begin{proof}
The trace inequality gives boundedness and the displayed lower
bound. The coercive Hilbert-space Riesz theorem proves existence,
uniqueness and the norm estimate. The original core Green identity
gives the stated sign of the distributional equation.
\end{proof}

For the annular example, with $\mu=Dh$, one has
$\int h^2\,d\mu=B\mathcal E(h,h)$. At $cB=1$ the nonzero vector
$h$ is even a weak kernel vector, since $h=-B$ at the support of
$\mu$ and $\mathcal E(h,\phi)=-\int\phi\,d\mu$. For $cB>1$
the form is negative on $h$. Thus the universal strict coercivity
threshold cannot be enlarged under only the stated barrier-height
hypothesis.

% END source: sections/dirichlet_energy_subsolution_hardy.tex


% BEGIN canonical R6 source: sections/capacity_energy_data_separation.tex
\section{A bounded continuous energy solution with no local classical potential}
\label{sec:r6-energy-data-separation}

Retain $n\ge3$, the full torus box $\Omega$, its smooth background
$S=\beta^{n-1}/(n-1)!$, the real compact set $A$, and the probability
$\vartheta$ from Section~\ref{sec:r6-capacity-energy-null-gap}.
Scale $A$ inside a real ball $B$ with $\overline B\Subset Q$ and put
\[
                 V=\operatorname{trop}^{-1}(B),\qquad
                 K=\sigma(A),\qquad \nu=\sigma_*\vartheta.
\]
The ordinary domain $V$ is relatively compact in $\Omega$.

\begin{lemma}[A bounded local potential can be extended near Gauss points]
\label{lem:r6-local-gauss-log-extension}
Let $x\in Q$, let $O$ be an ordinary neighborhood of $\sigma(x)$,
and let $v\in\mathcal C_b(O)$. There are a real box $Q_x\ni x$,
an ordinary neighborhood $N_x$ of $\sigma(x)$ inside
$O\cap\operatorname{trop}^{-1}(Q_x)$, and a bounded classical psh
function $U_x$ on the \emph{full} polyannulus
$\operatorname{trop}^{-1}(Q_x)$ such that
\[
                     U_x-v=M\log|P|-c\quad\hbox{on }N_x
\]
for a Laurent polynomial $P$ that is an analytic unit there, and
finite constants $M\ge0,c\in\mathbf R$. Consequently
$D(U_xS)=D(vS)$ on $N_x$.

If a family of such $v$ has one bound $|v|\le H$ on a fixed ordinary
neighborhood, the box, $N_x$, $P$, $M$, $c$, and the bound for $U_x$
can all be chosen independently of the member of that family.
\end{lemma}

\noindent\textbf{Status: \texttt{NEW\_PROVED}} elementary local
extension for this Gauss-section geometry. A general ordinary
neighborhood is not assumed to contain a full tropical box.

\begin{proof}
Shrink the given ordinary neighborhood so that $|v|\le H$ there.
The topology of the affine torus analytification is generated by
absolute values of Laurent polynomials. A finite basic neighborhood
therefore gives nonzero Laurent polynomials $F_1,\ldots,F_m$ and
strict lower and upper bounds around their values at $\sigma(x)$.
For every analytic point $y$ of the torus,
\[
                  |F_l(y)|\le|F_l(\sigma(\operatorname{trop}y))|
\]
by the ultrametric inequality on its Laurent monomials. The Gauss
norms on the right vary continuously with the real coordinates.
After shrinking a real box $Q_x$ around $x$, all upper bounds of
the basic neighborhood consequently hold on its entire tropical
inverse image. Thus for suitable $a_l<|F_l(\sigma(x))|$ the set
\[
 D=\{y:\operatorname{trop}y\in Q_x,\ |F_l(y)|>a_l\ (1\le l\le m)\}
\]
lies in the fixed ordinary neighborhood where the bound holds.
If no lower bound is needed, this is already a full polyannulus;
take $U_x=v+H$ and the assertion is immediate.

Otherwise set $c_l=\log|F_l(\sigma(x))|$,
$\delta=\min_l(c_l-\log a_l)>0$, $C=\sum_lc_l$, and
$P=\prod_lF_l$. Shrink $Q_x$ further so that
\[
 \big|\log|F_l(\sigma(s))|-c_l\big|<\delta/(4m)
                                                       \quad(s\in Q_x).
\]
On the Gauss section over $Q_x$ one has
$\log|P|\ge C-\delta/4$. At any relative boundary point of $D$
inside the full polyannulus, at least one lower bound is an equality;
the remaining Gauss upper bounds imply
$\log|P|\le C-3\delta/4$ there. Put $L=C-\delta/2$ and choose
$M$ with $M\delta/8>H+1$. Define
\[
 U_x=\max\{0,v+M(\log|P|-L)\}\quad\hbox{on }D,
                  \qquad U_x=0\quad\hbox{off }D.
\]
Near each relative boundary point the expression inside the maximum
is negative, by the uniform bound on $v$ and continuity of $\log|P|$.
Thus the two definitions agree with the constant zero on an ordinary
neighborhood of every gluing point. They define a bounded classical
psh function on the full polyannulus. Its bound is at most
$\max(0,H+3M\delta/4)$.

On the ordinary open subset of $D$ where
$\log|P|-L>(H+1)/M$, the second branch is strictly positive.
This subset contains the whole Gauss section over the small box.
There $U_x-v=M\log|P|-ML$, an original smooth pluriharmonic
function since $P$ is a unit. Separate closedness of $S$ gives
the current equality. All choices used only the fixed neighborhood
and the bound $H$, proving the uniform assertion.
\end{proof}

\begin{proposition}[Local outer-nullity for every bounded classical trace]
\label{prop:r6-local-gauss-nullity}
Let $O$ be any ordinary neighborhood of $K$ in the torus, with $S$
restricted there. Every $v\in\mathcal C_b(O)$ satisfies
\[
                              D(vS)(K)=0.
\]
More locally, on each sufficiently small ordinary neighborhood of a
point of $K$, the corresponding portion of $K$ has relative outer
BT capacity zero for normalized testers on that neighborhood.
The same nonconcentration holds for locally uniformly bounded
completed relative potentials.
\end{proposition}

\begin{proof}
At each point of $K$ apply the preceding lemma, and then restrict to
a smaller real box with compact closure in $Q_x$. The full-box
Gauss-section growth estimate of
Lemma~\ref{lem:r6-gauss-lipschitz-capacity}, after a fixed affine
normalization of the bounded $U_x$, implies that $D(U_xS)$ gives
zero mass to every Gauss subset of zero $\mathcal H^{n-1}$ measure
in that smaller box. On the protected ordinary neighborhood its
curvature equals $D(vS)$. Finitely many such neighborhoods cover
$K$, proving the first assertion.

For the capacity assertion start with normalized $0\le v\le1$ on
one fixed ordinary neighborhood. The extension lemma uses the same
data and the same bound for every such tester. Cover a compact
portion of $K$ in the protected neighborhood by small real balls
with arbitrarily small sum of $(n-1)$st radius powers. Intersect
their ordinary tropical inverse images with a fixed relatively
compact protected neighborhood. The full-box mass estimate gives
arbitrarily small capacity for these ordinary open neighborhoods,
uniformly over the original normalized class. This proves the local
outer-null assertion on its stated capacity domain.

For completed potentials use their locally uniformly bounded defining
rows on smaller ordinary neighborhoods. The same construction and
normalization give the same uniform open-neighborhood mass bound.
Ordinary weak convergence and the open-set inequality pass that
bound to the positive completed curvature measure. Letting the bound
tend to zero proves the last assertion. No restriction to $K$ is
passed directly through a weak limit.
\end{proof}

\begin{theorem}[Finite energy and bounded continuous solution without a psh trace]
\label{thm:r6-bounded-energy-no-local-psh}
The probability $\nu$ has finite core energy norm on $V$. There is
a bounded continuous actual function $a$ on $\overline V$, zero on
the ordinary boundary, such that
\[
                  a\in\mathscr D_0(V),\qquad
                  D(aS)=\nu\quad\hbox{on }V
\]
on every original compact smooth scalar test. This is the unique
energy-class solution.

Nevertheless, no ordinary neighborhood $O$ of $K$ has a locally
bounded actual classical psh potential $v$ with $D(vS)=\nu|_O$.
Nor can such an equation be solved by a locally uniformly bounded
completed relative potential. In particular there is no bounded
classical subsolution whose relative curvature dominates $\nu$ near
$K$.
\end{theorem}

\noindent\textbf{Status: \texttt{NEW\_PROVED}} separation of solution
classes and local representability, using the preceding null-set
construction. The energy Riesz theorem itself remains a direct input.

\begin{proof}
For an original compact smooth $f$ in $V$, the Newton trace inequality
and the all-original-test restriction-energy bound give
\[
 \left|\int f\,d\nu\right|^2
    =\left|\int f\circ\sigma\,d\vartheta\right|^2
    \le C_n I(\vartheta)\int|\nabla(f\circ\sigma)|^2
    \le C_n I(\vartheta)\mathcal E_S(f).
\]
Thus $\nu$ has finite energy.

We construct its bounded continuous solution explicitly. Let
$G(x)=c_n|x|^{2-n}$ be the positive Newton kernel with
$-\Delta G=\delta_0$, and put $P=G*\vartheta$.
The Frostman bound used in
Lemma~\ref{lem:r6-cantor-newton-set} implies uniformly in $x$
\[
       \int_{|x-y|<r}|x-y|^{2-n}\,d\vartheta(y)\le C r^d,
                         \qquad d=\log2/\log3>0.
\]
Splitting the kernel into this small tail and its continuous far
part proves that $P$ is bounded and continuous. Its weak gradient
is locally square integrable: the vector-kernel convolution bound
from the same lemma, applied first to smooth approximations of
$\vartheta$, bounds its gradient norm by a fixed multiple of
$I(\vartheta)^{1/2}$. Passing weakly and using local uniform
convergence identifies the gradient of $P$.

The measure is supported away from $\partial B$, so $P$ is smooth
near that sphere. Let $H$ be the ordinary real Poisson extension to
$B$ of its boundary values, and set $U=H-P$. Equivalently,
$-U$ is the positive Dirichlet Green potential of $\vartheta$ in
the ball. Hence
\[
 U\in C(\overline B)\cap H^1(B),\quad U<0\text{ on }B,
            \quad U|_{\partial B}=0,\quad\Delta U=\vartheta.
\]
The standard ball Poisson kernel gives the displayed harmonic
extension; positivity of its Green kernel gives the sign. These are
ordinary real integral identities, not an elliptic regularity claim
for a Berkovich distribution.

To verify the zero-boundary energy class directly, put
$U_\epsilon=\min(U+\epsilon,0)$. It is continuous with compact
support in $B$. The Sobolev contraction rule and $U<0$ show
$U_\epsilon\to U$ in $H^1(B)$: the gradient difference is supported
where $-\epsilon\le U<0$ and tends to zero by dominated convergence.
Mollification then gives real smooth compact functions converging to
$U$ in energy and uniformly on $B$. Their toric pullbacks are
original compact smooth functions in $V$, and
\eqref{eq:r6-torus-laplace-calibration} makes their energy differences
exactly the real ones. The frozen coercivity and uniform convergence
identify their part-space limit with the actual bounded continuous
function $a=U\circ\operatorname{trop}$. Thus $a\in\mathscr D_0(V)$.
Continuity on $\overline V$ and the ordinary zero boundary value
follow from the real boundary values and proper tropicalization.

It remains to check the equation on all original tests. Extend $U$
by zero beyond $B$ and convolve on the real box. The resulting smooth
$U_j$ converge locally uniformly to $U$. On each fixed compact
subset of $B$, for sufficiently small convolution radius,
$\Delta U_j=\vartheta*\kappa_j$; any boundary contribution is
outside that compact set. For an arbitrary original compact smooth
scalar test $\eta$ in $V$, the original torus calibration gives
\[
       \langle D((U_j\circ\operatorname{trop})S),\eta\rangle
           =\int_B(\eta\circ\sigma)\,\Delta U_j\,dx
           \longrightarrow\int_A\eta\circ\sigma\,d\vartheta.
\]
Here $\eta\circ\sigma$ is a continuous compact test, not assumed
smooth in one real chart. On the other hand, uniform coefficient
multiplication of the fixed $S$ gives convergence of the left-hand
currents to $D(aS)$. This proves the equation on the full original
test class. Energy uniqueness is the frozen Riesz uniqueness theorem.

Finally, Proposition~\ref{prop:r6-local-gauss-nullity} makes every
locally bounded classical or locally uniformly bounded completed
relative curvature give zero mass to $K$. Since $\nu(K)=1$, it
cannot equal or be dominated by such a curvature. This proves all
the asserted nonrepresentability statements.
\end{proof}

\paragraph{Scope.}
The background is smooth and fixed, the solution is actually bounded
and continuous, and the failure is local near its compact datum.
It is therefore stronger than a failure of one prescribed boundary
value problem. The successful solution belongs to the larger faithful
energy class. It supplies neither a new ambient psh solution nor a
counterexample to general relative compact no-gain or Borel outer
regularity.

% END source: sections/capacity_energy_data_separation.tex


% BEGIN canonical R6 source: sections/capacity_energy_separation_consequences.tex
\section{Consequences of the toric capacity and energy separation}
\label{sec:r6-energy-separation-consequences}

\noindent\textbf{Status: \texttt{DIRECT\_COROLLARY}}. Retain the
notation of Theorems~\ref{thm:r6-energy-BT-null-gap} and
\ref{thm:r6-bounded-energy-no-local-psh}; in particular $n\ge3$,
$K=\sigma(A)$, and $\nu(K)=1$.

\paragraph{No local difference of bounded classical psh potentials.}
There is no locally bounded actual function $b$ near $K$ which is
locally a difference of two members of $\mathcal C_b$ and satisfies
$D(bS)=\nu$. Indeed, on each ordinary neighborhood with
$b=u-v$, Proposition~\ref{prop:r6-local-gauss-nullity} makes both
positive measures $D(uS)$ and $D(vS)$ vanish on the corresponding
portion of $K$. Their signed difference therefore vanishes there.
A finite ordinary covering of the compact $K$ gives $D(bS)(K)=0$,
contradicting $\nu(K)=1$. The same argument applies to local
differences of locally uniformly bounded completed relative
potentials. In particular the explicit bounded continuous energy
solution in Theorem~\ref{thm:r6-bounded-energy-no-local-psh} lies
outside the local bounded DSH class near its datum. This statement
does not assert a failure of the frozen potential-quotient theorem,
whose hypotheses require an actual classical difference from the start.

\paragraph{Relative BT nullity does not mean ambient pluripolarity.}
Fix $x\in A$. Choose $M>0$ larger than every positive coordinate
increment $t_i-x_i$ for $t\in Q$, and put
\[
 w_x(t)=M^{-1}\max(0,t_1-x_1,\ldots,t_n-x_n).
\]
It is a normalized continuous convex function on $Q$, so its toric
pullback is an actual ambient classical tester on the whole $\Omega$.
The original torus slope-simplex calibration gives
\[
            \bigl(D(w_x\circ\operatorname{trop})\bigr)^n
                         =M^{-n}\delta_{\sigma(x)}.
\]
Indeed the only full-dimensional slope polytope is
$M^{-1}\operatorname{conv}(0,e_1,\ldots,e_n)$ at $x$; its volume
is $M^{-n}/n!$, canceled by the original top-degree factor $n!$.
The identity is valid on all original scalar tests, by the same
convex smoothing calibration used for the frozen radial examples.
Consequently
\[
             \operatorname{Cap}_{[1]}(K;\Omega)\ge M^{-n}>0,
                  \qquad \operatorname{Cap}_S^*(K;\Omega)=0.
\]
The tester can be restricted to any ordinary neighborhood of
$\sigma(x)$ while preserving its normalization and its atom.
Thus the frozen local ambient polar outer-zero theorem
(Corollary~\textup{R4--26.3}) excludes a local
pluripolar witness at every point $\sigma(x)$ of $K$. The discrepancy
already occurs for the smooth fixed $S$ used here; it does not depend
on a background carried by an ambient pluripolar set.

% END source: sections/capacity_energy_separation_consequences.tex


% BEGIN canonical R6 source: sections/dirichlet_toric_certificate.tex
\section{An explicit energy-dual certificate for the nonpolar obstruction}
\label{sec:r6-toric-energy-certificate}

\noindent\textbf{Status: \texttt{DIRECT\_COROLLARY}} of the new
separation criterion, with the following explicit real-variable
calculation. The underlying failure of a bounded-subsolution theorem is
the frozen R4 counterexample. The added conclusion is a strictly positive,
quantified distance from the datum to the feasible completed image in the
negative-energy norm, witnessed by an original compact smooth test.

Let $n\ge2$, $t_i=\log|T_i|$ on the original $n$-torus, and
\[
 B_1=\{x\in\mathbf R^n:|x|<1\},\qquad V=t^{-1}(B_1),\qquad
 \beta=\sum_i d't_i\wedge d''t_i,\qquad S=[\beta^{n-1}].
\]
Choose a larger ordinary norm-box domain $\Omega$ in the torus,
containing $\overline V$, and ordinary $W$ and compact $L$ with
$\overline V\subset L\Subset W\Subset\Omega$.
Put $h(t)=(|t|^2-1)/2$ and choose a common bound larger than its bound
on $\Omega$. Use the nonempty convex family
\[
 \mathscr P=\{u\in\mathcal C_b(\Omega):
       u=h\text{ on }\Omega\setminus V,\quad h\le u\le0\text{ on }V\}.
\]
All its members have one common global bound and collar. Its selected
completed limits retain ordinary boundary value zero.

Take a nonzero radial $f\in C_c^\infty(B_1)$ with $0\le f\le1$,
supported in $\{|x|<r_0\}$ for some $r_0<1$. Set
\[
                     \nu=f(t)\beta^n.
\]
The smooth subsolution $h$ has $D(hS)=\beta^n\ge\nu$ on $V$.
The signed datum $\delta=\nu-\mu_h|_V$ has finite energy: on the
compact $\overline V$, the smooth coefficient measure $\beta^n$ has a
bounded density against the faithful $\lambda$, so its density and that
of $\delta$ belong to $L^2(\lambda)$.

\begin{proposition}[A positive distance despite a bounded subsolution]
\label{prop:r6-toric-positive-separation-distance}
There are $\eta\in C_c^\infty(B_1)$, $\eta\ge0$, and $c>0$ such that
\begin{equation}
 \inf_{a\in\mathscr A}
       \|\nu-D(aS)|_V\|_{-1,V}
 \ \ge\sqrt{(n-1)!}\,
              \frac{c\int_{B_1}\eta\,dx}{\|\partial_1\eta\|_{L^2(B_1)}}
 \ >0.
 \label{eq:r6-toric-distance-lower-bound}
\end{equation}
In particular the smooth-test feasibility inequality fails for some
original compact smooth scalar on $V$. Here $\mathscr A$ is precisely the
bounded global capacity closure of the displayed family, not an
unspecified local sheaf.
\end{proposition}

\begin{proof}
Write $f(x)=f_0(|x|)$ and consider the explicit radial function
\[
 U_f(r)=-\int_r^1 s^{1-n}
                   \left(n\int_0^s f_0(t)t^{n-1}\,dt\right)ds.
\]
It is the smooth real solution of $\Delta U_f=nf$ with zero boundary
value; the integral is regular at zero because its inner integral is
$O(s^n)$. Moreover $h(x)\le U_f(x)\le0$ because $0\le f\le1$.
For $r>r_0$, with
\[
                  M=n\int_0^{r_0} f_0(t)t^{n-1}\,dt>0,
\]
one has $U_f'(r)=Mr^{1-n}$ and
$U_f''(r)=(1-n)Mr^{-n}<0$.
Choose $x_*=(r_*,0,\ldots,0)$ with $r_0<r_*<1$.
On a sufficiently small real neighborhood of $x_*$,
$\partial_1^2U_f\le-c$ for some $c>0$.
Choose a nonzero nonnegative smooth $\eta$ compactly supported there.

Let $\varphi\in H_0^1(B_1)$ be the real weak solution of
\[
                     \Delta\varphi=\partial_1^2\eta.
\]
It is obtained directly by Riesz representation for the gradient norm:
$\int\nabla\varphi\cdot\nabla v=\int\partial_1\eta\,\partial_1v$.
Consequently
\begin{equation}
                     \|\nabla\varphi\|_2\le\|\partial_1\eta\|_2.
 \label{eq:r6-real-test-energy}
\end{equation}
Choose $\varphi_m\in C_c^\infty(B_1)$ converging to $\varphi$ in that
norm and set $\Phi_m=\varphi_m(t)$. These are original compact smooth
scalars on $V$. The original torus calibration gives
\[
       \mathcal E_S(\Phi_m-\Phi_k)
            =(n-1)!\int_{B_1}|\nabla(\varphi_m-\varphi_k)|^2\,dx.
\]
Thus they converge to a vector $\Phi\in\mathscr D_{S,0}(V)$, with
$\|\Phi\|_H\le\sqrt{(n-1)!}\|\partial_1\eta\|_2$.
Only this form-domain vector and its original smooth approximations are
needed; no arbitrary energy representative is called a classical psh
function.

For $u\in\mathscr P$ let $U=u\circ\sigma$ be its actual Gauss trace.
The R4 full-annulus argument makes $U$ separately convex and continuous
on $B_1$, and $h\le U\le0$. These bounds also put $U-h$ in
$H_0^1(B_1)$. To see this without an unproved trace assertion, for
$x\in B_1$ set $d=1-|x|$. On each coordinate segment of radius $d/2$
about $x$, one has $|U|\le3d/2$. Convexity bounds each slope at $x$
in absolute value by $3$. On interior real boxes, Fubini therefore
gives uniformly bounded weak coordinate derivatives. The zero boundary
limit, squeezed between $h$ and zero, makes the zero extension a
Sobolev function with zero trace. Hence $U-h\in H_0^1(B_1)$.

For each original $\Phi_m$, calibration and Green give
\[
 \mathcal E_S(u-h,\Phi_m)
       =-(n-1)!\int_{B_1}(U-h)\Delta\varphi_m\,dx.
\]
Pass in the real gradient pairing and in the faithful form pairing to
obtain
\[
 \mathcal E_S(u-h,\Phi)
       =-(n-1)!\int_{B_1}(U-h)\partial_1^2\eta\,dx
       \le(n-1)!\int_{B_1}\eta\,dx.
\]
The inequality uses the positive distribution $\partial_1^2U$ and
$\partial_1^2h=1$. It is uniform over the entire classical family.
Thus its support functional satisfies
$\sigma_{C_V}(\Phi)\le(n-1)!\int\eta$.

Let $y$ be the Riesz vector of the datum. The same calibration, together
with $\Delta(U_f-h)=n(f-1)$ and zero real boundary values, gives
\begin{align*}
 \mathcal E_S(y,\Phi)
    &=-n!\int_{B_1}(f-1)\varphi\,dx\\
    &=-(n-1)!\int_{B_1}(U_f-h)\partial_1^2\eta\,dx\\
    &\ge(n-1)!\left(\int\eta+c\int\eta\right).
\end{align*}
Subtract the support bound. The gap is at least
$(n-1)!c\int\eta>0$. Normalize by $\|\Phi\|_H$ and use
\eqref{eq:r6-real-test-energy} and
Corollary~\ref{cor:r6-energy-data-projection} to obtain
\eqref{eq:r6-toric-distance-lower-bound}.
The support functional is energy-continuous on the bounded constraint,
so the strict gap persists for some original $\Phi_m$. This is the
promised original-test certificate, not only an abstract nonsmooth
separating direction.
\end{proof}

The pointwise lower and upper barriers are satisfied by $U_f$ itself.
The obstruction is its negative coordinate Hessian in the outer annulus,
which a Gauss trace of a globally defined classical psh row cannot have.
This calculation does not identify raw BT capacity with energy capacity,
and does not assert that all coefficient measures of $S$ live on the
Gauss skeleton. Only the specified toric top pairings are used.

% END source: sections/dirichlet_toric_certificate.tex


% BEGIN canonical R6 source: sections/dirichlet_radial_toric_criterion.tex
\section{A checkable radial criterion for a smooth toric background}
\label{sec:r6-radial-toric-dirichlet}

This qualified family supplies actual classical solutions as well as
an exact obstruction. It is a remaining-degree-one statement for the
specified smooth background, not a replacement for the arbitrary-$S$
criteria above.

Fix $n\ge2$, put $B=\{x\in\mathbf R^n:|x|<1\}$, and let
\[
 V=\operatorname{trop}^{-1}(B)\subset(\mathbf G_{m,K}^n)^{\an},
 \quad t_i=\log|T_i|,\quad
 \beta=\sum_i d't_i\wedge d''t_i,\quad
 S=\frac{\beta^{n-1}}{(n-1)!}\wedge[1].
\]
Write $\sigma:B\to V$ for the original Gauss section. The ordinary
boundary below is taken in the ambient analytic torus. Its Gauss
points include $\sigma(\partial B)$; it is not necessary to identify
the entire boundary with every fiber above $\partial B$.

\begin{proposition}[Radial classical solvability]
\label{prop:r6-radial-toric-dirichlet}
Let $\rho\in C([0,1])$ be nonnegative and set
\[
       \mu=\sigma_*(\rho(|x|)\,dx),\qquad
       M(r)=\int_0^r s^{n-1}\rho(s)\,ds.
\]
There exists a bounded actual classical psh $u$ on $V$ such that
$D(uS)=\mu$ on every original compact scalar test and
\[
          \lim_{\substack{y\in V\\y\to z}}u(y)=0
                                      \qquad(z\in\partial V)
\]
if and only if
\begin{equation}
                r^n\rho(r)\ge(n-1)M(r)
                                      \qquad(0<r<1).
 \label{eq:r6-radial-convexity-condition}
\end{equation}
When this condition holds, one such solution is the continuous
torus-invariant function
\begin{equation}
 u=U(|t|),\qquad
 U(r)=-\int_r^1 s^{1-n}M(s)\,ds.
 \label{eq:r6-radial-toric-solution}
\end{equation}
Every bounded classical solution with the stated boundary values
has this same actual Gauss-section trace $U(|x|)$. This assertion
does not claim uniqueness of all actual values away from that trace.
\end{proposition}

\noindent\textbf{Status: \texttt{DIRECT\_COROLLARY}} of the frozen
whole-domain Gauss-trace argument, original toric calibration, and
the elementary radial computation below. No general relative
subsolution theorem is assumed.

\begin{proof}
The formula for $M$ gives
\[
 0\le M(r)\le\|\rho\|_\infty r^n/n,
       \qquad -\|\rho\|_\infty/(2n)\le U(r)\le0.
\]
For $r>0$,
\[
 U'(r)=r^{1-n}M(r),\qquad
 U''(r)=\rho(r)-(n-1)r^{-n}M(r),\qquad
 U''(r)+(n-1)U'(r)/r=\rho(r).
\]
Both $U''(r)$ and $U'(r)/r$ tend to $\rho(0)/n$ at zero.
Thus $x\mapsto U(|x|)$ is $C^2$ on $B$, continuous on its closure,
solves $\Delta U=\rho(|x|)$, and has boundary value zero.

Suppose an actual bounded classical solution $u$ exists, and put
$g=u\circ\sigma$. On every coordinate segment contained in $B$,
the complete-extension annulus construction of
Lemma~\ref{lem:r6-gauss-lipschitz-capacity} applies inside this full
tropical inverse-image domain. Its two-endpoint Shilov comparison
shows that $g$ is separately convex. Its boundedness therefore
makes it locally Lipschitz on $B$.

For each $\varphi\in C_c^\infty(B)$, the pullback
$\varphi\circ\operatorname{trop}$ is an original compact smooth
scalar on $V$. Recursion and the original toric calibration give
\[
 \int_B g\,\Delta\varphi\,dx
  =\langle uS,D(\varphi\circ\operatorname{trop})\rangle
  =\int_B\varphi\rho(|x|)\,dx.
\]
Only this top trace is calibrated; no assertion about support of
every lower coefficient measure of $S$ is used.

Every Gauss point $\sigma(x_0)$ with $|x_0|=1$ belongs to the ordinary
boundary of $V$, by approaching $x_0$ through the real ball and using
continuity of $\sigma$. The specified ordinary boundary limit hence
extends $g$ continuously by zero to $\partial B$. Therefore
$g-U(|x|)$ is a continuous distributionally harmonic function on
$B$ with zero boundary values, and is zero. One can verify this last
uniqueness directly: mollify on successively smaller balls, apply
the ordinary smooth harmonic maximum principle there, and let their
boundaries approach $\partial B$, using uniform continuity on the
closed ball.

Consequently every such solution has trace $U(|x|)$. Separate
convexity on the first coordinate axis, at $(r,0,\ldots,0)$,
implies $U''(r)\ge0$. The displayed derivative formula is exactly
\eqref{eq:r6-radial-convexity-condition}, proving necessity.

Conversely, suppose that condition holds. Then $U''\ge0$ and
$U'\ge0$. The radial Hessian has eigenvalues $U''$ in the radial
direction and $U'/r$ in the tangential directions, with their stated
common limit at zero. It is positive semidefinite, so $U(|x|)$ is
convex. A continuous convex function of original logarithmic unit
moments is classical psh: on compact smaller moment domains one
may use convex smooth mollifications, or the defining curve test
with harmonic logarithmic unit coordinates. Thus the function in
\eqref{eq:r6-radial-toric-solution} is actual continuous classical
psh on $V$.

The calibration
$D(\varphi(t)S)=\sigma_*((\Delta\varphi)dx)$ first holds for smooth
real $\varphi$. Apply ordinary real mollification on compact smaller
balls to $U(|x|)$, with $C^2$ convergence, or pass the fixed
continuous multiplier through the original distributional recursion.
It gives $D(uS)=\mu$ on every original scalar test. Finally every
ordinary boundary point of $V$ has $|t|=1$, by continuity of $t$;
the continuous formula $U(1)=0$ gives the required boundary limits.
\end{proof}

\paragraph{Two immediate checks.}
If $\rho$ is nondecreasing, then
$M(r)\le r^n\rho(r)/n$, so the criterion holds. In particular
constant density $\rho=c$ gives
$u=c(|t|^2-1)/(2n)$. In contrast, if $\rho$ is nonzero and vanishes
on some terminal interval $(r_0,1)$, then $M>0$ there while
$\rho=0$, and the criterion fails. The latter recovers the frozen
compact radial-source obstruction while the former gives a direct
solvable family in every ambient dimension. No nonlinear
Gauss-pushforward formula is asserted for several moving slots.

% END source: sections/dirichlet_radial_toric_criterion.tex


% BEGIN canonical R6 source: sections/dirichlet_admissible_variations.tex
\section{Admissible variations in every remaining degree}
\label{sec:r6-dirichlet-admissible-variations}

Fix the original field and test sheaf, an arbitrary fixed original-positive
separately closed $(p,p)$-background $S$, and $q=n-p\ge1$. Retain the
norm-box domain $\Omega$, compact outer collar set $L$, smooth outer
potential $h$, and common bound $B$ of the R4 completed variational
construction. Let $V$ be ordinary open with $\overline V\subset L$.
As in Section~\textup{R4--81}, prescribe an actual
exterior function $e$ on $\Omega\setminus V$ and continuous barriers
$\ell\le r$ on $\overline V$ with the same ordinary boundary value $g$.
Assume that the classical family
\[
 \mathscr P_{\rm bd}=
 \{u\in\mathscr P_{h,L}^{B}: u=e\text{ on }\Omega\setminus V,
                                      \ \ell\le u\le r\text{ on }V\}
\]
is nonempty. Its bounded global capacity closure is denoted
$\mathscr K$. It is convex and capacity closed, and its selected actual
representatives retain the exterior, both barriers, and boundary value.

Let $\widetilde\nu\in\mathfrak N_{S,h,L}$ be admissible data in the
precise R4 sense, and write $\nu=\widetilde\nu|_V$. The completed
variational objective is
\[
 F(a)=\mathcal E_{\rm AM}(a)-\int_L(a-h)\,d\widetilde\nu,
       \qquad m=\max_{a\in\mathscr K}F(a).
\]
Its unique maximizing class is denoted $a_*$. Different admissible
extensions of $\nu$ outside $V$ change $F$ by a constant on
$\mathscr K$, because every selected representative has the same
actual exterior $e$. Thus they do not change $a_*$.

\subsection{What a verified two-sided variation gives quantitatively}

\begin{proposition}[A stationarity estimate with an explicit admissible radius]
\label{prop:r6-admissible-stationarity-estimate}
Let $\phi\in A_c^{0,0}(V;\mathbf R)$ be original smooth. There is a
finite $C_\phi$, depending only on the fixed collar data, $S$, the common
bound and this test, with the following property. If $a\in\mathscr K$
and $a\pm t_0\phi\in\mathscr K$ for some $t_0>0$, then for
$0<t\le t_0$,
\begin{equation}
 \left|\int_V\phi\,d(\mu_a-\nu)\right|
              \le \frac{m-F(a)}{t}+\frac{C_\phi t}{2}.
 \label{eq:r6-stationarity-defect}
\end{equation}
In particular, at the maximizer the displayed integral is zero whenever
some two-sided admissible radius exists for this test.
Moreover, for all $a\in\mathscr K$,
\begin{equation}
 I_S(a,a_*)\le\frac{m-F(a)}{2c_q},
                     \qquad c_q=\frac{q}{4^{q+1}}.
 \label{eq:r6-objective-energy-gap}
\end{equation}
The previously proved energy-dual and energy-capacity estimates turn
this into the corresponding quantitative errors for the fixed-background
measure and potential class.
\end{proposition}

\noindent\textbf{Status: \texttt{DIRECT\_COROLLARY}} of the proved
compact-collar energy polynomial, relative CLN and the R4 midpoint
inequality. The displayed radius is an actual admissibility hypothesis;
an obstacle gap alone does not provide it.

\begin{proof}
Convexity makes $a+s\phi$ admissible for $|s|\le t_0$.
The finite-polynomial energy identity and compact Green formula give
\[
 \frac{d}{ds}F(a+s\phi)
                  =\int\phi\,d(\mu_{a+s\phi}-\widetilde\nu),
\]
and
\[
 \frac{d^2}{ds^2}F(a+s\phi)
       =-q\int \Gamma_{T_{S,q-1}((a+s\phi)^{q-1})}(\phi).
\]
These identities for completed potentials follow from their defining
classical rows and the completed weighted calculus, exactly as in
Lemma~\textup{R4--82.2}. Equivalently the second
identity follows by polarization of the finite polynomial; the smooth
compact function $\phi$ is the only differentiated direction.
The last integrand is positive. Relative CLN applied to the fixed
positive test $d'\phi\wedge d''\phi$ bounds its total mass by
$C_\phi/q$, uniformly over the admissible segment. Every member has
the same bound $B$. Hence
\[
 F(a\pm t\phi)-F(a)
    \ge\pm t\int\phi\,d(\mu_a-\widetilde\nu)
                                      -\tfrac12C_\phi t^2.
\]
The left side is at most $m-F(a)$; use both signs and the support of
$\phi$ to obtain \eqref{eq:r6-stationarity-defect}.
At a maximizer let $t\downarrow0$.

For the last assertion, the midpoint belongs to $\mathscr K$ and the
linear datum term cancels from its midpoint difference. Thus
\[
 c_q I_S(a,a_*)
 \le F((a+a_*)/2)-\tfrac12F(a)-\tfrac12m
 \le\tfrac12(m-F(a)),
\]
which is \eqref{eq:r6-objective-energy-gap}.
\end{proof}

\subsection{A classical curvature reserve that supplies the missing directions}

For each original compact smooth $\phi$ in $V$, the explicit R4
compensator lemma supplies a bounded nonnegative original smooth
classical function $\rho_\phi$ on $\Omega$ such that
\[
              \rho_\phi+2\phi,\quad \rho_\phi-2\phi
                           \quad\text{are classical psh}.
\]
Its proof uses finite unit-moment presentations and a positive sum of
smooth convex functions of logarithms of affine regular functions.
In particular this assertion does not invoke a general equivalence
between original-current positivity and classical psh.

\begin{theorem}[A qualified interior Dirichlet criterion in all degrees]
\label{thm:r6-reserve-dirichlet-criterion}
Suppose there is a classical maximizing sequence
$u_j\in\mathscr P_{\rm bd}$, with $F(u_j)\to m$, having the following
property for every original $\phi\in A_c^{0,0}(V;\mathbf R)$.
There are an ordinary neighborhood $O_\phi$ of
$\operatorname{supp}\phi$, with $O_\phi\subset V$, constants
$\epsilon_\phi>0$, $b_\phi>0$, and an index $j_\phi$, such that
for all $j\ge j_\phi$:
\begin{enumerate}
\item the actual function $u_j-\epsilon_\phi\rho_\phi$ belongs to
      $\mathcal C_b(O_\phi)$;
\item on $\operatorname{supp}\phi$ one has the simultaneous margins
\[
  |u_j|\le B-b_\phi,
       \qquad \ell+b_\phi\le u_j\le r-b_\phi.
\]
\end{enumerate}
Then the maximizing completed class solves
\[
                   T_{S,q}(a_*^q)|_V=\nu
\]
on all original compact tests. Its selected representative retains the
actual exterior $e$ and ordinary boundary value $g$. It is the unique
relative class in $\mathscr K$ solving this interior equation.
\end{theorem}

\noindent\textbf{Status: \texttt{CONDITIONAL}} as a Dirichlet existence
criterion: the stated uniform reserve and margins are the additional
input. Their availability is not claimed for arbitrary data or from a
bounded subsolution. The implication itself is proved below for every
remaining degree and every allowed fixed background.

\begin{proof}
Fix $\phi$ and choose a positive $t_\phi$ with
\[
 t_\phi\le\epsilon_\phi,
       \qquad t_\phi\|\phi\|_\infty\le b_\phi.
\]
On $O_\phi$ the decomposition
\[
 u_j\pm t_\phi\phi
   =(u_j-\epsilon_\phi\rho_\phi)
      +\epsilon_\phi
            (\rho_\phi\pm(t_\phi/\epsilon_\phi)\phi)
\]
is a sum of classical psh functions. The second factor is classical by
convexity of the classical cone and the two compensator endpoints.
Off the support of $\phi$ it equals $u_j$. These local assertions glue
on the ordinary topology. The value margins preserve all bounds and
barrier inequalities, and the compact support inside $V$ preserves the
entire actual exterior. Therefore
\[
                 u_j\pm t_\phi\phi\in\mathscr P_{\rm bd}
                              \qquad(j\ge j_\phi).
\]

Every maximizing sequence converges to $a_*$ in all-domain local outer
capacity by the R4 variational theorem. Adding the same smooth function
does not change its positive gap sets. The two sequences just constructed
therefore define $a_*\pm t_\phi\phi$ in the same completed constraint
$\mathscr K$. The preceding proposition, or the exact R4 Euler
inequality applied to these two competitors, gives
\[
                     \int_V\phi\,d(\mu_{a_*}-\nu)=0.
\]
This holds for every original compact smooth scalar test in $V$, so it
is the required Radon-measure equality. No maximum locus has been
treated as an ordinary open set, and no contact-support formula is used.

For two solutions $a,b\in\mathscr K$, their selected actual difference
vanishes on $\Omega\setminus V$, including the ordinary boundary.
Consequently
\[
 I_S(a,b)=\int(a-b)\,d(\mu_b-\mu_a)
       =\int_V(a-b)\,d(\mu_b-\mu_a)=0.
\]
The fixed-collar energy-capacity estimate proves equality of their
relative classes. The preserved continuous barriers give the ordinary
boundary limit, as in the R4 interior theorem.
\end{proof}

\paragraph{How the reserve can be checked, and what it does not assert.}
Condition (1) is an actual classical decomposition on a prescribed
ordinary neighborhood. For potentials represented on one finite unit
chart, a uniform Hessian margin that dominates the compensator Hessian
on an open neighborhood of the compact moment image verifies it by
convex composition after every complete-extension curve pullback.
For a finite cover one verifies the condition on every member of that
cover. A merely positive difference of two original currents is not
silently substituted for this actual-function hypothesis.

The radius must be independent of the index for each fixed test.
Vanishing margins or reserves along a maximizing row do not justify the
limit argument. The smooth nonpolar counterexample of R4 shows why some
additional condition of this kind is necessary for this route: ordinary
boundary barriers and a bounded subsolution alone do not supply the
two-sided psh variations. No strict psh exhaustion of an arbitrary
analytic domain is assumed here.

% END source: sections/dirichlet_admissible_variations.tex


% BEGIN canonical R6 source: sections/dirichlet_approximate_variations.tex
\section{Approximate admissible directions and vanishing reserves}
\label{sec:r6-approximate-dirichlet-variations}

Retain the arbitrary fixed original-positive separately closed $S$, all
$q=n-p\ge1$, the common defining bound and collar, the convex completed
constraint $\mathscr K$, and the actual boundary data of
Section~\ref{sec:r6-dirichlet-admissible-variations}. Write
$\mu_a=T_{S,q}(a^q)$, let $\widetilde\nu$ be the same admissible datum,
and put
\[
 F(a)=\mathcal E(a)-\int_L(a-h)\,d\widetilde\nu,
 \qquad m=\max_{\mathscr K}F,
 \qquad \nu=\widetilde\nu|_V.
\]
All pairings below use actual selected representatives, or equivalently
their relative classes: each measure being used annihilates the relevant
relative-null sets. Measures are restricted to the fixed compact $L$
when an $L^1$ norm is written. No whole-domain finite mass is required.

\begin{lemma}[The completed first-order energy remainder]
\label{lem:r6-completed-energy-remainder}
For $a,b\in\mathscr K$ define
\[
 J_a(b)=\mathcal E(a)-\mathcal E(b)
                         +\int_L(b-a)\,d\mu_a.
\]
Then
\begin{equation}
 \frac{I_S(a,b)}{q+1}\le J_a(b)
                       \le\frac{qI_S(a,b)}{q+1},
 \qquad
 F(b)-F(a)=\int_L(b-a)\,d(\mu_a-\widetilde\nu)-J_a(b).
 \label{eq:r6-energy-remainder-identity}
\end{equation}
The map $b\mapsto J_a(b)$ is convex, $J_a(a)=0$, and for each fixed
$b$ one has $J_a(a+t(b-a))=O(t^2)$ as $t\downarrow0$.
\end{lemma}

\noindent\textbf{Status: \texttt{DIRECT\_COROLLARY}} of
Corollary~\textup{R4--38.4} and the completed weighted
calculus; it is not a new energy inequality.

\begin{proof}
For classical common-collar potentials the frozen polynomial calculation
gives
\[
 J_a(b)=\sum_{r=0}^{q-1}\frac{q-r}{q+1}
       \int\Gamma_{T_{S,q-1}(b^r,a^{q-1-r})}(b-a).
\]
Pass to bounded defining classical rows. Every term is evaluated with
one fixed compact original smooth cutoff equal to one near $L$, so the
completed weighted calculus passes each finite term. Positivity gives
the two stated bounds. The objective identity is its defining algebra.
Convexity follows from concavity of $\mathcal E$. Along a fixed segment
the energy is a polynomial, and the derivative at zero is precisely
$\int(b-a)\,d\mu_a$. Thus its constant and linear remainders cancel,
which proves the $O(t^2)$ assertion.
\end{proof}

\begin{theorem}[A criterion using approximate feasible variations]
\label{thm:r6-approximate-variation-dirichlet}
Let $a_j\in\mathscr K$ be a maximizing sequence and set
$\epsilon_j=m-F(a_j)\to0$; monotonicity of these numbers is
unnecessary. Suppose that for every original real
$\phi\in A_c^{0,0}(V)$ there are numbers $t_j>0$ and feasible
$b_j^+,b_j^-\in\mathscr K$ such that
\begin{align}
 &\epsilon_j/t_j\longrightarrow0,
       \qquad I_S(a_j,b_j^+)/t_j\longrightarrow0,
       \qquad I_S(a_j,b_j^-)/t_j\longrightarrow0,
       \label{eq:r6-variation-cost-condition}\\
 &\left\|(b_j^\pm-a_j)/t_j\mp\phi\right\|_
                 {L^1((\mu_{a_j}+\widetilde\nu)|_L)}
                         \longrightarrow0.
       \label{eq:r6-variation-direction-condition}
\end{align}
Then the maximizing completed class $a_*$ solves
\[
                          \mu_{a_*}|_V=\nu
\]
on all original compact smooth tests. It retains the prescribed actual
exterior and ordinary boundary trace, and is the unique solution class
in $\mathscr K$.

More precisely, for each sign and each index,
\begin{equation}
 \pm\int_V\phi\,d(\mu_{a_j}-\nu)
 \le \frac{\epsilon_j}{t_j}
      +\frac{qI_S(a_j,b_j^\pm)}{(q+1)t_j}
      +\left\|(b_j^\pm-a_j)/t_j\mp\phi\right\|_
                    {L^1((\mu_{a_j}+\widetilde\nu)|_L)}.
 \label{eq:r6-approximate-variation-defect}
\end{equation}
\end{theorem}

\noindent\textbf{Status: \texttt{CONDITIONAL}} as a solvability
criterion, with the explicitly stated feasible-direction conditions.
Its implication and error estimate are \texttt{DIRECT\_COROLLARY} of
the frozen energy identities. No availability of these competitors for
arbitrary data is asserted.

\begin{proof}
Since $F(b_j^\pm)\le m$, the exact identity
\eqref{eq:r6-energy-remainder-identity} gives
\[
 \int_L\frac{b_j^\pm-a_j}{t_j}
                 \,d(\mu_{a_j}-\widetilde\nu)
       \le \frac{\epsilon_j}{t_j}+\frac{J_{a_j}(b_j^\pm)}{t_j}.
\]
The signed residual is dominated in total variation by
$\mu_{a_j}+\widetilde\nu$. Subtract the indicated approximate
direction and use the upper bound for $J$ to obtain
\eqref{eq:r6-approximate-variation-defect}.

Every maximizing sequence converges to $a_*$ in all-domain local outer
relative capacity. The completed calculus therefore gives
$\mu_{a_j}\to\mu_{a_*}$ on the fixed original test $\phi$.
Let $j\to\infty$ in both signs of the estimate. The residual pairing
vanishes for every such $\phi$, proving equality of the Radon measures
on $V$. The boundary and uniqueness assertions are those of
Theorem~\ref{thm:r6-reserve-dirichlet-criterion}: actual differences of
two feasible representatives vanish outside $V$, including its
boundary, so the interior equation gives $I_S(a,b)=0$. No assertion
that boundary Monge--Ampere masses themselves agree is needed.
\end{proof}

\begin{corollary}[An admissible radius may tend to zero]
\label{cor:r6-vanishing-reserve-dirichlet}
Suppose for each original $\phi\in A_c^{0,0}(V;\mathbf R)$ there are
positive radii $r_{j,\phi}$ such that
\[
                  a_j\pm r_{j,\phi}\phi\in\mathscr K,
          \qquad \frac{m-F(a_j)}{r_{j,\phi}}\longrightarrow0.
\]
Then the conclusion of the preceding theorem holds. The radii need not
have a positive lower bound.

For classical $a_j$, these hypotheses are verified by the actual
curvature reserve and value margins of
Theorem~\ref{thm:r6-reserve-dirichlet-criterion}, with index-dependent
positive reserves $\varepsilon_{j,\phi}$ and margins $b_{j,\phi}$,
provided one can choose
\[
 0<r_{j,\phi}\le\varepsilon_{j,\phi},\qquad
 r_{j,\phi}\|\phi\|_\infty\le b_{j,\phi},\qquad
 (m-F(a_j))/r_{j,\phi}\longrightarrow0.
\]
\end{corollary}

\noindent\textbf{Status: \texttt{DIRECT\_COROLLARY}} as an implication;
the existence of such a rate-controlled maximizing row is an additional
\texttt{CONDITIONAL} hypothesis, not a consequence of a bounded
subsolution.

\begin{proof}
Write $\epsilon_j=m-F(a_j)$. If $\epsilon_j>0$, choose
$t_j=\min(r_{j,\phi},\sqrt{\epsilon_j})$. Then
$t_j\to0$ and $\epsilon_j/t_j\to0$. If $\epsilon_j=0$,
choose $t_j=\min(r_{j,\phi},1/j)$. Convexity preserves feasibility
at the smaller radius. Proposition
\ref{prop:r6-admissible-stationarity-estimate} gives
\[
 \left|\int_V\phi\,d(\mu_{a_j}-\nu)\right|
               \le\epsilon_j/t_j+C_\phi t_j/2\longrightarrow0.
\]
The maximizing-sequence limit proves the equation. For the final
assertion, the compensator decomposition and ordinary-open gluing in
the earlier reserve proof apply separately for each index, with the
displayed radius. No fixed-radius perturbation is passed to the limit.
\end{proof}

\subsection{The exact tangent cone that these tests must reach}

\begin{proposition}[Feasible tangent directions and the residual]
\label{prop:r6-energy-feasible-tangent-cone}
Fix $a\in\mathscr K$ and set
$\kappa_a=(\mu_a+\widetilde\nu)|_L$. In $L^1(\kappa_a)$ let
\[
 \mathfrak T_a=\overline{\{c(b-a):c\ge0,\ b\in\mathscr K\}}^{L^1}.
\]
This is a closed convex cone. A function $\xi$ belongs to it if and
only if there exist $t_j\downarrow0$ and $b_j\in\mathscr K$ with
\[
 (b_j-a)/t_j\longrightarrow\xi\quad\hbox{in }L^1(\kappa_a),
                       \qquad J_a(b_j)/t_j\longrightarrow0.
\]
At the variational maximizer the residual has the exact polar condition
\[
             \int_L\xi\,d(\mu_a-\widetilde\nu)\le0
                                    \qquad(\xi\in\mathfrak T_a).
\]
Conversely this condition characterizes the maximizer. The residual
annihilates the lineality space $\mathfrak T_a\cap(-\mathfrak T_a)$.
Consequently, if that space contains every original compact smooth
scalar test in $V$, the maximizer solves the interior equation.
\end{proposition}

\noindent\textbf{Status: \texttt{DIRECT\_COROLLARY}} of convexity and
the proved energy polynomial. This identifies a precise remaining
admissibility condition; it does not assert that this lineality space
contains all smooth tests for a general background.

\begin{proof}
The set before closure is a convex cone: the sum of two elements is
their total nonnegative coefficient times a convex combination of the
two competitors minus $a$. Any sequence in the displayed equivalent
description lies in this cone, proving one inclusion.

For the converse choose $\xi_j=c_j(v_j-a)\to\xi$ in $L^1$.
Along the fixed segment to $v_j$, the last assertion of
Lemma~\ref{lem:r6-completed-energy-remainder} implies
\[
 \frac{J_a(a+t c_j(v_j-a))}{t}\longrightarrow0
                       \quad(t\downarrow0)
\]
for each fixed $j$. Choose $t_j>0$ decreasing to zero, with
$t_jc_j\le1$, and with this quotient at most $1/j$.
Then $b_j=a+t_jc_j(v_j-a)$ is feasible, its normalized difference is
exactly $\xi_j$, and it has the required negligible first-order cost.

The functional $\xi\mapsto\int\xi\,d(\mu_a-\widetilde\nu)$ is
continuous on $L^1(\kappa_a)$ with norm at most one. The frozen Euler
inequality holds on $\mathscr K-a$, hence on its cone and closure.
Conversely the polar inequality includes every $b-a$ and therefore is
the frozen characterization of maximality. Using both $\xi$ and
$-\xi$ gives zero on the lineality space. For original tests this is
precisely the stated measure equality.
\end{proof}

The first-order energy requirement does not enlarge this closed
feasible tangent cone: the proposition proves equality with the usual
$L^1$ tangent cone. In particular energy estimates cannot silently
remove incompatible trace, curvature, or actual-boundary constraints.
Vanishing reserves alone give no conclusion; the objective-defect to
radius ratio in Corollary~\ref{cor:r6-vanishing-reserve-dirichlet} is the
quantitative condition that permits them.

% END source: sections/dirichlet_approximate_variations.tex


% BEGIN canonical R6 source: sections/dirichlet_actual_variation_radius.tex
\section{An exact curvature test for actual admissible variations}
\label{sec:r6-actual-variation-radius}

The new bounded-difference positivity theorem removes one previously
separate feasibility premise in the variational argument. It applies
to actual classical potentials and their smooth perturbations. It
does not identify an arbitrary relative completed representative with
an ambient classical function.

\begin{lemma}[Smooth directions belong locally to the difference class]
\label{lem:r6-smooth-direction-dsh}
Every real original smooth function $\phi$ is locally a difference
of locally bounded classical psh functions. Consequently, for actual
$a\in\mathcal C_b(\Omega)$ and $t\in\mathbf R$,
\[
 a+t\phi\in\mathcal C_b(\Omega)
       \quad\Longleftrightarrow\quad D[a]+tD\phi\ge0
\]
on all original compact complementary positive tests.
\end{lemma}

\noindent\textbf{Status: \texttt{DIRECT\_COROLLARY}} of
Theorem~\ref{thm:r6-positive-bounded-dsh-classical}.

\begin{proof}
Near a given point write $\phi=F(t_1,\ldots,t_r)$ in original
unit logarithms. Shrink to a relatively compact convex real ball
on which the Hessian of $F$ is bounded. For a sufficiently large
$M$, both $M\sum t_i^2$ and $M\sum t_i^2+F(t)$ are convex
there. Their pullbacks are bounded classical on a smaller ordinary
neighborhood: after every complete extension and curve pullback the
unit logarithms are harmonic, and convex composition is subharmonic.
Their difference is $\phi$. Thus $a+t\phi$ is locally a bounded
classical difference for either sign of $t$. The new positivity
criterion, applied locally and glued on ordinary opens, proves the
implication from the curvature condition. The converse
implication is the established positivity of classical curvature.
\end{proof}

Retain the actual constrained family $\mathscr P_{\rm bd}$ and its
data from Section~\ref{sec:r6-dirichlet-admissible-variations}.
In particular every member is bounded by $B$, equals the prescribed
actual exterior $e$ on $\Omega\setminus V$, and lies between
the actual barriers $\ell,r$ on $V$. Fix
$a\in\mathscr P_{\rm bd}$ and a real original
$\phi\in A_c^{0,0}(V)$. On $\operatorname{supp}\phi$ put
\[
 m_a(x)=\min\{B-a(x),B+a(x),a(x)-\ell(x),r(x)-a(x)\}\ge0.
\]
Define two extended nonnegative numbers by
\begin{align}
 r_{\rm val}(a,\phi)
   &=\inf_{\phi(x)\ne0}\frac{m_a(x)}{|\phi(x)|},\\
 r_{\rm curv}(a,\phi)
   &=\inf_{\substack{\eta\in A_c^{n-1,n-1}(\Omega),\ \eta\ge0\\
                     \langle D\phi,\eta\rangle\ne0}}
       \frac{\langle D[a],\eta\rangle}
            {|\langle D\phi,\eta\rangle|}.
       \label{eq:r6-original-curvature-radius}
\end{align}
The infimum of an empty family is $+\infty$. Tests in the second
line range over the full original positive cone, with arbitrary
compact support and moment refinements.

\begin{proposition}[Exact two-sided radius in the actual family]
\label{prop:r6-exact-actual-variation-radius}
For $t\ge0$,
\[
 a+t\phi,\ a-t\phi\in\mathscr P_{\rm bd}
   \quad\Longleftrightarrow\quad
  t\le\min\{r_{\rm val}(a,\phi),r_{\rm curv}(a,\phi)\}.
\]
In particular these two explicitly defined quantities give the
maximal admissible two-sided radius, including its finite endpoint.
The curvature radius is equivalently
\[
 \sup\{t\ge0:-D[a]\le tD\phi\le D[a]\}.
\]
It is unchanged by adding an actual pluriharmonic function to $a$;
the value radius need not be unchanged.
\end{proposition}

\noindent\textbf{Status: \texttt{DIRECT\_COROLLARY}}. This is
an actual feasibility criterion, not a Dirichlet existence assertion.

\begin{proof}
The two original-current inequalities $D[a]\pm tD\phi\ge0$
hold exactly when
\[
 t|\langle D\phi,\eta\rangle|
                     \le\langle D[a],\eta\rangle
\]
for every displayed positive test. Tests with zero denominator
give an automatic inequality. Thus the curvature inequalities are
equivalent to $t\le r_{\rm curv}$, including the infimum endpoint.
The preceding lemma identifies them with actual classicality of
both perturbed functions.

The simultaneous bounds $|a\pm t\phi|\le B$ and
$\ell\le a\pm t\phi\le r$ are exactly the inequalities
$t|\phi(x)|\le m_a(x)$ on the support of $\phi$.
They hold precisely when $t\le r_{\rm val}$. Elsewhere nothing
changes. Compact support in $V$ preserves the entire prescribed
actual exterior, hence also the fixed outer collar. These are all
the defining conditions of $\mathscr P_{\rm bd}$, proving the
equivalence. The curvature-radius assertion about pluriharmonic
functions follows from $D[a+H]=D[a]$; the value constraints refer
to the actual chosen function and do not have that invariance.
\end{proof}

\begin{corollary}[The variational reserve can be checked in current order]
\label{cor:r6-distributional-reserve-upgrade}
In Theorem~\ref{thm:r6-reserve-dirichlet-criterion}, replace its
first hypothesis by the original-current inequality
\[
           D[u_j]\ge\epsilon_\phi D[\rho_\phi]
                        \quad\text{on }O_\phi
                        \qquad(j\ge j_\phi).
\]
Retain every other hypothesis, in particular the actual classical
maximizing sequence, the uniform positive reserve for each fixed
test, the actual bound and barrier margins, and the specified datum
class. Then the same completed interior Dirichlet conclusion and
relative-class uniqueness hold for every remaining degree and every
allowed fixed positive closed $S$.

The vanishing-radius version in Corollary
\ref{cor:r6-vanishing-reserve-dirichlet} can likewise use these current
inequalities in place of its actual classical reserve, while retaining
the stated quantitative relation between objective deficit and radius.
\end{corollary}

\noindent\textbf{Status: \texttt{DIRECT\_COROLLARY}} for this
replacement. As a Dirichlet criterion it remains \texttt{CONDITIONAL}
on the displayed reserves, margins and maximizing sequence; their
existence is not proved for arbitrary data.

\begin{proof}
On $O_\phi$ the actual function
$u_j-\epsilon_\phi\rho_\phi$ is a locally bounded classical
difference. Its original curvature is nonnegative by the new
hypothesis, so Theorem~\ref{thm:r6-positive-bounded-dsh-classical}
makes that actual function classical. This is precisely the previous
first hypothesis. The remaining proof, including the capacity passage
of the two actual perturbed rows and preservation of the ordinary
boundary data, is unchanged. The same reasoning applies separately
at each positive radius in the vanishing-radius argument; it does
not supply the required rate automatically.
\end{proof}

\paragraph{What remains essential.}
Curvature positivity now decides whether an actual bounded classical
potential admits the specified smooth perturbation. It does not
provide slack in the pointwise obstacle or norm constraints, nor a
uniform radius along a maximizing row. For an arbitrary completed
potential, an actual ambient difference presentation has not been
supplied, so the positivity theorem cannot simply be applied to it.
The completed variational inequalities still require either verified
feasible directions or a separate argument producing them.

% END source: sections/dirichlet_actual_variation_radius.tex


% BEGIN canonical R6 source: logs/dirichlet_atomic_variation_obstruction.tex
\section{An atomic solution with no nonzero two-sided smooth direction}
\label{sec:r6-atomic-variation-obstruction}

\noindent\textbf{Status: \texttt{DIRECT\_COROLLARY}} of original
one-dimensional torus calibration, atomlessness of original smooth
top-form measures, and the curve maximum principle. This example
limits a sufficient variational criterion; it is not a Dirichlet
nonexistence example.

Let
\[
 \Omega=\{-2<t<2\}\subset\mathbf G_{m,K}^{\an},
       \qquad t=-\log|T|,\qquad S=[1],
       \qquad a=\frac12+\frac18|t|.
\]
The function $a$ is continuous classical, satisfies
$1/2\le a<3/4$, and has ordinary boundary limit $3/4$ at both
ends of the annulus. With $\zeta_1$ the original unit Gauss point,
the exact original curvature is
\[
                           D[a]=\frac14\delta_{\zeta_1}.
\]
Thus $a$ is an actual bounded classical solution of this atomic
Dirichlet datum. In the normalized constraints $0\le u\le1$,
it has pointwise obstacle margin at least $1/4$ everywhere.

Nevertheless, if $\phi\in A_c^{0,0}(\Omega;\mathbf R)$ is
original smooth and $s>0$, then
\[
       a+s\phi,\ a-s\phi\in\mathcal C_b(\Omega)
                         \quad\Longrightarrow\quad \phi=0.
\]
In particular the curvature radius in
\eqref{eq:r6-original-curvature-radius} is zero for every nonzero
compact original smooth direction, despite the uniform value slack.

\begin{proof}
The slope jump of $|t|/8$ is $1/4$. Original torus integration
therefore gives the displayed Dirac measure with no rescaling of
the absolute value. The boundary values and the stated range follow
from the formula for $a$.

Every original smooth $(1,1)$-form on a curve has an atomless
coefficient measure. Indeed, on a compact stage a finite original
moment presentation and its positive coefficient controls dominate
its variation by finitely many positive smooth top-form measures.
The pushforward of each such control by that presentation is
calibrated integration of a smooth density on finitely many real
one-dimensional polyhedra. It gives zero mass to a single real
point, hence the positive measure gives zero mass to the inverse
image of that point. In particular it gives zero mass to any
individual analytic point. If the presentation has rank zero, its
top form is zero. Finite localization proves the assertion without
claiming that arbitrary positive currents are atomless. Applied to
the compact smooth top form $D\phi$, this shows that its signed
Radon measure has no atom at $\zeta_1$.

Positivity of both proposed classical curvatures would imply
\[
     -\frac14\delta_{\zeta_1}\le sD\phi
                          \le\frac14\delta_{\zeta_1}.
\]
Thus $D\phi$ would be carried by the singleton $\{\zeta_1\}$.
Its atomlessness forces $D\phi=0$. The smooth function $\phi$
is locally a bounded classical difference by
Lemma~\ref{lem:r6-smooth-direction-dsh}; the frozen actual kernel
theorem therefore makes it harmonic on the curve. Equivalently,
apply the new bounded-difference positivity criterion to both
$\phi$ and $-\phi$.

Choose a closed subannulus compactly contained in $\Omega$ whose
ordinary interior contains $\operatorname{supp}\phi$.
The two Gauss boundary values of $\phi$ vanish. The maximum
principle for both harmonic signs gives $\phi=0$ on that
subannulus, and it is already zero outside its support. This proves
the implication. Hence no positive two-sided radius is available
for a nonzero direction.
\end{proof}

\paragraph{Meaning for the earlier criterion.}
An objective can attain the correct atomic Dirichlet solution at a
point where its actual psh constraint has no nonzero two-sided
original smooth direction. The reserve theorem is therefore a
sufficient route to its Euler equation, not a necessary property of
all solvable data. The pointwise obstacle margins and the curvature
reserve are independent requirements. This example does not alter
the exact radius formula or the conditional variational conclusions.

% END source: logs/dirichlet_atomic_variation_obstruction.tex


% BEGIN canonical R6 source: sections/dirichlet_relative_hessian_obstruction.tex
\section{Relative Hessian concavity and hyperbolicity can fail}
\label{sec:r6-relative-hessian-obstruction}

\noindent\textbf{Status of the proposed general assertions:
\texttt{FALSE}.} The following is an explicit construction over the
original finite extension $K/\mathbf Q_p$, using smooth bounded classical
potentials and a smooth strongly positive separately closed background.
It excludes two possible extensions of the relative Dirichlet theory:
unconditional root-concavity of the relative top density, and an
elliptic admissibility cone obtained merely by taking the positive
component of that density. It does not challenge the established BT
products, gradient Cauchy--Schwarz inequality, or compact-collar energy
concavity.

Fix $q\ge2$, let $n=q+2$, and work on an ordinary bounded full
polyannulus $\Omega$ in the original $n$-torus. Write
\[
 t_i=\log|T_i|,\qquad
 \beta_i=d't_i\wedge d''t_i,\qquad
 \gamma=\sum_{i=5}^{n}\beta_i,\qquad
 dV=\beta_1\wedge\cdots\wedge\beta_n.
\]
The empty sum is zero when $q=2$. Set
\[
                S=[\beta_1\wedge\beta_2+\beta_3\wedge\beta_4].
\]
Each summand is strongly positive. Pullback to every enlarged original
moment presentation preserves this cone, and its wedge with a weakly
positive complementary test is a positive top form. Thus $S$ is in
particular positive in the original middle CLD cone. It is symmetric
and separately closed because all the coordinate logarithmic
differentials are closed. Its bidegree is $(2,2)$, so its remaining
degree is exactly $q$. The brackets mean the ordinary current of the
displayed smooth form, not a current pushed forward from a skeleton.

Define
\[
 u=\tfrac12(t_1^2+t_2^2+\sum_{i=5}^nt_i^2),\qquad
 v=\tfrac12(t_3^2+t_4^2+\sum_{i=5}^nt_i^2).
\]
These functions are original smooth and bounded on $\Omega$. They are
classical psh after every complete extension: on any analytic curve
each unit logarithm is harmonic, and the displayed convex quadratic
compositions are subharmonic. Their curvatures are
\[
                    Du=\beta_1+\beta_2+\gamma,
       \qquad Dv=\beta_3+\beta_4+\gamma.
\]

\begin{proposition}[A smooth strongly positive counterexample in every remaining degree]
\label{prop:r6-relative-root-concavity-false}
For every $a,b\ge0$, the following equality holds on all original
compact smooth scalar tests:
\begin{equation}
 T_{S,q}((au+bv)^q)
       =q!(a^2+b^2)(a+b)^{q-2}\,dV.
 \label{eq:r6-relative-density-polynomial}
\end{equation}
Consequently the $q$th root of the relative density is not concave on
the cone of semipositive Hessians. For $q=2$ one also has
\[
        T_{S,2}(u,u)=T_{S,2}(v,v)=2dV,
                  \qquad T_{S,2}(u,v)=0.
\]
Thus a reverse Cauchy--Schwarz or Alexandrov--Fenchel inequality for
these relative mixed top densities is false for arbitrary backgrounds.
\end{proposition}

\begin{proof}
Every $\beta_i^2$ is zero, and the $(1,1)$-factors commute. The
diagonal curvature of $au+bv$ has coefficients $a,a,b,b$ in its first
four directions and $a+b$ in every remaining direction. Wedge with the
first summand of $S$ leaves the product with coefficients
$b^2(a+b)^{q-2}$, and the second leaves
$a^2(a+b)^{q-2}$. Each surviving top monomial occurs $q!$ times in
the $q$th power. This proves the equality as an identity of original
smooth forms and hence on every original test and refinement.

The measure $dV$ is nonzero and positive. More explicitly, on this
standard torus its original calibration is $dV=\sigma_*(dx)$, whereas
$(\sum_i\beta_i)^n=n!dV$. A nonnegative compact smooth scalar depending
on $t$ with positive real integral therefore has positive integral
against $dV$. No equality of the lower-degree coefficient measures of
$S$ with skeleton-supported measures is needed.

At $u$ and $v$ the density in \eqref{eq:r6-relative-density-polynomial}
equals $q!$, but at their midpoint it equals $q!/2$. Hence its $q$th
root at the midpoint is strictly smaller than the arithmetic mean of
its endpoint roots. When $q=2$, comparison of the coefficient of $ab$
gives the vanishing mixed term; the two diagonal terms are positive.
\end{proof}

The failure is not caused by a zero eigenvalue in the endpoint toric
Hessians. Add to both $u$ and $v$ the same small positive multiple of
$\frac12\sum_i t_i^2$. Their Hessians in the displayed $n$ coordinate
directions become positive definite. The strict failure of the root
inequality persists by continuity of the finite polynomial.

For $q>2$ the two-free-slot reverse Cauchy--Schwarz failure can also
be retained by freezing the remaining slots. Put
$w=\frac12\sum_{i=5}^nt_i^2$ and
$R=S\wedge(Dw)^{q-2}$. With
$u_0=(t_1^2+t_2^2)/2$ and $v_0=(t_3^2+t_4^2)/2$ one obtains
\[
 R\wedge(Du_0)^2=R\wedge(Dv_0)^2=2(q-2)!dV,
                 \qquad R\wedge Du_0\wedge Dv_0=0.
\]
This is a computation of legitimate mixed products over the original
fixed $S$, not an assertion about wedging arbitrary currents.

\begin{proposition}[The positive density component is not an elliptic branch]
\label{prop:r6-relative-hyperbolicity-false}
For this same background, let $P(H)$ be the coefficient of
$S\wedge(\sum_{i,j}H_{ij}d't_i\wedge d''t_j)^q$ relative to $dV$,
where $H$ is a real symmetric matrix. The homogeneous polynomial $P$
is not hyperbolic in the direction $I$. Moreover the connected
component of $\{P>0\}$ containing $I$ contains matrices at which
adding a positive semidefinite direction decreases $P$.
\end{proposition}

\begin{proof}
On diagonal matrices the exact polynomial is
\[
 P(\operatorname{diag}(\lambda_1,\ldots,\lambda_n))
        =q!(\lambda_1\lambda_2+\lambda_3\lambda_4)
                                          \prod_{i=5}^n\lambda_i.
\]
For $A=\operatorname{diag}(1,1,-1,-1,0,\ldots,0)$,
\[
                         P(A+tI)=2q!(t^2+1)t^{q-2}.
\]
It has the nonreal roots $i$ and $-i$, whereas $P(I)=2q!>0$.
This violates the defining real-root condition for hyperbolicity.

The path
$H_s=\operatorname{diag}(1,1,s,s,1,\ldots,1)$,
$-1\le s\le1$, has $P(H_s)=q!(1+s^2)>0$ and joins $I$ to
$A_*=\operatorname{diag}(1,1,-1,-1,1,\ldots,1)$.
But
\[
                     DP_{A_*}(E_{33})=-q!<0,
\]
although $E_{33}$ is positive semidefinite. This proves the asserted
failure of ellipticity on the positive component.
\end{proof}

The indefinite matrix in the last proof is used only to test a proposed
enlargement of the admissible class. It is not called a psh Hessian.
The first proposition already disproves root-concavity using genuinely
classical psh potentials.

\paragraph{Separation from the valid energy inequalities.}
The proved Cauchy--Schwarz inequality is for first-derivative energy
measures $\Gamma_R(f,g)$ over a fixed positive closed degree-one tail
$R$. The failed inequality here concerns second-derivative mixed top
densities; it is a different statement. Likewise the valid second
variation of compact-collar Aubin--Mabuchi energy is
\[
 \frac{d^2}{ds^2}\mathcal E_{\rm AM}(a+s\xi)
       =-q\int\Gamma_{T_{S,q-1}((a+s\xi)^{q-1})}(\xi)\le0
\]
along admissible segments. Its sign needs positivity of the indicated
tail and no pointwise root-concavity of $P$. Consequently all earlier
relative energy arguments remain compatible with this counterexample.

% END source: sections/dirichlet_relative_hessian_obstruction.tex


% BEGIN canonical R6 source: sections/dirichlet_nonlinear_image_nonconvexity.tex
\section{The actual relative Dirichlet image need not be convex}
\label{sec:r6-nonlinear-dirichlet-image}

R3 already disproved the unrestricted classical relative subsolution
theorem. The new assertion here is different: two data which each
have explicit continuous classical solutions on the same domain can
have an unsolvable midpoint. This occurs in every remaining degree
$q\ge2$. It obstructs a convex-image extension of the degree-one
Hilbert separation criterion.

Fix $q\ge2$, set $n=2q$, and write the coordinates as two blocks
$X=(X_1,\ldots,X_q)$ and $Y=(Y_1,\ldots,Y_q)$. On the ordinary
unit polydisc
\[
 \Omega=\{|X_i|<1,\ |Y_i|<1\}\subset(\mathbf A_K^{2q})^{\an}
\]
take the closed analytic coordinate polydiscs
\[
       Z_1=\{Y=0\}\cap\Omega,\qquad
       Z_2=\{X=0\}\cap\Omega,
       \qquad S=[Z_1]+[Z_2].
\]
Their only intersection is the rigid origin $o$. The fixed $S$ is
symmetric, original-positive, separately closed, and has bidegree
$(q,q)$. Thus its remaining degree is $q$. Put
$\lambda=-\log|\pi|>0$ and, on a block polydisc, define
\[
       g_t(Z)=\max\{\log|Z_1|,\ldots,\log|Z_q|,-t\}.
\]
The maximum is assigned its finite constant value at $Z=0$.
It is continuous classical psh, lies in $[-t,0]$, and has ordinary
boundary limit zero. Let $a_1$ be the radius-$|\pi|$ Gauss point
of $Z_1$, and let $b_1,b_2$ be the radius-$|\pi|$, respectively
radius-$|\pi|^2$, Gauss points of $Z_2$. These are original analytic
points, intrinsically divisorial on their carriers, and are distinct
from $o$ and from the other carrier.

\begin{theorem}[Nonconvexity with fixed continuous matching boundaries]
\label{thm:r6-relative-dirichlet-image-nonconvex}
Let $\mathscr I_S$ consist of the measures
$T_{S,q}(w,\ldots,w)$ for all bounded actual
$w\in\mathcal C_b(\Omega)$ having ordinary boundary limit zero.
The measures
\[
             \mu=\delta_{a_1}+2^{-q}\delta_{b_2},\qquad
             \nu=\delta_{a_1}+\delta_{b_1}
\]
belong to $\mathscr I_S$ and have continuous classical solutions.
Neither $\mu+\nu$ nor $(\mu+\nu)/2$ belongs to $\mathscr I_S$.
In particular this actual relative zero-boundary image is not convex.
Nevertheless each of these two unsolvable data has a bounded
continuous classical subsolution with the same zero boundary values.
\end{theorem}

\noindent\textbf{Status: \texttt{NEW\_PROVED}} for the nonconvexity
assertion and its explicit all-$q$ construction. The general failure
of the classical relative subsolution theorem is already an R3 result;
it is not counted again as a new conclusion.

\begin{proof}
We first specify the frozen intrinsic uniqueness input. Each $Z_i$
is a connected ordinary $q$-polydisc relatively compact in its smooth
affine $q$-space. The function $g_\lambda$ is a bounded continuous
negative classical exhaustion with ordinary boundary value zero.
It therefore satisfies condition \textup{R4--H}.
Theorem~\textup{R4--48.4} gives uniqueness among
\emph{all bounded actual classical potentials}, including nontoric
and discontinuous ones, for each finite divisorial datum below.
We use this exact intrinsic theorem, not uniqueness merely in an
energy or completed class.

The original radial calibration gives
\[
          (Dg_\lambda)^q=\delta_{\zeta_{|\pi|}},\qquad
          (D(\tfrac12 g_{2\lambda}))^q
                                  =2^{-q}\delta_{\zeta_{|\pi|^2}}.
\]
At the indicated Gauss point the slope simplex has volume $1/q!$,
which the original top-product factor $q!$ cancels; positive scaling
gives the second identity. These are all-original-test identities,
as in Lemma~\textup{R4--60.1} and its frozen
smoothing calibration.

Define on the whole ambient $\Omega$
\[
       u=\max\{g_\lambda(X),\tfrac12 g_{2\lambda}(Y)\},
               \qquad
       v=\max\{g_\lambda(X),g_\lambda(Y)\}.
\]
Both are continuous classical, take values in $[-\lambda,0]$, and
have ordinary boundary limit zero. Indeed any ordinary boundary
point is a limit of points of $\Omega$ with at least one block norm
tending to one, where its corresponding branch tends to zero.
Their restrictions are
\[
 \begin{array}{c|cc}
       &Z_1&Z_2\\ \hline
  u&g_\lambda&\tfrac12 g_{2\lambda}\\
  v&g_\lambda&g_\lambda.
 \end{array}
\]
For example the unused branch of $u$ is exactly $-\lambda$ on
each opposite carrier, so it does not change the indicated restriction.
Integration-current compatibility and background linearity give,
for every bounded actual classical $w$,
\begin{equation}
 T_{S,q}(w^q)
   =i_{1*}(D(w|_{Z_1}))^q+i_{2*}(D(w|_{Z_2}))^q.
 \label{eq:r6-nonconvex-carrier-splitting}
\end{equation}
This is the same all-original-test identity as
\textup{R4--64.4}; it uses the restrictions of
classical potentials, not a pullback of an arbitrary current.
Thus $T_{S,q}(u^q)=\mu$ and $T_{S,q}(v^q)=\nu$.

Suppose $w$ were a bounded actual classical solution for $\mu+\nu$.
Its analytic restrictions remain bounded actual classical. The two
positive summands in \eqref{eq:r6-nonconvex-carrier-splitting} are
carried by their respective $Z_i$. Restriction off $o$ separates
them, and the total measure has zero mass at $o$. Positivity hence
forces the intrinsic equations
\[
        (D(w|_{Z_1}))^q=2\delta_{a_1},\qquad
        (D(w|_{Z_2}))^q=\delta_{b_1}+2^{-q}\delta_{b_2}.
\]
There is no possible cancellation on the intersection. Moreover
every intrinsic ordinary boundary point of $Z_i$ is an ambient
ordinary boundary point of $\Omega$: it is a limit of points of
$Z_i\subset\Omega$ and lies outside $\Omega$ because a free block
coordinate has unit norm. Thus both restrictions have boundary
limit zero.

The first intrinsic solution is $2^{1/q}g_\lambda$, whose actual
origin value is $-\lambda2^{1/q}$. For the second, set
\[
          A=(1+2^{-q})^{1/q},\qquad B=\tfrac12,
          \qquad h=(A-B)g_\lambda+B g_{2\lambda}.
\]
This is continuous classical and has zero boundary limit. The exact
two-shell radial calculation in
\textup{R4--60.3}--\textup{R4--60.4}
gives outer mass $A^q-B^q=1$ and inner mass $B^q=2^{-q}$.
Explicitly, the slope shell at the outer point has volume
$(A^q-B^q)/q!$, and the inner slope simplex has volume $B^q/q!$;
the original factor $q!$ gives precisely those masses. Its actual
origin value is
\[
                         h(o)=-\lambda(A+B).
\]
The pinned intrinsic uniqueness theorem therefore forces these
two different values for the single actual number $w(o)$.
They are indeed different: strict convexity of the $\ell^q$ norm,
for $q>1$, gives
\[
 2^{1/q}=\|(1,1)\|_q
      <\|(1,\tfrac12)\|_q+\|(0,\tfrac12)\|_q
      =(1+2^{-q})^{1/q}+\tfrac12.
\]
The vectors are not proportional, so the inequality is strict.
This contradiction excludes $\mu+\nu$.

If a bounded actual classical zero-boundary potential solved the
midpoint datum, multiplication by $2^{1/q}$ would solve $\mu+\nu$,
by positive homogeneity. The midpoint is therefore also excluded.
Finally mixed positivity gives
\[
            T_{S,q}((u+v)^q)\ge T_{S,q}(u^q)+T_{S,q}(v^q)
                                      =\mu+\nu.
\]
Thus $u+v$ is the asserted continuous matching subsolution for the
sum, and $2^{-1/q}(u+v)$ is one for the midpoint. This uses only
the positive mixed expansion, not an invalid root-concavity
inequality for arbitrary relative backgrounds.
\end{proof}

\paragraph{Uniform capacity bounds do not repair the example.}
The normalized functions $1+u/\lambda$ and $1+v/\lambda$ are actual
testers on the whole $\Omega$. Every point in
$\{a_1,b_1,b_2\}$ therefore has relative singleton capacity at
least $c_0=2^{-q}\lambda^{-q}$. Each of $\mu$, $\nu$ and their
midpoint has mass at most two. Hence for every fixed $\alpha>0$
they all satisfy, on every Borel $E$,
\[
  \eta(E)\le2c_0^{-(1+\alpha)}
                    \operatorname{Cap}_S(E;\Omega)^{1+\alpha},
       \qquad\eta\in\{\mu,\nu,(\mu+\nu)/2\}.
\]
If $E$ misses the three support points the inequality is immediate;
otherwise the stated singleton lower bound proves it. The failure
therefore concerns actual carrier compatibility, even within a
fixed finite compact support and fixed common growth bounds.

\paragraph{Linear versus nonlinear image geometry.}
For $q=1$ the curvature map is linear, and addition of two classical
zero-boundary solutions solves the sum datum. The strict inequality
above becomes equality. The construction thus identifies a nonlinear
obstruction to extending the degree-one convex-image/Farkas method.
It does not disprove a conditional feasible-tangent theorem, and it
makes no nonexistence assertion for the larger completed or energy
solution classes.

% END source: sections/dirichlet_nonlinear_image_nonconvexity.tex


% BEGIN canonical R6 source: sections/dirichlet_radial_carrier_criterion.tex
\section{An explicit trace-compatible relative Dirichlet criterion}
\label{sec:r6-radial-carrier-criterion}

The following qualified theorem gives a numerical solvability test,
not a condition formulated as membership in an unknown operator image.
It concerns a specified family of analytic integration backgrounds and
finite radial data. It does not replace the arbitrary-background
energy alternative or assert sufficiency of a general bounded psh
subsolution.

Fix integers $q\ge1$ and $m\ge2$, and put $n=mq$. On the ordinary
unit polydisc $\Omega\subset(\mathbf A_K^n)^{\an}$ write the
coordinates as $m$ blocks $Z^{(i)}$ of length $q$. Let $Y_i$ be
the coordinate $q$-polydisc on which all blocks other than $i$ vanish.
Thus $Y_i\cap Y_l=\{o\}$ for $i\ne l$, where $o$ is the rigid
origin. Choose constants $c_i>0$ and set
\[
               S=\sum_{i=1}^m c_i[Y_i].
\]
This is an allowed original-positive separately closed background of
bidegree $(n-q,n-q)$, with remaining degree $q$.

Put $\lambda=-\log|\pi|$. For each $i$ fix an increasing finite
list $0<t_{i1}<\cdots<t_{iN_i}$ of positive integer multiples of
$\lambda$, and nonnegative real numbers $m_{ij}$. Write $z_{ij}$
for the radius-$e^{-t_{ij}}$ Gauss point of $Y_i$, viewed as a point
of the original ambient analytic space. These points are intrinsically
divisorial, lie off every other carrier, and are distinct from $o$.
Zero masses are allowed. Define
\begin{align}
 \nu&=\sum_{i=1}^m\sum_{j=1}^{N_i}m_{ij}\delta_{z_{ij}},
       \label{eq:r6-radial-carrier-data}\\
 A_{ij}&=\left(c_i^{-1}\sum_{k=j}^{N_i}m_{ik}\right)^{1/q},
       \qquad A_{i,N_i+1}=0,\qquad t_{i0}=0,\\
 d_i&=\sum_{j=1}^{N_i}(t_{ij}-t_{i,j-1})A_{ij}.
       \label{eq:r6-radial-carrier-depth}
\end{align}

\begin{theorem}[Finite radial data on coordinate carriers]
\label{thm:r6-radial-carrier-solvability}
There is a bounded actual $u\in\mathcal C_b(\Omega)$ with ordinary
boundary limit zero and
\[
                         T_{S,q}(u^q)=\nu
\]
if and only if
\[
                              d_1=\cdots=d_m.
\]
When these depths have common value $d$, a continuous classical
solution is given explicitly by
\begin{align}
 g_t(Z)&=\max\{\log|Z_1|,\ldots,\log|Z_q|,-t\},\\
 h_i(Z)&=\sum_{j=1}^{N_i}(A_{ij}-A_{i,j+1})g_{t_{ij}}(Z),
       \label{eq:r6-radial-carrier-intrinsic-solution}\\
 u(Z^{(1)},\ldots,Z^{(m)})&=\max_{1\le i\le m}h_i(Z^{(i)}).
       \label{eq:r6-radial-carrier-ambient-solution}
\end{align}
It has $\|u\|_{L^\infty(\Omega)}=d$. Every bounded actual classical
zero-boundary solution has restriction $h_i$ on $Y_i$, and its
supremum norm is at least $d$. Uniqueness off the union of the
carriers is not asserted.
\end{theorem}

\noindent\textbf{Status: \texttt{DIRECT\_COROLLARY}} of the
frozen all-original-test radial calibration, intrinsic bounded
classical uniqueness, and integration-current compatibility.
The explicit equal-depth criterion and extension formula record
their consequence for this family; no new general relative
Dirichlet existence theorem is being claimed.

\begin{proof}
Every $A_{ij}-A_{i,j+1}$ is nonnegative. The functions $g_t$ and
$h_i$ are therefore continuous classical on the whole intrinsic
polydisc and have ordinary boundary limit zero. Their finite values
at the rigid origin are assigned by the displayed maxima. The
frozen radial-shell calculation
\textup{R4--60.3}--\textup{R4--60.4}
gives
\[
 (Dh_i)^q
   =\sum_{j=1}^{N_i}(A_{ij}^q-A_{i,j+1}^q)\delta_{z_{ij}}
   =c_i^{-1}\sum_{j=1}^{N_i}m_{ij}\delta_{z_{ij}}.
\]
These are identities against all original compact tests. At a shell,
the slope polytope has volume
$(A_{ij}^q-A_{i,j+1}^q)/q!$; the original top-product
calibration multiplies it by $q!$. Thus no additional factorial or
residue-degree normalization occurs. The same calculation is valid
when two successive slopes agree, giving zero mass at that shell.

Summation by parts gives
\[
 h_i(o)=-\sum_j t_{ij}(A_{ij}-A_{i,j+1})
       =-d_i.
\]
Moreover $-t_{ij}\le g_{t_{ij}}\le0$, so
$-d_i\le h_i\le0$ everywhere and the lower bound is attained at
$o$. Each intrinsic polydisc satisfies
\textup{R4--H}, for example with the negative
exhaustion $g_\lambda$. The pinned intrinsic uniqueness theorem
\textup{R4--48.4} therefore applies to these
finite divisorial measures among all bounded actual classical
solutions, including nontoric and discontinuous ones.

Suppose an ambient bounded actual classical $w$ solves the stated
problem. Its analytic restrictions to $Y_i$ are actual bounded
classical, and the frozen integration-current identity gives
\[
           T_{S,q}(w^q)=\sum_i c_i\,i_{i*}(D(w|_{Y_i}))^q.
\]
All summands are positive. Off $o$ their carriers are disjoint,
whereas the prescribed total measure has zero mass at $o$.
Consequently each intrinsic measure is exactly the one displayed
for $h_i$. An intrinsic ordinary boundary point is a limit of
points of $Y_i\subset\Omega$ and is outside $\Omega$, because
a free block coordinate has unit norm. It is therefore an ambient
ordinary boundary point, and the restriction of $w$ has boundary
limit zero. Intrinsic uniqueness gives $w|_{Y_i}=h_i$ at every
actual point. In particular the single number $w(o)$ equals every
$-d_i$. This proves necessity, the asserted trace uniqueness, and
the lower bound for the norm of any solution.

Conversely suppose all depths equal $d$. Pulling the continuous
classical $h_i$ back by the block projections and taking their
finite maximum gives the continuous classical function in
\eqref{eq:r6-radial-carrier-ambient-solution}. On $Y_i$, every
other branch has the constant value $h_l(o)=-d$, while
$h_i\ge-d$. Hence $u|_{Y_i}=h_i$. The same integration-current
identity proves $T_{S,q}(u^q)=\nu$.

The bounds $-d\le u\le0$ and $u(o)=-d$ give its exact norm.
At any ordinary boundary point of $\Omega$, at least one block
norm tends to one. Along that block every $g_{t_{ij}}$ tends to
zero, so its $h_i$ tends to zero. Since $u$ is at least this branch
and at most zero, it has ordinary boundary limit zero. This also
handles the case $d=0$, when all data vanish and $u=0$.
\end{proof}

\paragraph{A numerical distinction between remaining degrees.}
For $q=1$, the depth is the linear expression
\[
                       d_i=c_i^{-1}\sum_j t_{ij}m_{ij}.
\]
For higher $q$ it is a sum of $q$th roots of tail masses. Thus
the actual compatibility condition is still a finite explicit
condition, but generally not a linear constraint on the data.
The nonconvex image in Theorem
\ref{thm:r6-relative-dirichlet-image-nonconvex} is one instance.
Equal depths prove sufficiency here because the specified radial
traces have a common attained global minimum and possess the
explicit ambient block extensions. No such extension assertion is
made for general intersecting analytic carriers or arbitrary
intrinsic trace data.

\paragraph{Completed solvability has no equal-depth obstruction.}
For every datum \eqref{eq:r6-radial-carrier-data}, whether or not the
$d_i$ agree, the frozen completed class already contains a bounded
zero-boundary solution. This is a \texttt{DIRECT\_COROLLARY} of
Proposition~\textup{R4--66.1}, not a new
completed existence theorem. To verify its bound explicitly, choose
$T\ge\max_{i,j}t_{ij}$ and $C\ge\max_i A_{i1}$, and put
\[
 G=\max\left(-T,\max_{1\le i\le m,\,1\le l\le q}
                                   \log|Z^{(i)}_l|\right).
\]
On each carrier,
\[
          CG|_{Y_i}=C g_T\le A_{i1}g_T\le h_i.
\]
The last inequality follows by adding the nonnegative terms in
\eqref{eq:r6-radial-carrier-intrinsic-solution} and using
$g_{t_{ij}}\ge g_T$. Apply the frozen coordinate separator
construction with branch functions $h_i\circ\operatorname{pr}_i$,
backgrounds $c_i[Y_i]$, and the common continuous floor $CG$.
It gives uniformly bounded continuous ambient classical functions
increasing to an actual completed potential $a$ with
\[
       CG\le a\le0,\qquad \mathcal T_{S,q}(a^q)=\nu.
\]
The same floor gives ordinary boundary limit zero. Multiplying a
carrier background by the positive constant $c_i$ preserves every
outer-null hypothesis of that construction. Thus the equal-depth
condition distinguishes actual classical compatibility at the
shared point from the already available completed solution class.

% END source: sections/dirichlet_radial_carrier_criterion.tex


% BEGIN canonical R6 source: sections/dirichlet_radial_measure_carriers.tex
\section{Compact radial measures on transverse coordinate carriers}
\label{sec:r6-radial-measure-carriers}

The equal-depth criterion extends from finite shell data to arbitrary
finite measures on a compact radial interval. This is an explicit
nonatomic family in every remaining degree; the geometry of the
carriers and their ambient extensions remains specified.

Retain the coordinate carriers $Y_i$, weights $c_i>0$, and background
$S=\sum_{i=1}^m c_i[Y_i]$ of
Theorem~\ref{thm:r6-radial-carrier-solvability}. For each $i$, let
$\tau_i$ be a finite positive Radon measure supported in a compact
interval $[a_i,b_i]\subset(0,\infty)$. Let
$\sigma_i(t)$ be the radius-$e^{-t}$ diagonal Gauss point of $Y_i$,
viewed in the original ambient space, and put
\[
 \nu=\sum_i(\sigma_i)_*\tau_i,\qquad
 A_i(t)=\bigl(c_i^{-1}\tau_i([t,\infty))\bigr)^{1/q}
       \quad(t>0),\qquad
 d_i=\int_0^\infty A_i(t)\,dt.
\]
All depths are finite. The convention at a jump of the tail function
does not affect these integrals.

\begin{proposition}[A radial-measure compatibility criterion]
\label{prop:r6-radial-measure-carrier-solvability}
There is a bounded actual classical $u$ on the whole $mq$-dimensional
unit polydisc, with ordinary boundary limit zero and
$T_{S,q}(u^q)=\nu$, if and only if $d_1=\cdots=d_m$.
When the depths agree, write their common value as $d$ and define
\[
 H_i(r)=-\int_0^r A_i(t)\,dt\quad(0\le r\le\infty),\qquad
 h_i(Z)=H_i\left(-\max_{1\le l\le q}\log|Z_l|\right),
\]
where $h_i(o)=H_i(\infty)=-d_i$. Then
\[
                    u=\max_i h_i(Z^{(i)})
\]
is a continuous classical solution with norm $d$. Every bounded
actual classical solution has exactly these restrictions to the
carriers. No uniqueness away from their union is asserted.
\end{proposition}

\noindent\textbf{Status: \texttt{DIRECT\_COROLLARY}} of the
finite-shell calibration, locally uniform continuous-potential
stability, and frozen pointwise ambient uniqueness. The measure
$\tau_i$ may be nonatomic, singular continuous, or mixed.

\begin{proof}
We first prove the intrinsic equation on one $q$-polydisc, suppressing
the index $i$ and retaining its positive weight $c$. Choose positive
finite atomic measures $\tau_N$ with the same total mass as $\tau$,
supported in one fixed compact interval $[a/2,2b]$, and converging
weakly to $\tau$. Such measures are obtained by finite interval
partitions and one representative per part; the representatives may
be chosen at rational multiples of $-\log|\pi|$. Their tails
converge to the tail of $\tau$ at every continuity point, hence at
Lebesgue-almost every $t$. They are bounded by the same total mass
and vanish for $t>2b$. Therefore
\[
 A_N(t)=\bigl(c^{-1}\tau_N([t,\infty))\bigr)^{1/q}
                     \longrightarrow A(t)
                \quad\hbox{in }L^1(0,\infty).
\]
For atomic $\tau_N$, summation by parts identifies the displayed
integral profile with the finite positive combination of $g_t$ in
\eqref{eq:r6-radial-carrier-intrinsic-solution}. The original
radial-shell calibration thus gives
\[
                     (Dh_N)^q=c^{-1}\sigma_*\tau_N.
\]
The same calibration is valid at any real radius: it is the
slope-polytope computation used in
\eqref{eq:r6-arbitrary-radius-capacity}, with shell masses given by
the difference of the $q$th powers of adjacent slopes. In particular
no strict-radius or new-field normalization is required here.

The estimate
\[
                   \|h_N-h\|_{L^\infty(Y)}
                         \le\|A_N-A\|_{L^1(0,\infty)}
\]
holds even at the rigid origin, since every profile is constant past
$2b$. Thus $h$ is a bounded continuous classical function on the
whole intrinsic polydisc. Locally uniform continuous-potential
stability, Proposition~\ref{prop:r6-continuous-joint-varying-background}
with fixed background $[1]$, gives
\[
                         (Dh)^q=c^{-1}\sigma_*\tau
\]
against every original compact smooth test. On the right, continuity
of the Gauss section on the fixed compact interval proves weak
convergence. The profile is between $-d$ and zero, attains $-d$ at
the origin, and tends to zero at every ordinary boundary point.

Intrinsic uniqueness for these possibly diffuse measures is the
frozen Corollary~\textup{R4--22.5}, rather than an
assumption that a finite-atomic existence theorem includes diffuse
data. The unit polydisc is an ordinary open set with compact closure
in smooth affine space, so it satisfies
\textup{R4--(G)}. Two bounded actual classical functions
with ordinary boundary limit zero satisfy the strict compact-gap
condition in both orders: a finite open cover of the compact ordinary
boundary excludes every fixed positive-margin gap. Equal intrinsic
Monge--Ampere measures therefore force actual pointwise equality.

Now let $w$ be any ambient actual bounded classical solution. Its
restrictions are bounded classical on each carrier, and original
integration-current compatibility decomposes its relative measure
as the sum of their positive intrinsic measures weighted by $c_i$.
The carrier intersections consist only of the origin, while $\nu$
has no mass there: its radial parameters are bounded above. Positivity
therefore identifies every individual intrinsic measure with
$c_i^{-1}(\sigma_i)_*\tau_i$. Intrinsic ordinary boundary points
are also ambient ordinary boundary points. The preceding uniqueness
argument gives $w|_{Y_i}=h_i$, so its single actual origin value
equals every $-d_i$. This proves necessity and trace uniqueness.

If the depths agree, each $h_i$ has common attained minimum $-d$.
The finite maximum of their block pullbacks is continuous classical;
on $Y_i$ all other branches are the constant $-d$, so its restriction
is $h_i$. Integration-current compatibility gives the required
relative measure. The bounds and origin value give norm $d$.
At any ordinary boundary point some block norm tends to one; the
corresponding $h_i$ tends to zero, and the maximum is squeezed between
this branch and zero. This proves the boundary condition.
\end{proof}

\paragraph{Completed potentials and scope.}
The frozen completed carrier-gluing theorem also applies to these
continuous intrinsic profiles for arbitrary depths. Indeed, with
$T\ge\max_i b_i$ and
$C\ge\max_i(c_i^{-1}\tau_i((0,\infty)))^{1/q}$, one has
$C g_T\le h_i$ on every carrier, since $A_i\le C$ and is supported
in $[0,T]$. The same global floor $CG$ as in the finite-shell theorem
gives a bounded zero-boundary completed solution. Equal depths are
precisely the additional actual point-value compatibility for the
specified ambient classical construction. This statement neither
extends arbitrary intrinsic psh data to ambient classical functions
nor asserts a general relative subsolution theorem.

% END source: sections/dirichlet_radial_measure_carriers.tex


% BEGIN canonical R6 source: logs/dirichlet_radial_tail_depth.tex
\section{A finite-depth criterion for noncompact radial data}
\label{sec:r6-radial-tail-depth}

\begin{corollary}[Exact radial tail integrability]
\label{cor:r6-radial-tail-depth}
Retain the coordinate carriers and weights of
Proposition~\ref{prop:r6-radial-measure-carrier-solvability}, but let
each $\tau_i$ be any finite positive Radon measure on $(0,\infty)$.
Set
\[
 \nu=\sum_i(\sigma_i)_*\tau_i,
 \qquad d_i=\int_0^\infty
       \bigl(c_i^{-1}\tau_i([t,\infty))\bigr)^{1/q}\,dt
                  \in[0,\infty].
\]
The datum has a bounded actual classical solution on the whole
ambient unit polydisc with ordinary boundary limit zero if and
only if all $d_i$ are finite and equal. In that event the same
integral profiles and finite maximum from the cited proposition
give a continuous classical solution. The measures may accumulate
at the rigid origin or towards the ordinary boundary; no mass is
assigned to the origin itself.
\end{corollary}

\noindent\textbf{Status: \texttt{DIRECT\_COROLLARY}} of the
compact radial-measure construction, frozen actual measure order,
and uniform continuous-potential stability. This is a criterion for
the specified radial data, not for arbitrary measures.

\begin{proof}
On one intrinsic carrier put $\tau_N=\mathbf1_{[1/N,N]}\tau$.
The corresponding tail roots $A_N$ increase pointwise to
$A(t)=(c^{-1}\tau([t,\infty)))^{1/q}$, up to the irrelevant
choice of values at endpoints, and $d_N=\int A_N\uparrow d$.
The intrinsic construction in the compact radial-measure proposition gives continuous
classical zero-boundary $h_N$ with
\[
          (Dh_N)^q=c^{-1}\sigma_*\tau_N,
          \qquad h_N(o)=-d_N.
\]
If an actual bounded classical zero-boundary $h$ has intrinsic
measure $c^{-1}\sigma_*\tau$, then
$(Dh)^q\ge(Dh_N)^q$. The frozen measure-order lemma
\textup{R4--48.5} gives $h\le h_N$ at every actual
point. Thus $h(o)\le-d_N$ for all $N$, excluding $d=\infty$.
This use of order does not assume the existence of a solution for
the limiting datum.

If $d<\infty$, the integral profile
$H(r)=-\int_0^r A(t)\,dt$, with $H(\infty)=-d$, satisfies
\[
                         \|h_N-h\|_\infty\le d-d_N\longrightarrow0.
\]
Hence its radial pullback $h$ is bounded continuous classical on
the entire intrinsic polydisc, including the rigid origin. Uniform
BT stability and monotone convergence of the finite measures give
$(Dh)^q=c^{-1}\sigma_*\tau$ on all original compact tests.
Since $A(t)\le(c^{-1}\tau((0,\infty)))^{1/q}$, the estimate
$|H(r)|\le r(c^{-1}\tau((0,\infty)))^{1/q}$ proves the ordinary
boundary limit zero. Intrinsic actual uniqueness follows from
Corollary~\textup{R4--22.5} and its two compact-gap
conditions, exactly as in the compact radial-measure proposition.

For an ambient solution, positivity of the carrier decomposition
identifies the individual intrinsic measures. The datum gives no
mass to the shared origin, and the other carrier parts are disjoint.
The preceding argument forces each depth to be finite, and actual
intrinsic uniqueness forces its value at the common origin to be
every $-d_i$. Conversely finite equal depths give the same
max-of-blocks solution as before. Boundary convergence follows from
the displayed linear bound near radial parameter zero. No compact
support of the datum is needed for these steps.
\end{proof}

\paragraph{A bounded completed solution when the depths differ.}
If every $d_i$ is finite, bounded completed solvability remains a
direct consequence of the frozen carrier-gluing theorem, even when
the depths differ. Choose
\[
 C>0,\qquad
 C\ge\max_i\bigl(c_i^{-1}\tau_i((0,\infty))\bigr)^{1/q},
 \qquad T>0,\qquad CT\ge\max_i d_i.
\]
For the
intrinsic radial parameter $r$,
\[
 h_i(r)\ge-Cr\quad(r\le T),\qquad
 h_i(r)\ge-d_i\ge-CT\quad(r\ge T).
\]
Thus $C g_T\le h_i$ on every carrier, and the same global
continuous floor $CG$ supplies bounded completed gluing and the
ordinary zero-boundary squeeze. This paragraph asserts sufficiency
only in the completed class; it makes no boundary-representative
claim for an arbitrary completed solution.

% END source: logs/dirichlet_radial_tail_depth.tex


\part{Failed mechanisms and the next proof gates}

% BEGIN canonical R6 source: logs/analytic_translation_smoothing_obstruction.tex
\section[Small analytic translations and transverse Cantor currents]{Small analytic translations do not regularize a transverse Cantor current}
\label{sec:r6-translation-smoothing-obstruction}

\noindent\textbf{Status: \texttt{FAILED\_ATTEMPT}} for a proposed
regularization mechanism; the construction and its calculation are
\texttt{DIRECT\_COROLLARY} of torus calibration and the finite
polyhedral description of actual delta-forms. This is not a counterexample
to weak delta approximation or to the conjectural equality of all allowed
backgrounds with the weak approximation class.

Fix the original finite extension $K/\mathbf Q_p$, choose a uniformizer
$\pi$, and put $\lambda=-\log|\pi|>0$. On the ordinary bidannulus
\[
 V=\{e^{-2\lambda}<|T_i|<e^{2\lambda}:i=1,2\},
 \qquad t_i=\log|T_i|,
\]
let $\nu$ be a nonzero finite atomless singular measure supported on a
middle-thirds Cantor subset of $(-\lambda/2,\lambda/2)$. Define
\[
             g(s)=\int(s-r)_+\,d\nu(r),\qquad
             S=D[g(t_1)].
\]
The function $g$ is convex and $C^1$: its derivative is the continuous
distribution function of the atomless measure $\nu$. Its distributional
second derivative is exactly $\nu$. The potential $g(t_1)$ is continuous,
locally bounded and classical psh. Indeed convex piecewise affine
approximants are finite sums of maxima of affine functions of the
pluriharmonic function $t_1$, and converge locally uniformly to it.
Thus $S$ is an allowed symmetric positive separately closed original
$(1,1)$-current on the smooth two-dimensional $K$-space $V$.

\begin{proposition}
The following assertions hold.
\begin{enumerate}
\item For every $a_1,a_2\in K$ with $|a_i|<e^{-2\lambda}$, the
analytic automorphism
\[
             \tau_a(T_1,T_2)=(T_1+a_1,T_2+a_2)
\]
of $V$ satisfies $\tau_a^*S=S$. Consequently averaging such translates,
with any probability law for which the scalar averages are defined,
leaves $S$ unchanged.
\item On any smaller bidannulus whose $t_1$-interval contains the support
of $\nu$, the current $S$ is not the ordinary current of an actual
delta-form.
\item Nevertheless $S$ belongs to the fixed-domain weak closure of
positive separately closed actual delta-forms.
\end{enumerate}
\end{proposition}

\begin{proof}
At every Berkovich point of $V$, the ultrametric inequality and
$|a_i|<|T_i|$ give $|T_i+a_i|=|T_i|$. The inverse translation has
the same property, so $\tau_a$ is an automorphism and
$g(t_1)\circ\tau_a=g(t_1)$. Naturality of $D$ under an analytic
isomorphism proves the first assertion. The scalar average of a fixed
current is that same current.

For the second assertion, choose nonnegative smooth real cutoffs
$\chi_1,\chi_2$ compactly supported in the smaller radius intervals,
with $\chi_1=1$ on $\operatorname{supp}\nu$ and
$\int\chi_2(s)\,ds=1$. The smooth original form
\[
                \eta=\chi_1(t_1)\chi_2(t_2)
                       d't_2\wedge d''t_2
\]
has compact analytic support. Original torus calibration, first for
smooth convex $g$ and then by locally uniform convergence of its
convex approximants, shows that the pushforward of the positive top
measure $S\wedge\eta$ under $t_1$ is exactly $\nu$.
All constants in this assertion use the original absolute value; the
two coordinate moments have calibration one.

Suppose instead that $S=[\sigma]$ for an actual delta-form $\sigma$
near this compact support. Then $S\wedge\eta=[\sigma\wedge\eta]$.
Take a finite common moment presentation of this compactly supported
top delta-form, adjoining $T_1$ to the moments. The integration
definition expresses its measure as a finite sum of pushforwards of
integrable piecewise smooth densities on calibrated polyhedral cells.
On each cell, the coordinate $t_1$ is affine. If its linear rank is
one, pushing the density to the real line is absolutely continuous
by Fubini. If its rank is zero, the pushforward is supported at one
point. Hence the total pushforward is the sum of an absolutely
continuous signed measure and finitely many atoms. This remains true
after every original refinement and after adding the finitely many
local presentations. It cannot equal the nonzero singular continuous
measure $\nu$. This proves the contradiction without restricting a
delta-form along an arbitrary analytic embedding.

Finally, take convex polygonal interpolants $g_N$ on nested finite
partitions containing the endpoints of an interval which contains
$\operatorname{supp}\nu$. Extend them by the unchanged affine tails.
They are finite convex piecewise affine functions, and $g_N\to g$
uniformly on compact real intervals. Each
$D_\delta(g_N(t_1))$ is an actual delta-form whose ordinary current
is positive and separately closed. Distributional differentiation
of locally uniform functions gives
\[
                     [D_\delta(g_N(t_1))]\longrightarrow S
\]
on every original compact test of $V$. Thus $S$ lies in the same-domain
weak approximation class, despite being fixed by all the indicated
analytic translations.
\end{proof}

The failed inference is therefore specific: analytic automorphisms
which are sufficiently close to the identity in the non-Archimedean
additive coordinates do not move tropical radii at all. A complex-style
average over these automorphisms supplies no smoothing in the real
transverse variable. Convex polygonal approximation succeeds here by
changing that transverse distribution directly; the example supplies
no such operation for a general positive closed background.

% END source: logs/analytic_translation_smoothing_obstruction.tex


% BEGIN canonical R6 source: logs/higher_delta_relation_stratum_gate.tex
\section{Why one positive stratum decomposition does not settle higher delta relations}
\label{sec:r6-higher-delta-stratum-gate}

\noindent\textbf{Status: \texttt{FAILED\_ATTEMPT}} for a proposed
reduction. The explicit nonclosed restriction below is a proved scope
counterexample, not a counterexample to a full delta action.

The positive bounded-difference theorem makes every continuous scalar
relation among represented $(1,1)$-curvatures act independently of its
presentation over any allowed background. It does not alone identify
the kernel of the map from polynomials in those curvatures to actual
higher-codimension delta-forms. In particular, freezing scalar
coefficients in a relation between products still leaves a relation
between products. It does not turn that relation into a degree-one
curvature relation.

A tempting next step is to stratify the actual moment image, restrict
the positive background to each Borel stratum, and apply the closed
background cycle calculus to those restrictions. Level tangency does
hold after Borel restriction. Separate closedness, however, need not
survive, even in a very simple positive actual delta example.

\begin{example}[A wall restriction of a positive closed current is not closed]
\label{ex:r6-nonclosed-real-wall-restriction}
Let $K$ be the original field, and work on the ordinary tropical box
\[
     V=\{|\log|T_1||<1/4,\ |\log|T_2||<1/4\}
         \subset (\mathbf G_{m,K}^{2})^{\rm an}.
\]
Write $x=\log|T_1|$ and $y=\log|T_2|$, and set
\[
 u=x^2+y^2+(1+y)x_+,
 \qquad S=D[u],\qquad H=\{x=0\},\qquad R=\mathbf1_H S.
\]
Then $u$ is continuous classical psh, $S$ is a symmetric positive
separately closed original current, and $R$ is coefficient-Radon and
positive, but neither $d'R$ nor $d''R$ is zero.
\end{example}

\begin{proof}
On this box $1+y>0$, so
\[
 u=\max\{x^2+y^2, x^2+y^2+x+xy\}.
\]
The two smooth functions have real Hessian matrices
$\operatorname{diag}(2,2)$ and
$\left(\begin{smallmatrix}2&1\\1&2\end{smallmatrix}\right)$.
They are original smooth psh, and their maximum is continuous
classical psh. Its piecewise smooth expression also defines an
actual delta curvature. Therefore $S$ has all the stated positivity
and closedness properties. Borel multiplication defines $R$.

Choose $\phi,\psi\in C_c^\infty((-1/4,1/4))$, with
$\phi(0)=1$, $\psi\ge0$, and $\int\psi(y)\,dy>0$. Put
\[
                 \eta=\phi(x)\psi(y)d''y.
\]
This is an original compact smooth $(0,1)$-test. On the toric test
subspace, the smooth Hessian pieces of $S$ have two-dimensional
Lebesgue coefficient measures, so restricting their pairing with
$d'\eta$ to $H$ gives zero. The singular wall contribution is
\[
                    (1+y)\delta_{x=0}\,d'x\wedge d''x.
\]
The torus moments have calibration one. Since the term of
$d'\eta$ containing $d'x$ is annihilated by this wall factor,
\[
 \langle R,d'\eta\rangle
       =\phi(0)\int(1+y)\psi'(y)\,dy
       =-\int\psi(y)\,dy.
\]
With the original total-degree current convention,
\[
                \langle d'R,\eta\rangle
                       =\int\psi(y)\,dy>0.
\]
Likewise, using $\eta'=\phi(x)\psi(y)d'y$ gives
$\langle d''R,\eta'\rangle=-\int\psi(y)\,dy\ne0$.
Thus neither separate derivative vanishes. Every test in this
calculation is an original compact smooth test on the actual
Berkovich space.

The wall expression was used only to evaluate this toric test
subspace. It is not asserted to be a full-current formula for $R$
after all moment refinements. For example, a refined moment can
detect smooth coefficient mass in branches on which $|T_1|=1$.
Such mass contributes zero to the displayed toric tests, so it
cannot change the nonzero derivative calculation. This distinction
avoids an independent false assertion about all refined coefficient
measures.
\end{proof}

\paragraph{Exact remaining gate.}
Let $\mathscr P$ denote a finite local presentation algebra generated
by continuous piecewise linear curvature operators and compatible
piecewise smooth coefficient forms. The established constructions
annihilate ambient cycle relations, pointwise piecewise-polynomial
zero relations, and actual degree-one curvature relations. For a
general analytic moment presentation $t$, the remaining assertion is
that the additional kernel of
\[
        \mathscr P\longrightarrow B_X^{\bullet,\bullet},
        \qquad P\longmapsto t^*P,
\]
annihilates every allowed positive separately closed background.
The statement must be proved with the full original refinement
equivalence, not merely on the tangent spaces of one selected image.
The example explains why the following argument does not prove it:
``restrict the background to each face and reapply closed-background
cycle insertion.'' Those restricted backgrounds need not be closed.

This leaves two mathematically distinct requirements: either an actual
local relation theorem reducing the higher kernel to relations already
annihilated by the operator calculus, or a separate approximation or
continuity theorem controlling the corresponding current actions.
No such general relation theorem is asserted here. The already proved
matroidal saturated case supplies one sufficient relation theorem;
the weak delta approximation class supplies one sufficient continuity
mechanism. The example itself is in that approximation class and
therefore cannot refute either construction.

The wall $H=\{\log|T_1|=0\}$ is a real moment-level set, not the
complete analytic pole set used in the frozen polar-restriction
theorem. The latter theorem's separate closedness is unaffected.

% END source: logs/higher_delta_relation_stratum_gate.tex


% BEGIN canonical R6 source: sections/round6_status.tex
\section{Status change from R5 and the remaining minimal gates}
\label{sec:r6-status}

\subsection{What is already settled before this record}

The seven bounded BT properties over an arbitrary fixed original
positive separately closed background remain \texttt{PROVED} in
the frozen calculus. The arbitrary-background potential quotient
is also \texttt{PROVED} by R5. The present record neither reopens
that quotient nor claims those previous proofs as new results.
Its main further advance is the ordered potential quotient: positivity
of the original curvature of an actual bounded classical difference
forces that actual difference to be classical psh. This is stronger
than preservation of the zero kernel. Further results concern the
actual delta-test category, quantitative relative capacities, compact
represented curvature families, and verifiable Dirichlet constraints.

\begingroup
\setlength{\tabcolsep}{3pt}
\begin{longtable}{@{}p{.24\textwidth}p{.22\textwidth}p{.45\textwidth}@{}}
\toprule
Precise assertion & Before R6 & R6 conclusion\\
\midrule
\endhead
Ordinary fixed-$S$ bounded calculus and differentiated potential
quotient & \texttt{PROVED} & Unchanged input. The positive-cone
question explicitly left aside in R5's signed-current section is
closed below only for its actual bounded classical difference class.\\[4pt]
Positive curvature in the actual bounded classical difference class &
R5 settled the zero kernel only & \texttt{NEW\_PROVED}: $D[a-b]\ge0$
if and only if the actual $a-b$ is classical psh; all complete-field
curve tests and ambient upper semicontinuity are proved.\\[4pt]
Represented curvature order over arbitrary $S$ & No general order
transfer for discontinuous differences & \texttt{DIRECT\_COROLLARY}:
$D[a]\ge D[b]$ implies $D(aS)\ge D(bS)$ in all differentiated slots;
explicit mixed order-error and top-degree variation bounds follow.\\[4pt]
Continuous scalar relations among bounded-potential curvatures &
Q handled constant differentiated-potential relations &
\texttt{NEW\_PROVED}: positive order freezing transfers all finite
continuous-coefficient relations to every $S$, even with discontinuous
classical slots. Positive products of this module are a direct
consequence; arbitrary Borel coefficient descent is false.\\[4pt]
Actual total-degree $(1,1)$ delta-tests over arbitrary $S$ & No
general actual relation theorem & \texttt{NEW\_PROVED}: canonical
ordinary local action, all refinements and ordinary supports; in
remaining degree one this gives a full closed complementary-test
extension and the specified positivity-class uniqueness.\\[4pt]
Ambient real cycle relations & Generator formulas did not remove
all relations & \texttt{NEW\_PROVED}: a piecewise polynomial
operator descends through the full ambient cycle ring over arbitrary
$S$, including real directions.\\[4pt]
Generation of every real delta-form & Constant-cycle generation
alone did not cover all smooth coefficient arrays &
\texttt{NEW\_PROVED}: exact finite anchored-face generation by
PS coefficients and constant real cycles, in every tridegree.\\[4pt]
Full actual delta-tests in arbitrary remaining degree & Ordinary
equality alone was insufficient & \texttt{NEW\_PROVED} for
$\mathcal W_{\rm loc}$, including all its bounded BT tails;
weighted joint decreasing convergence and both derivatives hold.
All-background higher-degree relation annihilation remains
\texttt{OPEN}.\\[4pt]
Geometric and diffuse examples in $\mathcal W$ & No R6 closure
description & \texttt{DIRECT\_COROLLARY}: smooth regular flags, singular effective
Cartier divisors, curvatures of coefficient-integrable extended
classical potentials, same-domain barycenters, and diffuse
polar-supported fiber currents.\\[4pt]
Saturated torus/matroidal charts & No all-background relation
mechanism at this scope & \texttt{NEW\_PROVED} under the exact
fixed-chart and calibration assumptions; locality is saturated
unless another ordinary-local theorem applies.\\[4pt]
Represented curvature compactness & No common-domain compactness
mechanism for these families & \texttt{NEW\_PROVED}: normalized
continuous order intervals are locally uniformly compact; the positive
DSH criterion removes the initial continuity premise. Compact upper-
potential families also work without any common Radon dominator.\\[4pt]
Weak curvature and varying-background stability & Unrestricted weak
stability is false & \texttt{DIRECT\_COROLLARY} in the proved compact
represented families. Explicit finite Shilov gauges give weighted
stability for those chosen local representatives.\\[4pt]
Relative energy versus normalized degree-one BT measures & A
square-root capacity upper bound was available &
\texttt{NEW\_PROVED}: Hardy constant one and strong capacitary
constant four, both optimal. The energy-capacity bound is linear.\\[4pt]
Capacity-growth data with finite degree-one energy & A weaker
sufficient growth condition & \texttt{PROVED}: powers
$\alpha>1/2$ suffice and the endpoint fails in general. The
concave-modulus theorem gives its precise universal extension.\\[4pt]
Reverse equality of energy and BT capacity null sets & Not
established & \texttt{FALSE}: a compact Gauss-section
Cantor-product set has zero BT outer capacity and positive faithful
energy capacity in every $n\ge3$.\\[4pt]
General relative C1 and C2 & \texttt{OPEN} outside stated
subclasses & Still \texttt{OPEN}; compact upper-potential families give C1, but
Borel C2 and outer continuity from below for their restricted capacities
are \texttt{FALSE} for every $p,q\ge1$: explicit smooth-torus examples have raw value
$1/2$ and outer value $1$. Common
Radon domination, uniformly $L^r$ relative density bounds obtained
from variable continuous curvature coefficients, and, in remaining degree one, energy-null
relative boundaries give qualified C2.
On full torus boxes, exact convex reduction gives C2 for finite
Hausdorff $(n-1)$-measure Borel sets, Borel coordinate sheets,
and relatively compact open sets with compact additions,
and null equivalence under sigma-finiteness. Exact diffuse rectangle/graph formulas also remain. A new exact
continuous psh interpolation construction on the fast pole Cantor set
extends the graph formula to arbitrary bounded Borel radial thresholds
and to Borel sets between those graphs and their filled radial tubes.\\[4pt]
Capacity domination characterizes the normalized image & Not
proved & \texttt{FALSE}: one original smooth square separates a
capacity-dominated annular measure from the weak normalized image.
This is a stronger inverse/image requirement, not a counterexample to C1.\\[4pt]
Completed degree-one Dirichlet feasibility & A general usable
criterion remained open & \texttt{NEW\_PROVED}: exact energy-dual
support inequalities, a separating smooth test, and a distance
formula for finite-energy data. Evaluating those inequalities in
new geometries can still be difficult.\\[4pt]
Relative energy-class Dirichlet solutions & Faithful form available &
\texttt{DIRECT\_COROLLARY}: finite-energy measures and dominated
bounded energy subsolutions give solutions in that energy class.
Classical psh feasibility is a distinct constraint.\\[4pt]
Actual nonlinear Dirichlet image & Earlier subsolution obstructions
existed; image convexity was not settled & \texttt{FALSE} in general:
for every $q\ge2$, two solvable zero-boundary data have unsolvable sum
and midpoint, despite a continuous classical subsolution and all finite
power-capacity bounds.\\[4pt]
Higher-degree stationary solutions & Feasible variations not
automatic & \texttt{CONDITIONAL} on the stated reserves and margins.
R6 computes the exact two-sided actual variation radius from value
margins and original current order; the new positive-DSH theorem
removes the formerly separate classicality premise for that reserve.\\[4pt]
Unrestricted relative Hessian root concavity and hyperbolicity &
No justified general hypothesis & \texttt{FALSE}: explicit smooth
strongly positive backgrounds give the displayed polynomial
counterexamples.\\
\bottomrule
\end{longtable}
\endgroup

\subsection{The eleven classical directions, with current qualifications}

\begin{longtable}{@{}p{.29\textwidth}p{.65\textwidth}@{}}
\toprule
Direction & Current mathematical scope\\
\midrule
\endhead
1. Locally bounded psh & B's actual class and regularization are
inputs. R5's zero kernel is continuous PH; R6 now identifies the
positive cone within the actual bounded classical difference class.
Every actual member under a continuous curvature upper bound is
continuous. Arbitrary bounded USC/Borel representatives are excluded.\\
2. Positive closed currents and mixed products & Ordinary arbitrary
fixed-$S$ calculus and Q are proved. R6 adds the order-preserving
quotient, continuous coefficient products, all-background degree-one
actual delta action and full higher-degree testing on the stated
weak approximation class.\\
3. CLN & The earlier finite-control relative CLN estimates remain.
They provide uniform masses in the new fixed-test and capacity
arguments; no arbitrary moving-background theorem is inferred.\\
4. Symmetry and weighted monotonicity & The ordinary fixed-$S$ theorem remains. R6 adds explicit order-error
bounds, qualified joint weak curvature/background convergence, and
weighted convergence after a fixed local normalization. Full delta-test
convergence retains the stated background/degree scopes.\\
5. Locality and maximum principles & Frozen ordinary and Borel
maximum locality remain. The new degree-one action has ordinary
locality; saturated chart theorems retain their more limited scope.
Strict gaps of two discontinuous psh functions need not be ordinary
open.\\
6. Capacities & General C1/C2 remain open. Sharp energy bounds, compact upper-potential
families, exact full-torus convex reduction, finite-Hausdorff Borel sets
and explicit diffuse polar carriers add distinct qualified mechanisms.\\
7. Quasicontinuity and capacity convergence & Frozen bounded and
energy-representative results retain their exact hypotheses. New
energy traces use these representatives, rather than arbitrary
coefficient-null modifications.\\
8. Pluripolar and negligible sets & Arbitrary backgrounds can charge
ambient complete-polar sets. Diffuse examples in this record are
not countable sums of fiber subcurrents. The BT and energy capacity
null ideals can themselves differ.\\
9. Comparison and domination & The qualified compact-gap and
relative-q.e. principles remain. Energy order comparison and
Dirichlet distance statements do not imply unrestricted actual
pointwise classical uniqueness.\\
10. Extremal functions & The frozen constrained and qualified
ambient results remain. The upper-contact counterexample rules
out a simple ambient-envelope repair of relative compact mass.\\
11. Monge--Ampere and Dirichlet & Linear energy solutions, completed-class dual feasibility and actual
classical nonlinear feasibility are distinguished. The latter image
is explicitly nonconvex in every $q\ge2$; scalar energy solutions and
capacity growth do not resolve the carrier compatibility constraint.
General existence retains the false/open scopes recorded here.\\
\bottomrule
\end{longtable}

\subsection{Primary comparisons that do not close the remaining gates}

Mihatsch's actual sheaf is faithful in ordinary currents
\cite[Proposition 4.11]{Mih24}. This resolves equality for two
actual delta-forms and is not an arbitrary-current extension theorem.
The real balanced definition and its product are
\cite[Theorem 3.3 and Theorem 4.1]{Mih23}. Esterov's ring theorem
\cite[Theorem 3.1]{Est12} removes ambient real fan relations;
Gross--Shokrieh \cite[Lemma 4.4 and Theorem 4.5]{GS21} supplies
the distinct matroidal relative-duality input. None of these inputs
was silently extended to arbitrary singular moment images.

The focused capacity-source check also preserves a genuine
distinction. Dhouib--Elkhadhra's discussion of arbitrary-background
capacity \cite[Section 3]{DE15} does not provide the missing general
compact continuity theorem. In Elaini--Elkhadhra,
\cite[Proposition 3]{EE21} gives basic relative-capacity properties,
whereas the extremal-capacity results used in Section 4.2 have the
explicit $T=1$ specialization. Their subsolution theorem has an
additional current class and approximation/tail hypotheses. It
cannot be cited as unconditional arbitrary-background classical
psh Dirichlet existence in this project. Exact printed locations
and hypothesis pins are retained in the source register.

\subsection{Minimal proof gates after the proved advances}

\begin{enumerate}
\item \textbf{Higher-degree actual relation annihilation.}
Given an actual zero delta-presentation of total complementary
degree at least two, show its prescribed finite PL/PS BT expression
vanishes over every allowed fixed $S$. Exact generation is now
proved, so surjectivity of the test presentation is no longer the
gap. The missing assertion is its analytic kernel. Weak positive
closed delta approximation is sufficient, not yet proved universal.
The degree-one order argument does not automatically lift to a
higher-degree relation theorem. The explicit real-wall example
\ref{ex:r6-nonclosed-real-wall-restriction} shows that Borel stratum
restrictions need not be separately closed. They therefore cannot
simply be reused as admissible backgrounds, despite level tangency.
The Cantor example also rules out small analytic translations as a
universal smoothing mechanism; it does not refute weak delta density.

\item \textbf{C1 without a uniform domination mechanism.}
On the same whole normalized domain, recover the compact mass of
a weak-limit or completed measure from actual normalized testers.
The new compact upper-potential argument solves this for its entire
represented family, even without one common measure dominator.
The general problem still needs lower compact-mass approximation,
not merely weak density or a capacitary upper bound on the datum.

\item \textbf{C2 beyond the established set and family classes.}
Control the simultaneous compact extraction and normalized-tester
pasting losses needed for arbitrary Borel ordinary-open regularity.
Already on a full torus box this reduces exactly to the upper envelope
of Laplacian masses of normalized convex functions. Finite Hausdorff
measure, Borel coordinate sheets, open sets with compact additions,
and sigma-finite nullity are solved
there; positive-capacity Borel sets outside those hypotheses remain a
concrete gate. Even compact normalized curvature families can have
restricted raw/outer values $1/2$ and $1$, so their compactness alone
cannot supply the missing simultaneous extraction.
On the explicit pole Cantor lamination, exact scalar interpolation
and Lusin localization now handle arbitrary bounded Borel radial
graphs. Arbitrary Borel subsets with general fiber geometry still
require a joint extension or selection mechanism; independent
measurable fiber optimizers do not supply that mechanism.

\item \textbf{Classical feasibility of nonlinear Dirichlet data.}
Supply verifiable admissible directions or a separate construction
inside the actual classical constraint. The exact actual variation
radius separates current-order slack from pointwise obstacle slack.
For intersecting coordinate carriers, the depths of intrinsic unique solutions must agree at their
common actual point; the nonconvexity example violates this compatibility.
The radial coordinate-carrier theorems give a necessary and sufficient
finite-equal-depth condition for arbitrary finite radial data, with
depth given by the integral of the $q$th root of the weighted tail
mass and an explicit max extension when it holds. When every individual depth is finite, a bounded completed solution
exists even if the depths differ; no necessity for arbitrary completed
solutions is asserted. The relative energy solution operator
solves a larger linear problem. Hessian root concavity and
hyperbolicity cannot be assumed for every positive background.
\end{enumerate}

\subsection{Excluded shortcuts}

The following implications have deliberately not been used: ordinary
current equality implies equality for arbitrary delta-test functionals;
ambient cycle relations remove every analytic pullback relation;
generator formulas alone cover smooth flat coefficient remainders;
energy capacity and BT capacity have identical null sets; capacity
domination characterizes the normalized measure image; a pointwise
obstacle gap supplies a classical psh variation; or an arbitrary
positive background has the usual relative Hessian hyperbolicity.
Each failure or missing hypothesis is recorded with its own example
or proof gate above. The project still defines no wedge of two
completely arbitrary positive closed currents and claims no blanket
replacement of GK's nonclosed, nonpositive, off-diagonal calculus.

% END source: sections/round6_status.tex


% BEGIN canonical R6 source: bibliography.tex
\begin{thebibliography}{CLD25}
\bibitem[CLD12]{CLD12}
A.~Chambert-Loir and A.~Ducros,
\emph{Formes diff\'erentielles r\'eelles et courants sur les espaces
de Berkovich}, arXiv:1204.6277v1 (2012).
\url{https://arxiv.org/pdf/1204.6277v1}.

\bibitem[CLD25]{CLD25}
A.~Chambert-Loir and A.~Ducros,
\emph{Formes diff\'erentielles r\'eelles et courants sur les espaces
de Berkovich}, arXiv:1204.6277v2 (2025), 608 pp.
\url{https://arxiv.org/pdf/1204.6277v2}.

\bibitem[DE15]{DE15}
A.~Dhouib and F.~Elkhadhra,
\emph{Complex Hessian operator, $m$-capacity, Cegrell's classes
and $m$-potential associated to a positive closed current},
checked preprint arXiv:1504.03519v1 (2015).
\url{https://arxiv.org/pdf/1504.03519v1}.

\bibitem[EE21]{EE21}
H.~Elaini and F.~Elkhadhra,
\emph{Complex Hessian operator associated to an $m$-positive
closed current and weighted $m$-capacity},
checked preprint arXiv:1912.00763v3 (2021).
\url{https://arxiv.org/pdf/1912.00763v3}.

\bibitem[Est12]{Est12}
A.~Esterov,
\emph{Tropical varieties with polynomial weights and corner loci
of piecewise polynomials}, Moscow Math. J. \textbf{12} (2012),
no.~1, 55--76. Checked preprint arXiv:1012.5800v3.
\url{https://arxiv.org/pdf/1012.5800v3}.

\bibitem[GK17]{GK17}
W.~Gubler and K.~K\"unnemann,
\emph{A tropical approach to non-Archimedean Arakelov geometry},
Algebra Number Theory \textbf{11} (2017), no.~1, 77--180.
Checked preprint arXiv:1406.7637v3.
\url{https://arxiv.org/pdf/1406.7637v3}.

\bibitem[GS21]{GS21}
A.~Gross and F.~Shokrieh,
\emph{Cycles, cocycles, and duality on tropical manifolds},
Proc. Amer. Math. Soc. \textbf{149} (2021), 2429--2444.
Checked preprint arXiv:1910.04805v2.
\url{https://arxiv.org/pdf/1910.04805v2}.

\bibitem[GR25]{GR25}
W.~Gubler and J.~Rabinoff,
\emph{Classically psh and pluriharmonic functions on Berkovich spaces},
checked preprint arXiv:2506.13548v2 (17 September 2025).
\url{https://arxiv.org/pdf/2506.13548v2}.

\bibitem[JY16]{JY16}
A.~Jensen and J.~Yu,
\emph{Stable intersections of tropical varieties},
J. Algebraic Combin. \textbf{43} (2016), no.~1, 101--128.
Checked preprint arXiv:1309.7064v3.
\url{https://arxiv.org/pdf/1309.7064v3}.

\bibitem[Mih23]{Mih23}
A.~Mihatsch, \emph{On tropical intersection theory},
Selecta Math. (N.S.) \textbf{29} (2023), article 17.
Checked preprint arXiv:2107.12067v3 (2022).
\url{https://doi.org/10.1007/s00029-022-00818-7}.

\bibitem[Mih24]{Mih24}
A.~Mihatsch, \emph{$\delta$-Forms on Lubin--Tate Spaces},
arXiv:2112.10018v4, September 7, 2024.
\url{https://arxiv.org/pdf/2112.10018v4}.

\bibitem[Wan19]{Wan19}
V.~Wanner,
\emph{Comparison of two notions of subharmonicity on non-archimedean
curves}, Math. Z. \textbf{293} (2019), 443--474.
Checked preprint arXiv:1801.04713v3 (8 November 2018).
\url{https://arxiv.org/pdf/1801.04713v3}.

\bibitem[Wan22]{Wan22}
V.~Wanner,
\emph{Energy Minimization Principle for non-archimedean curves},
J. Th\'eor. Nombres Bordeaux \textbf{34} (2022), no.~1, 1--39.
\url{https://www.numdam.org/item/JTNB_2022__34_1_1_0.pdf}.
\end{thebibliography}

% END source: bibliography.tex

\clearpage
\appendix

% BEGIN canonical R6 source: logs/capacity_convex_compact_barycenter_failure.tex
\section{The capacity gap for convex compact potential families}
\label{sec:r6-convex-compact-capacity-gap}

This is a supplemental consequence of the accepted restricted-capacity
example. Convexity of the entire normalized potential family does not
repair its Borel outer regularity. The unrestricted relative BT capacity
is not determined by this construction.

\begin{proposition}[Convex compact potential families with a Borel gap]
\label{prop:r6-convex-compact-capacity-gap}
In the setting of Proposition~\ref{prop:r6-compact-dominator-all-degrees},
for every $p,q\ge1$ one may take the upper-potential family and the
entire associated normalized competitor family to be convex and compact,
while their restricted capacity still satisfies
\[
                         C(E)=\tfrac12,\qquad C^*(E)=1
\]
on the same $F_\sigma$ set in a fixed compact core. Compactness of the
competitor family is in local uniform topology; its actual degree-$q$
measure image is weakly compact.
\end{proposition}

\noindent\textbf{Status: \texttt{DIRECT\_COROLLARY}} of the
accepted product example, its mixed order theorem, and compact
upper-potential compactness. The single varying Hessian direction is
essential to the linear barycenter formula proved below.

\begin{proof}
Use exactly the original field, domain, background and coordinates
of Proposition~\ref{prop:r6-compact-dominator-all-degrees}:
\[
 \begin{gathered}
 \Omega=\{-1<x_i,y_a<2\}\subset
                 (\mathbf G_{m,K}^{p+q})^{\rm an},\qquad
 A=\{0\}\cup\{1/j:j\ge1\},\\
 g_s=\frac{(x_1-s)_++\sum_{i=2}^q(x_i)_+}{3q},\qquad
 S=\frac{(3q)^q}{q!}
       \left[\bigwedge_{a=1}^p d'y_a\wedge d''y_a\right].
 \end{gathered}
\]
For a probability measure $\lambda$ on the compact real set $A$, put
\[
 g_\lambda(z)=\int_A g_s(z)\,d\lambda(s),\qquad
 \overline{\mathscr G}=
               \{g_\lambda:\lambda\in\operatorname{Prob}(A)\}.
\]
These integrals are ordinary real-valued integrals. They do not
require an integration operation on psh classes over another field.

We first verify actual classicality and compactness. Collapse the
tail $\{1/j:j>N\}$ of a probability $\lambda$ onto $0$, retaining
the first $N$ masses, and call the resulting finite probability
$\lambda_N$. The uniform Lipschitz estimate of the product example
gives, on the entire $\Omega$,
\[
       \|g_{\lambda_N}-g_\lambda\|_\infty
                    \le\frac{1}{3q(N+1)}.
\]
Each $g_{\lambda_N}$ is a finite convex combination of continuous
actual classical functions. The uniform classical-closure property
therefore makes $g_\lambda$ continuous classical. The same bound
$0\le g_\lambda<2/3$ holds globally.

The map $\lambda\mapsto g_\lambda$ is continuous from weak
probability convergence to the global uniform norm. Indeed, the
first $N$ isolated singleton masses are continuous weak-probability
coordinates, because their indicator functions are continuous on
$A$. Thus the finite collapsed barycenters converge uniformly for
fixed $N$, and the displayed uniform tail estimate is independent
of the probability. Compactness of $\operatorname{Prob}(A)$ now
proves compactness of $\overline{\mathscr G}$ in the global uniform
norm. Its convexity follows from
$g_{t\lambda+(1-t)\kappa}=t g_\lambda+(1-t)g_\kappa$.

Let $\sigma$ be the original full Gauss section. The exact top
measure is
\begin{equation}
 T_{S,q}(g_\lambda,\ldots,g_\lambda)
   =\nu_\lambda
   :=\sigma_*\bigl(\lambda\otimes\delta_0^{\otimes(q-1)}
                                      \otimes dy_1\cdots dy_p\bigr)
   =\int_A\nu_s\,d\lambda(s).
 \label{eq:r6-convex-barycenter-top-measure}
\end{equation}
Here every passive coordinate has range $(-1,2)$. To justify this
identity without a false general nonlinear barycenter rule, write
\[
 a_\lambda(x_1)=\int_A(x_1-s)_+\,d\lambda(s).
\]
Only this first-coordinate term varies with $\lambda$; the other
$q-1$ hinge terms remain fixed. Smooth convex one-variable
approximations show that two curvature factors in the same
coordinate wedge to zero. This holds for bounded continuous BT
products by their local uniform continuity. Consequently, in the
$q$th power of the sum only the terms using each of the $q$
coordinates once remain. There are $q!$ such terms, and the
coefficient is exactly
\[
                     \frac{(3q)^q}{q!}\frac{q!}{(3q)^q}=1.
\]
For finite $\lambda_N$, linearity of the surviving first-coordinate
factor and the accepted product calibration give
\eqref{eq:r6-convex-barycenter-top-measure}. Passing $N\to\infty$
uses the global uniform convergence just proved and fixed-background
continuous BT convergence. On the measure side it uses weak
probability convergence and continuity of the original Gauss section;
pairings with every original compact smooth scalar test converge.
Thus the identity is on all original tests, not only on tests which
factor through the displayed radius coordinates. The product measures
are finite Radon, and the same identity gives the displayed Borel
barycenter formula.

Now take the entire normalized family
\[
 \mathscr P(\overline{\mathscr G})=
 \{u\in\mathcal C_b(\Omega):0\le u\le1,
               \quad D[u]\le D[g_\lambda]
                         \text{ for some }\lambda\in\operatorname{Prob}(A)\}.
\]
It is convex: if $u,v$ have upper probabilities $\lambda,\kappa$,
then $tu+(1-t)v$ has upper probability $t\lambda+(1-t)\kappa$.
Every member is continuous by
Corollary~\ref{cor:r6-continuity-from-continuous-curvature-domination}.
Proposition~\ref{prop:r6-compact-upper-potential-family} therefore
proves local-uniform compactness of this whole family. Its degree-$q$
measure image is weakly compact by
Corollary~\ref{cor:r6-compact-upper-potential-stability}.

For each such $u$, mixed order replacement in all slots gives
\[
              0\le T_{S,q}(u,\ldots,u)\le\nu_\lambda.
\]
For every Borel $B\subset\Omega$, it follows that
\[
 \begin{split}
 C_{\overline{\mathscr G}}(B)
   &\le\sup_{\lambda\in\operatorname{Prob}(A)}\nu_\lambda(B)\\
   &=\sup_{s\in A}\nu_s(B)
    =C_{\mathscr G}(B).
 \end{split}
\]
The scalar function $s\mapsto\nu_s(B)$ is Borel by the product
kernel formula, so the middle equality follows from averaging and
Dirac probabilities. Conversely every $g_s$ is a member of the
enlarged family and attains $\nu_s$. Hence equality holds throughout
for all Borel $B$, in particular for every ordinary open set. The
raw and outer capacities are therefore identical to those in the
accepted all-degree example, which gives the claimed $1/2$ and $1$.
\end{proof}

\paragraph{Exact limitation.}
Convexity here refers to the actual normalized potential family and
its upper family. For $q=1$ its measure image is also convex, by
linearity; no such nonlinear image-convexity assertion is needed or
made for $q>1$. Closure under arbitrary maxima is not asserted.
The example still concerns a restricted competitor class and does
not establish failure of unrestricted relative BT C1 or C2.

% END source: logs/capacity_convex_compact_barycenter_failure.tex

\end{document}
